%\documentclass[a4paper,9pt]{article}
\documentclass[12pt, a4paper]{article}
\usepackage{amsthm}
\usepackage{amsmath}
\usepackage{amssymb}
\usepackage[T1,T2A,T8K,T8M]{fontenc}
\usepackage[utf8]{inputenc}
\usepackage{tkz-fct}
\usepackage{tikz-cd}
\usepackage{tikz}
\usetikzlibrary{math}
\usepackage{pgfplots}
\pgfplotsset{compat=1.17}
\usepackage[english,russian,georgian]{babel}
%\usepackage{index}
\usepackage{imakeidx}
\usepackage[columns=3,totoc=true]{idxlayout}
\usepackage{nomencl}
\usepackage[plainpages=false,bookmarksopen,pdfpagelabels,bookmarksnumbered,unicode=true,hyperindex=true,pagebackref=true]{hyperref}
\usepackage{graphicx}
\usepackage{nameref}
\newcommand{\Index}[1]{\index{#1}#1}
\hyphenation{ჯამ-ში მუხ-ტი-სათ-ვის გან-სა-კუთ-რე-ბუ-ლი}
\makeindex[intoc]
\usepackage[nottoc]{tocbibind}
\makenomenclature
%\renewindex{default}{idx}{ind}{საძიებელი}
\renewcommand\nomname{აღნიშვნები}
\newtheorem{theorem}{თეორემა}
\newtheorem{lemma}{ლემა}
\newtheorem{corollary}{შედეგი}
\newtheorem{remark}{შენიშვნა}
\newtheorem{ex}{მაგალითი }
\newtheorem{cons}{კონსტრუქცია}
\newtheorem{convention}{შეთანხმება}
\newtheorem{definition}{განსაზღვრება}
\newtheorem{proposition}{წინადადება}
\newtheorem{example}{}
\newtheorem{Example}{მაგალითი}
\newtheorem{exercise}{სავარჯიშო}
\definecolor{boundarycolor}{cmyk}{1,0,0,0}%Cyan
\definecolor{thirdcolor}{cmyk}{0.82,0.00,0.04,0.30}%LightSeaGreen
\definecolor{fillcolor}{cmyk}{0.12,0,0,0}%LightCyan
\DeclareMathOperator{\rank}{rank}
\DeclareMathOperator{\Hom}{Hom}
\DeclareMathOperator{\id}{id}


\begin{document}
	
	
	
	
	\thispagestyle{empty} 
	
	\begin{tikzpicture}[remember picture, overlay]
		% ფონი იხატება მთელი გვერდის ზომაზე
		\shade[shading=axis, top color=blue!20, bottom color=blue!80]
		(current page.south west) rectangle (current page.north east);
		
		% --- 1. ჯერ იხატება 3D ტორი, რომ ის ფონზე იყოს ---
		\node at (current page.center) {
			\begin{axis}[
				width=0.8\paperwidth,
				height=0.8\paperwidth,
				view={120}{30},
				hide axis,
				colormap/viridis,
				]
				\addplot3[
				surf,
				shader=flat,
				domain=0:360,
				y domain=0:360,
				samples=40,
				variable=\u,
				variable y=\v,
				point meta=u
				]
				({(2 + 0.7*cos(u)) * cos(v)}, 
				{(2 + 0.7*cos(u)) * sin(v)}, 
				{0.7*sin(u)});
			\end{axis}
		};
		
		% --- 2. შემდეგ ტექსტი იხატება ტორის გამოსახულებაზე ზემოდან ---
		\node[font=\Huge\bfseries\sffamily, white] at ([yshift=5cm]current page.center) {დიფერენციალური გეომეტრია};
		\node[font=\large\sffamily, white] at ([yshift=3.5cm]current page.center) {პირველი კურსი};
		\node[font=\sffamily\large, white] at (current page.south) [yshift=3cm] {მალხაზ ბაკურაძე};	
	\end{tikzpicture}
	
	\newpage
	
	
	
	
	
	
	
	\sloppy
	\title{დიფერენციალური  გეომეტრია\\
	პირველი კურსი}
	% repeat info for each author  .
	\author{მალხაზ ბაკურაძე\\
		\small ალგებრის და გეომეტრიის  კათედრა\\[-0.8ex]
		\small ზუსტ და საბუნებისმეტყველო მეცნიერებათა ფაკულტეტი\\[-0.8ex]
		\small ივ. ჯავახიშვილის თბილისის სახელმწიფო   უნივერსიტეტი   \\[-0.8ex]
		%\small \textsf{malkhaz.bakuradze@tsu.ge}\\
	}
	\date{თბილისი  2025}
	%\netaddress{levan.shoshiashvili (at) tsu.ge}
	\maketitle
	
	\ 
	
	\
	
	\
	
	\
	
	\
	
	\
	
	\
	
%	\newpage
	\vfill
	\begin{center}
		თეონას
	\end{center}
	\vfill
	
	
	
	\newpage
	
	\tableofcontents
	
	
	
	\newpage

		
		%======================================================================
		\part{კლასიკური დიფერენციალური გეომეტრია}
		%=====================================================================
		%\addcontentsline{toc}{section}{შესავალი და ისტორიული მიმოხილვა}
		\setcounter{page}{1}
		
		\subsection*{რას შეისწავლის დიფერენციალური   გეომეტრია?}
		დიფერენციალური გეომეტრია არის მათემატიკის დარგი, რომელიც იყენებს კალკულუსს, წრფივ ალგებრასა და ანალიზს გეომეტრიული ობიექტების — წირების, ზედაპირებისა და მათი მაღალგანზომილებიანი ანალოგების (მრავალნაირობების) — შესასწავლად. კლასიკური დიფერენციალური გეომეტრია, რომელსაც ამ კურსში განვიხილავთ, ძირითადად ფოკუსირებულია სამგანზომილებიან ევკლიდურ სივრცეში ($ \mathbb{R}^3 $) არსებული ობიექტების თვისებებზე.
		
		ჩვენი კვლევის მთავარი ობიექტი იქნება \textbf{ზედაპირი}. ჩვენ ვუპასუხებთ ისეთ კითხვებს, როგორიცაა:
		\begin{itemize}
			\item როგორ განვსაზღვროთ და გავზომოთ „სიმრუდე“?
			\item რა განსხვავებაა ზედაპირის შინაგან თვისებებსა (რომელთა გაზომვაც ზედაპირზე მცხოვრებ დამკვირვებელს შეუძლია) და გარეგან თვისებებს (რომლებიც დამოკიდებულია იმაზე, თუ როგორაა ზედაპირი ჩადგმული სივრცეში) შორის?
			\item რა არის „სწორი ხაზის“ ანალოგი მრუდე ზედაპირზე?
			\item როგორ უკავშირდება ზედაპირის ლოკალური გეომეტრია მის გლობალურ, ტოპოლოგიურ ფორმას?
		\end{itemize}
		ამ კითხვებზე პასუხის გასაცემად, ჩვენ შემოვიტანთ ორ ფუნდამენტურ ინსტრუმენტს: \textbf{პირველ} და \textbf{მეორე ფუნდამენტურ ფორმებს}.
		
		\subsection*{მოკლე  ისტორიული   მიმოხილვა}
		დიფერენციალური გეომეტრიის ფესვები ანტიკურ პერიოდშია, თუმცა მისი თანამედროვე ფორმით განვითარება მე-17 საუკუნეში, კალკულუსის აღმოჩენასთან ერთად დაიწყო.
		\begin{itemize}
			\item \textbf{ლეონარდ ეილერი (1707-1783)} ითვლება ზედაპირთა თეორიის ერთ-ერთ ფუძემდებლად. მან შემოიტანა ზედაპირის პარამეტრიზაციის იდეა და პირველმა შეისწავლა ნორმალური სიმრუდე.
			\item \textbf{გასპარ მონჟი (1746-1818)} ითვლება დიფერენციალური გეომეტრიის მამად. მან სისტემატურად გამოიყენა კალკულუსი სივრცული წირებისა და ზედაპირების შესასწავლად.
			\item \textbf{კარლ ფრიდრიხ გაუსი (1777-1855)} არის ცენტრალური ფიგურა. მისმა ნაშრომმა \textit{Disquisitiones Generales circa Superficies Curvas} (1827) რევოლუცია მოახდინა საგანში. მან შემოიტანა შინაგანი გეომეტრიის ცნება და დაამტკიცა თავისი \textit{Theorema Egregium}, რითაც გზა გაუხსნა გეომეტრიის უფრო აბსტრაქტულ გაგებას.
			\item \textbf{ბერნჰარდ რიმანი (1826-1866)}, გაუსის მოსწავლემ, განაზოგადა ეს იდეები $n$-განზომილებიან სივრცეებზე (რიმანის მრავალნაირობები), რითაც საფუძველი ჩაუყარა თანამედროვე დიფერენციალურ გეომეტრიას და იმ მათემატიკურ აპარატს, რომელიც მოგვიანებით აინშტაინმა ფარდობითობის ზოგად თეორიაში გამოიყენა.
		\end{itemize}
		
		\subsection*{კურსის სტრუქტურა}		
		კურსი დაყოფილია რამდენიმე ლოგიკურ ნაწილად. თავდაპირველად, ჩვენ განვიხილავთ წირების თეორიას, რადგან ის უფრო მარტივ, ერთგანზომილებიან ანალოგიას წარმოადგენს ზედაპირების შესწავლისთვის. შემდეგ, დეტალურად შევისწავლით პირველ და მეორე ფუნდამენტურ ფორმებს, რომლებიც ზედაპირის გეომეტრიის სრულ აღწერას გვაძლევს. ამის შემდეგ, განვიხილავთ კლასიკური თეორიის უმნიშვნელოვანეს თეორემებს. კურსის მსვლელობისას, ჩვენ დეტალურად გავაანალიზებთ კონკრეტულ მაგალითებს,  რათა თეორიული კონცეფციები უფრო  თვალსაჩინო  გახდეს.
		
		\newpage
		
		
			
		\section{ბრტყელი  წირები. \\
			პარამეტრული  აღწერა, რკალის  სიგრძე \\და ნატურალური პარამეტრი}
		
		
	ანალიზურ გეომეტრიაში ხშირად ვიყენებთ წირის განტოლებას $y=f(x)$ სახით. თუმცა, ეს ფორმა ყოველთვის მოსახერხებელი არ არის, განსაკუთრებით შეკრული ფორმის  მქონე წირებისთვის (მაგალითად, წრეწირი ან ელიფსი).
			
	უფრო ზოგადი და მძლავრი მიდგომაა წირის პარამეტრული  აღწერა. ამ დროს, სიბრტყეზე წერტილის $x$ და $y$ კოორდინატები გამოისახება მესამე, დამხმარე ცვლადით, რომელსაც  პარამეტრი ეწოდება და აღინიშნება, როგორც წესი, $t$ ასოთი.
			
	წარმოვიდგინოთ, რომ წირი არის სიბრტყეზე მოძრავი წერტილის ტრაექტორია, ხოლო პარამეტრი $t$ — დრო. დროის ყოველ $t$ მომენტში წერტილი იმყოფება სიბრტყის კონკრეტულ $(x(t), y(t))$ ადგილას.
			
	\subsection{განსაზღვრება}
			ბრტყელი წირის პარამეტრული განტოლება არის ფუნქციათა სისტემა:
			\begin{equation}
				\begin{cases}
					x = x(t) \\
					y = y(t)
				\end{cases}
				\quad t \in [a, b]
			\end{equation}
			სადაც $t$ არის პარამეტრი, რომელიც იცვლება $[a, b]$ შუალედში. ხშირად, ამას ვექტორული სახით წერენ:
			$$ \vec{r}(t) = (x(t), y(t)) \quad \text{ან} \quad \vec{r}(t) = x(t)\vec{i} + y(t)\vec{j} $$
			სადაც $\vec{r}(t)$ არის წერტილის რადიუს-ვექტორი.
			
			\subsection*{მაგალითი: წრეწირი}
			განვიხილოთ $R$ რადიუსიანი წრეწირი, რომლის ცენტრი კოორდინატთა სათავეშია. მისი პარამეტრული განტოლებაა:
			\begin{equation}
				\begin{cases}
					x(t) = R \cos(t) \\
					y(t) = R \sin(t)
				\end{cases}
				\quad t \in [0, 2\pi]
			\end{equation}
			აქ პარამეტრი $t$ არის კუთხე, რომელსაც რადიუს-ვექტორი ქმნის $Ox$ ღერძის დადებით მიმართულებასთან.
			
			\begin{center}
				\begin{tikzpicture}[scale=1.5]
					\draw[->, gray] (-2,0) -- (2,0) node[right] {$x$};
					\draw[->, gray] (0,-2) -- (0,2) node[above] {$y$};
					\draw[blue, thick] (0,0) circle (1.5); % წრეწირი R=1.5
					
					\def\t{40} % კუთხე t
					\draw[->, thick, red] (0,0) -- ({1.5*cos(\t)}, {1.5*sin(\t)}) node[midway, above right] {$\vec{r}(t)$};
					\draw[dashed] ({1.5*cos(\t)}, 0) -- ({1.5*cos(\t)}, {1.5*sin(\t)}) node[midway, right] {$y(t)$};
					\draw[dashed] (0, {1.5*sin(\t)}) -- ({1.5*cos(\t)}, {1.5*sin(\t)});
					\draw[->] (0.5,0) arc (0:\t:0.5);
					\node at ({0.7*cos(\t/2)}, {0.7*sin(\t/2)}) {$t$};
					\node at ({1.5*cos(\t)}, -0.2) {$x(t)$};
				\end{tikzpicture}
			\end{center}
			
			
			
			\subsection{რკალის  სიგრძე}
			
			როგორ გამოვთვალოთ წირის ნაწილის, ანუ რკალის, სიგრძე, თუ ის პარამეტრულად არის მოცემული? იდეა მდგომარეობს იმაში, რომ წირი დავყოთ ძალიან მცირე მონაკვეთებად, რომლებიც თითქმის სწორხაზოვანია.
			
			პარამეტრის მცირე ცვლილებას, $dt$-ს, შეესაბამება კოორდინატების მცირე ცვლილებები $dx$ და $dy$. რკალის უსასრულოდ მცირე ელემენტის, $ds$-ის, სიგრძე პითაგორას თეორემით იქნება:
			$$ (ds)^2 = (dx)^2 + (dy)^2 $$
			გავყოთ ეს გამოსახულება $(dt)^2$-ზე:
			$$ \left(\frac{ds}{dt}\right)^2 = \left(\frac{dx}{dt}\right)^2 + \left(\frac{dy}{dt}\right)^2 $$
			აქედან გამომდინარე:
			$$ \frac{ds}{dt} = \sqrt{x'(t)^2 + y'(t)^2} $$
			ეს გამოსახულება გვიჩვენებს "სიჩქარეს", რომლითაც წერტილი მოძრაობს წირზე. ვექტორული ფორმით, ეს არის $\vec{r}'(t) = (x'(t), y'(t))$ ვექტორის სიგრძე: $\|\vec{r}'(t)\|$.
			
			რკალის სრული სიგრძე $L$, როცა პარამეტრი $t$ იცვლება $a$-დან $b$-მდე, მიიღება $ds$ ელემენტების შეჯამებით, ანუ ინტეგრებით:
			\begin{equation}
				L = \int_{a}^{b} \sqrt{x'(t)^2 + y'(t)^2} \, dt = \int_{a}^{b} \|\vec{r}'(t)\| \, dt
			\end{equation}
			
			\subsubsection*{მაგალითი: წრეწირის სიგრძე}
			გამოვთვალოთ $R$ რადიუსიანი წრეწირის სრული სიგრძე. ჩვენ ვიცით, რომ:
			$$ x(t) = R \cos(t) \quad \implies \quad x'(t) = -R \sin(t) $$
			$$ y(t) = R \sin(t) \quad \implies \quad y'(t) = R \cos(t) $$
			ჩავსვათ ეს წარმოებულები ფორმულაში:
			$$ L = \int_{0}^{2\pi} \sqrt{(-R \sin t)^2 + (R \cos t)^2} \, dt $$
			$$ L = \int_{0}^{2\pi} \sqrt{R^2 \sin^2 t + R^2 \cos^2 t} \, dt = \int_{0}^{2\pi} \sqrt{R^2(\sin^2 t + \cos^2 t)} \, dt $$
			რადგან $\sin^2 t + \cos^2 t = 1$, მივიღებთ:
			$$ L = \int_{0}^{2\pi} \sqrt{R^2} \, dt = \int_{0}^{2\pi} R \, dt = R [t]_{0}^{2\pi} = R(2\pi - 0) = 2\pi R $$
			შედეგი, როგორც მოსალოდნელი იყო, დაემთხვა წრეწირის სიგრძის ცნობილ ფორმულას.
			
			\subsection{ნატურალური   პარამეტრი}
			
			პარამეტრი $t$ ხშირად თვითნებურია. ერთი და იგივე წირი შეიძლება სხვადასხვაგვარად იყოს პარამეტრიზებული. მაგალითად, წერტილმა შეიძლება წრეწირს ორჯერ უფრო სწრაფად შემოუაროს.
			
			არსებობს პარამეტრი, რომელიც წირის შინაგანი, გეომეტრიული თვისებაა და არ არის დამოკიდებული იმაზე, თუ "რა სისწრაფით" ვივლით მასზე. ასეთ პარამეტრს ნატურალური პარამეტრი ეწოდება და მას, როგორც წესი, $s$ ასოთი აღნიშნავენ.
			
			\textbf{ნატურალური  პარამეტრი $s$ არის რკალის სიგრძე, რომელიც ათვლილია  წირის რაღაც ფიქსირებული  საწყისი წერტილიდა ნ}.
			
			მისი განმსაზღვრელი ფუნქციაა:
			\begin{equation}
				s(t) = \int_{a}^{t} \sqrt{x'(\tau)^2 + y'(\tau)^2} \, d\tau = \int_{a}^{t} \|\vec{r}'(\tau)\| \, d\tau
			\end{equation}
			აქ $\tau$ ინტეგრების ცვლადია, რათა არ აგვერიოს ზედა ზღვარში, $t$-ში.
			
			\subsubsection{ნატურალური პარამეტრის მთავარი თვისება}
			თუ წირი პარამეტრიზებულია რკალის სიგრძით $s$, ანუ $\vec{r} = \vec{r}(s)$, მაშინ მისი სიჩქარის ვექტორის სიგრძე ყოველთვის 1-ის ტოლია:
			$$ \left\| \frac{d\vec{r}}{ds} \right\| = 1 $$
			ეს ნიშნავს, რომ ნატურალური პარამეტრით მოძრაობისას, წერტილი ყოველთვის "ერთეულოვანი სიჩქარით" გადაადგილდება. ეს თვისება ძალიან ამარტივებს დიფერენციალური გეომეტრიის ბევრ ფორმულას (მაგალითად, სიმრუდის ფორმულას).
			
			\subsubsection{რეპარამეტრიზაცია}
			წირის ხელახლა პარამეტრიზაცია ნატურალური პარამეტრით ნიშნავს, რომ $t$ პარამეტრის ნაცვლად უნდა გამოვიყენოთ $s$. ამისათვის საჭიროა:
			\begin{enumerate}
				\item ვიპოვოთ $s(t)$ ფუნქცია.
				\item ამ განტოლებიდან ვიპოვოთ $t$, როგორც $s$-ის ფუნქცია: $t = t(s)$.
				\item ჩავსვათ მიღებული $t(s)$ გამოსახულება საწყის პარამეტრულ განტოლებებში: $\vec{r}(t(s))$.
			\end{enumerate}
			მაგალითად, ჩვენი წრეწირისთვის:
			$$ s(t) = \int_{0}^{t} R \, d\tau = R t $$
			აქედან, $t = s/R$. ჩავსვათ ეს საწყის განტოლებებში:
			$$
			\begin{cases}
				x(s) = R \cos(s/R) \\
				y(s) = R \sin(s/R)
			\end{cases}
			$$
			ეს არის წრეწირის განტოლება, რომელიც პარამეტრიზებულია ნატურალური პარამეტრით.
			

		
		
		

		
		
		\begin{figure}[h!]
			\centering
			\begin{tikzpicture}
				\draw[gray, thick, ->] (-4,0) -- (7,0) node[anchor=west]{X};
				\draw[gray, thick, ->] (0,-3) -- (0,5) node[anchor=south]{Y};
				\draw [line width=1pt] (-1,-1) .. controls (1,1) and (2,6) .. (5,-1);
				\draw [line width=1pt] (-1,-1) -- (-1.3,-1.2);
				\draw[gray, thick, ->] (-1,-1) --(-0.2,-0.35); 	
				\draw[gray, thick, ->] (-1,-1) --(-2.35,0.2);
				\filldraw[black] (-1,-1) circle (2pt) node[anchor=north west]{P};
				\filldraw[black] (5,-1) circle (2pt) node[anchor=south west]{Q};
				\draw[gray, thick] (5,-1) --(3.5,2.5);
			\end{tikzpicture}
			\caption{წირის მხები და ნორმალური ვექტორები $P$ წერტილში. მხები წრფე $Q$ წერტილში}
			\label{fig:tangent-p}
		\end{figure}
		
		
		\bigskip
		
		
		შევთანხმდეთ, რომ $s$ წირის ნატურალური  პარამეტრია, ხოლო  $t$ ნებისმიერი პარამეტრი.
		
		თუ $c(t)$ არის ბრტყელი  გლუვი რეგულარული   წირი, ე.ი. თუ $c(t)$ ორჯერ უწყვეტად დიფერენცირებადია  და ყოველი  პარამეტრისთვის $|c'(t)|\neq 0$, მაშინ $c(s)$-ს ერთეულოვანი  მხები ვექტორია $T:=c'(s)$, ხოლო  ერთეულოვანი  ნორმალური  ვექტორი $N$ მიიღება $T$ ვექტორის მობრუნებით  საათის ისრის საწინააღმდეგო მიმართულებით $\pi/2$ მართი კუთხით.
		\[
		|c'|=1\implies c'\cdot c'=1\implies 0=(c'\cdot c')'=2c'\cdot c''\implies c''\perp T.
		\]
		ე.ი $c''$ პარალელურია $N$ ვექტორის, ანუ $c''(s)=k(s)N(s)$, სადაც $k$ რაღაც ფუნქციაა. ამ $k$ ფუნქციას ეწოდება $c$ წირის ორიენტირებული სიმრუდე. მისი ნიშანი მიუთითებს თუ მოცემულ წერტილში საით უხვევს წირი (ან მისი მხები). თუ $k>0$, მაშინ მხები მიდის მარცხნივ, მაშინ როცა $k<0$ მიუთითებს, რომ მხები ბრუნავს მარჯვნივ. წერტილს, სადაც $k$ იცვლის ნიშანს, ეწოდება გადაღუნვის წერტილი.
		
		\begin{definition}
			\label{plane-curvature}
			ლიტერატურაში უფრო ხშირად განისაზღვრება $c(s)$ ბრტყელი წირის \textbf{არაორიენტირებული სიმრუდე}, როგორც
			\[k(s)=|T'(s)|=|c''(s)|\]
			და \textbf{სიმრუდის რადიუსი} $R=1/k$.
		\end{definition}
		
		\medskip
		
		\begin{example}
			\label{ex:1}
			წრფე
		\end{example}
		${\bf q}=(q_1,q_2,q_3)\neq 0 $ ვექტორის მიმართულებით ${\bf p}=(p_1,p_2,p_3)$
		წერტილზე გამავალი წრფის განტოლებაა
		\[\alpha(t)={\bf p}+t {\bf q}=(p_1+tq_1,p_2+tq_2,p_3+tq_3).\]
		
		
		\medskip
		
		\begin{example}
			სურ. \ref{*1} მიხედვით დაწერეთ ერთეულოვანი წრეწირიდან $(-1,0)$ წერტილის
			ამოშლით მიღებული $\tilde{S}^1$ წირის პარამეტრული განტოლება $h$ სიმაღლით. 	
		\end{example}
		
		\begin{figure}[h!]
			\centering
			\begin{tikzpicture}
				\begin{axis}[
					trig format plots=rad,
					samples=100,
					axis equal image,
					xtick=\empty, ytick=\empty,
					axis lines=middle, enlargelimits=true
					]
					\addplot [
					no markers,
					red,
					domain=0:2*pi
					] ( {cos(x)}, { sin(x) } );
					
					\draw (0.2,0.2) node[below] {$t$};
					
					\draw (-1,0) node[below] {$(-1,0)$};
					\draw[dashed] (0,0) -- (0.5,0.866) node[right] {$P$};
					\draw [red, thick](-1,0) -- (0.5,0.866);
					\draw [line width=2pt](0,0) -- (0,0.6);
					\draw (0,0.3) node[left] {$h$};
					\draw (-0.8,0.1) node[right] {$t/2$};
				\end{axis}
			\end{tikzpicture}
			\caption{$\tilde{S}^1$ წირი პარამეტრით $h$ სიმაღლე: $r(h)=(\frac{1-h^2}{1+h^2},\frac{2h}{1+h^2})$.}
			\label{*1}
		\end{figure}
		
		\medskip	
		
		\begin{example}
			\label{ex:3}
			დიფერენცირებადი, მაგრამ არარეგულარული ბრტყელი წირი (იხ. ფიგურა \ref{t3t2})
		\end{example}
		\[
		\alpha(t)=(t^3,t^2)
		\]
		
		\bigskip
		
		\begin{theorem}
			რეგულარული ბრტყელი $c(s)$ წირის სიმრუდე მუდმივია მაშინ და მხოლოდ მაშინ, როცა
			წირი არის წრეწირის ნაწილი ($k>0$), ან წრფის ნაწილი ($k=0$).
		\end{theorem}
		\begin{proof}
			თუ $k(s_0)= 0$, მაშინ $T$ მუდმივია და გვაქვს წრფე. ვთქვათ $k(s_0)\neq 0$. $c(s)$ წრეწირის ნაწილია $\iff$ გამოსახულება $c(s)+\frac{1}{k(s_0)}N(s)$ მუდმივია (როგორც წრეწირის ცენტრი) $\iff$ $k=k(s_0)$ ყველგან, რადგან $\left(c(s)+\frac{1}{k(s_0)}N(s)\right)'=c'(s)+\frac{1}{k(s_0)}N'(s)=T(s)-\frac{k(s)}{k(s_0)}T(s)=0$. რ.დ.გ.
		\end{proof}
		
		
		

		\bigskip
		
		
		
		\bigskip
		
		
		
		\begin{figure}[h!]
			\centering
			\begin{tikzpicture}
				\begin{axis}[
					axis x line=middle,
					axis y line=middle,
					axis line style={<->,color=blue},
					xlabel={$x$},
					ylabel={$y$},
					xmin=-4,xmax=4,
					ymin=-2,ymax=4,
					]
					\addplot [domain=-2:2, samples=100, smooth]({x^3},{x^2});
				\end{axis}
			\end{tikzpicture}
			\caption{$\alpha(t)=(t^3,t^2)$}
			\label{t3t2}
		\end{figure}
		
		\medskip
		
		\begin{figure}[h!]
			\centering
			\begin{tikzpicture}[scale=2.5]
				\draw [color=red, thick, domain=0:2.7, samples=100, smooth, variable=\t] plot({1/cosh(\t)},{\t-tanh(\t)});
				\draw [color=red, thick, domain=-2.7:0, samples=100, smooth, variable=\t] plot({1/cosh(\t)},{\t-tanh(\t)});
				\draw[-{Latex},thick] (0,-2) -- (0,2) node[left] {$y$};
				\draw[-{Latex},thick] (0,0) -- (1.5,0) node[above] {$x$};
				\filldraw[blue,very thin]({1/cosh(0.8)},{0.8-tanh(0.8)}) circle[radius = 0.4pt]--++(138.39:1);
				\filldraw[blue,very thin]({1/cosh(1.3)},{1.3-tanh(1.3)}) circle[radius = 0.4pt]--++(120.48:1);
				\filldraw[blue,very thin]({1/cosh(1.8)},{1.8-tanh(1.8)}) circle[radius = 0.4pt]--++(108.77:1);
			\end{tikzpicture}
			\caption{ტრაქტრიქსი $(1/\cosh t, t-\tanh t)$}
			\label{tractrix}
		\end{figure}
		
		\medskip
		
		\begin{figure}[h!]
			\centering
			\begin{tikzpicture}
				\begin{axis}[
					axis x line=middle,
					axis y line=middle,
					axis line style={<->,color=blue},
					xlabel={$x$},
					ylabel={$y$},
					xmin=-8,xmax=8,
					ymin=-10,ymax=10,
					]
					\addplot [domain=-3.5:3.5, samples=100, smooth]({x^3-3*x},{3*x^2-9});
				\end{axis}
			\end{tikzpicture}
			\caption{$\alpha(t)=(t^3-3t, 3t^2-9)$}
			\label{t3-3t,t2}
		\end{figure}
		
		\begin{figure}[h!]
			\centering
			\begin{tikzpicture}
				\begin{axis}[
					axis x line=middle,
					axis y line=middle,
					axis line style={<->,color=blue},
					xlabel={$x$},
					ylabel={$y$},
					xmin=-pi,xmax=pi,
					ymin=0,ymax=4,
					axis equal,
					]
					\addplot [color=black, thick, domain=-pi:pi, samples=100, smooth]({x-tanh(x)},{1/cosh(x)});
					\addplot [color=red, thick, domain=-2:2, samples=100, smooth]({x},{cosh(x)});
				\end{axis}
			\end{tikzpicture}
			\caption{წითელი კატენარი $(t,\cosh t)$ და შავი ტრაქტრიქსი $(t-\tanh t, 1/\cosh t)$}
			\label{cathenar-tractrix}
		\end{figure}
		

	
		
		
		\bigskip
		
		\bigskip
		
		
		
		\subsection{წირის  ევოლუტა}
		
		ამ პარაგრაფში ჩვენ განვიხილავთ ბრტყელ წირთა სპეციალურ წყვილს, ინვოლუტა-ევოლუტას.
		
		\begin{example}
			\label{example-involute}
			სურ. \ref{involute-of-cycle} მოცემული წრეწირი კომპლექსური სახით ჩაწერილი განტოლებით
			$r(t)=e^{it}$. წრეწირის \textbf{ინვოლუტას (ან ევოლვენტას)} ყოველი $C$ წერტილი მიიღება წრეწირის $B$ წერტილიდან მხებზე $t$ კუთხის ტოლი $|BC|= \text{arc}(B,1)$ რკალის გადაზომვით.
		\end{example}
		
		\begin{figure}[h!]
			\centering
			\begin{tikzpicture}
				\begin{axis}[
					trig format plots=rad,
					samples=100,
					axis equal image,
					xtick=\empty, ytick=\empty,
					axis lines=middle, enlargelimits=true
					]
					\addplot [
					no markers, black, thick,
					domain=0:2*pi
					] ( {cos(x) + x*sin(x)}, {sin(x) - x*cos(x)} );
					
					\addplot [
					no markers, red, thick,
					domain=0:2*pi
					] ( {cos(x)}, {sin(x)} );
					
					\draw (0.3,0.3) node[below] {$\angle t$};
					\draw (1,0) node[below left] {$1$};
					\draw[dashed] (0,0) -- ({cos(pi/3)},{sin(pi/3)}) node[above] {$B=e^{it}$};
					\draw[dashed] ({cos(pi/3)},{sin(pi/3)})-- ({cos(pi/3) + (pi/3)*sin(pi/3)}, {sin(pi/3) - (pi/3)*cos(pi/3)}) node[right] {$C$};
				\end{axis}
			\end{tikzpicture}
			\caption{წრეწირის ინვოლუტა $C(t)=B+BC=e^{it}+te^{i(t-\frac{\pi}{2})}=(1-it)e^{it}$.}
			\label{involute-of-cycle}
		\end{figure}
		
		\begin{figure}[h!]
			\centering
			\begin{tikzpicture}
				\begin{axis}[
					trig format plots=rad,
					samples=100,
					axis equal image,
					xtick=\empty, ytick=\empty,
					axis lines=middle, enlargelimits=true
					]
					\addplot [no markers, black, thick, domain=0:2*pi] ( {cos(x) +x*sin(x)}, { sin(x) - x*cos(x)} );
					\addplot [no markers, green, thick, domain=0:2*pi] ( {cos(x+0.5*pi) +x*sin(x+0.5*pi)}, { sin(x+0.5*pi) - x*cos(x+0.5*pi)} );
					\addplot [no markers, purple, thick, domain=0:2*pi] ( {cos(x+0.8*pi) +x*sin(x+0.8*pi)}, { sin(x+0.8*pi) - x*cos(x+0.8*pi)} );
					\addplot [no markers, red, thick, domain=0:2*pi] ( {cos(x)}, { sin(x) } );
				\end{axis}
			\end{tikzpicture}
			\caption{წრეწირის რამდენიმე ინვოლუტა.}
			\label{3involutes-of-cycle}
		\end{figure}
		
		საზოგადოდ $r(s)$ წირის ინვოლუტის განსაზღვრება ზემოთ მოყვანილი მაგალითის ანალოგიურია, იხ. სურ. \ref{involute}
		
		\begin{figure}[h!]
			\centering
			\begin{tikzpicture}
				\begin{axis}[
					trig format plots=rad,
					samples=100,
					axis equal image,
					xtick=\empty, ytick=\empty,
					axis lines=middle, enlargelimits=true
					]
					\addplot [no markers, black, thick, domain=0:0.7*pi] ( {cos(x) +x*sin(x)}, { sin(x) - x*cos(x)} );
					\addplot [no markers, red, thick, domain=-0.25*pi:0.6*pi] ( {cos(x)}, { sin(x) } );
					
					\draw (0.3,0.3) node[below] {$s$};
					\draw (1,0) node[below] {$r(0)$};
					\draw[dashed] (0,0) -- (0.5,0.866) node[right] {$r(s)$};
					\draw[dashed] (0.5,0.866)-- (1.4,0.2) node[above] {$\mathrm{in}(s)$};
				\end{axis}
			\end{tikzpicture}
			\caption{$r(s)$ წირის ინვოლუტა $\mathrm{in}(s)=r(s)- s r'(s)$.}
			\label{involute}
		\end{figure}
		
		$r$ წირის მიმხები წრეწირების ცენტრები, ანუ სიმრუდის ცენტრები, ქმნიან წირს, რომელსაც $r$ წირის \textbf{ევოლუტა} ეწოდება და მოიცემა ფორმულით
		\[ \mathrm{ev}(u)=r(u)+\frac{1}{k(u)} N(u). \]
		
		\subsection*{კატენარი და ტრაქტრიქსი}
		
		წირთა ეს წყვილი ინვოლუტა-ევოლუტას კარგი მაგალითია.
		
		\begin{example}
			\textbf{კატენარი, ანუ ჯაჭვის წირი (catenary)}. სახელწოდება მოდის ბოლოებით ჰორიზონტალურად დაკიდული ჯაჭვიდან. კატენარის ცხადი განტოლება მოიცემა ჰიპერბოლური კოსინუსით
			\begin{equation}
				\label{catenary}
				y=a\cosh \frac{x}{a}, \quad a=\text{const}.
			\end{equation}
		\end{example}
		
		\begin{example}
			ვაჩვენოთ, რომ ტრაქტრიქსი-კატენარი არის წყვილი ინვოლუტა-ევოლუტა.
		\end{example}
		
		ტრაქტრიქსი (tractrix) ლათინურად ნიშნავს "მიუშვი-მოქაჩე".
		
		\begin{figure}[h!]
			\centering
			\begin{tikzpicture}
				\draw[gray, thick] (-4,0) -- (4,0);
				\draw [line width=1pt] (0,1) .. controls (0.3,0.6) and (0.6,0.4) .. (2,0.1);
				\draw [line width=1pt] (0,1) .. controls (-0.3,0.6) and (-0.6,0.4) .. (-2,0.1);
				\draw [line width=1pt] (0,1) .. controls (-1,1.3) and (-1.4,2.1) .. (-2, 3.3);
				\draw [line width=1pt] (0,1) .. controls (1,1.3) and (1.4,2.1) .. (2, 3.3);
				\draw[gray, thick] (1.4,2.1) -- (1.4,0);
				\draw[gray, thick] (1.4,2.1) -- (0.2,0);
				\draw[gray, thick] (1.4,2.1) -- (2.4,2.1) node[anchor=south east]{t};
				\draw[gray, thick] (1.4,2.1) -- (1.9,2.7) ;
				\filldraw[black] (1.4,0) circle (2pt) node[anchor=north ]{S};
				\draw[red, thick] (0.2,0.7) -- (1.4,0);
				\draw[red, thick] (0,1) -- (0,0);
				\filldraw[black] (0,1) circle (2pt) node[anchor=south ]{O};
				\filldraw[black] (1.4,2.1) circle (2pt) node[anchor=south east]{P};
				\filldraw[black] (0.5,0.5) circle (2pt) node[anchor=south west]{K};
			\end{tikzpicture}
			\caption{ტრაქტრიქსი და მისი ევოლუტა კატენარი. $|OK|=|PS|=a=1$}
			\label{tractrix-catenary}
		\end{figure}
		
		
		
		
		
		
		
		
		
		
	
		\medskip
		
		\begin{example}
			\label{ex:cycloid}
			\textbf{ციკლოიდა} არის წრეწირის წერტილის ტრაექტორია სიბრტყეზე, როცა წრეწირი გორავს რაიმე წრფეზე. (იხ. სურ. \ref{fig:cycloid}).
		\end{example}
		სურ. \ref{fig:cycloid} მოცემულია ციკლოიდა და მისი პარამეტრული განტოლება. $a$ რადიუსიანი წრეწირის აბცისთა ღერძთან კონტაქტის წერტილიდან მანძილი სათავემდე $\theta$ კუთხის შესაბამისი $a\theta$ რკალის ტოლია. ამიტომ ციკლოიდას $(x,y)$ წერტილის რადიუს ვექტორი არის სამი ვექტორის ჯამი:
		\[(x,y)=(a\theta,0)+(0,a)+(-a\sin \theta,-a \cos \theta).\]
		
		ციკლოიდას სინგულარულ წერტილებს (უკუქცევის წერტილები) შორის რკალის სიგრძე ტოლია:
		\begin{align*}
			S &=\int_0^{2\pi}\sqrt{\left(\frac{dx}{d\theta}\right)^2+
				\left(\frac{dy}{d\theta}\right)^2} d\theta \\
			&=\int_0^{2\pi} a\sqrt{(1-\cos\theta)^2 + (\sin\theta)^2} d\theta \\
			&=\int_0^{2\pi} a\sqrt{1-2\cos \theta + \cos^2\theta + \sin^2\theta} d\theta \\
			&=\int_0^{2\pi} a\sqrt{2-2\cos \theta} d\theta\\
			&=2a\int_0^{2\pi}\sin \frac{\theta}{2}d\theta\\
			&=8a.
		\end{align*}
		
		\begin{figure}[h!]
			\centering
			\begin{tikzpicture}
				\coordinate (O) at (0,0);
				\coordinate (A) at (0,3);
				\def\r{1} % radius
				\def\c{1.4} % center x-coordinate
				\coordinate (C) at (\c, \r);
				
				\draw[-latex] (O) -- (A) node[anchor=south] {$y$};
				\draw[-latex] (O) -- (2.6*pi,0) node[anchor=west] {$x$};
				\draw[red,domain=0:2.5*pi,samples=100, line width=1]
				plot ({\x - sin(\x r)},{1 - cos(\x r)});
				\draw[blue, line width=1] (C) circle (\r);
				
				% Point P on the cycloid
				\def\thetaVal{1.4} % Angle for the point P
				\coordinate (P) at ({\thetaVal - sin(\thetaVal r)}, {1 - cos(\thetaVal r)});
				\coordinate (Center) at (\thetaVal, 1);
				
				% draw center of circle
				\draw[fill=blue] (Center) circle (1pt);
				
				% draw radius of the circle
				\draw[] (Center) -- (P) node[midway, anchor=south east] {$a$};
				
				% bottom of circle, radius to the bottom
				\coordinate (B) at (\thetaVal, 0);
				\draw[dashed] (Center) -- (B);
				\node[anchor=north] at (B) {$a\theta$};
				
				% projections of point P
				\draw[dotted] (P) -- (P |- O) node[anchor=north] {$x$};
				\draw[dotted] (P) -- (P -| O) node[anchor=east] {$y$};
				
				% arc theta
				\draw[->] (Center) ++(-90:0.4) arc (-90:{-90-rad(\thetaVal)}:0.4);
				\node[xshift=-3mm, yshift=-3mm] at (Center) {\scriptsize $\theta$};
				
				% P dot
				\draw[fill=black] (P) circle (2pt);
				
				% top label
				\coordinate (T) at (pi, 2);
				\node[anchor=south] at (T) {$(\pi a, 2a)$};
				\draw[fill=black] (T) circle (1pt);
				
				% equations
				\coordinate (E) at ( 4,1.2);
				\coordinate (F) at ( 4,0.9);
				\node[] at (E) {\scriptsize $x=a(\theta - \sin \theta)$};
				\node[] at (F) {\scriptsize $y=a(1 - \cos \theta)$};
				
				% label 2pi a
				\coordinate (TPA) at (2*pi, 0);
				\node[anchor=north] at (TPA) {$2 \pi a$};
			\end{tikzpicture}
			\caption{ციკლოიდა.}
			\label{fig:cycloid}
			%\end{center}
		\end{figure}
		
		ციკლოიდას ბევრი შესანიშნავი თვისება აქვს. მაგალითად ბერნულიმ (Bernoulli 1696) აღმოაჩინა (იხ. სურ. \ref{bernuli}), რომ არა $AC$ მონაკვეთი, ან $ABC$ ტეხილი, არამედ ამობრუნებული ციკლოიდა არის ისეთი წირი, რომლის უკუქცევის $A$ წერტილიდან ჩამოსრიალებული სხეული ყველაზე მოკლე დროში მიაღწევს $C$ ფინიშს (brachistochrone curve). ციკლოიდას კიდევ უფრო საკვირველი თვისება აღმოაჩინა ჰიუგენსმა (Huygens 1659). კერძოდ,
		ციკლოიდას $AC$ რკალის ნებისმიერი წერტილიდან ჩამოსრიალებული სხეული $C$ ფინიშს
		ერთი და იგივე ფიქსირებულ დროში მიაღწევს (tautochrone curve). ეს მუდმივი დროა $T=\pi\sqrt{\frac{a}{g}}$, სადაც $a$ არის ციკლოიდას წარმომქმნელი მგორავი წრეწირის რადიუსი, ხოლო $g$ არის თავისუფალი ვარდნის აჩქარება.
		
		შესაბამისი გამოთვლები დამოუკიდებლად შეიძლება ვცადოთ. მითითება: ისარგებლეთ სურ. \ref{bernuli}-ით. გვაქვს $T=\int_{\text{start}}^{\text{finish}}\frac{ds}{v}$, აქ $v$ სიჩქარე, როგორც სხეულის კოორდინატების ფუნქცია გამოითვლება $y_0$ სიმაღლიდან ჩამოსრიალებული სხეულისთვის $y$ სიმაღლეზე მისი კინეტიკური და პოტენციური ენერგიების ტოლობით
		\begin{equation}
			\label{pot=kinet}
			\frac{1}{2}mv^2=mg(y_0-y).
		\end{equation}
		გარდა ამისა,
		\[ ds=\sqrt{(dx)^2+(dy)^2}. \]
		ბოლოს უნდა გადავიდეთ $\theta$ პარამეტრზე ამოტრიალებული ციკლოიდას პარამეტრულ განტოლებაში. $C$ წერტილს შეესაბამება $y=0$ ანუ $\theta=\pi$. გამოთვლით უნდა დავრწმუნდეთ, რომ მიღებული ინტეგრალი $T$ არაა დამოკიდებული სტარტზე, ანუ $y_0$-ზე და ამიტომ $\theta_0$-ზე. პასუხი უკვე ვთქვით: $T=\pi\sqrt{\frac{a}{g}}$.
		
		\bigskip
		
		\begin{example} 
			აბელის მექანიკური ამოცანა
		\end{example}
		
		არსებობს წინა ამოცანის განზოგადება. კერძოდ, როგორ ვიპოვოთ წირის განტოლება, თუ მოცემულია ფუნქცია $T(y)$, როგორც სიმაღლის ფუნქცია, რომელიც განისაზღვრება როგორც წირზე მოძრავი ნაწილაკის ნულოვან დონემდე ჩამოსრიალების დრო. წინა ამოცანა კერძო მაგალითია, რადგან იქ $T(y)$ მუდმივია.
		
		\[ ds = \sqrt{1 + (x'(y))^2} dy = \frac{ds}{dy} dy. \]
		ამიტომ \eqref{pot=kinet} გადავწეროთ
		\[ \frac{1}{2}m\left(\frac{ds}{dt}\right)^2=mg(y_0-y). \]
		\[ dt=\pm\frac{ds}{\sqrt{2g(y_0-y)}}, \quad \text{ანუ} \quad dt=\pm\frac{1}{\sqrt{2g(y_0-y)}}\frac{ds}{dy}dy, \]
		სადაც ბოლო ტოლობაში $ds$ როგორც $y$-ის ფუნქცია ასე ჩავწერეთ $ds=\frac{ds}{dy}dy$.
		
		ეხლა ვაინტეგროთ $y=y_0$-დან $y=0$-მდე
		\begin{equation}
			\label{eq:convolution}	
			T(y_0)=\int_{0}^{y_0} dt=
			\frac{1}{\sqrt{2g}}\int_{0}^{y_0}\frac{1}{\sqrt{y_0-y}}\frac{ds}{dy}dy.
		\end{equation}
		ეს არის აბელის ინტეგრალური განტოლება.
		
		აბელის მექანიკური ამოცანა ამის შებრუნებულ პროცესს მოითხოვს, ჩვენ გვინდა ვიპოვოთ $f(y)=\frac{ds}{dy}$, საიდანაც აღვადგენთ წირის განტოლებას.
		
		ეს საინტერესო ამოცანა მოითხოვს წავიკითხოთ: ორი ფუნქციის კონვოლუცია $f \star g$, ფუნქციის ლაპლასის გარდაქმნა $\mathcal{L}[f]$ და მათი თვისებები. აქ მხოლოდ მოკლე მითითებით შემოვიფარგლებით.
		\begin{align*}
			& (f\star g)(x) = \int_0^x f(x-t)g(t) dt; \\
			& \mathcal{L}[f \star g]=\mathcal{L}[f] \cdot \mathcal{L}[g];\\
			& \mathcal{L}[y^{-1/2}](\xi)=\sqrt{\pi/\xi}.
		\end{align*}
		კერძოდ, \eqref{eq:convolution} განტოლებაში ინტეგრალი არის კონვოლუცია,
		\[ \sqrt{2g}T(y_0) = \left(y^{-1/2} \star \frac{ds}{dy}\right)(y_0). \]
		ამიტომ, თუ \eqref{eq:convolution}-ს ლაპლასის გარდაქმნას მოვდებთ, გვექნება:
		\[ \sqrt{2g}\mathcal{L}[T(y_0)](\xi) = \mathcal{L}[y^{-1/2}](\xi) \cdot \mathcal{L}\left[\frac{ds}{dy}\right](\xi) = \sqrt{\frac{\pi}{\xi}} \mathcal{L}\left[\frac{ds}{dy}\right](\xi). \]
		\[ \mathcal{L}\left[\frac{ds}{dy}\right](\xi)=\sqrt{\frac{2g}{\pi}}\sqrt{\xi}\mathcal{L}[T(y_0)](\xi). \]
		
		წინა ამოცანაში $T(y_0)=T_0$ მუდმივია, ამიტომ ციკლოიდა შეესაბამება ასეთ პასუხს
		\[ \frac{ds}{dy}=T_0\frac{\sqrt{2g}}{\pi}\frac{1}{\sqrt{y}}, \]
		და დაგვრჩება მხოლოდ ინტეგრება.
		
		\bigskip
		\bigskip
		
		\begin{example}
			აჩვენეთ, რომ:
			
			ა. წრფის ევოლუტა არის უსასრულოდ დაშორებული წერტილი.
			
			ბ. წრეწირის ევოლუტა მისი ცენტრია.
			
			გ. ციკლოიდის ევოლუტა არის მისი კონგრუენტული ციკლოიდა.
		\end{example}
		
		\begin{example}
			იპოვეთ:
			
			ა. ელიფსის ევოლუტა.
			
			ბ. პარაბოლას ევოლუტა.
			
			გ. ჰიპერბოლას ევოლუტა.
		\end{example}
	
	
	
	
	
	
	
	
	
	
	
	
	
	
		\textbf{ელიფსის ევოლუტა} ელიფსი, მოცემულია პარამეტრული განტოლებებით:
		\begin{equation*}
			x = a \cos(t), \quad y = b \sin(t)
		\end{equation*}
		სადაც \(a\) და \(b\) ელიფსის ნახევარღერძებია. მისი ევოლუტა (ასტროიდა ) აღიწერება შემდეგი განტოლებებით:
		\begin{equation}
			X = \frac{a^2 - b^2}{a} \cos^3(t), \quad Y = \frac{b^2 - a^2}{b} \sin^3(t)
		\end{equation}
	
	
	%\documentclass[tikz, border=5mm]{standalone}
	
		
		\begin{tikzpicture}[
			scale=1.5, % ნახატის მასშტაბი
			declare function={
				a = 2; % ელიფსის დიდი ნახევარღერძი
				b = 1; % ელიფსის მცირე ნახევარღერძი
			}
			]
			
			% კოორდინატთა ღერძები
			\draw[->, gray] (-3,0) -- (3,0) node[right] {$x$};
			\draw[->, gray] (0,-2) -- (0,2) node[above] {$y$};
			
			% ელიფსის დახატვა
			% ვიყენებთ პარამეტრულ ფორმას, \t არის კუთხე გრადუსებში
			\draw[blue, thick, variable=\t, domain=0:360, samples=200]
			plot ({a*cos(\t)}, {b*sin(\t)});
			
			% ევოლუტის დახატვა
			\draw[red, thick, variable=\t, domain=0:360, samples=200]
			plot ({( (a^2-b^2)/a ) * (cos(\t))^3}, { ( (b^2-a^2)/b ) * (sin(\t))^3 });
			
			% ლეგენდა
			\node[draw, rounded corners, fill=white, fill opacity=0.8, text opacity=1] at (2.5, 1.8) {
				\begin{tabular}{l}
					\textcolor{blue}{\rule{0.5cm}{2pt}} ელიფსი \\
					\textcolor{red}{\rule{0.5cm}{2pt}} ევოლუტა \\
				\end{tabular}
			};
			
		\end{tikzpicture}
		
	
	
	
	
	
	
		
		\textbf{ჰიპერბოლის ევოლუტა}
		ჰიპერბოლა მოცემულია  პარამეტრული განტოლებებით:
		\begin{equation*}
			x = a \cosh(t), \quad y = b \sinh(t)
		\end{equation*}
		სადაც \(a\) და \(b\) ჰიპერბოლის ნახევარღერძებია. მისი ევოლუტა აღიწერება შემდეგი განტოლებებით:
		\begin{equation*}
			X = \frac{a^2 + b^2}{a} \cosh^3(t), \quad Y = -\frac{a^2 + b^2}{b} \sinh^3(t)
		\end{equation*}
	
	
	
	
		
		 % მათემატიკური ფუნქციებისთვის, როგორიცაა sqrt
		
		
			
			\begin{tikzpicture}[
				scale=0.8,
				declare function={
					a = 2; % ჰიპერბოლის რეალური ნახევარღერძი
					b = 1.5; % ჰიპერბოლის წარმოსახვითი ნახევარღერძი
				}
				]
				
				% ვიანგარიშებთ ფოკუსის მდებარეობას: c = sqrt(a^2 + b^2)
				\tikzmath{
					\c = sqrt(a*a + b*b);
				}
				
				% სურათის ცენტრირებისთვის ვიყენებთ უხილავ ჩარჩოს
				\useasboundingbox (-8,-7) rectangle (8,7);
				
				% კოორდინატთა ღერძები
				\draw[->, gray] (-8,0) -- (8,0) node[right] {$x$};
				\draw[->, gray] (0,-7) -- (0,7) node[above] {$y$};
				
				% --- ფოკუსების დამატება ---
				\fill[black] (\c,0) circle (2.5pt);
				\node[above=2pt, font=\small] at (\c,0) {$F_1$};
				\fill[black] (-\c,0) circle (2.5pt);
				\node[above=2pt, font=\small] at (-\c,0) {$F_2$};
				
				% --- ჰიპერბოლის დახატვა (4 ნაწილად, მხოლოდ t > 0 გამოყენებით) ---
				\def\hyperdomain{1.3}
				% მარჯვენა შტოს ზედა ნაწილი (x > 0, y > 0)
				\draw[blue, thick, variable=\t, domain=0:\hyperdomain, samples=100] plot ({a*cosh(\t)}, {b*sinh(\t)});
				% მარჯვენა შტოს ქვედა ნაწილი (x > 0, y < 0)
				\draw[blue, thick, variable=\t, domain=0:\hyperdomain, samples=100] plot ({a*cosh(\t)}, {-b*sinh(\t)});
				% მარცხენა შტოს ზედა ნაწილი (x < 0, y > 0)
				\draw[blue, thick, variable=\t, domain=0:\hyperdomain, samples=100] plot ({-a*cosh(\t)}, {b*sinh(\t)});
				% მარცხენა შტოს ქვედა ნაწილი (x < 0, y < 0)
				\draw[blue, thick, variable=\t, domain=0:\hyperdomain, samples=100] plot ({-a*cosh(\t)}, {-b*sinh(\t)});
				
				
				% --- ევოლუტის დახატვა (4 ნაწილად, მხოლოდ t > 0 გამოყენებით) ---
				\def\evoldomain{1.0}
				\pgfmathsetmacro{\Ca}{(a*a+b*b)/a}
				\pgfmathsetmacro{\Cb}{(a*a+b*b)/b}
				% მარჯვენა შტოს ქვედა ნაწილი
				\draw[red, thick, variable=\t, domain=0.05:\evoldomain, samples=100] plot ({\Ca*cosh(\t)^3}, {-\Cb*sinh(\t)^3});
				% მარჯვენა შტოს ზედა ნაწილი
				\draw[red, thick, variable=\t, domain=0.05:\evoldomain, samples=100] plot ({\Ca*cosh(\t)^3}, {\Cb*sinh(\t)^3});
				% მარცხენა შტოს ქვედა ნაწილი
				\draw[red, thick, variable=\t, domain=0.05:\evoldomain, samples=100] plot ({-\Ca*cosh(\t)^3}, {-\Cb*sinh(\t)^3});
				% მარცხენა შტოს ზედა ნაწილი
				\draw[red, thick, variable=\t, domain=0.05:\evoldomain, samples=100] plot ({-\Ca*cosh(\t)^3}, {\Cb*sinh(\t)^3});
				
				
				% ლეგენდა
				\node[draw, rounded corners, font=\small, fill=white, fill opacity=0.8, text opacity=1] at (6, 6) {
					\begin{tabular}{l}
						\textcolor{blue}{\rule{0.4cm}{2pt}} ჰიპერბოლა \\
						\textcolor{red}{\rule{0.4cm}{2pt}} ევოლუტა \\
						\textbf{F} \hspace{0.2cm} ფოკუსები
					\end{tabular}
				};
				
			\end{tikzpicture}
			
		
		
		
		
		
		
		
		\textbf{პარაბოლის ევოლუტა}
		პარაბოლა მოცემულია  განტოლებით \(y^2 = 4ax\), და რომლის  პარამეტრული  ფორმაა:
		\begin{equation*}
			x = at^2, \quad y = 2at
		\end{equation*}
		მისი ევოლუტა (ნახევრადკუბური პარაბოლა) აღიწერება პარამეტრულად:
		\begin{equation*}
			X = 2a + 3at^2, \quad Y = -2at^3
		\end{equation*}
		ან არაცხადი ფორმით:
		\begin{equation}
			27aY^2 = 4(X - 2a)^3
		\end{equation}
	




	\begin{tikzpicture}[
		scale=0.8, 
		declare function={
			a = 1; % პარაბოლის პარამეტრი (y^2 = 4ax)
		}
		]
		
		% სურათის ცენტრირებისთვის ვამატებთ უხილავ ჩარჩოს
		\useasboundingbox (-2,-6.5) rectangle (10.5,6.5);
		
		
		% კოორდინატთა ღერძები
		\draw[->, gray] (-1,0) -- (10,0) node[right] {$x$};
		\draw[->, gray] (0,-6) -- (0,6) node[above] {$y$};
		
		% --- ფოკუსის დამატება ---
		\fill[black] (a,0) circle (2.5pt); % ფოკუსის წერტილი
		\node[above=2pt, font=\small] at (a,0) {$F$}; % ფოკუსის აღნიშვნა
		
		
		% --- პარაბოლა (გასწორებული მეთოდით) ---
		% ზედა შტო
		\draw[blue, thick, variable=\t, domain=0:2.4, samples=100]
		plot ({a*\t^2}, {2*a*\t});
		% ქვედა შტო
		\draw[blue, thick, variable=\t, domain=0:2.4, samples=100]
		plot ({a*\t^2}, {-2*a*\t});
		
		
		% --- ევოლუტა (გასწორებული მეთოდით) ---
		\def\evoldomain{1.5}
		
		% ქვედა შტო
		\draw[red, thick, variable=\t, domain=0:\evoldomain, samples=100]
		plot ({2*a + 3*a*\t^2}, {-2*a*\t^3});
		% ზედა შტო
		\draw[red, thick, variable=\t, domain=0:\evoldomain, samples=100]
		plot ({2*a + 3*a*\t^2}, {2*a*\t^3});
		
		
		% ლეგენდა
		\node[draw, rounded corners, fill=white, fill opacity=0.8, text opacity=1] at (8, 5) {
			\begin{tabular}{l}
				\textcolor{blue}{\rule{0.5cm}{2pt}} პარაბოლა \\
				\textcolor{red}{\rule{0.5cm}{2pt}} ევოლუტა \\
				\textbf{F} \hspace{0.2cm} ფოკუსი
			\end{tabular}
		};
		
	\end{tikzpicture}

	
	











		
		\begin{figure}[h!]
			\centering
			\begin{tikzpicture}
				\draw[gray, thick, ->] (0,0) -- (7,0) node[anchor=west]{X};
				\draw[gray, thick, ->] (0,-2) -- (0,5) node[anchor=south]{Y};
				\filldraw[black] (0,0) circle (2pt) node[anchor=north east]{B};
				\filldraw[black] (0,4) circle (2pt) node[anchor=south east]{A};
				\filldraw[black] (6,0) circle (2pt) node[anchor=south west]{C};
				\draw [black, thick] (0,4) -- (0,0) -- (6,0);	
				\draw [color=red, line width=2pt] (0,4) .. controls (0,1) and (3,0) .. (6,0);
				\draw [line width=1pt] (0,4) -- (6,0);
			\end{tikzpicture}
			\caption{ამობრუნებული ციკლოიდას რკალი}
			\label{bernuli}
		\end{figure}
		
		\bigskip
	
	
	

		
	
	
	
	
	
	
	
	
	
	
	
	
	
	
	
		\newpage
		
		\subsection{წირთა ოჯახის მომვლები}
		
		განვიხილოთ განტოლება
		\begin{equation}
			\label{family}
			F(x, y, t) = 0,
		\end{equation}
		სადაც $F(x, y, t)$ არის სამი ცვლადის გლუვი (ანუ $C^{\infty}$ კლასის) ფუნქცია. $t$ ცვლადს ეწოდება პარამეტრი. თუ ჩვენ დავაფიქსირებთ $t$ პარამეტრის მნიშვნელობას განტოლებაში, მაშინ ვიღებთ წირის განტოლებას.
		ამრიგად \eqref{family} განსაზღვრავს წირთა ერთპარამეტრიან
		ოჯახს. ოჯახის წირი, რომელიც შეესაბამება $t$ პარამეტრს,
		აღინიშნება $C_t$-თი და მთელი ოჯახი კი $C$-თი. ჩვენ განვიხილავთ მარტივ წირებს, ანუ ჯერადი წერტილების გარეშე.
		
		\begin{definition}
			\label{def:discriminant}
			ვთქვათ მარტივ წირთა $C$ ოჯახი მოცემულია \eqref{family} განტოლებით $(x_0,y_0, t_0)$ წერტილის მიდამოში, ე.ი.,
			\[|x-x_0|<\epsilon_1,\quad |y-y_0|<\epsilon_2,\quad |t-t_0|<\epsilon \]
			და $F_t(x,y,t)$ იგივურად ნული არაა იმავე მიდამოში.
			
			თუ ამ მიდამოში არსებობს რეგულარული გლუვი
			\[ l \text{ წირი: } x=\phi(t), \quad y=\psi(t),\]
			ისეთი, რომ ნებისმიერი $t'$ პარამეტრისთვის $F(x,y,t')=0$ წირი ეხება $l$ წირს $(\phi(t'), \psi(t'))$ წერტილში, მაშინ $l$-ს ეწოდება $C$ ოჯახის მომვლები.
			
			მოკლედ, $l$ წირს ეწოდება $C$ ოჯახის მომვლები, თუ $l$ მის თითოეულ წერტილში
			ეხება ამ ოჯახის რომელიმე ერთ წირს (განსხვავებულ წერტილებში განსხვავებულ წირებს).
		\end{definition}
		
		\textbf{მარტივ წირთა ოჯახის მომვლების აუცილებელი პირობა.}
		
		\begin{theorem}
			\label{thm:discriminant}
			სიბრტყის წერტილთა სიმრავლეს, რომელიც აკმაყოფილებს განტოლებათა სისტემას
			\begin{equation}
				\label{eq:envelope}
				\begin{cases}
					F(x,y,t)=0,\\
					F_t(x,y,t)=0
				\end{cases}
			\end{equation}
			ეწოდება \eqref{family} წირთა ოჯახის დისკრიმინანტული წირი.
		\end{theorem}
		
		\begin{theorem}
			\label{Thm-envelope}
			თუ $C$ ოჯახის მომვლები არსებობს, ის მოიცემა \eqref{eq:envelope} სისტემით, ანუ თუ მომვლები არსებობს, ის არის დისკრიმინანტული წირის ქვესიმრავლე.
		\end{theorem}
		\begin{proof}
			ვაჩვენოთ, რომ მომვლები $x=\phi(t), y=\psi(t)$ აკმაყოფილებს თეორემის განტოლებებს. მომვლების წერტილი, რომელიც შეესაბამება პარამეტრის ნებისმიერ მნიშვნელობას, არის აგრეთვე ოჯახის შესაბამისი წირის წერტილი, ამიტომ სრულდება თეორემის პირველი ტოლობა. შემდეგ, ამ საერთო წერტილში მომვლების და ოჯახის წირის მხებიც საერთოა, ამიტომ მხებების მიმმართველი ვექტორები პარალელურია, ანუ $(F_x,F_y) \perp (\phi',\psi')$, ანუ
			\[ F_x\phi'+F_y\psi'=0. \]
			თეორემის მეორე ტოლობას მივიღებთ, თუ ამ ტოლობას გამოვაკლებთ $t$-თი გაწარმოებულ პირველ ტოლობას, ანუ ტოლობას:
			\[\frac{d}{dt}F(\phi(t),\psi(t),t) = F_x\phi'(t)+F_y\psi'(t)+F_t = 0.\]
		\end{proof}
		
		\begin{figure}[h!]
			\centering
			\begin{tikzpicture}
				\tikzset{
					dot/.style 2 args={fill, circle, inner sep=1pt, label={#1:#2}}
				}
				
				\newcommand\Radius{6}
				\begin{scope}[rotate=90,yscale=-1,xscale=1]
					\draw[name path=quartarc, thick] (\Radius,0) arc (0:90:\Radius) coordinate[pos=0] (C0);
					
					\foreach \angle [
					count=\n,
					evaluate=\n as \xn using int(5-\n)
					] in
					{20,45,60,80,90}{%
						\path[name path={tan\n}] (0,0) -- (\angle:{\Radius+1});
						\path[name intersections={of={tan\n} and quartarc,by={P\n}}];
						\pgfmathsetmacro\pointA{cos(\angle)*\Radius}
						\pgfmathsetmacro\pointB{sin(\angle)*\Radius}
						\pgfmathsetmacro\ArcLen{(2*pi*\Radius/360)*\angle}
						\pgfmathsetmacro\pointX{\pointA+sin(\angle)*\ArcLen}
						\pgfmathsetmacro\pointY{\pointB-cos(\angle)*\ArcLen}
						\coordinate (C\n) at (\pointX,\pointY);
					}
				\end{scope}
				
				\foreach \j [remember=\j as \lastj (initially 0)] in {0,...,5}{
					\draw[thick] (C\lastj) edge[bend left=7.5] (C\j);
				}
				
				\node[dot={below}{$N_0$}] at (C0) {};
				\foreach \x in {1,...,5}{
					\draw[thick] (P\x) -- (C\x) coordinate[pos=0,dot={}{}] coordinate[pos=1,dot={}{}] (a\x);
					\ifnum\x<4\relax
					\node[dot={180-(15*\x)}{$M_\x$}] at (C\x) {};
					\node[dot={255-(15*\x)}{$N_\x$}] at (P\x) {};
					\fi
				}
				
				\draw (40:\Radius) --++ (225:3cm) node[fill=white] {$E$};
				\draw ($(a3)!.5!(a4)+(110:1mm)$) --++ (110:2cm) node[fill=white] {$I$};
			\end{tikzpicture}
			\caption{$E$ ევოლუტა, $I$ ინვოლუტა. $E$ მომვლებია $(M_c N_c)$ წრფეთა ოჯახის.}
			\label{fig:envelpe}
		\end{figure}
		
		\bigskip
		
		ამრიგად, \eqref{eq:envelope} მხოლოდ მომვლების არსებობის აუცილებელი პირობაა. საკმარისი პირობა დამატებით მოითხოვს, რომ შესრულდეს შემდეგი \textbf{მარტივ წირთა ოჯახის მომვლების საკმარისი პირობა}.
		
		\begin{theorem}
			\label{suf}
			თუ $F(x,y,t)$ სამი ცვლადის ანალიზური ფუნქციაა $(x_0,y_0,t_0)$ წერტილის რაიმე $U$ მიდამოში და თუ $(x_0,y_0,t_0)$ წერტილში
			\begin{enumerate}
				\item $F(x_0,y_0,t_0)=0$ და $F_t(x_0,y_0,t_0)=0$;
				\item $\det \begin{bmatrix}
					\frac{\partial F}{\partial x} & \frac{\partial F}{\partial y} \\
					\frac{\partial F_t}{\partial x} & \frac{\partial F_t}{\partial y}
				\end{bmatrix}\neq 0$ $(x_0,y_0,t_0)$-ში;
				\item $\frac{\partial^2 F}{\partial t^2}(x_0,y_0,t_0)\neq 0$.
			\end{enumerate}
			მაშინ $F(x,y,t)=0$ წირთა ოჯახს აქვს ერთადერთი $x=\phi(t), y=\psi(t)$ მომვლები ანალიზური წირი $(x_0,y_0,t_0)$ წერტილის რაიმე $V \subset U$ მიდამოში.
		\end{theorem}
		
		\begin{proof}
			თეორემის (2) პირობა საჭიროა, რომ (1) პირობიდან არაცხადი ფუნქციის შესახებ თეორემით $x,y$ ცვლადები გამოვსახოთ $t$-თი და მივიღოთ წირი
			\[x=\phi(t),\quad y=\psi(t).\]
		


როცა  $\frac{\partial^2 F}{\partial t^2}(x_0,y_0,t_0)= 0$, ამ შემთხვევაში $t_0$ პარამეტრის შესაბამისი წირის გრაფიკს $(x_0,y_0)$ წერტილში  აქვს გადაღუნვის წერტილი. ის არ იქცევა პარაბოლასავით. ეს ნიშნავს, რომ მეზობელი  წირების გადაკვეთა არ არის სტანდარტული. ასეთ წერტილების  ზღვრული  მდებარეობით   მომვლების  წერტილის  ნაცვლად შეიძლება მივიღოთ  წვეტი ან სხვა სინგულარული  წერტილი.


ამის მაგალითია  $L$ სიგრძის მონაკვეთით მოღებული  ოჯახი, როცა მონაკვეთის ბოლოები  $x$ და $y$ საკოორდინატო  ღერძებზე სრიალებს  და პარამეტრი მონაკვეთსა და $x$ ღერძს  შორის კუთხეა. ამ ოჯახის განტოლებაა  
$$
F(x,y,t)=x \sin(t) + y \cos(t) - L \sin(t) \cos(t) = 0,
$$
ხოლო  ოჯახის მომვლები არის ასტროიდა, თავისი ოთხი წვეტით.

\end{proof}


		
		\begin{figure}[h!]
			\centering
			\begin{tikzpicture}[scale=0.6]
				\draw [blue] circle [radius=4.3];
				\draw [blue] circle [radius=3.3];
				\draw [blue] circle [radius=2.3];
				\draw [blue] circle [radius=1.3];
				\draw [blue] circle [radius=0.3];
			\end{tikzpicture}
			\caption{კონცენტრული წრეწირების ოჯახს არა აქვს მომვლები.}
			\label{contr-envelope}
		\end{figure}
		
		\bigskip
		
		
		
		\subsection{სავარჯიშოები}
		
		\begin{enumerate}
			\item \textbf{წრეწირი:} იპოვეთ $\mathbf{r}(t) = \langle a\cos(t), a\sin(t) \rangle$ წრეწირის სიმრუდე, სადაც $a > 0$.
			
			\textbf{ამოხსნა:} $\kappa(t) = 1/a$.
			
			
			
			\item \textbf{პარაბოლა:} იპოვეთ $\mathbf{r}(t) = \langle t, t^2 \rangle$ პარაბოლის სიმრუდე.
			
			\textbf{ამოხსნა:} $\kappa(t) = \frac{2}{(1 + 4t^2)^{3/2}}$.
			
			\item \textbf{კუბური წირი:} იპოვეთ $\mathbf{r}(t) = \langle t, t^3 \rangle$ წირის სიმრუდე.
			
			\textbf{ამოხსნა:} $\kappa(t) = \frac{6|t|}{(1 + 9t^4)^{3/2}}$.
			
			\item \textbf{ექსპონენციალური წირი:} იპოვეთ $\mathbf{r}(t) = \langle e^t, e^{-t} \rangle$ წირის სიმრუდე.
			
			\textbf{ამოხსნა:} $\kappa(t) = \frac{2}{(e^{2t} + e^{-2t})^{3/2}}$.
			
			
			
			\item \textbf{ლოგარითმული წირი:} იპოვეთ $\mathbf{r}(t) = \langle t, \ln(t) \rangle$ წირის სიმრუდე.
			
			\textbf{ამოხსნა:} $\kappa(t) = \frac{1}{\sqrt{(t^2+1)^3}}$.
			
			\item \textbf{ტრიგონომეტრიული   წირი:} იპოვეთ $\mathbf{r}(t) = \langle \cos(2t), \sin(2t) \rangle$ წირის სიმრუდე.
			
			\textbf{ამოხსნა:} $\kappa(t) = 1$.
			
			\
			
			\item \textbf{მარტივი წირი:} იპოვეთ $\mathbf{r}(t) = \langle t^3, t^2 \rangle$ წირის სიმრუდე.
			
			\textbf{ამოხსნა:} $\kappa(t) = \frac{6}{t(9t^2+4)^{3/2}}$.
		\end{enumerate}
		
		\subsection*{მოკლე კითხვა-პასუხი წირების სიმრუდეზე}
		
		\begin{enumerate}
			\item \textbf{კითხვა:} რას ზომავს სიმრუდე?
			
			\textbf{პასუხი:} სიმრუდე ზომავს, თუ რამდენად იხრება წირი, ანუ რა სიჩქარით ბრუნავს წირის მხები, მოცემულ  წერტილში.
			
			\item \textbf{კითხვა:} რისი ტოლია წრფის სიმრუდე?
			
			\textbf{პასუხი:} ნულის.
			
			\item \textbf{კითხვა:} რისი ტოლია $r$ რადიუსიანი წრეწირის სიმრუდე?
			
			\textbf{პასუხი:} $\frac{1}{r}$.
			
			\item \textbf{კითხვა:} რა არის სიმრუდის, $\kappa$-ს, ფორმულა წირის $\mathbf{r}(t)$ ვექტორული ფუნქციის გამოყენებით?
			
			\textbf{პასუხი:} $\kappa(t) = \frac{|\mathbf{r}'(t) \times \mathbf{r}''(t)|}{|\mathbf{r}'(t)|^3}$.
			
			\item \textbf{კითხვა:} რა არის სიმრუდის, $\kappa$-ს, ფორმულა რკალის სიგრძის $s$ პარამეტრის გამოყენებით?
			
			\textbf{პასუხი:} $\kappa(s) = |\mathbf{r}''(s)|$.
			
			\item \textbf{კითხვა:} რაზე მიუთითებს სიმრუდის მაღალი მნიშვნელობა?
			
			\textbf{პასუხი:} წირი მკვეთრად იხრება.
			
			\item \textbf{კითხვა:} რაზე მიუთითებს სიმრუდის დაბალი მნიშვნელობა?
			
			\textbf{პასუხი:} წირი თითქმის სწორია.
			
			\item \textbf{კითხვა:} რა კავშირია სიმრუდესა და მხები წრეწირის რადიუსს შორის?
			
			\textbf{პასუხი:} მხები წრეწირის რადიუსი სიმრუდის შებრუნებული სიდიდეა (რადიუსი = $\frac{1}{\kappa}$).
			
			\item \textbf{კითხვა:} რა არის მხები წრეწირი?
			
			\textbf{პასუხი:} ეს არის წრეწირი, რომელიც საუკეთესოდ მიახლოებს წირს მოცემულ წერტილში და ამ წერტილში წირის იდენტური მხები და სიმრუდე აქვს.
			
			\item \textbf{კითხვა:} მოქმედებს თუ არა პარამეტრიზაციის სიჩქარე სიმრუდეზე?
			
			\textbf{პასუხი:} დიახ, თუ იყენებთ ფორმულას $t$ პარამეტრით, მაგრამ არა, თუ იყენებთ რკალის სიგრძის $s$ პარამეტრს.
			
			\item \textbf{კითხვა:} რისი ტოლია $y=f(x)$ გრაფიკის სიმრუდე?
			
			\textbf{პასუხი:} $\kappa(x) = \frac{|f''(x)|}{(1 + [f'(x)]^2)^{3/2}}$.
			
			\item \textbf{კითხვა:} რა კავშირია $\mathbf{r}(t)$ წირის სიმრუდესა და ერთეულოვან  ნორმალურ  ვექტორს, $\mathbf{N}$-ს, შორის?
			
			\textbf{პასუხი:} $\kappa(t)=\frac{|\mathbf{T}'(t)|}{|\mathbf{r}'(t)|}$, და $\mathbf{T}'(t)=|\mathbf{r}'(t)|\kappa(t)\mathbf{N}(t)$.
		\end{enumerate}
		
		
		
		
		
		
		
		
		
		\newpage
		
		
		%======================================================================
		\section{წირების ლოკალური თეორია: ფრენე-სერეს აპარატი}
		%======================================================================
		
		
		\begin{definition}[პარამეტრიზებული წირი]
			$I \subset \mathbb{R}$ ინტერვალიდან $\mathbb{R}^3$-ში ასახვას $\gamma: I \to \mathbb{R}^3$ ეწოდება \textbf{პარამეტრიზებული წირი}. წირი არის \textbf{რეგულარული}, თუ $\gamma'(t) \neq 0$ ყველა $t \in I$-თვის.
		\end{definition}
		
		\subsection{ნატურალური პარამეტრიზაცია}
		\begin{definition}[რკალის სიგრძე]
			რეგულარული წირის $\gamma(t)$ რკალის სიგრძე $t_0$-დან $t$-მდე არის:
			\[ s(t) = \int_{t_0}^t \|\gamma'(\tau)\| \, d\tau \]
		\end{definition}
		რადგან $\frac{ds}{dt} = \|\gamma'(t)\| > 0$, $s(t)$ მონოტონურად ზრდადი ფუნქციაა და მას გააჩნია შებრუნებული ფუნქცია $t(s)$. შეგვიძლია წირის ხელახალი პარამეტრიზაცია $s$ პარამეტრით: $\beta(s) = \gamma(t(s))$. ამ პარამეტრიზაციას \textbf{ნატურალური} ეწოდება.
		\begin{proposition}
			ნატურალური პარამეტრიზაციისთვის $\|\beta'(s)\| = 1$.
		\end{proposition}
		\begin{proof}
			ჯაჭვური წესით: $\beta'(s) = \gamma'(t(s)) \frac{dt}{ds}$. რადგან $\frac{dt}{ds} = 1 / \frac{ds}{dt} = 1 / \|\gamma'(t)\|$, მივიღებთ:
			\[ \|\beta'(s)\| = \|\gamma'(t)\| \cdot \left| \frac{dt}{ds} \right| = \|\gamma'(t)\| \cdot \frac{1}{\|\gamma'(t)\|} = 1 \]
		\end{proof}
		შემდგომში, თუ სხვა რამ არ არის ნათქვამი, ვიგულისხმებთ, რომ წირი ნატურალურად არის პარამეტრიზებული და აღვნიშნავთ მას $\gamma(s)$-ით.
		
		\subsection{ფრენეს რეპერი}
		
		იმისთვის, რომ აღვწეროთ წირის ლოკალური გეომეტრია — ანუ, როგორ იხრება და იგრეხება ის სივრცეში ნებისმიერ წერტილთან ახლოს — სტანდარტული, ფიქსირებული  $\{ \mathbf{i}, \mathbf{j}, \mathbf{k} \}$ ბაზისი მოუხერხებელია. წირი გამუდმებით იცვლის მიმართულებას ამ ფიქსირებულ ვექტორებთან მიმართებით, რაც ანალიზს ართულებს.
		
		ბევრად ბუნებრივი და მძლავრი მიდგომაა, შევქმნათ კოორდინატთა სისტემა, რომელიც „მიბმულია“ თავად წირზე და მასთან ერთად მოძრაობს. ასეთი „მოძრავი რეპერი“ (moving frame) საშუალებას მოგვცემს, წირის თვისებები აღვწეროთ თავად წირის პერსპექტივიდან.
		
		როგორ უნდა ავაგოთ ასეთი რეპერი ბუნებრივად?
		\begin{enumerate}
			\item \textbf{პირველი მიმართულება:} ყველაზე ბუნებრივი არჩევანია წირის მოძრაობის მიმართულება. ეს არის \textbf{მხები ვექტორი}.
			\item \textbf{მეორე მიმართულება:} შემდეგი კითხვაა — საით იხრება წირი? წირის მოხრა ნიშნავს, რომ მისი მხები ვექტორი იცვლის მიმართულებას. შესაბამისად, მხები ვექტორის წარმოებული გვიჩვენებს ზუსტად იმ მიმართულებას, საითაც წირი იხრება. ეს არის \textbf{მთავარი ნორმალის} მიმართულება.
			\item \textbf{მესამე მიმართულება:} ამ ორი, უკვე ნაპოვნი ორთოგონალური მიმართულების შემდეგ, მესამე ვექტორი (ამ ორი ვექტორის ვექტორული  ნამრავლი )უბრალოდ ავსებს ჩვენს სისტემას სამგანზომილებიან, მარჯვენა ორიენტაციის ორთონორმირებულ ბაზისამდე.
		\end{enumerate}
		სწორედ ამ სამი, ურთიერთორთოგონალური ვექტორისგან შემდგარ მოძრავ ბაზისს ეწოდება \textbf{ფრენეს რეპერი}. ის იდეალური ინსტრუმენტია წირის სიმრუდისა და გრეხის ზუსტად გასაზომად. ახლა კი, განვსაზღვროთ მისი კომპონენტები ფორმალურად.
		
		\vspace{1em} % მცირე დაშორება ტექსტსა და განსაზღვრებებს შორის
		
		ვთქვათ, $\gamma(s)$ არის ნატურალურად პარამეტრიზებული წირი.
		\begin{definition}[მხები ვექტორი]
			$\mathbf{T}(s) = \gamma'(s)$. რადგან $\|\mathbf{T}(s)\| = 1$, ვექტორი $\mathbf{T}'(s)$ ორთოგონალურია $\mathbf{T}(s)$-ის.
		\end{definition}
		\begin{definition}[სიმრუდე და მთავარი ნორმალი]
			წირის \textbf{სიმრუდე} არის $\kappa(s) = \|\mathbf{T}'(s)\|$. თუ $\kappa(s) \neq 0$, მაშინ \textbf{მთავარი ნორმალური ვექტორი} არის $\mathbf{N}(s) = \frac{\mathbf{T}'(s)}{\kappa(s)}$.
		\end{definition}
		\begin{definition}[ბინორმალი და გრეხა]
			\textbf{ბინორმალური ვექტორი} არის $\mathbf{B}(s) = \mathbf{T}(s) \times \mathbf{N}(s)$. ვექტორთა სიმრავლე $\{\mathbf{T}, \mathbf{N}, \mathbf{B}\}$ ქმნის მოძრავ ორთონორმირებულ რეპერს. \textbf{გრეხა} $\tau(s)$ განისაზღვრება განტოლებით $\mathbf{B}'(s) = -\tau(s)\mathbf{N}(s)$. 
		\end{definition}
	ეს ფორმულები   მიიღება ასე:
		
		
		\subsection{ფრენე-სერეს ფორმულები}
		\begin{theorem}[ფრენე-სერეს ფორმულები]
			ნატურალურად პარამეტრიზებული წირისთვის, რომლის $\kappa > 0$, სამართლიანია შემდეგი განტოლებები:
			\begin{align*}
				\mathbf{T}' &= \phantom{-}\kappa \mathbf{N} \\
				\mathbf{N}' &= -\kappa \mathbf{T} + \tau \mathbf{B} \\
				\mathbf{B}' &= \phantom{- \kappa \mathbf{T}} - \tau \mathbf{N}.
			\end{align*}
		\end{theorem}
	მატრიცული განტოლება ასე გამოიყურება:
	\[
	\begin{pmatrix} \mathbf{T}' \\ \mathbf{N}' \\ \mathbf{B}' \end{pmatrix}
	=
	\begin{pmatrix} 0 & \kappa & 0 \\ -\kappa & 0 & \tau \\ 0 & -\tau & 0 \end{pmatrix}
	\begin{pmatrix} \mathbf{T} \\ \mathbf{N} \\ \mathbf{B} \end{pmatrix}
	\].
	
		\begin{proof}
			(1) პირველი ფორმულა გამომდინარეობს $\mathbf{N}$-ის განსაზღვრებიდან. (2) $\mathbf{B}' = (\mathbf{T} \times \mathbf{N})' = \mathbf{T}' \times \mathbf{N} + \mathbf{T} \times \mathbf{N}' = (\kappa \mathbf{N}) \times \mathbf{N} + \mathbf{T} \times \mathbf{N}' = \mathbf{T} \times \mathbf{N}'$. რადგან $\mathbf{B}'$ ორთოგონალურია $\mathbf{B}$-ს, ის უნდა იყოს $\mathbf{T}$ და $\mathbf{N}$ ვექტორების წრფივი კომბინაცია. მაგრამ $\mathbf{B}' \cdot \mathbf{T} = -\mathbf{B} \cdot \mathbf{T}' = -\mathbf{B} \cdot (\kappa\mathbf{N}) = 0$, ამიტომ $\mathbf{B}'$ უნდა იყოს $\mathbf{N}$-ის პარალელური. შესაბამისად, $\mathbf{B}' = -\tau \mathbf{N}$ რომელიღაც $\tau$ სკალარისთვის. (3) $\mathbf{N}' = (\mathbf{B} \times \mathbf{T})' = \mathbf{B}' \times \mathbf{T} + \mathbf{B} \times \mathbf{T}' = (-\tau\mathbf{N}) \times \mathbf{T} + \mathbf{B} \times (\kappa\mathbf{N}) = \tau\mathbf{B} - \kappa\mathbf{T}$.
		\end{proof}
		
		
			
			
	\bigskip
			
			
			\begin{theorem}[ფრენეს ფორმულები ზოგადი პარამეტრიზაციისთვის] \label{thm:frenet_general}
				ვთქვათ, $\alpha: I \to \mathbb{R}^3$ არის რეგულარული წირი, რომელიც, ზოგადად, არ არის ნატურალურად პარამეტრიზებული, ანუ $\|\alpha'(t)\| \neq 1$. თუ აღვნიშნავთ სიჩქარეს $v(t) = \|\alpha'(t)\|$-თი, მაშინ ფრენეს ფორმულები იღებს შემდეგ სახეს:
				\[
				\begin{cases} 
					\mathbf{T}'(t) = v(t)k(t)\mathbf{N}(t) \\ 
					\mathbf{N}'(t) = -v(t)k(t)\mathbf{T}(t) + v(t)\tau(t)\mathbf{B}(t) \\ 
					\mathbf{B}'(t) = -v(t)\tau(t)\mathbf{N}(t) 
				\end{cases}
				\]
			\end{theorem}
			
			\begin{proof}
				ვთქვათ, $\bar{\alpha}(s)$ არის $\alpha(t)$ წირის რეპარამეტრიზაცია ნატურალური პარამეტრით $s$ (რკალის სიგრძე), ანუ $\alpha(t) = \bar{\alpha}(s(t))$.
				
				მხები ვექტორი $\mathbf{T}$ არ არის დამოკიდებული პარამეტრიზაციაზე. ჯაჭვური წესის გამოყენებით, გვაქვს:
				\[ \frac{d\mathbf{T}}{dt} = \frac{d\mathbf{T}}{ds} \cdot \frac{ds}{dt} \]
				ჩვენ ვიცით, რომ ნატურალური პარამეტრიზაციისთვის $\frac{d\mathbf{T}}{ds} = k(s)\mathbf{N}(s)$, ხოლო სიჩქარეა $v(t) = \frac{ds}{dt}$. ამასთან, სიმრუდე $k$ და ნორმალი $\mathbf{N}$ გეომეტრიული სიდიდეებია და არ არის დამოკიდებული პარამეტრიზაციაზე, ანუ $k(s) = k(t)$ და $\mathbf{N}(s) = \mathbf{N}(t)$. ამ ყველაფრის გათვალისწინებით მივიღებთ პირველ ფორმულას:
				\[ \mathbf{T}'(t) = k(t)\mathbf{N}(t) \cdot v(t) = v(t)k(t)\mathbf{N}(t) \]
				დანარჩენი ორი ფორმულა ანალოგიურად მტკიცდება $\mathbf{N}$ და $\mathbf{B}$ ვექტორების გაწარმოებით ჯაჭვური წესის გამოყენებით.
			\end{proof}
			
			
			\begin{theorem}[სიმრუდისა და გრეხის პრაქტიკული ფორმულები] \label{thm:k_tau_practical}
				ვთქვათ, $\alpha: I \to \mathbb{R}^3$ არის ნებისმიერი რეგულარული პარამეტრიზებული წირი. მაშინ ფრენეს რეპერი, სიმრუდე და გრეხა გამოითვლება შემდეგი ფორმულებით:
				\begin{align*}
					\mathbf{T}(t) &:= \frac{\alpha'(t)}{\|\alpha'(t)\|} \\
					\mathbf{B}(t) &:= \frac{\alpha'(t) \times \alpha''(t)}{\|\alpha'(t) \times \alpha''(t)\|} \\
					\mathbf{N}(t) &:= \mathbf{B}(t) \times \mathbf{T}(t) \\
					k(t) &= \frac{\|\alpha'(t) \times \alpha''(t)\|}{\|\alpha'(t)\|^3} \\
					\tau(t) &= \frac{(\alpha'(t) \times \alpha''(t)) \cdot \alpha'''(t)}{\|\alpha'(t) \times \alpha''(t)\|^2}
				\end{align*}
			\end{theorem}
			
			\begin{proof}
				აღვნიშნოთ $v = v(t) = \|\alpha'(t)\|$. მაშინ, განმარტების თანახმად, $\alpha'(t) = v\mathbf{T}$.
				გავაწარმოოთ ეს გამოსახულება $t$-თი და გამოვიყენოთ წინა თეორემის პირველი ფორმულა $\mathbf{T}' = vk\mathbf{N}$:
				\begin{equation} \label{eq:alpha_pp}
					\alpha''(t) = (v\mathbf{T})' = v'\mathbf{T} + v\mathbf{T}' = v'\mathbf{T} + v(vk\mathbf{N}) = v'\mathbf{T} + v^2k\mathbf{N}
				\end{equation}
				ახლა გამოვთვალოთ ვექტორული ნამრავლი $\alpha' \times \alpha''$:
				\begin{align*}
					\alpha' \times \alpha'' &= (v\mathbf{T}) \times (v'\mathbf{T} + v^2k\mathbf{N}) \\
					&= (v\mathbf{T} \times v'\mathbf{T}) + (v\mathbf{T} \times v^2k\mathbf{N}) \\
					&= \mathbf{0} + v^3k(\mathbf{T} \times \mathbf{N}) \quad (\text{რადგან } \mathbf{T} \times \mathbf{T} = \mathbf{0}) \\
					&= v^3k\mathbf{B}
				\end{align*}
				ამრიგად, მივიღეთ საკვანძო ტოლობა:
				\begin{equation} \label{eq:cross_product}
					\alpha' \times \alpha'' = v^3k\mathbf{B}
				\end{equation}
				თუ ამ ტოლობის ორივე მხარეს ავიღებთ ნორმას (სიგრძეს), მივიღებთ:
				\[ \|\alpha' \times \alpha''\| = \|v^3k\mathbf{B}\| = v^3k\|\mathbf{B}\| = v^3k \]
				აქედან პირდაპირ გამომდინარეობს \textbf{სიმრუდის} ფორმულა:
				\[ k = \frac{\|\alpha' \times \alpha''\|}{v^3} = \frac{\|\alpha' \times \alpha''\|}{\|\alpha'\|^3} \]
				ასევე, \eqref{eq:cross_product} ტოლობიდან გამომდინარეობს \textbf{ბინორმალის} ფორმულა, რადგან $\mathbf{B}$ არის $\alpha' \times \alpha''$ ვექტორის მიმართულების ერთეულოვანი ვექტორი.
				
				\textbf{გრეხის} ფორმულის გამოსაყვანად, გავაწარმოოთ \eqref{eq:alpha_pp} ტოლობა:
				\[ \alpha''' = (v'\mathbf{T} + v^2k\mathbf{N})' = (v''\mathbf{T} + v'\mathbf{T}') + ((v^2k)'\mathbf{N} + v^2k\mathbf{N}') \]
				ახლა გამოვთვალოთ შერეული ნამრავლი $(\alpha' \times \alpha'') \cdot \alpha'''$. \eqref{eq:cross_product} ტოლობიდან ვიცით, რომ $\alpha' \times \alpha''$ ვექტორი $\mathbf{B}$-ს პარალელურია. რადგან $\mathbf{B}$ ორთოგონალურია $\mathbf{T}$-სა და $\mathbf{N}$-ის, სკალარული ნამრავლისას $\alpha'''$-ის ის კომპონენტები, რომლებიც $\mathbf{T}$-სა და $\mathbf{N}$-ს შეიცავს, განულდება. ჩვენ გვჭირდება მხოლოდ $\alpha'''$-ის $\mathbf{B}$ კომპონენტი.
				\[ \alpha''' = \dots + v^2k\mathbf{N}' = \dots + v^2k(-vk\mathbf{T} + v\tau\mathbf{B}) = (\dots)\mathbf{T} + (\dots)\mathbf{N} + v^3k\tau\mathbf{B} \]
				მაშასადამე,
				\begin{align*}
					(\alpha' \times \alpha'') \cdot \alpha''' &= (v^3k\mathbf{B}) \cdot ((\dots)\mathbf{T} + (\dots)\mathbf{N} + v^3k\tau\mathbf{B}) \\
					&= (v^3k\mathbf{B}) \cdot (v^3k\tau\mathbf{B}) \\
					&= v^6k^2\tau (\mathbf{B} \cdot \mathbf{B}) = v^6k^2\tau
				\end{align*}
				საიდანაც:
				\[ \tau = \frac{(\alpha' \times \alpha'') \cdot \alpha'''}{v^6k^2} = \frac{(\alpha' \times \alpha'') \cdot \alpha'''}{(v^3k)^2} = \frac{(\alpha' \times \alpha'') \cdot \alpha'''}{\|\alpha' \times \alpha''\|^2} \]
				რითაც დამტკიცდა გრეხის ფორმულაც.
			\end{proof}
			
		
	
	
	
		
		
		\bigskip
		
		
		\subsection{წირების კლასიფიკაცია სიმრუდისა და გრეხის მიხედვით}
		
		შემდეგი წინადადება წარმოადგენს წირთა ლოკალური თეორიის კლასიკურ შედეგებს, რომლებიც აკავშირებს წირის სიმრუდე-გრეხას მის გეომეტრიულ ფორმასთან.
		
		\begin{proposition}[წირების კლასიფიკაცია სიმრუდისა და გრეხის მიხედვით]
			ვთქვათ, $\gamma(s)$ არის ნატურალური პარამეტრით მოცემული რეგულარული წირი, $k(s)$ მისი სიმრუდით და $\tau(s)$ მისი გრეხით. მაშინ სამართლიანია შემდეგი დებულებები:
			\begin{enumerate}
				\item $\gamma$ არის წრფის ნაწილი მაშინ და მხოლოდ მაშინ, როდესაც $k(s) = 0$ ყველა $s$-თვის.
				\item $\gamma$ არის ბრტყელი წირი (ძევს სიბრტყეში) მაშინ და მხოლოდ მაშინ, როდესაც $\tau(s) = 0$ ყველა $s$-თვის (იმ წერტილებში, სადაც $k(s) \neq 0$).
				\item $\gamma$ არის წრიული  ჰელიქსი (circular helix) მაშინ და მხოლოდ მაშინ, როდესაც $k$ და $\tau$ არის არანულოვანი მუდმივები.
				\item $\gamma$ არის ზოგადი ცილინდრული   ჰელიქსი (cylindrical helix) მაშინ და მხოლოდ მაშინ, როდესაც $\tau/k$ არის მუდმივი.
			\end{enumerate}
		\end{proposition}
		
		\begin{proof}
			
			\subsubsection*{1-ლი პუნქტის დამტკიცება ($k=0 \iff$ წრფე)}
			($\Rightarrow$) ვთქვათ, $k(s) = \|\mathbf{T}'(s)\| = 0$ ყველა $s$-თვის. ეს ნიშნავს, რომ $\mathbf{T}'(s) = \mathbf{0}$. ამ განტოლების ინტეგრებით მივიღებთ, რომ მხები ვექტორი $\mathbf{T}(s)$ მუდმივია: $\mathbf{T}(s) = \mathbf{T}_0$. რადგან $\gamma'(s) = \mathbf{T}(s)$, მივიღებთ $\gamma'(s) = \mathbf{T}_0$. ამის კიდევ ერთხელ ინტეგრებით კი მივიღებთ:
			\[ \gamma(s) = s\mathbf{T}_0 + \mathbf{C} \]
			სადაც $\mathbf{C}$ მუდმივი ვექტორია. ეს არის წრფის პარამეტრული განტოლება.
			
			($\Leftarrow$) ვთქვათ, $\gamma(s)$ არის წრფე. მაშინ მისი პარამეტრული განტოლებაა $\gamma(s) = s\mathbf{a} + \mathbf{b}$, სადაც $\mathbf{a}$ მუდმივი ერთეულოვანი ვექტორია. მაშინ $\mathbf{T}(s) = \gamma'(s) = \mathbf{a}$. რადგან $\mathbf{T}(s)$ მუდმივია, მისი წარმოებული $\mathbf{T}'(s) = \mathbf{0}$, და შესაბამისად, სიმრუდე $k(s) = \|\mathbf{T}'(s)\| = 0$.
			
			\vspace{1em} % <-- დამატებულია მცირე დაშორება
			
			\subsubsection*{2-ლი პუნქტის დამტკიცება ($\tau=0 \iff$ ბრტყელი წირი)}
			($\Rightarrow$) ვთქვათ, $\tau(s) = 0$ ყველა $s$-თვის. ფრენეს ფორმულიდან $\mathbf{B}'(s) = -\tau(s)\mathbf{N}(s) = \mathbf{0}$. ეს ნიშნავს, რომ ბინორმალური ვექტორი $\mathbf{B}(s)$ მუდმივია: $\mathbf{B}(s) = \mathbf{B}_0$. განვიხილოთ ფუნქცია $f(s) = (\gamma(s) - \gamma(s_0)) \cdot \mathbf{B}_0$. მისი წარმოებული არის $f'(s) = \gamma'(s) \cdot \mathbf{B}_0 = \mathbf{T}(s) \cdot \mathbf{B}_0$. რადგან $\mathbf{T}$ ყოველთვის ორთოგონალურია $\mathbf{B}$-ს, $f'(s)=0$. მაშასადამე, $f(s)$ მუდმივია. რადგან $f(s_0)=0$, $f(s)=0$ ყველა $s$-თვის. ეს ნიშნავს, რომ ვექტორი $\gamma(s) - \gamma(s_0)$ ყოველთვის $\mathbf{B}_0$-ის პერპენდიკულარულია, ანუ წირი $\gamma(s)$ ძევს სიბრტყეში.
			
			($\Leftarrow$) ვთქვათ, წირი ძევს სიბრტყეში, რომლის ნორმალია მუდმივი ვექტორი $\mathbf{n}$. მაშინ $\mathbf{T}(s)$ და $\mathbf{N}(s)$ ყოველთვის ამ სიბრტყის პარალელურია და, შესაბამისად, $\mathbf{n}$-ის ორთოგონალური. აქედან გამომდინარე, ბინორმალი $\mathbf{B} = \mathbf{T} \times \mathbf{N}$ უნდა იყოს $\mathbf{n}$-ის პარალელური. რადგან $\mathbf{B}$ ერთეულოვანი ვექტორია, $\mathbf{B} = \pm\mathbf{n}$. ორივე შემთხვევაში, $\mathbf{B}$ მუდმივი ვექტორია, მისი წარმოებული $\mathbf{B}'(s)=\mathbf{0}$ და, შესაბამისად, გრეხა $\tau(s)=0$.
			
			\vspace{1em} % <-- დამატებულია მცირე დაშორება
			
			\subsubsection*{3-ლი პუნქტის დამტკიცება ($k, \tau$ მუდმივია $\iff$ წრიული ჰელიქსი)}
			($\Rightarrow$) ვთქვათ, $k>0$ და $\tau \neq 0$ მუდმივებია. ფრენეს ფორმულების გაწარმოებით, შეგვიძლია მივიღოთ დიფერენციალური განტოლება:
			\[ \mathbf{T}''' + (k^2+\tau^2)\mathbf{T}' = \mathbf{0} \]
			ამ განტოლების ამონახსნი, შესაბამისი ინტეგრებით, იძლევა წრიული ჰელიქსის პარამეტრულ ფორმას:
			\[ \gamma(s) = \left( a\cos\left(\frac{s}{c}\right), a\sin\left(\frac{s}{c}\right), b\left(\frac{s}{c}\right) \right) \]
			სადაც $a, b, c$ მუდმივებია, რომლებიც $k$ და $\tau$-ზეა დამოკიდებული: $k = \frac{a}{a^2+b^2}$ და $\tau = \frac{b}{a^2+b^2}$.
			
			($\Leftarrow$) თუ დავიწყებთ წრიული ჰელიქსის პარამეტრული განტოლებით, პირდაპირი გამოთვლებით შეგვიძლია ვაჩვენოთ, რომ მისი სიმრუდე და გრეხა მუდმივი სიდიდეებია.
			
			\vspace{1em} % <-- დამატებულია მცირე დაშორება
			
			\subsubsection*{4-ლი პუნქტის დამტკიცება ($\tau/k$ მუდმივია $\iff$ ცილინდრული ჰელიქსი)}
			ცილინდრული ჰელიქსის განმარტება ისაა, რომ მისი მხები ვექტორი $\mathbf{T}(s)$ ქმნის მუდმივ $\theta$ კუთხეს ფიქსირებულ $\mathbf{a}$ მიმართულებასთან (ჰელიქსის ღერძი). ანუ, $\mathbf{T}(s) \cdot \mathbf{a} = \cos\theta = \text{const}$.
			
			($\Rightarrow$) ვთქვათ, $\gamma$ ჰელიქსია. გავაწარმოოთ $\mathbf{T}(s) \cdot \mathbf{a} = \cos\theta$:
			\[ \mathbf{T}'(s) \cdot \mathbf{a} = 0 \implies k(s)\mathbf{N}(s) \cdot \mathbf{a} = 0 \]
			რადგან $k \neq 0$, მივიღებთ $\mathbf{N}(s) \cdot \mathbf{a} = 0$. ეს ნიშნავს, რომ ჰელიქსის ღერძი ყოველთვის მთავარი ნორმალის პერპენდიკულარულია. გავაწარმოოთ კიდევ ერთხელ:
			\[ \mathbf{N}'(s) \cdot \mathbf{a} = 0 \implies (-k(s)\mathbf{T}(s) + \tau(s)\mathbf{B}(s)) \cdot \mathbf{a} = 0 \]
			\[ -k(s)(\mathbf{T} \cdot \mathbf{a}) + \tau(s)(\mathbf{B} \cdot \mathbf{a}) = 0 \]
			რადგან $\mathbf{a}$ პერპენდიკულარულია $\mathbf{N}$-ის, ის ძევს $\mathbf{T}$ და $\mathbf{B}$ ვექტორებით შედგენილ სიბრტყეში, ამიტომ $\mathbf{B} \cdot \mathbf{a} = \pm \sin\theta$. ჩავსვათ ეს მნიშვნელობები:
			\[ -k(s)\cos\theta \pm \tau(s)\sin\theta = 0 \implies \frac{\tau(s)}{k(s)} = \pm \frac{\cos\theta}{\sin\theta} = \pm \cot\theta = \text{const} \]
			
			($\Leftarrow$) თუ $\tau/k = C = \text{const}$, მაშინ შეგვიძლია ვიპოვოთ ისეთი მუდმივი კუთხე $\theta$, რომ $C = \cot\theta$. ავაგოთ ვექტორი $\mathbf{a}(s) = \cos\theta\mathbf{T}(s) + \sin\theta\mathbf{B}(s)$. ვაჩვენოთ, რომ ეს ვექტორი მუდმივია:
			\[ \mathbf{a}'(s) = \cos\theta\mathbf{T}'(s) + \sin\theta\mathbf{B}'(s) = \cos\theta(k\mathbf{N}) + \sin\theta(-\tau\mathbf{N}) = (k\cos\theta - \tau\sin\theta)\mathbf{N} \]
			ჩვენი პირობიდან $\tau = k\cot\theta = k\frac{\cos\theta}{\sin\theta}$, ამიტომ ფრჩხილებში მივიღებთ $k\cos\theta - (k\frac{\cos\theta}{\sin\theta})\sin\theta = 0$. მაშასადამე, $\mathbf{a}'(s) = \mathbf{0}$, ანუ $\mathbf{a}$ მუდმივი ვექტორია. ამასთან, $\mathbf{T}(s) \cdot \mathbf{a} = \cos\theta$, რაც ნიშნავს, რომ წირი მუდმივ კუთხეს ქმნის $\mathbf{a}$ ღერძთან, ანუ ის ჰელიქსია.
		\end{proof}
		
		
		
		
		\bigskip
		
		
		
		
		\begin{figure}[h!]
			\centering
			\includegraphics[width=0.8\textwidth]{helix_figure.pdf}
			\caption{ჰელიქსის პარამეტრული გამოსახულება და მისი პროექცია.}
			\label{helix}
		\end{figure}
		
	
	
	
	\bigskip
	
			
			\subsection{მაგალითები}
			
			\begin{example} \label{ex:45}
				%oneil p.73
				დაწერეთ $\alpha(t) = \langle 3t-t^3, 3t^2, 3t+t^3 \rangle$ წირის ფრენეს ბაზისი, სიმრუდე და გრეხა.
			\end{example}
			\begin{proof}[ამოხსნა]
				პირველ რიგში, ვიპოვოთ პირველი სამი წარმოებული:
				\begin{align*}
					\alpha'(t) &= \langle 3-3t^2, 6t, 3+3t^2 \rangle = 3\langle 1-t^2, 2t, 1+t^2 \rangle \\
					\alpha''(t) &= \langle -6t, 6, 6t \rangle = 6\langle -t, 1, t \rangle \\
					\alpha'''(t) &= \langle -6, 0, 6 \rangle = 6\langle -1, 0, 1 \rangle
				\end{align*}
				
				ვიპოვოთ სიჩქარის ვექტორის ნორმა:
				\begin{align*}
					\|\alpha'(t)\| &= 3\sqrt{(1-t^2)^2 + (2t)^2 + (1+t^2)^2} \\
					&= 3\sqrt{(1-2t^2+t^4) + 4t^2 + (1+2t^2+t^4)} \\
					&= 3\sqrt{2t^4 + 4t^2 + 2} = 3\sqrt{2(t^4 + 2t^2 + 1)} = 3\sqrt{2(t^2+1)^2} \\
					&= 3\sqrt{2}(1+t^2)
				\end{align*}
				
				გამოვთვალოთ ვექტორული ნამრავლი $\alpha'(t) \times \alpha''(t)$:
				\begin{align*}
					\alpha'(t) \times \alpha''(t) &= 3 \cdot 6 \cdot
					\begin{vmatrix}
						\mathbf{i} & \mathbf{j} & \mathbf{k} \\
						1-t^2 & 2t & 1+t^2 \\
						-t & 1 & t
					\end{vmatrix} \\
					&= 18 \left[ \mathbf{i}(2t^2 - (1+t^2)) - \mathbf{j}(t(1-t^2) - (-t)(1+t^2)) + \mathbf{k}(1-t^2 - (-2t^2)) \right] \\
					&= 18 \left[ \mathbf{i}(t^2-1) - \mathbf{j}(t-t^3+t+t^3) + \mathbf{k}(1+t^2) \right] \\
					&= 18 \langle t^2-1, -2t, 1+t^2 \rangle
				\end{align*}
				
				ვიპოვოთ ამ ვექტორის ნორმა:
				\begin{align*}
					\|\alpha'(t) \times \alpha''(t)\| &= 18\sqrt{(t^2-1)^2 + (-2t)^2 + (1+t^2)^2} \\
					&= 18\sqrt{2(t^2+1)^2} = 18\sqrt{2}(1+t^2)
				\end{align*}
				
				ახლა შეგვიძლია დავწეროთ ფრენეს ბაზისის ვექტორები:
				\begin{align*}
					\mathbf{T}(t) &= \frac{\alpha'(t)}{\|\alpha'(t)\|} = \frac{3\langle 1-t^2, 2t, 1+t^2 \rangle}{3\sqrt{2}(1+t^2)} = \frac{1}{\sqrt{2}(1+t^2)}\langle 1-t^2, 2t, 1+t^2 \rangle \\
					\mathbf{B}(t) &= \frac{\alpha'(t) \times \alpha''(t)}{\|\alpha'(t) \times \alpha''(t)\|} = \frac{18\langle t^2-1, -2t, 1+t^2 \rangle}{18\sqrt{2}(1+t^2)} = \frac{1}{\sqrt{2}(1+t^2)}\langle t^2-1, -2t, 1+t^2 \rangle \\
					\mathbf{N}(t) &= \mathbf{B}(t) \times \mathbf{T}(t) = \frac{1}{2(1+t^2)^2} 
					\begin{vmatrix}
						\mathbf{i} & \mathbf{j} & \mathbf{k} \\
						t^2-1 & -2t & 1+t^2 \\
						1-t^2 & 2t & 1+t^2
					\end{vmatrix} \\
					%&= \frac{1}{2(1+t^2)^2} \langle -2t(1+t^2)-2t(1+t^2), %(1+t^2)(1-t^2)-(t^2-1)(1+t^2), 2t(t^2-1)-(-2t)(1-t^2) \rangle \\
					&= \frac{1}{2(1+t^2)^2} \langle -4t(1+t^2), 2(1-t^2)(1+t^2), 0 \rangle \\
					&= \frac{1}{1+t^2} \langle -2t, 1-t^2, 0 \rangle
				\end{align*}
				
				და ბოლოს, სიმრუდე და გრეხა:
				\begin{align*}
					k(t) &= \frac{\|\alpha'(t) \times \alpha''(t)\|}{\|\alpha'(t)\|^3} = \frac{18\sqrt{2}(1+t^2)}{\left(3\sqrt{2}(1+t^2)\right)^3} = \frac{18\sqrt{2}(1+t^2)}{54\sqrt{2}(1+t^2)^3} = \frac{1}{3(1+t^2)^2} \\
					\tau(t) &= \frac{(\alpha'(t) \times \alpha''(t)) \cdot \alpha'''(t)}{\|\alpha'(t) \times \alpha''(t)\|^2} = \frac{18\langle t^2-1, -2t, 1+t^2 \rangle \cdot 6\langle -1, 0, 1 \rangle}{(18\sqrt{2}(1+t^2))^2} \\
					&= \frac{108(-(t^2-1) + 0 + (1+t^2))}{648(1+t^2)^2} = \frac{108(2)}{648(1+t^2)^2} = \frac{216}{648(1+t^2)^2} = \frac{1}{3(1+t^2)^2}
				\end{align*}
			\end{proof}
			
			\begin{example} \label{ex:47}
				იპოვეთ $\alpha(t) = \langle t-\sin t, 1-\cos t, 4\sin(t/2) \rangle$ წირის სიმრუდე და გრეხა $t=\pi$ წერტილში.	
			\end{example}
			\begin{proof}[ამოხსნა]
				პირველ რიგში, ვიპოვოთ წარმოებულები ზოგადი $t$-თვის:
				\begin{align*}
					\alpha'(t) &= \langle 1-\cos t, \sin t, 2\cos(t/2) \rangle \\
					\alpha''(t) &= \langle \sin t, \cos t, -\sin(t/2) \rangle \\
					\alpha'''(t) &= \langle \cos t, -\sin t, -1/2 \cos(t/2) \rangle
				\end{align*}
				ახლა ჩავსვათ $t=\pi$:
				\begin{align*}
					\alpha'(\pi) &= \langle 1-(-1), 0, 0 \rangle = \langle 2, 0, 0 \rangle \\
					\alpha''(\pi) &= \langle 0, -1, -1 \rangle \\
					\alpha'''(\pi) &= \langle -1, 0, 0 \rangle
				\end{align*}
				გამოვთვალოთ საჭირო სიდიდეები $t=\pi$ წერტილში:
				\begin{align*}
					\|\alpha'(\pi)\| &= \sqrt{2^2+0^2+0^2} = 2 \\
					\alpha'(\pi) \times \alpha''(\pi) &= 
					\begin{vmatrix}
						\mathbf{i} & \mathbf{j} & \mathbf{k} \\
						2 & 0 & 0 \\
						0 & -1 & -1
					\end{vmatrix}
					= \langle 0, 2, -2 \rangle \\
					\|\alpha'(\pi) \times \alpha''(\pi)\| &= \sqrt{0^2+2^2+(-2)^2} = \sqrt{8} = 2\sqrt{2} \\
					(\alpha'(\pi) \times \alpha''(\pi)) \cdot \alpha'''(\pi) &= \langle 0, 2, -2 \rangle \cdot \langle -1, 0, 0 \rangle = 0
				\end{align*}
				საბოლოოდ, ვწერთ პასუხს:
				\begin{align*}
					k(\pi) &= \frac{\|\alpha'(\pi) \times \alpha''(\pi)\|}{\|\alpha'(\pi)\|^3} = \frac{2\sqrt{2}}{2^3} = \frac{2\sqrt{2}}{8} = \frac{\sqrt{2}}{4} \\
					\tau(\pi) &= \frac{(\alpha'(\pi) \times \alpha''(\pi)) \cdot \alpha'''(\pi)}{\|\alpha'(\pi) \times \alpha''(\pi)\|^2} = \frac{0}{(2\sqrt{2})^2} = 0
				\end{align*}
			\end{proof}
			%--- დანარჩენი მაგალითები ფორმატირებულია, მაგრამ ამოხსნის გარეშე ---
			\begin{example} \label{ex:46}
				იპოვეთ $\alpha(t) = \langle 2t, t^2, t^3/3 \rangle$ წირის ფრენეს ბაზისი, სიმრუდე და გრეხა.	
			\end{example}
			
			\begin{example} \label{ex:48}
				იპოვეთ $\alpha(t) = \langle t+t^3, t, -2t+t^3 \rangle$ წირის სიმრუდე და გრეხა $t=-2$ წერტილში.	
			\end{example}
			
			\begin{example} \label{ex:49}
				იპოვეთ $x^2=3y$, $2xy=9z$ განტოლებებით მოცემული წირის სიმრუდე და გრეხა.	
			\end{example}
			
			\begin{example} \label{ex:50}
				იპოვეთ $x^2-y^2+z^2=1$, $y^2-2x+z=0$ განტოლებებით მოცემული წირის სიმრუდე და გრეხა $(1,1,1)$ წერტილში.	
			\end{example}
			
			\begin{example} \label{ex:51}
				იპოვეთ $y^2=x$, $x^2=z$ განტოლებებით მოცემული წირის სიმრუდე და გრეხა.	
			\end{example}
			
			\begin{example} \label{ex:52}
				იპოვეთ $x^2=2az$, $y^2=2bz$ განტოლებებით მოცემული წირის სიმრუდე და გრეხა.	
			\end{example}
			
			\begin{example} \label{ex:53}
				იპოვეთ შემდეგი წირების სიმრუდე და გრეხა:
				\begin{enumerate}
					\item[ა)] $\alpha(t) = \langle t-\sin t, 1-\cos t, 4\sin(t/2) \rangle$;
					\item[ბ)] $\alpha(t) = \langle e^t, e^{-t}, t\sqrt{2} \rangle$;
					\item[გ)] $\alpha(t) = \langle 2t, \ln t, t^2 \rangle$;
					\item[დ)] $\alpha(t) = \langle e^t\sin t, e^{t}\cos t, e^t \rangle$;
					\item[ე)] $\alpha(t) = \langle 3t-t^3, 3t^2, 3t+t^3 \rangle$;
					\item[ვ)] $\alpha(t) = \langle \cos^3 t, \sin^3 t, \cos 2t \rangle$;
				\end{enumerate}
			\end{example}
			
			\begin{example} \label{ex:54}
				ცნობილია, რომ $\alpha(t)$ წირის მიმხები სიბრტყე ნებისმიერ წერტილში  ერთი და იგივე წრფის პარალელურია. დაამტკიცეთ, რომ ეს წირი ბრტყელია.
			\end{example}
			
		
		
		
			\subsection{შერეული სავარჯიშოები}
		
		\begin{example}
			რას წარმოადგენს კომპლექსური სახით ჩაწერილი ბრტყელი წირი?
			
			ა) $r(u)=1+iu$;
			% წრფე
			
			ბ) $r(u)=e^{iu}$;
			% წრეწირი
			
			გ) $r(u)=\sqrt{1+iu}$;
			% ჰიპერბოლა
			% x+iy=\sqrt{1+iu};
			% x^2-y^2=1+u; x^2-y^2=1; 2xy=u; მეორე განტოლება პარამეტრის სკალაა.
			
			დ) $r(u)=\frac{1}{1-iu}$.
			% წრეწირი
			% r=\frac{1+iu}{1+u^2}=x+iy=\frac{1}{1+u^2}+i\frac{u}{1+u^2}
			% $(x-1/2)^2+y^2=1/4.
		\end{example}
		
		\medskip
		\medskip
		
		\begin{example}
			ყველა წირს აქვს ევოლუტა? მეორე ევოლუტა?
			% არა
			რა არის წრეწირის ევოლუტა? მეორე ევოლუტა?
			% წერტილი, მისი ევოლუტა არ არსებობს.
		\end{example}
		
		\medskip
		
		\begin{example}
			დაწერეთ წირის ინვოლუტას ინვოლუტას განტოლება.
			
			დაწერეთ წირის $n$-ური ინვოლუტას განტოლება.
		\end{example}
		
		\medskip
		\medskip
		
		\begin{example}
			განსაზღვრეთ ფუნქცია $f(u)$ ისე, რომ $r=r(u)$ წირი იყოს ბრტყელი
			\[ r=(a\cos u, a\sin u, f(u)). \]
		\end{example}
		
		\medskip
		
		\begin{example}
			იპოვეთ $r=r(u)$ წირის სიმრუდე და გრეხა
			\begin{align*}
				&r=(a(3u-u^3),3au^2, a(3u+u^3)); \\
				&r=(a(u-\sin u), a(1-\cos u), bu); \\
				&r=(a(1+\cos u), a\sin u, 2a\sin \frac{u}{2}).
			\end{align*}
		\end{example}
		
		\medskip
		
		\begin{example}
			იპოვეთ $r=r(u)$ წირის სფერული  სიმრუდის ცენტრის კოორდინატები
			\[ r=(a\cos u, a\sin u, a\cos 2u). \]
		\end{example}
		
		\medskip
		
		\begin{example}
			დაამტკიცეთ, რომ შემდეგი წირი ძევს სფეროზე
			\[ x=a\sin^2 u; \quad y=a\sin u \cos u; \quad z=a\cos u. \]
		\end{example}
		
		
		
		\begin{example}
			\label{ex:4}
			წირი $\alpha(t)=(e^t,e^{-t},\sqrt{2}t)$ ჰელიქსის მსგავსად  მუდმივად მაღლდება  პარამეტრის ზრდისას, მაგრამ ჰელიქსისგან განსხვავებით $Oxy$ სიბრტყეზე პროექტირდება $xy=1$ ჰიპერბოლაზე.
		\end{example}
		
		
		
		
		
		
		
		
		
		
		
		
  \newpage
		
		
		
		
	\subsection{წირთა  თეორიის  ფუნდამენტური  თეორემა}
	
	\begin{theorem}[წირთა თეორიის ფუნდამენტური თეორემა]
		ვთქვათ, $\kappa(s) > 0$ და $\tau(s)$ არის ნებისმიერი ორი გლუვი ფუნქცია $s \in I$-თვის. მაშინ არსებობს რეგულარული წირი $\gamma(s)$, რომლისთვისაც $s$ არის ნატურალური პარამეტრი, $\kappa(s)$ არის მისი სიმრუდე და $\tau(s)$ არის მისი გრეხა. ეს წირი ერთადერთია ევკლიდური სივრცის მოძრაობის (გადაადგილება და მობრუნება) სიზუსტით.
	\end{theorem}
	
	
	
	თეორემის დამტკიცებამდე განვიხილოთ  {\bf  თეორემის  მნიშვნელობა და მოტივაცია}
	
	ეს თეორემა ბევრად მეტია, ვიდრე უბრალოდ  ტექნიკური შედეგი. აი, რამდენიმე ფრაზა და იდეა, რომელიც მის ღრმა მნიშვნელობასა და სილამაზეს უსვამს ხაზს:
	
	\begin{itemize}
		\item \textbf{წირის გენეტიკური კოდი:} ეს თეორემა, არსებითად, სივრცული წირების „გენეტიკური კოდია“. ის გვეუბნება, რომ წირის მთელი უნიკალური სამგანზომილებიანი ფორმა სრულად კოდირებულია მხოლოდ ორ ლოკალურ, ერთგანზომილებიან ფუნქციაში — მის სიმრუდეში ($\kappa$) და გრეხაში ($\tau$). ეს არის გეომეტრიული ინფორმაციის საოცარი შეკუმშვა.
		
		\item \textbf{ლოკალურიდან გლობალურისკენ:} ეს თეორემა ლოკალური ანალიზის ტრიუმფია. ის გვიჩვენებს, რომ თუ ჩვენ ვიცით წირის ქცევა მის თითოეულ უსასრულოდ მცირე ნაწილზე (როგორ იხრება და იგრეხება), ჩვენ შეგვიძლია სრულად აღვადგინოთ მისი მთლიანი, გლობალური ფორმა.
		
		\item \textbf{გეომეტრიული დეტერმინიზმი:} თეორემა ამჟღავნებს საოცარ დეტერმინიზმს გეომეტრიაში. თუ თქვენ იცით, როგორ იწყება წირი (საწყისი წერტილი და საწყისი რეპერი) და რა არის მისი „მოძრაობის წესები“ (სიმრუდისა და გრეხის ფუნქციები), მაშინ მისი მთელი მომავალი ტრაექტორია ცალსახად და უნიკალურად არის განსაზღვრული. არ არსებობს ორი სხვადასხვა გზა, რომელიც ერთსა და იმავე წესებს ემორჩილება.
		
		\item \textbf{კონსტრუქციის ძალა:} ეს თეორემა არა მხოლოდ აღწერითია, არამედ კონსტრუქციულიც. ის გვაძლევს რეცეპტს, თუ როგორ ავაგოთ ნებისმიერი სირთულის წირი, თუკი ხელთ გვაქვს მხოლოდ მისი სიმრუდისა და გრეხის „ინსტრუქციები“. წარმოიდგინეთ, რომ შეგიძლიათ  აღწეროთ  ნებისმიერი ფორმის  ტრაექტორია, თუკი იცით მხოლოდ ორი ინსტრუქცია ყოველ ნაბიჯზე: „რამდენად მკვეთრად მოვუხვიო“ ($\kappa$) და „რამდენად მკვეთრად ავიხარო ან დავიხარო სიბრტყიდან“ ($\tau$).
		
		\item \textbf{პასუხი ფუნდამენტურ კითხვაზე:} თეორემა ფუნდამენტურად პასუხობს კითხვას: „რა განსაზღვრავს წირის ფორმას?“ პასუხი გასაოცრად მარტივი და ელეგანტურია: მხოლოდ მისი სიმრუდე და გრეხა. ყველა სხვა თვისება ამ ორი მახასიათებლიდან გამომდინარეობს. სწორედ ამიტომ ეწოდება მას „ფუნდამენტური“.
	\end{itemize}
	
	
	\begin{proof}
		დამტკიცება შედგება ორი ნაწილისგან: \textbf{არსებობა} და \textbf{ერთადერთობა}.
		
		\subsection*{ნაწილი 1: არსებობა}
		
		ჩვენი მიზანია, მოცემული $\kappa(s)$ და $\tau(s)$ ფუნქციების გამოყენებით ავაგოთ (შევქმნათ) შესაბამისი წირი $\gamma(s)$.
		
		\begin{enumerate}
			\item \textbf{ფრენეს სისტემის აგება:} განვიხილოთ ფრენე-სერეს ფორმულების მატრიცული ჩანაწერი, როგორც წრფივ დიფერენციალურ განტოლებათა სისტემა სამ უცნობ ვექტორ-ფუნქციაზე: $\mathbf{T}(s), \mathbf{N}(s), \mathbf{B}(s)$.
			\[
			\frac{d}{ds}
			\begin{pmatrix} \mathbf{T} \\ \mathbf{N} \\ \mathbf{B} \end{pmatrix}
			=
			\begin{pmatrix} 0 & \kappa(s) & 0 \\ -\kappa(s) & 0 & \tau(s) \\ 0 & -\tau(s) & 0 \end{pmatrix}
			\begin{pmatrix} \mathbf{T} \\ \mathbf{N} \\ \mathbf{B} \end{pmatrix}
			\]
			აღვნიშნოთ კოეფიციენტთა მატრიცა $A(s)$-ით.
			
			\item \textbf{საწყისი პირობების არჩევა:} ავირჩიოთ საწყისი წერტილი $s_0 \in I$. ამ წერტილში ჩვენ თავისუფლად შეგვიძლია ავირჩიოთ საწყისი რეპერი. ყველაზე მარტივია, ავირჩიოთ სტანდარტული ორთონორმირებული ბაზისი:
			\begin{align*}
				\mathbf{T}(s_0) &= (1, 0, 0) = \mathbf{i} \\
				\mathbf{N}(s_0) &= (0, 1, 0) = \mathbf{j} \\
				\mathbf{B}(s_0) &= (0, 0, 1) = \mathbf{k}
			\end{align*}
			
			\item \textbf{დიფერენციალურ განტოლებათა სისტემის ამოხსნა:} წრფივ დიფერენციალურ განტოლებათა თეორიის ფუნდამენტური თეორემის თანახმად, ზემოთ მოცემულ სისტემას, მოცემული საწყისი პირობებით, გააჩნია ერთადერთი გლუვი ამონახსნი $(\mathbf{T}(s), \mathbf{N}(s), \mathbf{B}(s))$ მთელ $I$ ინტერვალზე.
			
			\item \textbf{ამონახსნის ორთონორმირებულობის ჩვენება:} უნდა დავამტკიცოთ, რომ მიღებული ვექტორები $\{\mathbf{T}(s), \mathbf{N}(s), \mathbf{B}(s)\}$ ქმნიან ორთონორმირებულ ბაზისს ყველა $s \in I$-თვის. განვიხილოთ $3 \times 3$ მატრიცა $M(s)$, რომლის სვეტებია $\mathbf{T}(s), \mathbf{N}(s), \mathbf{B}(s)$. გვინდა ვაჩვენოთ, რომ $M(s)$ ორთოგონალური მატრიცაა, ანუ $M(s)^T M(s) = I_{3\times3}$. განვიხილოთ $(M^T M)' = (M')^T M + M^T M'$. რადგან $M' = A M$, მივიღებთ:
			\[ (M^T M)' = (AM)^T M + M^T (AM) = M^T A^T M + M^T A M = M^T (A^T + A) M \]
			ფრენეს მატრიცა $A(s)$ ანტისიმეტრიულია, ანუ $A^T = -A$. შესაბამისად, $A^T + A = 0$. აქედან გამომდინარე, $(M^T M)' = 0$, რაც ნიშნავს, რომ $M^T M$ მუდმივი მატრიცაა. რადგან $s_0$ წერტილში $M(s_0)$-ის სვეტებია $\mathbf{i}, \mathbf{j}, \mathbf{k}$, $M(s_0)^T M(s_0) = I_{3\times3}$. საბოლოოდ, $M(s)^T M(s) = I_{3\times3}$ ყველა $s$-თვის, რაც რეპერის ორთონორმირებულობას ნიშნავს.
			
			\item \textbf{წირის კონსტრუირება:} ახლა, როცა გვაქვს მხები ვექტორი $\mathbf{T}(s)$, შეგვიძლია ავაგოთ თავად წირი ინტეგრებით. ავირჩიოთ ნებისმიერი საწყისი წერტილი $\mathbf{p}_0 \in \mathbb{R}^3$ და განვსაზღვროთ:
			\[ \gamma(s) = \mathbf{p}_0 + \int_{s_0}^{s} \mathbf{T}(u) \, du \]
			
			\item \textbf{თვისებების შემოწმება:}
			\begin{itemize}
				\item $\gamma'(s) = \mathbf{T}(s)$. რადგან $\|\mathbf{T}(s)\|=1$, $s$ არის ნატურალური პარამეტრი.
				\item წირის სიმრუდეა $\|\gamma''(s)\| = \|\mathbf{T}'(s)\|$. რადგან ჩვენი სისტემიდან $\mathbf{T}'=\kappa\mathbf{N}$ და $\|\mathbf{N}\|=1$, $\kappa>0$, მივიღებთ $\|\mathbf{T}'(s)\| = \kappa(s)$.
				\item წირის გრეხა განისაზღვრება $\mathbf{B}' = -\tau\mathbf{N}$-ით, რაც ზუსტად ემთხვევა ჩვენს სისტემას.
			\end{itemize}
			ამრიგად, ჩვენ ავაგეთ წირი, რომელიც აკმაყოფილებს თეორემის ყველა პირობას.
		\end{enumerate}
		
		\subsection*{ნაწილი 2: ერთადერთობა}
		
		ახლა ვაჩვენოთ, რომ ნებისმიერი ორი წირი, $\gamma_1(s)$ და $\gamma_2(s)$, რომლებსაც ერთი და იგივე $\kappa(s)$ და $\tau(s)$ აქვთ, ერთმანეთისგან მხოლოდ ევკლიდური მოძრაობით (მყარი სხეულის გადაადგილება) განსხვავდება.
		\begin{enumerate}
			\item ვთქვათ, $\{\mathbf{T}_1, \mathbf{N}_1, \mathbf{B}_1\}$ და $\{\mathbf{T}_2, \mathbf{N}_2, \mathbf{B}_2\}$ არის ამ ორი წირის შესაბამისი ფრენეს რეპერები. ორივე რეპერი აკმაყოფილებს ფრენე-სერეს ერთსა და იმავე განტოლებათა სისტემას, რადგან $\kappa(s)$ და $\tau(s)$ ემთხვევა.
			
			\item ავირჩიოთ $s_0 \in I$. რადგან ევკლიდურ მოძრაობას შეუძლია ნებისმიერი წერტილი გადაიტანოს ნებისმიერ სხვა წერტილში და ნებისმიერი ორთონორმირებული რეპერი — ნებისმიერ სხვა რეპერში, ჩვენ შეგვიძლია ვიპოვოთ ისეთი უნიკალური ევკლიდური მოძრაობა $M(\mathbf{x}) = R\mathbf{x} + \mathbf{v}$ (სადაც $R$ მობრუნების მატრიცაა, $\mathbf{v}$ — ტრანსლაციის ვექტორი), რომელიც ერთდროულად ასწორებს საწყის წერტილებსა და საწყის რეპერებს:
			\[ M(\gamma_2(s_0)) = \gamma_1(s_0) \quad \text{და} \quad R(\mathbf{T}_2(s_0)) = \mathbf{T}_1(s_0), \quad \text{და ა.შ. რეპერისთვის.} \]
			
			\item განვიხილოთ ახალი წირი $\tilde{\gamma}_2(s) = M(\gamma_2(s))$. ევკლიდური მოძრაობები ინახავენ სკალარულ ნამრავლს და ვექტორულ ნამრავლს, ამიტომ $\tilde{\gamma}_2$ წირსაც ექნება იგივე სიმრუდე $\kappa(s)$ და გრეხა $\tau(s)$. მისი ფრენეს რეპერი $\{\tilde{\mathbf{T}}_2, \tilde{\mathbf{N}}_2, \tilde{\mathbf{B}}_2\}$ იქნება $R\{\mathbf{T}_2, \mathbf{N}_2, \mathbf{B}_2\}$.
			
			\item ახლა ჩვენ გვაქვს ორი წირი, $\gamma_1$ და $\tilde{\gamma}_2$, რომლებსაც აქვთ:
			\begin{itemize}
				\item ერთი და იგივე $\kappa(s)$ და $\tau(s)$.
				\item ერთი და იგივე საწყისი წერტილი: $\gamma_1(s_0) = \tilde{\gamma}_2(s_0)$.
				\item ერთი და იგივე საწყისი რეპერი: $\mathbf{T}_1(s_0) = \tilde{\mathbf{T}}_2(s_0)$, და ა.შ.
			\end{itemize}
			
			\item ორივე წირის რეპერი აკმაყოფილებს ფრენე-სერეს ერთსა და იმავე დიფერენციალურ განტოლებათა სისტემას ერთი და იგივე საწყისი პირობებით. დიფერენციალური განტოლებების ამოხსნის ერთადერთობის თეორემის თანახმად, მათი ამონახსნები უნდა ემთხვეოდეს, ანუ:
			\[ \mathbf{T}_1(s) = \tilde{\mathbf{T}}_2(s), \quad \mathbf{N}_1(s) = \tilde{\mathbf{N}}_2(s), \quad \mathbf{B}_1(s) = \tilde{\mathbf{B}}_2(s) \quad \text{ყველა } s \in I\text{-თვის.} \]
			
			\item რადგან წირების მხები ვექტორები ემთხვევა ($\gamma_1'(s) = \tilde{\gamma}_2'(s)$) და ისინი იწყებიან ერთი და იგივე წერტილიდან ($\gamma_1(s_0) = \tilde{\gamma}_2(s_0)$), მათი ინტეგრალებიც უნდა ემთხვეოდეს. მაშასადამე, $\gamma_1(s) = \tilde{\gamma}_2(s)$ ყველა $s \in I$-თვის.
			
			\item ეს ნიშნავს, რომ $\gamma_1(s) = M(\gamma_2(s))$, ანუ $\gamma_1$ და $\gamma_2$ წირები ერთმანეთისგან მხოლოდ ერთი ევკლიდური მოძრაობით განსხვავდებიან. ერთადერთობის ნაწილიც დამტკიცებულია.
		\end{enumerate}
	\end{proof}
	
	
	\bigskip
	
	
	
	
	
	
	
	
	\subsection{სივრცული წირის ევოლუტა, როგორც   წირის მიმხები სფეროების  ცენტრების სიმრავლე}
	
	
	აქ ვნახავთ, რომ წირის მიმხები სფეროს $\mathbf{c}(s)$ ცენტრის ფორმულაა:
	$$
	\mathbf{c}(s) = \mathbf{r}(s) + \frac{1}{\kappa}\mathbf{N} + \frac{1}{\tau}\frac{d(1/\kappa)}{ds}\mathbf{B}
	$$
	სადაც $\mathbf{r}(s)$ არის $s$ რკალის სიგრძით პარამეტრიზებული წირის რადიუს-ვექტორი, ხოლო $\mathbf{T}$, $\mathbf{N}$, $\mathbf{B}$, $\kappa$ და $\tau$ არის, შესაბამისად, ფრენეს რეპერის ვექტორები, სიმრუდე და გრეხა.
	
	\subsection*{1. მიმხები სფეროს  განსაზღვრება}
	მიმხები სფერო არის ერთადერთი სფერო, რომელსაც  მოცემულ წერტილში  წირთან მესამე რიგის შეხება აქვს. ეს ნიშნავს, რომ მანძილს წირის ნებისმიერი წერტილიდან სფეროს ცენტრამდე, $d(s) = |\mathbf{c} - \mathbf{r}(s)|$, აქვს მესამე რიგის მინიმუმი შეხების წერტილში. აქედან გამომდინარეობს, რომ მანძილის კვადრატის, $D(s) = d(s)^2$, პირველი სამი წარმოებული $s$ რკალის სიგრძით ნულის ტოლია შეხების წერტილში.
	
	\subsection*{2. პირველი და მეორე წარმოებულების პირობების გამოყვანა}
	მანძილის კვადრატის პირველი წარმოებულია:
	\begin{align*}
		\frac{dD}{ds} &= \frac{d}{ds}((\mathbf{c} - \mathbf{r}) \cdot (\mathbf{c} - \mathbf{r})) = 2(\mathbf{c}' - \mathbf{r}') \cdot (\mathbf{c} - \mathbf{r}) \\
		&= 2(-\mathbf{T}) \cdot (\mathbf{c} - \mathbf{r}) = 0
	\end{align*}
	ეს ნიშნავს, რომ ვექტორი წირის წერტილიდან სფეროს ცენტრამდე, $\mathbf{c}-\mathbf{r}$, მხები ვექტორის $\mathbf{T}$-ს პერპენდიკულარულია. მაშასადამე, ის ძევს მთავარი ნორმალისა ($\mathbf{N}$) და ბინორმალის ($\mathbf{B}$) ვექტორებით შედგენილ სიბრტყეში:
	$$
	\mathbf{c} - \mathbf{r}(s) = a\mathbf{N} + b\mathbf{B}
	$$
	მანძილის კვადრატის მეორე წარმოებულია:
	\begin{align*}
		\frac{d^2D}{ds^2} &= \frac{d}{ds}(2(\mathbf{c} - \mathbf{r})\cdot(-\mathbf{T})) \\
		&= 2(\mathbf{c}' - \mathbf{r}')\cdot(-\mathbf{T}) + 2(\mathbf{c} - \mathbf{r})\cdot(-\mathbf{T}') \\
		&= 2(-\mathbf{T})\cdot(-\mathbf{T}) - 2(\mathbf{c}-\mathbf{r})\cdot(\kappa\mathbf{N}) \\
		&= 2 - 2(a\mathbf{N} + b\mathbf{B})\cdot(\kappa\mathbf{N}) = 2 - 2a\kappa = 0
	\end{align*}
	აქედან მივიღებთ კოეფიციენტს $\mathbf{N}$ კომპონენტისთვის:
	$$
	a = \frac{1}{\kappa}
	$$
	
	\subsection*{3. მესამე წარმოებულის  პირობის გამოყვანა}
	
	მანძილის კვადრატის მესამე წარმოებული:
	$$
	\frac{d^3D}{ds^3} = \frac{d}{ds}\left(2-2a\kappa\right) = -2(a'\kappa + a\kappa') = 0
	$$
	რადგან $\mathbf{c} = \mathbf{r}(s) + a\mathbf{N} + b\mathbf{B}$ და $\mathbf{c}$ მუდმივია, გვაქვს $\mathbf{c}' = \mathbf{T} + a'\mathbf{N} + a\mathbf{N}' + b'\mathbf{B} + b\mathbf{B}'=0$.
	ფრენეს ფორმულების გამოყენებით:
	$(1-a\kappa)\mathbf{T} + (a'-b\tau)\mathbf{N} + (a\tau+b')\mathbf{B}=0$.
	ეს გვაძლევს $1-a\kappa=0 \implies a=1/\kappa$, რაც უკვე ვიცით. ასევე, $a'-b\tau = 0$ და $a\tau+b' = 0$.
	პირობა $\frac{d^3D}{ds^3}=0$ ეკვივალენტურია $(\mathbf{c}-\mathbf{r}) \cdot \frac{d^3\mathbf{r}}{ds^3}=0$ პირობისა.
	$\mathbf{r}''' = (\mathbf{T}')' = (\kappa\mathbf{N})' = \kappa'\mathbf{N}+\kappa\mathbf{N}' = \kappa'\mathbf{N}+\kappa(-\kappa\mathbf{T}+\tau\mathbf{B}) = -\kappa^2\mathbf{T}+\kappa'\mathbf{N}+\kappa\tau\mathbf{B}$.
	$(\frac{1}{\kappa}\mathbf{N}+b\mathbf{B})\cdot(-\kappa^2\mathbf{T}+\kappa'\mathbf{N}+\kappa\tau\mathbf{B})=0$.
	სკალარული ნამრავლის გამოთვლით მივიღებთ:
	$\frac{\kappa'}{\kappa}+b\kappa\tau=0$.
	აქედან, $b=-\frac{\kappa'}{\kappa^2\tau}$. თუ გამოვიყენებთ $\frac{d(1/\kappa)}{ds}=-\frac{\kappa'}{\kappa^2}$, მივიღებთ:
	$$
	b = \frac{1}{\tau} \frac{d(1/\kappa)}{ds}
	$$
	
	\subsection*{დასკვნა}
	მესამე რიგის შეხების პირობის დაკმაყოფილებით, ჩვენ განვსაზღვრეთ მიმხები სფეროს რადიუს-ვექტორის კომპონენტები. მაშასადამე, სფეროს ცენტრი მოცემულია ფორმულით:
	$$
	\mathbf{c}(s) = \mathbf{r}(s) + \frac{1}{\kappa}\mathbf{N} + \frac{1}{\tau}\frac{d(1/\kappa)}{ds}\mathbf{B}
	$$
	
	\newpage
	
	\section*{მაგალითი 1: წრიული  ჰელიქსი}
	
	მოცემულია წრიული ჰელიქსი $\mathbf{r}(t) = \langle a\cos t, a\sin t, bt \rangle$. იპოვეთ მისი ევოლუტა.
	
	\subsection*{ამოხსნა}
	\textbf{1. წირი ნატურალური  პარამეტრით ($s$)}:
	პირველ რიგში, ვიპოვოთ რკალის სიგრძე $s$. $\mathbf{r}(t)$-ის წარმოებულია $\mathbf{r}'(t) = \langle -a\sin t, a\cos t, b \rangle$.
	მისი ნორმაა $\|\mathbf{r}'(t)\| = \sqrt{a^2\sin^2t + a^2\cos^2t + b^2} = \sqrt{a^2+b^2} = c$.
	რადგან სიჩქარე მუდმივია, $s = ct \implies t = s/c$.
	წირის განტოლება ნატურალური პარამეტრით არის $\mathbf{r}(s) = \langle a\cos(s/c), a\sin(s/c), bs/c \rangle$.
	
	\textbf{2. ფრენეს აპარატი, სიმრუდე და გრეხა}:
	\begin{itemize}
		\item მხები ვექტორი ($\mathbf{T}$): $\mathbf{T}(s) = \mathbf{r}'(s) = \langle -\frac{a}{c}\sin(s/c), \frac{a}{c}\cos(s/c), \frac{b}{c} \rangle$.
		\item სიმრუდე ($\kappa$): $\kappa = \|\mathbf{T}'(s)\| = \|\langle -\frac{a}{c^2}\cos(s/c), -\frac{a}{c^2}\sin(s/c), 0 \rangle\| = \frac{a}{c^2} = \frac{a}{a^2+b^2}$.
		\item მთავარი ნორმალი ($\mathbf{N}$): $\mathbf{N} = \frac{1}{\kappa}\mathbf{T}'(s) = \langle -\cos(s/c), -\sin(s/c), 0 \rangle$.
		\item ბინორმალი ($\mathbf{B}$): $\mathbf{B} = \mathbf{T} \times \mathbf{N} = \langle \frac{b}{c}\sin(s/c), -\frac{b}{c}\cos(s/c), \frac{a}{c} \rangle$.
		\item გრეხა ($\tau$): $\tau = \mathbf{N} \cdot \mathbf{B}'(s) = \frac{b}{c^2} = \frac{b}{a^2+b^2}$.
	\end{itemize}
	
	\textbf{3. მიმხები სფეროს ცენტრი}:
	ცენტრის რადიუს-ვექტორი მოცემულია ფორმულით:
	$$
	\mathbf{c}(s) = \mathbf{r}(s) + \frac{1}{\kappa}\mathbf{N} + \frac{1}{\tau}\frac{d(1/\kappa)}{ds}\mathbf{B}
	$$
	რადგან $\kappa$ მუდმივია, $\frac{d(1/\kappa)}{ds} = 0$. ფორმულა მარტივდება:
	\begin{align*}
		\mathbf{c}(s) &= \mathbf{r}(s) + \frac{1}{\kappa}\mathbf{N} \\
		&= \langle a\cos\frac{s}{c}, a\sin\frac{s}{c}, \frac{bs}{c} \rangle + \frac{a^2+b^2}{a}\langle -\cos\frac{s}{c}, -\sin\frac{s}{c}, 0 \rangle \\
		&= \left\langle \left(a - \frac{a^2+b^2}{a}\right)\cos\frac{s}{c}, \left(a - \frac{a^2+b^2}{a}\right)\sin\frac{s}{c}, \frac{bs}{c} \right\rangle \\
		&= \left( -\frac{b^2}{a}\cos\frac{s}{c}, -\frac{b^2}{a}\sin\frac{s}{c}, \frac{bs}{c} \right)
	\end{align*}
	
	\textbf{4. ევოლუტის  პარამეტრული  განტოლება}:
	წირის ევოლუტა არის წირის მიმხები სფეროების ცენტრების სიმრავლე. წირის ევოლუტის  პარამეტრული  განტოლება იგივეა, რაც $\mathbf{c}(s)$-ის გამოსახულება.
	$$
	\mathbf{E}(s) = \left( -\frac{b^2}{a}\cos\frac{s}{c}, -\frac{b^2}{a}\sin\frac{s}{c}, \frac{bs}{c} \right)
	$$
	ეს ასევე წრიული  ჰელიქსია, რომელიც  საწყისი ჰელიქსის თანაღერძულია.
	
	
	
	\bigskip
	
		
	
	
	\bigskip
			
			\subsection{წირები სამყაროში }
			
			მზის სისტემაში პლანეტების, კომეტებისა და მეტეორების ორბიტები შემთხვევითი ტრაექტორიები კი არა, ბრტყელი  მეორე რიგის წირებია. ამის მიზეზი ფიზიკის ერთ-ერთი ფუნდამენტური  პრინციპი — კუთხოვანი იმპულსის შენახვის კანონია.
			
			ორბიტის სიბრტყეში მოქცევის ასახსნელად, განვიხილოთ სამი ძირითადი ცნება:
			
			\begin{itemize}
				\item \textbf{კუთხოვანი იმპულსი ($\mathbf{L}$):} ეს არის სხეულის ბრუნვითი მოძრაობის საზომი. ვექტორულად, ის განისაზღვრება, როგორც სხეულის პოზიცია-ვექტორის ($\mathbf{r}$) და მისი წრფივი იმპულსის ($\mathbf{p}=m\mathbf{v}$) ვექტორული ნამრავლი.
				\[ \mathbf{L} = \mathbf{r} \times \mathbf{p} \]
				
				\item \textbf{ბრუნვითი მომენტი ($\boldsymbol{\tau}$):} ეს არის ძალის ბრუნვითი ეფექტის საზომი.
				\[ \boldsymbol{\tau} = \mathbf{r} \times \mathbf{F} \]
				
				\item \textbf{ბრუნვის კანონი:} ნიუტონის მეორე კანონის ანალოგი ბრუნვითი მოძრაობისთვის, რომელიც ამბობს, რომ კუთხოვანი იმპულსის ცვლილების სიჩქარე მოქმედი ბრუნვითი მომენტის ტოლია.
				\[ \frac{d\mathbf{L}}{dt} = \boldsymbol{\tau} \]
			\end{itemize}
			
	
			
\textbf{გრავიტაცია ცენტრალური ძალაა:} ორსხეულიან სისტემაში (მაგ. პლანეტა და მზე), გრავიტაციული ძალა ($\mathbf{F}$) ყოველთვის მიმართულია სხეულების დამაკავშირებელი ხაზის გასწვრივ. ეს ნიშნავს, რომ ძალის ვექტორი $\mathbf{F}$ ყოველთვის პოზიცია-ვექტორ $\mathbf{r}$-ის პარალელურია.
				
\textbf{ბრუნვითი მომენტი ნულის ტოლია:} რადგან $\mathbf{F}$ და $\mathbf{r}$ პარალელური ვექტორებია, მათი ვექტორული ნამრავლი ნულია. შესაბამისად, სისტემაზე არ მოქმედებს გარეგანი ბრუნვითი მომენტი.
				\[ \boldsymbol{\tau} = \mathbf{r} \times \mathbf{F} = 0 \]
				
\textbf{კუთხოვანი იმპულსი ინახება:} ბრუნვის კანონის თანახმად, თუ $\boldsymbol{\tau}=0$, მაშინ $\frac{d\mathbf{L}}{dt}=0$. ეს ნიშნავს, რომ კუთხოვანი იმპულსის ვექტორი $\mathbf{L}$ დროში არ იცვლება — ის მუდმივია როგორც სიდიდით, ისე მიმართულებით.

			
ვექტორული ნამრავლის განმარტების თანახმად, $\mathbf{L}$ ვექტორი ყოველთვის ორთოგონალურია როგორც $\mathbf{r}$ ვექტორის, ასევე $\mathbf{p}$ ვექტორის მიმართ. თუ $\mathbf{L}$ ვექტორი დროში არ იცვლის მიმართულებას, ეს იმას ნიშნავს, რომ $\mathbf{r}$ და $\mathbf{p}$ ვექტორები მუდამ უნდა დარჩნენ ერთ, ფიქსირებულ სიბრტყეში, რომელიც $\mathbf{L}$-ის პერპენდიკულარულია. სწორედ ეს არის \textbf{ორბიტალური სიბრტყე}.
			
			
უფრო  დიდ მასშტაბებზე გადასვლისას, მოძრაობის სურათი რთულდება, თუმცა ფუნდამენტური პრინციპები კვლავ მოქმედებს.
			
			
			
ჩვენი მზე არ მოძრაობს იდეალურად  ბრტყელ  ორბიტაზე. მისი ტრაექტორია ორი მოძრაობის კომბინაციაა:
			
\begin{enumerate}
				\item \textbf{წრიული ორბიტა:} მზე დაახლოებით  230 კმ/წმ სიჩქარით მოძრაობს გალაქტიკის ცენტრის გარშემო, რომელსაც დაახლოებით 230 მილიონ წელიწადში (ერთ „კოსმოსურ წელში“) ერთხელ უვლის გარს.
				\item \textbf{ვერტიკალური რხევა:} გალაქტიკის დისკის გრავიტაციული ველი მზეს აიძულებს, ირხეოდეს ზევით და ქვევით დისკის ცენტრალური სიბრტყის მიმართ, ქანქარის მსგავსად.
\end{enumerate}
			ამ ორი მოძრაობის შეერთება ქმნის მზის გრძელვადიან, ძალიან დიდი მასშტაბის სპირალურ ტრაექტორიას გალაქტიკაში.
			
			
	
		
		\newpage
		
		%======================================================================
		\section{ზედაპირები}
		%======================================================================
		%\addcontentsline{toc}{section}{ზედაპირები და პირველი ფუნდამენტური ფორმა}
		
	
		\textbf{რეგულარული პარამეტრიზებული ზედაპირი} არის ზედაპირი სამგანზომილებიან სივრცეში ($\mathbb{R}^3$), რომელიც ყოველ წერტილში გლუვი და არაგადაგვარებულია. ეს ფორმალურად განისაზღვრება \textbf{რეგულარული კოორდინატთა არეს} ცნებით.
		
		\begin{definition}[რეგულარული არე]
			დავუშვათ $D \subset \mathbb{R}^2$ ღია სიმრავლეა და განვიხილოთ ვექტორ-ფუნქცია $\mathbf{r}: D \to \mathbb{R}^3$, რომელიც მოიცემა როგორც $\mathbf{r}(u,v)=(x(u,v),y(u,v),z(u,v))$.
			
			ამ პარამეტრიზაციას ეწოდება \textbf{რეგულარული არე} (ან კოორდინატთა არე), თუ იგი აკმაყოფილებს შემდეგ სამ პირობას:
			\begin{enumerate}
				\item \textbf{სირგლუვე:} ფუნქცია $\mathbf{r}(u, v)$ არის უწყვეტად დიფერენცირებადი ($C^1$ ან უფრო მაღალი კლასის).
				
				\item \textbf{ინექციურობა:} ფუნქცია $\mathbf{r}(u, v)$ არის ცალსახა (ჩადგმა). ეს ნიშნავს, რომ მას არ აქვს თვითგადაკვეთები.
				
				\item \textbf{რეგულარულობის პირობა:} კერძო წარმოებულების ვექტორები
				\[
				\mathbf{r}_u = \frac{\partial \mathbf{r}}{\partial u} \quad \text{და} \quad \mathbf{r}_v = \frac{\partial \mathbf{r}}{\partial v}
				\]
				წრფივად დამოუკიდებელია ყოველი $(u, v) \in D$ წერტილისთვის. ეს ეკვივალენტურია იმისა, რომ მათი ვექტორული ნამრავლი არასდროსაა ნულოვანი ვექტორი:
				\[
				\mathbf{r}_u \times \mathbf{r}_v \neq \mathbf{0}
				\]
			\end{enumerate}
		\end{definition}
		
		$S \subset \mathbb{R}^3$ ზედაპირს ეწოდება \textbf{რეგულარული ზედაპირი}, თუ $S$-ის ყოველ წერტილს გააჩნია მიდამო, რომელიც წარმოადგენს რეგულარული არეს ანასახს.
		
		\subsection{გეომეტრიული მნიშვნელობა}
		რეგულარულობის პირობა გადამწყვეტია, რადგან $\mathbf{r}_u$ და $\mathbf{r}_v$ ვექტორები წარმოქმნიან ზედაპირის \textbf{მხებ სიბრტყეს} $\mathbf{r}(u, v)$ წერტილში. პირობა $\mathbf{r}_u \times \mathbf{r}_v \neq \mathbf{0}$ უზრუნველყოფს, რომ ეს მხები სიბრტყე ყოველთვის კარგად განმარტებული, ორგანზომილებიანი სიბრტყეა. ეს ხელს უშლის ზედაპირზე წამახვილებული წერტილების, კუსპების ან წიბოების არსებობას, სადაც მხები სიბრტყე გადაგვარებული ან განუსაზღვრელი იქნებოდა.
		
		
		
			
		\subsection*{ზედაპირის მხები სიბრტყის  განსაზღვრება}
		
		
		
		
		
		\begin{definition}
			ვთქვათ, გვაქვს რეგულარული   პარამეტრიზებული  $S$ ზედაპირი, რომელიც მოცემულია ვექტორ-ფუნქციით $\mathbf{r}(u, v)$ და განვიხილოთ  ამ ზედაპირზე მდებარე  $P_0 = \mathbf{r}(u_0, v_0)$ წერტილი .
			 $\mathbf{v}$ ვექტორს ეწოდება $S$ ზედაპირის \textbf{მხები ვექტორი} $P_0$ წერტილში, თუ არსებობს რაიმე დიფერენცირებადი წირი $\alpha(t)$, რომელიც:
			\begin{enumerate}
				\item მთლიანად ძევს $S$ ზედაპირზე.
				\item გადის $P_0$ წერტილზე (მაგალითად, $\alpha(0) = P_0$).
				\item მისი სიჩქარის ვექტორი $t=0$ მომენტში არის $\mathbf{v}$ (ანუ, $\alpha'(0) = \mathbf{v}$).
			\end{enumerate}
		\end{definition}
		
		
		ზედაპირის \textbf{მხები სიბრტყე} $P_0$ წერტილში არის ისეთი სიბრტყე, რომელიც:
		\begin{enumerate}
			\item გადის $P_0$ წერტილზე.
			\item შეიცავს ზედაპირის ყველა შესაძლო მხებ ვექტორს ამ წერტილში.
		\end{enumerate}
	
	 ზედაპირის ნორმალური  ვექტორი $\mathbf{N}$ არის $\mathbf{r}_u$ და $\mathbf{r}_v$ ვექტორების ვექტორული  ნამრავლი: $\mathbf{N} = \mathbf{r}_u \times \mathbf{r}_v$ (გამოთვლილი $(u_0, v_0)$ წერტილში).
		
	მხები სიბრტყის ზოგადი განტოლება $P_0 = (x_0, y_0, z_0)$ წერტილში, სადაც ნორმალია $\mathbf{N} = \langle A, B, C \rangle$, მოიცემა ფორმულით:
		\[
		A(x - x_0) + B(y - y_0) + C(z - z_0) = 0.
		\]
	
	
		
		\subsection*{საკოორდინატო წირები და მათი მხები ვექტორები}
		
		
		
		წარმოვიდგინოთ პარამეტრების არე $D$ როგორც ბრტყელი რუკა. ამ რუკაზე $v=v_0$ განტოლებით მოცემული სიმრავლე არის ჰორიზონტალური \textbf{წრფე}, ხოლო $u=u_0$ განტოლებით მოცემული სიმრავლე — ვერტიკალური \textbf{წრფე}.
		
		ფუნქცია $\mathbf{r}(u,v)$ ამ ბრტყელ წრფეებს \textbf{ასახავს} სამგანზომილებიან სივრცეში და აქცევს მათ \textbf{წირებად}, რომლებიც უკვე ზედაპირზე მდებარეობენ.
		\begin{itemize}
			\item ჰორიზონტალური წრფე $(u, v_0)$ გარდაიქმნება \textbf{$u$-საკოორდინატო წირად} ზედაპირზე.
			\item ვერტიკალური წრფე $(u_0, v)$ გარდაიქმნება \textbf{$v$-საკოორდინატო წირად} ზედაპირზე.
		\end{itemize}
		სწორედ ამიტომ, ამ წირებს საკოორდინატო წირები ეწოდება. ისინი ქმნიან ბადეს ზედაპირზე, რომელიც წარმოადგენს საწყისი, ბრტყელი საკოორდინატო ბადის ანასახს.
		
		\subsection*{რატომ  არიან $\mathbf{r}_u$ და $\mathbf{r}_v$ მხები ვექტორები?}
		ზემოთქმულიდან გამომდინარე, $\mathbf{r}_u$ და $\mathbf{r}_v$ ვექტორები არიან სწორედ ამ საკოორდინატო წირების სიჩქარის ვექტორები. ეს შეგვიძლია ვაჩვენოთ ზოგადი სანსაზღვრების გამოყენებით:
		
		\begin{itemize}
			\item \textbf{$\mathbf{r}_u$ ვექტორი:} განვიხილოთ $u$-საკოორდინატო წირი $\alpha(t) = \mathbf{r}(u_0 + t, v_0)$. ეს წირი ძევს ზედაპირზე და $t=0$ მომენტში გადის $P_0$ წერტილს. მისი სიჩქარის ვექტორი $t=0$ მომენტში არის $\alpha'(0) = \mathbf{r}_u(u_0, v_0)$. მაშასადამე, $\mathbf{r}_u$ არის მხები ვექტორი.
			
			\item \textbf{$\mathbf{r}_v$ ვექტორი:} ანალოგიურად, $v$-საკოორდინატო წირის, $\beta(t) = \mathbf{r}(u_0, v_0 + t)$, სიჩქარის ვექტორი $t=0$ მომენტში არის $\beta'(0) = \mathbf{r}_v(u_0, v_0)$. მაშასადამე, $\mathbf{r}_v$ ასევე მხები ვექტორია.
		\end{itemize}
		
		\subsection*{რატომ წარმოქმნიან $\mathbf{r}_u$ და $\mathbf{r}_v$ მხებ სიბრტყეს?}
		\begin{enumerate}
			\item ჩვენ უკვე ვიცით, რომ $\mathbf{r}_u$ და $\mathbf{r}_v$ ორივე მხები ვექტორია $P_0$ წერტილში.
			
			\item რადგან ზედაპირი \textbf{რეგულარულია},  $\mathbf{r}_u \times \mathbf{r}_v \neq \mathbf{0}$. ეს ნიშნავს, რომ ეს ორი ვექტორი \textbf{არ არის პარალელური} (წრფივად დამოუკიდებელია). 
			
			\item სამგანზომილებიან სივრცეში ნებისმიერი ორი არაპარალელური ვექტორი, რომლებსაც საერთო სათავე აქვთ, ცალსახად განსაზღვრავს სიბრტყეს და არის მისი ბაზისი.
		\end{enumerate}
		სწორედ ეს სიბრტყე, რომელიც გადის $P_0$ წერტილზე და შეიცავს $\mathbf{r}_u$ და $\mathbf{r}_v$ ვექტორებს, არის განსაზღვრების თანახმად ზედაპირის მხები სიბრტყე.
		
		\subsection*{მაგალითი}
		ვიპოვოთ პარაბოლოიდის $z = x^2 + y^2$ მხები სიბრტყე $(1, 2, 5)$ წერტილში.
		\begin{enumerate}
			\item \textbf{პარამეტრიზაცია:} ზედაპირის ბუნებრივი პარამეტრიზაციაა $x = u$, $y = v$. მაშინ $z = u^2 + v^2$.
			\[
			\mathbf{r}(u, v) = \langle u, v, u^2 + v^2 \rangle
			\]
			წერტილი $(1, 2, 5)$ შეესაბამება პარამეტრებს $(u, v) = (1, 2)$.
			
			\item \textbf{კერძო წარმოებულების გამოთვლა:}
			\[
			\mathbf{r}_u = \frac{\partial \mathbf{r}}{\partial u} = \langle 1, 0, 2u \rangle, \quad
			\mathbf{r}_v = \frac{\partial \mathbf{r}}{\partial v} = \langle 0, 1, 2v \rangle
			\]
			
			\item \textbf{მხები ვექტორები წერტილში (1, 2):}
			\[
			\mathbf{r}_u(1, 2) = \langle 1, 0, 2(1) \rangle = \langle 1, 0, 2 \rangle
			\]
			\[
			\mathbf{r}_v(1, 2) = \langle 0, 1, 2(2) \rangle = \langle 0, 1, 4 \rangle
			\]
			
			\item \textbf{ნორმალური ვექტორის გამოთვლა:}
			\[
			\mathbf{N} = \mathbf{r}_u \times \mathbf{r}_v = 
			\begin{vmatrix} 
				\mathbf{i} & \mathbf{j} & \mathbf{k} \\ 
				1 & 0 & 2 \\ 
				0 & 1 & 4 
			\end{vmatrix} 
			= \mathbf{i}(0-2) - \mathbf{j}(4-0) + \mathbf{k}(1-0) = \langle -2, -4, 1 \rangle
			\]
			
			\item \textbf{მხები სიბრტყის განტოლება:} ვიყენებთ წერტილს $P_0 = (1, 2, 5)$ და ნორმალს $\mathbf{N} = \langle -2, -4, 1 \rangle$.
			\begin{align*}
				-2(x - 1) - 4(y - 2) + 1(z - 5) &= 0 \\
				-2x + 2 - 4y + 8 + z - 5 &= 0 \\
				-2x - 4y + z + 5 &= 0
			\end{align*}
			და, საბოლოოდ
			\[
			2x + 4y - z = 5
			\]
		\end{enumerate}
		
		
		
		
		
		
		
		\subsection{მაგალითები}
		
		\subsubsection{სფერო (რეგულარული)}
		ერთეულოვანი სფეროს სტანდარტული პარამეტრიზაციაა 
		\[
		\mathbf{r}(\theta, \phi) = \langle \sin(\phi)\cos(\theta), \sin(\phi)\sin(\theta), \cos(\phi) \rangle.
		\]
		ეს კოორდინატთა არე რეგულარულია $D = \{(\theta, \phi) \mid 0 < \theta < 2\pi, 0 < \phi < \pi\}$ არეზე. სფერო რეგულარული ზედაპირია, რადგან ჩვენ შეგვიძლია დავფაროთ პოლუსები (სადაც ეს არე არ არის რეგულარული) სხვა, შებრუნებული რეგულარული არეებით.
		
		\subsubsection{კონუსი (არარეგულარული წვეროში)}
		პარამეტრიზაცია $\mathbf{r}(u, v) = \langle v \cos(u), v \sin(u), v \rangle$ აღწერს კონუსს. წვეროში, სადაც $v=0$, გვაქვს $\mathbf{r}(u, 0) = \langle 0, 0, 0 \rangle$. კერძო წარმოებულებია $\mathbf{r}_u = \langle -v\sin(u), v\cos(u), 0 \rangle$ და $\mathbf{r}_v = \langle \cos(u), \sin(u), 1 \rangle$.
		როდესაც $v=0$, მაშინ $\mathbf{r}_u = \langle 0, 0, 0 \rangle$, შესაბამისად ვექტორული ნამრავლი $\mathbf{r}_u \times \mathbf{r}_v = \mathbf{0}$. ზედაპირი არ არის რეგულარული წვეროში, რადგან მას არ გააჩნია კარგად განმარტებული მხები სიბრტყე.
			
			
			
			
		
		\subsection{პირველი ფუნდამენტური ფორმა}
		\begin{definition}[პირველი ფუნდამენტური ფორმა]
			$p \in S$ წერტილში \textbf{პირველი ფუნდამენტური ფორმა} $I_p$ არის ბილინიარული, სიმეტრიული ფორმა $I_p: T_pS \times T_pS \to \mathbb{R}$, რომელიც განისაზღვრება, როგორც სკალარული ნამრავლი $\mathbb{R}^3$-ში:
			\[ I_p(\mathbf{w}_1, \mathbf{w}_2) = \mathbf{w}_1 \cdot \mathbf{w}_2 \]
		\end{definition}
		შესაბამისი კვადრატული ფორმაა $I_p(\mathbf{w}) = \mathbf{w} \cdot \mathbf{w} = \|\mathbf{w}\|^2$.
		$\{\mathbf{r}_u, \mathbf{r}_v\}$ ბაზისში, ამ ფორმის მატრიცას აქვს კომპონენტები:
		\begin{align*}
			g_{11} &= I(\mathbf{r}_u, \mathbf{r}_u) = \mathbf{r}_u \cdot \mathbf{r}_u = E \\
			g_{12} &= I(\mathbf{r}_u, \mathbf{r}_v) = \mathbf{r}_u \cdot \mathbf{r}_v = F \\
			g_{21} &= I(\mathbf{r}_v, \mathbf{r}_u) = \mathbf{r}_v \cdot \mathbf{r}_u = F \\
			g_{22} &= I(\mathbf{r}_v, \mathbf{r}_v) = \mathbf{r}_v \cdot \mathbf{r}_v = G
		\end{align*}
		ამ მატრიცას $\begin{pmatrix} E & F \\ F & G \end{pmatrix}$ ეწოდება \textbf{მეტრიკული ტენზორი} მოცემულ კოორდინატებში.
		უსასრულოდ მცირე გადაადგილების $d\mathbf{r} = \mathbf{r}_u du + \mathbf{r}_v dv$ სიგრძის კვადრატი, $ds^2$, არის:
		\[
		ds^2 = E \, du^2 + 2F \, du \, dv + G \, dv^2
		\]
		
		\subsection{შინაგანი გეომეტრიის  ელემენტები}
		\begin{itemize}
			\item \textbf{წირის სიგრძე:} $$ L(\gamma) = \int_{a}^{b} \sqrt{E(u')^2 + 2F u'v' + G(v')^2} \, dt.$$
			\item \textbf{კუთხე ორ ვექტორს შორის:} $$\cos(\theta) = \frac{E a_1 a_2 + F(a_1 b_2 + a_2 b_1) + G b_1 b_2}{\sqrt{E a_1^2 + 2F a_1 b_1 + G b_1^2}\sqrt{E a_2^2 + 2F a_2 b_2 + G b_2^2}}.$$
			\item \textbf{ზედაპირის ფართობი:} $A(R) = \iint_{R} \sqrt{EG - F^2} \, du \, dv$.
		\end{itemize}
	
	
	
	\subsection*{დეტალური   განმარტება}
	
	ზემოთ მოცემული ფორმულები არის პირველი ფუნდამენტური ფორმის უშუალო გამოყენება ზედაპირის გეომეტრიული თვისებების გასაზომად.
	
	\subsubsection*{1. წირის სიგრძე}
	
	ეს ფორმულა გამოიყენება  $\mathbf{r}=\mathbf{r}(u, v)$ \textbf{ზედაპირზე მდებარე} $\gamma$ წირის სიგრძის გამოსათვლელად.
	\[ L(\gamma) = \int_{a}^{b} \sqrt{E(u')^2 + 2F u'v' + G(v')^2} \, dt \]
	აქ:
	\begin{itemize}
		\item $\gamma$ არის წირი, რომელიც მთლიანად ძევს ჩვენს $S$ ზედაპირზე. მისი პარამეტრული განტოლებაა:
		\[ \gamma(t) = \mathbf{r}(u(t), v(t)), \quad t \in [a, b] \]
		
		\item $u(t)$ და $v(t)$ არის ფუნქციები, რომლებიც გვიჩვენებს, თუ ზედაპირის რომელი $(u,v)$ წერტილი შეესაბამება $t$ მომენტს.
		
		\item $u'$ და $v'$ არის ამ ფუნქციების წარმოებულები $t$-თი: $u' = \frac{du}{dt}$ და $v' = \frac{dv}{dt}$. ამ სიდიდეების გეომეტრიული შინაარსის გასაგებად, გავაწარმოოთ $\gamma(t)$ ჯაჭვური წესით, რომ ვიპოვოთ მისი სიჩქარის ვექტორი $\gamma'(t)$:
		\[ \gamma'(t) = \frac{d\mathbf{r}}{dt} = \frac{\partial \mathbf{r}}{\partial u}\frac{du}{dt} + \frac{\partial \mathbf{r}}{\partial v}\frac{dv}{dt} = u'(t)\mathbf{r}_u + v'(t)\mathbf{r}_v \]
		ეს განტოლება ცხადად გვიჩვენებს, რომ სიდიდეები $u'$ და $v'$ არის წირის  სიჩქარის ვექტორის, ანუ $\gamma'(t)$-ის \textbf{კოორდინატები} $\{\mathbf{r}_u, \mathbf{r}_v\}$ ბაზისში. ისინი გვიჩვენებენ, თუ წირის მოძრაობის რა ნაწილი მოდის $u$-მიმართულებაზე და რა ნაწილი — $v$-მიმართულებაზე.
	\end{itemize}
	
	
	
	\subsubsection*{2. კუთხე ორ ვექტორს შორის}
	ეს ფორმულა  ზომავს კუთხეს ზედაპირის ერთსა და იმავე წერტილის მხებ სიბრტყეში მდებარე \textbf{ორ $\mathbf{w}_1$ და $\mathbf{w}_2$ ვექტორს} შორის.
	\[ \cos(\theta) = \frac{E a_1 a_2 + F(a_1 b_2 + a_2 b_1) + G b_1 b_2}{\sqrt{E a_1^2 + 2F a_1 b_1 + G b_1^2}\sqrt{E a_2^2 + 2F a_2 b_2 + G b_2^2}} \]
	აქ:
	\begin{itemize}
		\item $\mathbf{w}_1$ და $\mathbf{w}_2$ არის ორი ვექტორი, რომლებიც მიეკუთვნება $p$ წერტილში არსებულ მხებ სიბრტყეს, $T_pS$.
		\item ამ მხებ სიბრტყეში ბაზისს ქმნის ვექტორული წყვილი $\{\mathbf{r}_u, \mathbf{r}_v\}$. შესაბამისად, ნებისმიერი მხები ვექტორი შეგვიძლია ჩავწეროთ ამ ბაზისში:
		\begin{align*}
			\mathbf{w}_1 &= a_1 \mathbf{r}_u + b_1 \mathbf{r}_v \\
			\mathbf{w}_2 &= a_2 \mathbf{r}_u + b_2 \mathbf{r}_v
		\end{align*}
		\item $(a_1, b_1)$ და $(a_2, b_2)$ არის ამ ვექტორების \textbf{კოორდინატები} $\{\mathbf{r}_u, \mathbf{r}_v\}$ ბაზისში. ფორმულა საშუალებას გვაძლევს, ვიპოვოთ კუთხე ამ ვექტორებს შორის მხოლოდ მათი ლოკალური კოორდინატებისა და ზედაპირის $E,F,G$ კოეფიციენტების ცოდნით.
	\end{itemize}
	
	\subsubsection*{3. ზედაპირის ფართობი}
	ეს ფორმულა გვაძლევს \textbf{ზედაპირის $S_D$ ნაწილის} ფართობს.
	\[ A(S_D) = \iint_{D} \sqrt{EG - F^2} \, du \, dv \]
	აქ:
	\begin{itemize}
		\item $S_D$ არის კონკრეტული რეგიონი ან „ნაჭერი“, რომელიც მდებარეობს ჩვენს $S$ \textbf{ზედაპირზე} და რომლის ფართობის გამოთვლაც გვსურს.
		\item $D$ არის შესაბამისი რეგიონი \textbf{პარამეტრების $(u,v)$ სიბრტყეზე}. პარამეტრიზაცია $\mathbf{r}(u,v)$ ასახავს ბრტყელ $D$ რეგიონს გამრუდებულ $S_D$ რეგიონზე ზედაპირზე.
		\item ინტეგრება ხდება ბრტყელ $D$ რეგიონზე $u$ და $v$ პარამეტრების მიხედვით.
		\item $\sqrt{EG - F^2}$ წევრი არის „გამრუდების ფაქტორი“ ანუ იაკობიანი. ის გვიჩვენებს, თუ რა ფართობისაა ზედაპირზე ის უსასრულოდ მცირე პარალელოგრამი, რომელიც შეესაბამება $(u,v)$ სიბრტყეზე არსებულ $du \times dv$ ფართობის მქონე მართკუთხედს.
	\end{itemize}
		
		\bigskip
		

		განვიხილოთ  ამ ფორმულების  გამოყენება.
		
		\subsubsection*{მაგალითი 1: წირის სიგრძის გამოთვლა} 
		მოცემულია  ჰელიქსი(სპირალი)
		
		\[
		\gamma(t) = \left( r\cos(t), r\sin(t), \frac{h}{2\pi}t \right)
		\]
		
		სადაც:
		\begin{itemize}
			\item \textbf{$(r\cos(t), r\sin(t))$} კომპონენტები აღწერს წრიულ მოძრაობას $xy$-სიბრტყეში $r$ რადიუსით.
			\item \textbf{$\frac{h}{2\pi}t$} კომპონენტი აღწერს წრფივ, თანაბარ მოძრაობას $z$-ღერძის გასწვრივ.
		\end{itemize}
		
		პარამეტრები კი შემდეგნაირად განიმარტება:
		\begin{itemize}
			\item \textbf{$r$} — არის ცილინდრის რადიუსი, რომელზეც „დახვეულია“ სპირალი.
			\item \textbf{$h$} — არის სიმაღლე, რომელსაც  ჰელიქსი აღწევს ერთი სრული ბრუნის ($t=2\pi$) შემდეგ. ამ სიდიდეს ჰელიქსის „ბიჯი“ ეწოდება.
			\item \textbf{$t$} — არის პარამეტრი, რომელიც ამ შემთხვევაში ემთხვევა $\theta$ კუთხეს რადიანებში.
		\end{itemize}
		
		ვიპოვოთ $R$ რადიუსიან ცილინდრზე  არსებული (ჰელიქსის) ერთი სრული  ბიჯის სიგრზე.
		\vspace{1em}
		
		ჰელიქსი აღიწერება პარამეტრებით: $\theta(t) = t$ და $z(t) = \frac{h}{2\pi}t$, სადაც $t$ იცვლება $0$-დან $2\pi$-მდე.
		ამ შემთხვევაში, $u=\theta$ და $v=z$, ამიტომ:
		\[ u'(t) = \frac{d\theta}{dt} = 1 \quad \text{და} \quad v'(t) = \frac{dz}{dt} = \frac{h}{2\pi} \]
		ჩავსვათ ეს მნიშვნელობები და $E=r^2, F=0, G=1$ კოეფიციენტები სიგრძის ფორმულაში:
		\begin{align*}
			L(\gamma) &= \int_{0}^{2\pi} \sqrt{E(u')^2 + 2F u'v' + G(v')^2} \, dt \\
			&= \int_{0}^{2\pi} \sqrt{r^2(1)^2 + 2(0)(1)(\frac{h}{2\pi}) + 1(\frac{h}{2\pi})^2} \, dt \\
			&= \int_{0}^{2\pi} \sqrt{r^2 + \frac{h^2}{4\pi^2}} \, dt
		\end{align*}
		რადგან ფესვის ქვეშ მყოფი გამოსახულება მუდმივია, ინტეგრალი მარტივდება:
		\[ L(\gamma) = \sqrt{r^2 + \frac{h^2}{4\pi^2}} \cdot \int_{0}^{2\pi} dt = \sqrt{r^2 + \frac{h^2}{4\pi^2}} \cdot [t]_0^{2\pi} = 2\pi \sqrt{r^2 + \frac{h^2}{4\pi^2}} \]
		\[ L(\gamma) = \sqrt{4\pi^2 r^2 + h^2} \]
		ეს არის გაშლილი ცილინდრის  (მართკუთხედის) დიაგონალის სიგრძე, რომლის გვერდებია $2\pi r$ (წრეწირის სიგრძე) და $h$ (სიმაღლე).
		
		\subsubsection*{მაგალითი 2: კუთხის გამოთვლა}
		 ვიპოვოთ კუთხე ცილინდრის ზედაპირზე არსებულ ჰორიზონტალურ წრეწირსა და ვერტიკალურ წარმომქმნელ წრფეს შორის.
		\vspace{1em}
		
		\begin{itemize}
			\item \textbf{ვექტორი 1 ($\mathbf{w}_1$):} ჰორიზონტალური წრეწირის მხები ვექტორი. ეს არის $\theta$ პარამეტრის ცვლილების მიმართულება, ანუ $\mathbf{r}_\theta$. ამ ვექტორისთვის $a_1=1, b_1=0$.
			\item \textbf{ვექტორი 2 ($\mathbf{w}_2$):} ვერტიკალური წრფის მხები ვექტორი. ეს არის $z$ პარამეტრის ცვლილების მიმართულება, ანუ $\mathbf{r}_z$. ამ ვექტორისთვის $a_2=0, b_2=1$.
		\end{itemize}
		გამოვიყენოთ კუთხის ფორმულა $E=r^2, F=0, G=1$ კოეფიციენტებით:
		\begin{align*}
			\cos(\theta) &= \frac{E a_1 a_2 + F(a_1 b_2 + a_2 b_1) + G b_1 b_2}{\sqrt{I(\mathbf{w}_1)}\sqrt{I(\mathbf{w}_2)}} \\
			&= \frac{r^2(1)(0) + 0((1)(1) + (0)(0)) + 1(0)(1)}{\sqrt{I(\mathbf{w}_1)}\sqrt{I(\mathbf{w}_2)}} \\
			&= \frac{0}{\sqrt{I(\mathbf{w}_1)}\sqrt{I(\mathbf{w}_2)}} = 0
		\end{align*}
		რადგან $\cos(\theta) = 0$, კუთხე $\theta = \frac{\pi}{2}$ (ანუ $90^\circ$). ეს შედეგი ემთხვევა ჩვენს ინტუიციას: ცილინდრის ჰორიზონტალური და ვერტიკალური ხაზები ურთიერთმართობულია.
		
		\subsubsection*{მაგალითი 3: ზედაპირის ფართობის გამოთვლა}
		 ვიპოვოთ $R$ რადიუსიანი და $H$ სიმაღლის მქონე ცილინდრის გვერდითი ზედაპირის ფართობი.
		\vspace{1em}
		
		ეს შეესაბამება პარამეტრების ცვლილების შემდეგ დიაპაზონს: $\theta$ იცვლება $0$-დან $2\pi$-მდე, ხოლო $z$ იცვლება $0$-დან $H$-მდე.
		ფართობის ელემენტია $\sqrt{EG - F^2}$. ჩვენს შემთხვევაში:
		\[ \sqrt{EG - F^2} = \sqrt{r^2 \cdot 1 - 0^2} = \sqrt{r^2} = R \]
		ახლა გამოვთვალოთ ორმაგი ინტეგრალი:
		\begin{align*}
			A(R) &= \int_{0}^{H} \int_{0}^{2\pi} \sqrt{EG - F^2} \, d\theta \, dz \\
			&= \int_{0}^{H} \int_{0}^{2\pi} r \, d\theta \, dz \\
			&= r \int_{0}^{H} \left[ \theta \right]_0^{2\pi} \, dz \\
			&= r \int_{0}^{H} 2\pi \, dz \\
			&= 2\pi r \left[ z \right]_0^{H} = 2\pi r H
		\end{align*}
		ეს არის ცილინდრის  გვერდითი ზედაპირის ფართობის კარგად ცნობილი ფორმულა.
		
		
		
		
		
		
		
		
		
		
		
		
		
		

		
		\bigskip
		
		
		
		\subsection{იზომეტრია და ზედაპირების შინაგანი გეომეტრია}
		
		\begin{definition}[ლოკალური იზომეტრია]
			$S_1$ და $S_2$ ზედაპირებს შორის ასახვას $\phi: S_1 \to S_2$ ეწოდება \textbf{ლოკალური იზომეტრია}, თუ ის ინახავს პირველ ფუნდამენტურ ფორმას. ეს ნიშნავს, რომ ყოველი $p \in S_1$-თვის და ყოველი $\mathbf{v} \in T_pS_1$-თვის, $d\phi_p$ (ასახვის დიფერენციალი) ინახავს ვექტორის ნორმას:
			\[ \|d\phi_p(\mathbf{v})\|_{\phi(p)} = \|\mathbf{v}\|_p \]
			სადაც ნორმა იზომება შესაბამისი ზედაპირის პირველი ფუნდამენტური ფორმით.
		\end{definition}
		
		იზომეტრია არის ზედაპირების „მყარი“ (rigid) გარდაქმნა. თუ ორი ზედაპირი ლოკალურად იზომეტრიულია, მათი შინაგანი გეომეტრია აბსოლუტურად იდენტურია.
		
		\subsubsection{იზომეტრიის ინტუიციური მნიშვნელობა}
		
		წარმოიდგინეთ, რომ ზედაპირი დამზადებულია იდეალურად ელასტიური, მაგრამ გაუჭიმავი მასალისგან, მაგალითად, ქაღალდის ფურცლისგან. იზომეტრია არის ამ ფურცლის ნებისმიერი დეფორმაცია სივრცეში, რომელიც არ იწვევს მის გაჭიმვას, შეკუმშვას ან გახევას.
		
		\begin{itemize}
			\item \textbf{ანალოგია:} წარმოიდგინეთ ჭიანჭველა, რომელიც ცხოვრობს ქაღალდის ფურცელზე. მისთვის სამყარო ორგანზომილებიანია. ის ზომავს მანძილებს და კუთხეებს თავის სამყაროში. თუ ჩვენ ამ ფურცელს ავიღებთ და ცილინდრის ფორმაზე დავახვევთ, ჭიანჭველასთვის არაფერი შეიცვლება. მის მიერ გაზომილი ნებისმიერი მანძილი ორ წერტილს შორის ან კუთხე ორ გზას შორის იგივე დარჩება. ჭიანჭველამ არ იცის, რომ მისი სამყარო „გამრუდდა“ სამგანზომილებიან სივრცეში.
		\end{itemize}
		ეს იმიტომ ხდება, რომ სიბრტყე და ცილინდრი \textbf{ლოკალურად იზომეტრიულია}. მათი \textit{შინაგანი} გეომეტრია ერთნაირია. ცილინდრის გამრუდება, რომელსაც ჩვენ ვხედავთ, \textit{გარეგანი} (extrinsic) თვისებაა.
		
		\subsubsection*{იზომეტრიის შედეგები}
		რადგან იზომეტრია ინახავს პირველ ფუნდამენტურ ფორმას, ის ავტომატურად ინახავს ყველა იმ სიდიდეს, რომელიც მხოლოდ ამ ფორმით გამოითვლება. კერძოდ:
		\begin{enumerate}
			\item \textbf{წირების სიგრძეები ინახება:} თუ $\gamma$ არის წირი $S_1$-ზე, მაშინ $L(\gamma) = L(\phi(\gamma))$.
			\item \textbf{კუთხეები ინახება:} თუ ორი წირი $S_1$-ზე $\theta$ კუთხით იკვეთება, მათი სახეები $S_2$-ზეც იმავე $\theta$ კუთხით იკვეთება.
			\item \textbf{ფართობები ინახება:} $S_1$ ზედაპირის ნაწილის ფართობი მისი სახის ფართობის ტოლია $S_2$-ზე.
		\end{enumerate}
		
		\subsubsection*{იზომეტრია და გაუსის სიმრუდე: Theorema Egregium-ის როლი}
		ჩნდება კითხვა: როგორ შეგვიძლია დავადგინოთ, არის თუ არა ორი ზედაპირი იზომეტრიული? აქ გადამწყვეტ როლს თამაშობს გაუსის „შესანიშნავი თეორემა“.
		
		\begin{proposition}
			თუ ორ ზედაპირს, $S_1$ და $S_2$, შორის არსებობს ლოკალური იზომეტრია $\phi: S_1 \to S_2$, მაშინ შესაბამის წერტილებში მათი გაუსის სიმრუდეები ტოლია:
			\[ K_{S_1}(p) = K_{S_2}(\phi(p)) \quad \text{ყველა } p \in S_1\text{-თვის.} \]
		\end{proposition}
		\begin{proof}[დამტკიცების იდეა]
			გაუსის \textit{Theorema Egregium}-ის თანახმად, გაუსის სიმრუდე $K$ შეიძლება გამოითვალოს მხოლოდ პირველი ფუნდამენტური ფორმის კოეფიციენტებით ($E,F,G$) და მათი წარმოებულებით. რადგან იზომეტრია ინახავს $E,F,G$ კოეფიციენტებს, ის აუცილებლად შეინახავს ნებისმიერ სიდიდეს, რომელიც მათგან გამომდინარეობს. შესაბამისად, გაუსის სიმრუდეც ინახება.
		\end{proof}
		
		ეს მძლავრი შედეგი გვაძლევს საშუალებას, მარტივად დავამტკიცოთ, რომ ზოგიერთ ზედაპირს შორის იზომეტრია შეუძლებელია.
		
		\begin{example}
			\textbf{ცილინდრი და სიბრტყე:} ცილინდრის გაუსის სიმრუდე $K=0$. სიბრტყის გაუსის სიმრუდეც $K=0$. რადგან სიმრუდეები ემთხვევა, მათ შორის იზომეტრია \textit{შესაძლებელია}. სწორედ ამიტომ შეგვიძლია ცილინდრის „გაშლა“ სიბრტყეზე.
			
			\textbf{სფერო და სიბრტყე:} $R$ რადიუსიანი სფეროს გაუსის სიმრუდეა $K = 1/R^2 > 0$, ხოლო სიბრტყისა $K=0$. რადგან $K_{სფერო} \neq K_{სიბრტყე}$, მათ შორის \textbf{იზომეტრია შეუძლებელია}. ეს არის მათემატიკური მიზეზი, თუ რატომ არ შეიძლება დედამიწის რუკის შექმნა ბრტყელ ფურცელზე მანძილების დაუმახინჯებლად.
		\end{example}
		
		\subsubsection{უფრო ზოგადი ცნება: კონფორმული ასახვა}
		იზომეტრიასთან ახლოს მდგომი, მაგრამ უფრო ზოგადი ცნებაა კონფორმული ასახვა.
		
		\begin{definition}[კონფორმული ასახვა]
			ასახვას $\phi: S_1 \to S_2$ ეწოდება \textbf{კონფორმული}, თუ ის ინახავს კუთხეებს, მაგრამ არაა აუცილებელი, ინახავდეს სიგრძეებს. ეს მათემატიკურად ნიშნავს, რომ პირველი ფუნდამენტური ფორმა ინახება გარკვეული დადებითი ფუნქციის ($e^{2\sigma}$) სიზუსტით:
			\[ I_p(\mathbf{v})_1 = e^{2\sigma(p)} \cdot I_{\phi(p)}(d\phi_p(\mathbf{v}))_2 \]
		\end{definition}
		კონფორმული ასახვები ინახავენ „ფორმას“ ლოკალურად, მაგრამ ცვლიან „ზომას“. კარტოგრაფიაში (რუკების შედგენისას) კონფორმული ასახვები ძალიან მნიშვნელოვანია.
		\begin{itemize}
			\item \textbf{მერკატორის პროექცია:} ეს არის კონფორმული ასახვა სფეროდან ცილინდრზე (რომელიც შემდეგ იშლება სიბრტყედ). ის ინახავს კუთხეებს, რის გამოც ნავიგაციისთვისაა მოსახერხებელი, მაგრამ საშინლად ამახინჯებს ფართობებს პოლუსებთან ახლოს (მაგ. გრენლანდია აფრიკის ზომის ჩანს).
			\item \textbf{სტერეოგრაფიული პროექცია:} ეს არის კონფორმული ასახვა სფეროდან სიბრტყეზე.
		\end{itemize}
		
		ეს ასახვა ასე ხება: წარმოიდგინეთ  სფერო (მაგალითად, გლობუსი) სიბრტყეზე (მაგალითად, მაგიდის ზედაპირზე), ისე რომ სფეროს ერთ წერტილში, სამხრეთ პოლუსით ეხება მაგიდას.
		
		ახლა, სფეროს ჩრდილოეთ   პოლუსი, ავირჩიოთ ჩვენს „პროექტორად“ ან „სინათლის წყაროდ“.
		
		ასახვის წესი: სფეროს ზედაპირზე არსებული ნებისმიერი სხვა წერტილის  ასასახად, ჩვენ ვავლებთ სხივს ჩრდილოეთ   პოლუსიდან  ამ წერტილის  გავლით. წერტილი, სადაც ეს სხივი გადაკვეთს ჩვენს სიბრტყეს (მაგიდის ზედაპირს), არის ამოსავალი წერტილის სტერეოგრაფიული   პროექცია.
		
		რა ხდება ამ პროცესში?
		სამხრეთ ნახევარსფერო აისახება სიბრტყის იმ წრეში, რომლის საზღვარიც ეკვატორის პროექციაა. სამხრეთ პოლუსი რჩება თავის ადგილზე.
		
		ჩრდილოეთ ნახევარსფერო კი აისახება ამ წრის გარეთ არსებულ მთელ უსასრულო სიბრტყეზე. რაც უფრო ახლოსაა წერტილი ჩრდილოეთ პოლუსთან, მით უფრო შორს აისახება ის სიბრტყეზე.
		
		თავად ჩრდილოეთ პოლუსი, საიდანაც პროექცია ხდება, აისახება „უსასრულოდ შორეულ წერტილში“, რადგან მისგან გავლებულ მხებ ხაზს სიბრტყე არასდროს გადაკვეთს.
		
		ძირითადი თვისებები:
		სტერეოგრაფიულ პროექციას ორი უმნიშვნელოვანესი თვისება აქვს:
		
		კონფორმულობა (კუთხეების შენახვა): ეს მისი ყველაზე ფასეული თვისებაა. თუ სფეროზე ორი წირი რაღაც კუთხით იკვეთება, მათი პროექციები სიბრტყეზე ზუსტად იგივე კუთხით გადაიკვეთება. ამის გამო, ის ინახავს პატარა ფიგურების ფორმას, მაგრამ არა მათ ზომას.
		
		წრეწირების ასახვა: სფეროზე არსებული ნებისმიერი წრეწირი სიბრტყეზე აისახება ან წრეწირად, ან წრფედ.
		
		წრეწირები, რომლებიც არ გადიან ჩრდილოეთ პოლუსზე, რჩებიან წრეწირებად (თუმცა მათი ცენტრი და რადიუსი იცვლება).
		
		წრეწირები, რომლებიც გადიან ჩრდილოეთ პოლუსზე (მაგალითად, მერიდიანები), გარდაიქმნებიან სიბრტყეზე გამავალ წრფეებად.
		
		ამ თვისებების გამო, სტერეოგრაფიული პროექცია ფართოდ გამოიყენება მათემატიკის სხვადასხვა დარგში, როგორიცაა კომპლექსური ანალიზი (რიმანის სფერო), კარტოგრაფია (პოლარული რეგიონების რუკები) და კრისტალოგრაფია.
		
		
		
		
		\subsection{სავარჯიშოები}
		
		\begin{enumerate}
			\item \textbf{სიბრტყე:} მოცემულია $\mathbf{r}(u, v) = \langle u, v, 0 \rangle$. იპოვეთ პირველი ფუნდამენტური ფორმა და $\mathbf{c}(t) = \mathbf{r}(t, t)$ წირის სიგრძე, როცა $0 \le t \le 1$.
			
			ამოხსნა
			
			1.  გამოთვალეთ კერძო წარმოებულები:
			$$ \mathbf{r}_u = \langle 1, 0, 0 \rangle, \quad \mathbf{r}_v = \langle 0, 1, 0 \rangle $$
			
			2.  გამოთვალეთ პირველი ფუნდამენტური ფორმის კოეფიციენტები:
			$$ E = 1, \quad F = 0, \quad G = 1 $$
			
			3.  ჩაწერეთ პირველი ფუნდამენტური ფორმა:
			$$ ds^2 = du^2 + dv^2 $$
			
			4.  გააპარამეტრეთ წირი $\mathbf{c}(t)$:
			$$ u(t) = t, \quad v(t) = t $$
			
			5.  გამოთვალეთ $u(t)$-სა და $v(t)$-ს წარმოებულები:
			$$ \frac{du}{dt} = 1, \quad \frac{dv}{dt} = 1 $$
			
			6.  გამოთვალეთ წირის სიგრძე:
			$$ L = \int_0^1 \sqrt{1 (1)^2 + 2(0)(1)(1) + 1 (1)^2} \, dt = \sqrt{2} $$
			
			მაშასადამე, წირის სიგრძეა $\sqrt{2}$.
			
			
			\item \textbf{ცილინდრი:} მოცემულია $\mathbf{r}(u, v) = \langle \cos(u), \sin(u), v \rangle$. იპოვეთ პირველი ფუნდამენტური ფორმა და $\mathbf{c}(t) = \mathbf{r}(t, t)$ წირის სიგრძე, როცა $0 \le t \le 2\pi$.
			
			ამოხსნა.
			
			1.  გამოთვალეთ კერძო წარმოებულები:
			$$ \mathbf{r}_u = \langle -\sin(u), \cos(u), 0 \rangle, \quad \mathbf{r}_v = \langle 0, 0, 1 \rangle $$
			
			2.  გამოთვალეთ პირველი ფუნდამენტური ფორმის კოეფიციენტები:
			$$ E = 1, \quad F = 0, \quad G = 1 $$
			
			3.  ჩაწერეთ პირველი ფუნდამენტური ფორმა:
			$$ ds^2 = du^2 + dv^2 $$
			
			4.  გააპარამეტრეთ წირი $\mathbf{c}(t)$:
			$$ u(t) = t, \quad v(t) = t $$
			
			5.  გამოთვალეთ $u(t)$-სა და $v(t)$-ს წარმოებულები:
			$$ \frac{du}{dt} = 1, \quad \frac{dv}{dt} = 1 $$
			
			6.  გამოთვალეთ წირის სიგრძე:
			$$ L = \int_0^{2\pi} \sqrt{1 (1)^2 + 2(0)(1)(1) + 1 (1)^2} \, dt = 2\pi\sqrt{2} $$
			
			მაშასადამე, წირის სიგრძეა $2\pi\sqrt{2}$.
			
			\item \textbf{სფერო:} მოცემულია $\mathbf{r}(u, v) = \langle \sin(u)\cos(v), \sin(u)\sin(v), \cos(u) \rangle$. იპოვეთ პირველი ფუნდამენტური ფორმა და $\mathbf{c}(t) = \mathbf{r}(t, \pi/4)$ წირის სიგრძე, როცა $0 \le t \le \pi$.
			
			\item \textbf{კონუსი:} მოცემულია $\mathbf{r}(u, v) = \langle u\cos(v), u\sin(v), u \rangle$. იპოვეთ პირველი ფუნდამენტური ფორმა და $\mathbf{c}(t) = \mathbf{r}(t, t)$ წირის სიგრძე, როცა $1 \le t \le 2$.
			
			\item \textbf{პარაბოლოიდი:} მოცემულია $\mathbf{r}(u, v) = \langle u, v, u^2 + v^2 \rangle$. იპოვეთ პირველი ფუნდამენტური ფორმა და $\mathbf{c}(t) = \mathbf{r}(t, 0)$ წირის სიგრძე, როცა $0 \le t \le 1$.
			
			\item \textbf{ჰელიკოიდი:} მოცემულია $\mathbf{r}(u, v) = \langle u\cos(v), u\sin(v), v \rangle$. იპოვეთ პირველი ფუნდამენტური ფორმა და $\mathbf{c}(t) = \mathbf{r}(t, t)$ წირის სიგრძე, როცა $0 \le t \le \pi$.
			
			\item \textbf{ტორი:} მოცემულია $\mathbf{r}(u, v) = \langle (2 + \cos(u))\cos(v), (2 + \cos(u))\sin(v), \sin(u) \rangle$. იპოვეთ პირველი ფუნდამენტური ფორმა და $\mathbf{c}(t) = \mathbf{r}(0, t)$ წირის სიგრძე, როცა $0 \le t \le 2\pi$.
			
			\item \textbf{ერთფურცელა ჰიპერბოლოიდი:} მოცემულია $\mathbf{r}(u, v) = \langle \cosh(u)\cos(v), \cosh(u)\sin(v), \sinh(u) \rangle$. იპოვეთ პირველი ფუნდამენტური ფორმა და $\mathbf{c}(t) = \mathbf{r}(0, t)$ წირის სიგრძე, როცა $0 \le t \le \pi$.
			
			\item \textbf{ელიფსოიდი:} მოცემულია $\mathbf{r}(u, v) = \langle 2\sin(u)\cos(v), \sin(u)\sin(v), \cos(u) \rangle$. იპოვეთ პირველი ფუნდამენტური ფორმა და $\mathbf{c}(t) = \mathbf{r}(t, 0)$ წირის სიგრძე, როცა $0 \le t \le \pi/2$.
			
			\item \textbf{მაიმუნის უნაგირი:} მოცემულია $\mathbf{r}(u, v) = \langle u, v, u^3 - 3uv^2 \rangle$. იპოვეთ პირველი ფუნდამენტური ფორმა და $\mathbf{c}(t) = \mathbf{r}(t, 0)$ წირის სიგრძე, როცა $0 \le t \le 1$.
			
			\item \textbf{კატენოიდი:} მოცემულია $\mathbf{r}(u,v) = \langle \cosh(v)\cos(u), \cosh(v)\sin(u), v \rangle$. იპოვეთ პირველი ფუნდამენტური ფორმა და $\mathbf{c}(t) = \mathbf{r}(t,0)$ წირის სიგრძე, როცა $0\le t\le 2\pi$.
			
			\item \textbf{ბრუნვის ზედაპირი:} მოცემულია $\mathbf{r}(u,v) = \langle u \cos(v), u \sin(v), u^2 \rangle$. იპოვეთ პირველი ფუნდამენტური ფორმა და $\mathbf{c}(t) = \mathbf{r}(t, \pi)$ წირის სიგრძე, როცა $1 \le t \le 2$.
		\end{enumerate}
		
		
		
		
		\subsubsection*{დამატებითი  სავარჯიშოები ზედაპირის ფართობის  გამოთვლაზე}
		\label{SA}
		
		\begin{enumerate}
			\item \textbf{პარაბოლოიდი:} იპოვეთ 
			$$\mathbf{r}(u, v) = \langle u, v, u^2 + v^2 \rangle$$ 
			პარაბოლოიდის  ფართობი  $D = \{(u, v) \mid 0 \le u \le 1, 0 \le v \le 1\}$ არეზე.
			\item \textbf{კონუსი:} იპოვეთ 
			$$\mathbf{r}(u, v) = \langle u\cos(v), u\sin(v), u \rangle$$ 
			კონუსის ფართობი $D = \{(u, v) \mid 0 \le u \le 2, 0 \le v \le 2\pi\}$ არეზე.
			\item \textbf{ცილინდრის ნაწილი:} იპოვეთ 
			$$\mathbf{r}(u, v) = \langle 2\cos(u), 2\sin(u), v \rangle$$ 
			ცილინდრის ფართობი $D = \{(u, v) \mid 0 \le u \le \pi/2, 0 \le v \le 4\}$ არეზე.
			\item \textbf{ჰელიკოიდი:} იპოვეთ 
			$$\mathbf{r}(u, v) = \langle u\cos(v), u\sin(v), v \rangle$$ 
			ჰელიკოიდის ფართობი $D = \{(u, v) \mid 0 \le u \le 1, 0 \le v \le \pi\}$ არეზე.
			\item \textbf{ტორის ნაწილი:} იპოვეთ 
			$$\mathbf{r}(u, v) = \langle (3 + \cos(u))\cos(v), (3 + \cos(u))\sin(v), \sin(u) \rangle$$ 
			ტორის ფართობი $D = \{(u, v) \mid 0 \le u \le \pi, 0 \le v \le \pi\}$ არეზე.
			\item \textbf{ერთფურცელა  ჰიპერბოლოიდის  ნაწილი:} იპოვეთ 
			$$\mathbf{r}(u, v) = \langle \cosh(u)\cos(v), \cosh(u)\sin(v), \sinh(u) \rangle$$ 
			ჰიპერბოლოიდის ფართობი $D = \{(u, v) \mid 0 \le u \le 1, 0 \le v \le \pi/2\}$ არეზე.
			\item \textbf{ელიფსოიდის  ნაწილი:} იპოვეთ 
			$$\mathbf{r}(u, v) = \langle 2\sin(u)\cos(v), \sin(u)\sin(v), 3\cos(u) \rangle$$ 
			ელიფსოიდის ფართობი $D = \{(u, v) \mid 0 \le u \le \pi/2, 0 \le v \le \pi/2\}$ არეზე.
			\item \textbf{მაიმუნის უნაგირის ნაწილი:} იპოვეთ 
			$$\mathbf{r}(u, v) = \langle u, v, u^3 - 3uv^2 \rangle$$ 
			ზედაპირის ფართობი $D = \{(u, v) \mid 0 \le u \le 1, 0 \le v \le 1\}$ არეზე.
			\item \textbf{კატენოიდის ნაწილი:} იპოვეთ 
			$$\mathbf{r}(u,v) = \langle \cosh(v)\cos(u), \cosh(v)\sin(u), v \rangle$$ კატენოიდის ფართობი $D=\{(u,v) \mid 0 \le u \le \pi, 0 \le v \le 1\}$ არეზე.
		\end{enumerate}
		
		
		
		
		
		
		
		\newpage
		%======================================================================
		\section{მეორე ფუნდამენტური ფორმა და ზედაპირის სიმრუდე}
		%======================================================================
		%\addcontentsline{toc}{section}{მეორე ფუნდამენტური ფორმა და ზედაპირის სიმრუდე}
		
		\subsection{გაუსის ასახვა და ვაინგარტენის ოპერატორი}
		\begin{definition}[გაუსის ასახვა]
			ვთქვათ $S$ ორიენტირებადი ზედაპირია. \textbf{გაუსის ასახვა} $N: S \to S^2$ არის ასახვა, რომელიც ყოველ $p \in S$ წერტილს შეუსაბამებს მის ერთეულოვან ნორმალურ ვექტორს $\mathbf{n}(p)$, რომელიც განხილულია როგორც $S^2$ ერთეულოვანი სფეროს წერტილი.
		\end{definition}
		გაუსის ასახვა გვიჩვენებს, თუ როგორ "იხვევა" ზედაპირი. რაც უფრო სწრაფად იცვლება ნორმალი, მით უფრო "მრუდეა" ზედაპირი.
		\begin{definition}[ვაინგარტენის ოპერატორი]
			გაუსის ასახვის დიფერენციალს $dN_p: T_pS \to T_{\mathbf{n}(p)}S^2$ ეწოდება \textbf{ვაინგარტენის ოპერატორი} $S_p$. რადგან $T_pS$ და $T_{\mathbf{n}(p)}S^2$ პარალელური სიბრტყეებია, შეგვიძლია $S_p$ განვიხილოთ, როგორც წრფივი ოპერატორი $S_p: T_pS \to T_pS$.
		\end{definition}
		$S_p(\mathbf{v}) = -d\mathbf{n}_p(\mathbf{v})$ ზომავს ნორმალური ვექტორის ცვლილების სიჩქარეს $\mathbf{v}$ მიმართულებით.
		
		\begin{proposition}
			ვაინგარტენის ოპერატორი $S_p$ არის თვითდართული (სიმეტრიული) ოპერატორი პირველი ფუნდამენტური ფორმით განმარტებულ სკალარულ ნამრავლთან მიმართებით: $I_p(S_p(\mathbf{v}), \mathbf{w}) = I_p(\mathbf{v}, S_p(\mathbf{w}))$.
		\end{proposition}
		\begin{proof}
			(ეს დამტკიცება მოითხოვს კოვარიანტული წარმოებულის ცნებას, თუმცა მისი შედეგი ფუნდამენტურია).
		\end{proof}
		
		\subsection{მეორე ფუნდამენტური ფორმა}
		\begin{definition}[მეორე ფუნდამენტური ფორმა]
			$p \in S$ წერტილში \textbf{მეორე ფუნდამენტური ფორმა} $II_p$ არის ბილინიარული, სიმეტრიული ფორმა, რომელიც განისაზღვრება ვაინგარტენის ოპერატორის საშუალებით:
			\[ II_p(\mathbf{v}, \mathbf{w}) = I_p(S_p(\mathbf{v}), \mathbf{w}) = S_p(\mathbf{v}) \cdot \mathbf{w} \]
		\end{definition}
		შესაბამისი კვადრატული ფორმაა $II_p(\mathbf{v}) = S_p(\mathbf{v}) \cdot \mathbf{v}$. $\{\mathbf{r}_u, \mathbf{r}_v\}$ ბაზისში, ამ ფორმის კოეფიციენტებია:
		\begin{align*}
			L &= II(\mathbf{r}_u, \mathbf{r}_u) = S(\mathbf{r}_u) \cdot \mathbf{r}_u = -d\mathbf{n}(\mathbf{r}_u) \cdot \mathbf{r}_u = -\mathbf{n}_u \cdot \mathbf{r}_u = \mathbf{n} \cdot \mathbf{r}_{uu} \\
			M &= II(\mathbf{r}_u, \mathbf{r}_v) = S(\mathbf{r}_u) \cdot \mathbf{r}_v = -\mathbf{n}_u \cdot \mathbf{r}_v = \mathbf{n} \cdot \mathbf{r}_{uv} \\
			N &= II(\mathbf{r}_v, \mathbf{r}_v) = S(\mathbf{r}_v) \cdot \mathbf{r}_v = -\mathbf{n}_v \cdot \mathbf{r}_v = \mathbf{n} \cdot \mathbf{r}_{vv}
		\end{align*}
		ეს გვიჩვენებს კოეფიციენტების $L, M, N$ ეკვივალენტურ განსაზღვრებას.
		
		\subsection{სიმრუდეები}
		\begin{definition}[ნორმალური სიმრუდე]
			$k_n(\mathbf{v}) = \frac{II(\mathbf{v})}{I(\mathbf{v})}$ არის ნორმალური სიმრუდე $\mathbf{v}$ მიმართულებით.
		\end{definition}
		\begin{theorem}[ეილერის თეორემა]
			თუ $k_1$ და $k_2$ მთავარი სიმრუდეებია, ხოლო $\mathbf{e}_1, \mathbf{e}_2$ შესაბამისი მთავარი მიმართულებები, მაშინ ნებისმიერი $\mathbf{v} = \cos\theta \mathbf{e}_1 + \sin\theta \mathbf{e}_2$ მიმართულებით ნორმალური სიმრუდე გამოითვლება ფორმულით:
			\[ k_n(\theta) = k_1 \cos^2\theta + k_2 \sin^2\theta \]
		\end{theorem}
		\begin{proof}
			$II(\mathbf{v}) = I(S(\mathbf{v}), \mathbf{v}) = S(\cos\theta \mathbf{e}_1 + \sin\theta \mathbf{e}_2) \cdot (\cos\theta \mathbf{e}_1 + \sin\theta \mathbf{e}_2) = (k_1\cos\theta \mathbf{e}_1 + k_2\sin\theta \mathbf{e}_2) \cdot (\cos\theta \mathbf{e}_1 + \sin\theta \mathbf{e}_2) = k_1\cos^2\theta + k_2\sin^2\theta$. რადგან $I(\mathbf{v})=1$, $k_n = II(\mathbf{v})$.
		\end{proof}
		\begin{definition}[გაუსის და საშუალო სიმრუდე]
			$K = k_1 k_2 = \det(S)$ და $H = \frac{1}{2}(k_1+k_2) = \frac{1}{2}\text{tr}(S)$.
		\end{definition}
		
		\subsection{ზედაპირის წერტილების კლასიფიკაცია}
		\begin{itemize}
			\item \textbf{ელიფსური წერტილი:} $K > 0$. $k_1, k_2$ ერთი ნიშნისაა. ზედაპირი ლოკალურად გუმბათს ჰგავს.
			\item \textbf{ჰიპერბოლური წერტილი:} $K < 0$. $k_1, k_2$ სხვადასხვა ნიშნისაა. ზედაპირი ლოკალურად უნაგირს ჰგავს.
			\item \textbf{პარაბოლური წერტილი:} $K=0$, მაგრამ $H \neq 0$. ერთ-ერთი მთავარი სიმრუდე ნულია. ზედაპირი ლოკალურად ცილინდრს ჰგავს.
			\item \textbf{ბრტყელი წერტილი:} $K=H=0$. ორივე მთავარი სიმრუდე ნულია.
			\item \textbf{ომბილური წერტილი:} $k_1=k_2$. ამ წერტილში ნორმალური სიმრუდე ყველა მიმართულებით ერთნაირია ($k_n=k_1=k_2$). ელიფსური ომბილური წერტილი სფეროს წერტილის მსგავსია, ბრტყელი ომბილური წერტილი კი სიბრტყის წერტილია.
		\end{itemize}
		
		\subsection{სავარჯიშოები}
		\begin{exercise}
			დაამტკიცეთ, რომ თუ ზედაპირის ყველა წერტილი ომბილურია, მაშინ ის ან სფეროს ნაწილია, ან სიბრტყის ნაწილი.
		\end{exercise}
		\begin{exercise}
			იპოვეთ ჰიპერბოლური პარაბოლოიდის $z = y^2 - x^2$ მთავარი სიმრუდეები სათავეში.
		\end{exercise}
		
		\newpage
		%======================================================================
		\section{ზედაპირთა თეორიის ფუნდამენტური განტოლებები}
		%======================================================================
		%\addcontentsline{toc}{section}{ზედაპირთა თეორიის ფუნდამენტური განტოლებები}
		
		\subsection{კოვარიანტული წარმოებული და კრისტოფელის სიმბოლოები}
		მრუდე ზედაპირზე ვექტორული ველის წარმოებულის აღება რთული საკითხია, რადგან ვექტორები სხვადასხვა მხებ სიბრტყეშია. \textbf{კოვარიანტული წარმოებული} არის ამ პრობლემის გადაჭრის საშუალება.
		გაუსის ფორმულა გვიჩვენებს, რომ მეორე წარმოებული $\mathbf{r}_{ij}$ იშლება მხებ და ნორმალურ კომპონენტებად:
		\[ \mathbf{r}_{ij} = \sum_k \Gamma^k_{ij} \mathbf{r}_k + L_{ij}\mathbf{n} \]
		სადაც $L_{ij}$ არის მეორე ფუნდამენტური ფორმის კომპონენტები ($L_{11}=L$, $L_{12}=M$ და ა.შ.).
		\begin{definition}[კრისტოფელის სიმბოლოები]
			$\Gamma^k_{ij}$ სიდიდეებს ეწოდება \textbf{კრისტოფელის სიმბოლოები} და ისინი გამოითვლება მხოლოდ პირველი ფუნდამენტური ფორმით:
			\[ \Gamma^k_{ij} = \frac{1}{2} \sum_{m=1}^2 g^{km} \left( \frac{\partial g_{mi}}{\partial u^j} + \frac{\partial g_{mj}}{\partial u^i} - \frac{\partial g_{ij}}{\partial u^m} \right) \]
		\end{definition}
		
		\subsection{გაუს-ვაინგარტენის განტოლებები}
		ეს განტოლებათა სისტემა აღწერს, თუ როგორ იცვლება მოძრავი რეპერი $\{\mathbf{r}_u, \mathbf{r}_v, \mathbf{n}\}$:
		\begin{itemize}
			\item \textbf{გაუსის ფორმულები:}
			\begin{align*}
				\mathbf{r}_{uu} &= \Gamma^1_{11}\mathbf{r}_u + \Gamma^2_{11}\mathbf{r}_v + L\mathbf{n} \\
				\mathbf{r}_{uv} &= \Gamma^1_{12}\mathbf{r}_u + \Gamma^2_{12}\mathbf{r}_v + M\mathbf{n} \\
				\mathbf{r}_{vv} &= \Gamma^1_{22}\mathbf{r}_u + \Gamma^2_{22}\mathbf{r}_v + N\mathbf{n}
			\end{align*}
			\item \textbf{ვაინგარტენის ფორმულები:}
			\begin{align*}
				\mathbf{n}_u &= \frac{FM-GL}{EG-F^2}\mathbf{r}_u + \frac{FL-EM}{EG-F^2}\mathbf{r}_v \\
				\mathbf{n}_v &= \frac{FN-GM}{EG-F^2}\mathbf{r}_u + \frac{FM-EN}{EG-F^2}\mathbf{r}_v
			\end{align*}
		\end{itemize}
		
		\subsection{თავსებადობის განტოლებები და Theorema Egregium-ის დამტკიცება}
		შერეული წარმოებულების ტოლობის პირობიდან $(\mathbf{r}_{uu})_v = (\mathbf{r}_{uv})_u$ გამომდინარეობს თავსებადობის განტოლებები. ამ ვექტორული ტოლობის დაშლით მხებ და ნორმალურ კომპონენტებად, მივიღებთ:
		\begin{itemize}
			\item \textbf{გაუსის განტოლება:} ნორმალური კომპონენტების გატოლებით მიიღება:
			\[ LN - M^2 = E(\Gamma^2_{11})_v - \dots \]
			ეს რთული გამოსახულება აკავშირებს $L,M,N$-ს მხოლოდ $E,F,G$-სთან და მათ წარმოებულებთან, რაც ამტკიცებს გაუსის \textbf{შესანიშნავ თეორემას (Theorema Egregium)}. ანუ, $K = \frac{LN-M^2}{EG-F^2}$ შეიძლება გამოისახოს მხოლოდ შინაგანი სიდიდეებით.
			\item \textbf{კოდაცი-მაინარდის განტოლებები:} მხები კომპონენტების გატოლებით მიიღება:
			\begin{align*}
				L_v - M_u &= L\Gamma^1_{12} + M(\Gamma^2_{12} - \Gamma^1_{11}) - N\Gamma^2_{11} \\
				M_v - N_u &= L\Gamma^1_{22} + M(\Gamma^2_{22} - \Gamma^1_{12}) - N\Gamma^2_{12}
			\end{align*}
		\end{itemize}
		
		\subsection{ბონეს თეორემა}
		\begin{theorem}[ზედაპირთა თეორიის ფუნდამენტური თეორემა]
			ვთქვათ მოცემულია $E,F,G$ და $L,M,N$ გლუვი ფუნქციები, სადაც $\begin{pmatrix} E & F \\ F & G \end{pmatrix}$ დადებითად განსაზღვრული მატრიცაა. თუ ეს ექვსი ფუნქცია აკმაყოფილებს გაუსისა და კოდაცი-მაინარდის თავსებადობის განტოლებებს, მაშინ არსებობს (ლოკალურად და უნიკალურად, ევკლიდური მოძრაობის სიზუსტით) ზედაპირი, რომლისთვისაც ეს ფუნქციები პირველი და მეორე ფუნდამენტური ფორმების კოეფიციენტებია.
		\end{theorem}
		ეს თეორემა ამბობს, რომ მთელი ლოკალური გეომეტრია კოდირებულია ამ 6 ფუნქციასა და 3 განტოლებაში.
		
		\subsection{სავარჯიშოები}
		\begin{exercise}
			გამოთვალეთ კრისტოფელის სიმბოლოები სიბრტყისთვის პოლარულ კოორდინატებში.
		\end{exercise}
		\begin{exercise}
			შეამოწმეთ კოდაცი-მაინარდის განტოლებების შესრულება სფეროსთვის.
		\end{exercise}
	
	
	
	
	\newpage
	
	\section{დიფერენციალური ფორმები}
	
	განვიხილოთ ევკლიდური სივრცე $R^3$ საკოორდინაატო ღერძებით $x,y,z$. 
	
	სივრცის ყოველი წერტილი ამავე დროს არის ვექტორი. ყოველი $p$ წერტილისთვის და $v$ ვექტორისთვის $v$ ვექტორი შეგვიძლია სივრვის მხებ ვექტორად ჩავთვალოთ $p$-წერტილში. ყოველი ფიქსირებული $p$-თვის ელემენტარული 1-ფორმა $dx$ (ფორმები $dy$ ან $dz$ ) არის წრფივი ასახვა
	$$
	R^3\to R
	$$
	რომელიც $p$-ში მხები სივრცის ყოველ ვექტორს (ანუ $R^3$-ის ნებისმიერ ვექტორს) მიაწერს მის პირველ კოორდინატს ($dy$ მეორე კოორდინატს, ხოლო $dz$ მესამეს), ანუ თუ $v=(v_1,v_2,v_3)$, ანუ
	
	\begin{equation*}  
		dx(v)=v_1,\,\,\,dy(v)=v_2,\,\,\,\,dz(v)=v_3.
	\end{equation*}
	
	ელემენტარული 2-ფორმებს $dx\land dy$, $dx\land dz$, $dy\land dz$ მოკლედ ასე ავღნიშვნავთ $dxdy$, $dxdz$, $dydz$ და განსაზღვრებით ვექტორთა წყვილისათვის $u=(u_1,u_2,u_3)$, $v=(v_1,v_2,v_3)$ გვაქვს
	
	\begin{equation*}
		dxdy(u,v)=\left|
		\begin{array}{cc}
			u_1 & v_1 \\ 
			u_2 & v_2 \\
		\end{array}
		\right|,
	\end{equation*}
	\begin{equation*}
		dxdz(u,v)=\left|
		\begin{array}{cc}
			u_1 & v_1 \\ 
			u_3 & v_3 \\
		\end{array}
		\right|,
	\end{equation*}
	\begin{equation*}
		dydz(u,v)=\left|
		\begin{array}{cc}
			u_2 & v_2 \\ 
			u_3 & v_3 \\
		\end{array}
		\right|
	\end{equation*}
	
	სხვა სიტყვებით, $dxdy$ ელემენტარული 2-ფორმა ყოველ $u,v$ ვექტორთა წყვილს შეუსაბამებს $oxy$ სიბრტყეზე $u,v$ ვექტორების პროექციებით მიღებულ პარალელოგრამის ორიენტირებულ (ნიშნიან) ფართობს. ანალოგიურად $dxdz$ და $dydz$.
	
	ელემენტარული 3-ფორმა $dxdydz$ ვექტორთა სამეულს $u,v,w$ შეუსაბამებს მათგან მიღებულ პარალელოპიპედის ნიშნიან მოცულობას, ანუ
	
	\begin{equation*}
		dxdydz(u,v,w)=\left|
		\begin{array}{ccc}
			u_1 & v_1 & w_1\\ 
			u_2 & v_2 & w_2 \\
			u_3 & v_3 & w_3\\
		\end{array}
		\right|.
	\end{equation*}
	
	დიფერენციალური 0- ფორმა არის უბრალოდ $R^3$-ში განსაზღვრული ფუნქცია $f=f(x,y,z)$. 
	
	დიფერენციალური 1- ფორმა ეწოდება გამოსახულებას 
	$$
	\phi=f_1 dx+f_2 dy+f_3 dz, \text{ სადაც } f_i=f_i(x,y,z) \text{ რაღაც ფუნქციებია}.
	$$
	
	დიფერენციალური 2- ფორმა ეწოდება გამოსახულებას 
	$$
	\psi=f_1 dxdy+f_2 dydz+f_3 dxdz, \text{ სადაც } f_i=f_i(x,y,z) \text{ რაღაც ფუნქციებია}.
	$$
	
	დიფერენციალური 3- ფორმა ეწოდება გამოსახულებას 
	$$
	\rho=f dxdydz, \text{ სადაც } f=f(x,y,z) \text{ რაღაც ფუნქციაა}.
	$$
	
	1-დიფერენციალური ფორმების \textbf{გარე $\land$}-ნამრავლი არის 2-ფორმა 
	$$\phi_1 \land \phi_2$$ 
	რომელიც მიიღება ასე: ვთქვათ 
	$$
	\phi_1=f_1 dx+f_2 dy+f_3 dz,\,\,\,\phi_2=g_1 dx+g_2 dy+g_3 dz.
	$$
	
	ეს 1-ფორმები ფორმალურად, როგორც მრავალწევრები გადავამრავლოთ, ისე რომ უნდა დავიცვათ წესები:
	$$
	dxdx=0, dydy=0, dzdz=0, dxdy=-dydx, dzdx=-dxdz, dzdy=-dydz.
	$$
	და ბოლოს დავაჯგიფოთ მსგავსი წევრები, $dxdy$, $dxdz$ და $dydz$.
	მაგალითად
	\begin{align*}
		(2dx+fdy)\land (gdx+4dz) &= \\
		& 2gdx\land dx+8dx\land dz+fgdy \land dx+4fdy\land dz=\\
		& -fgdx \land dy+8dx\land dz+4fdy\land dz=\\
		& -fgdxdy+8dxdz+4fdydz.
	\end{align*}
	
	1-დიფერენციალური ფორმის გარე ნამრავლი 2-ფორმასთან არის 3-ფორმა 
	$$\phi \land \psi$$ 
	რომელიც მიიღება ასე: ვთქვათ 
	$$
	\phi=f_1 dx+f_2 dy+f_3 dz,\,\,\,\psi=g_1 dxdy+g_2 dxdz+g_3 dydz.
	$$
	
	ეს ფორმები ფორმალურად გადავამრავლოთ, ისე რომ უნდა დავიცვათ ისევ ძველი წესები
	$$
	dxdx=0, dydy=0, dzdz=0, dxdy=-dydx, dzdx=-dxdz, dzdy=-dydz.
	$$
	ამასთან თუ გადასმა ორჯერ მოხდა მაშინ ნიშანი იგივე რჩება, მაგალითად $dzdydx=-dydxdz=dxdydz$.
	ბოლოს ისევ უნდა დავაჯგიფოთ წევრები $dxdydz$ ელემენტარულ ფორმასთან.
	
	მაგალითად
	
	\begin{align*}
		(2dydz+fdxdz)\land (gdx+4dy) &= \\
		& 2gdydz\land dx+4fdxdz\land dy=\\
		& 2gdxdydz-4fdxdydz=\\
		& (2g-4f)dxdydz.
	\end{align*}
	
	გადამრავლების წესები გვიჩვენებს რომ \textbf{$R^3$-ში 2-ფორმების გარე ნამრავლი 0-ია}.
	
	თურმე ფორმის დიფერენცირებაც შეიძლება: 0-ფორმის ანუ ფუნქციის გარე დიფერენციალი
	არის უბრალოდ დიფერენციალი
	$$
	df=\frac{\partial f}{\partial x}dx+\frac{\partial f}{\partial y}dy+\frac{\partial f}{\partial z}dz.
	$$
	
	\bigskip
	
	1-ფორმის გარე დიფერენციალი არის 2- ფორმა: 
	$$d\phi=d(fdx+gdy+hdz)=df\land dx+dg\land dy+dh\land dz,$$
	შემდეგ უნდა დავასრულოთ ამ 1-ფორმების გადამრავლება და შევკრიბოთ.
	
	\begin{example} \label{27}
		გამოთვალეთ შემდეგი 1-ფორმის დიფერენციალი $\phi= xy \,\,dx+ x^2 \,\,dz$.
	\end{example}
	
	ამოხსნა: $d\phi=d(xy)\land dx+ d(x^2)\land dz=(y\,\,dx+x\,\,dy)\land dx+(2x\,\,dx)\land dz=-x\,\,dx\,\,dy+2x\,\,dx\,\,dz$
	
	\bigskip
	
	2-ფორმის გარე დიფერენციალი არის 3-ფორმა. დიფერენციალი განისაზღვრება ასე:
	$$d\psi=d(fdxdy+gdxdz+hdydz)=df\land fdxdy+dg\land dxdz+dh\land dydz,$$
	შემდეგ უნდა გამოვთვალოთ 1-და 2-ფორმების გარე ნამრავლი და შევკრიბოთ.
	
	ამ წესებით ცხადია,
	რომ $R^3$-ში 3-ფორმის დიფერენციალი 0-ია.
	
	\begin{example} \label{26}
		გამოთვალეთ შემდეგი 1-ფორმების გარე ნამრავლი 
		
		$\phi= x \,\,dx- y \,\,dy$ და $\psi= z\,\, dx+x\,\,dz$.
	\end{example}
	
	\begin{example} \label{28}მოცემულია
		
		$\phi= yz \,\,dx+ dz$;\,\,\, $\psi=\sin z\,\, dx+\cos z\,\,dy$; \,\,\, $\xi= xydxdy+ z\,dxdz$. 
		
		იპოვეთ: 
		
		ა) $\phi \land \psi$ ,\,\,\,\, $ \psi \land \xi$,\,\,\, $\xi \land \phi$ \,\,\,\,\,\,\,\,\,\,\,\,\,\,\,\,\,\, ბ) $d\phi,\,\,\,\,\,\, d\psi,\,\,\,\,\,\, d\xi$.
	\end{example}
	
	\begin{example}
		აჩვენეთ რომ ნებისმიერი $f,g$ ფუნქციებისთვის 
		$$d(df)=0, \text{ და გამოიყვანეთ აქედან, რომ } d(fdg)=df \land dg.$$
	\end{example}
	
	\begin{example}
		გაამარტივეთ
		
		\medskip
		
		\,\,\,\,\,\,\,\,\,\,\,\,
		ა) $d(fdg+dfg)$.\,\,\,\,\,\,\,\,\,\,\,\,\,\,\,\,\,\,\, ბ) $d((f-g)(df+dg))$.
		
		\,\,\,\,\,\,\,\,\,\,\,\,
		გ) $d(fdg\land dfg)$.\,\,\,\,\,\,\,\,\,\,\,\,\,\,\,\,\,\,\, დ) $d(gfdf)+d(fdg)$.
		
	\end{example}
	
	\bigskip
	
	\begin{theorem}
		ვთქვათ $\land^k(R^n)$ არის $R^n$-ში განსაზღვრული დიფერენციალური $k$-ფორმების სიმრავლე $R^n$-ის ნებისმიერ ფიქსირებულ წერტილში. მაშინ $\land^k(R^n)$ არის $\binom{n}{k}$ განზომილებიანი ვექტორული სივრცეს ბაზისით $dx_{i_1}\land \cdots \land dx_{i_k}$. 
	\end{theorem}
	სიმარტივისთვის განვიხილოთ $n=3$. 0-ფორმები არის ფუნქციები $R^3$-ში. ფიქსირებული წერტილში ფუნქციების შეზღუდვების სიმრავლე არის $R$, ანუ $\binom{3}{0}=1$ ერთგანზომილებიანი ნამდვილი ვექტორული სივრცე ბაზისით 1. შემდეგ, $\land^1(R^3)$-ის ყოველი 1-ფორმა სივრცის $p=(x_0,y_0,z_0)$ წერტილში ჩაიწერება $dx,dy,dz$ ბაზისში 
	$$f(p)dx+g(p)dy+h(p)dz,$$
	ანუ აქვს კოორდინატები $f(p),g(p),h(p)$. ყოველ წერტილში მოცემული 1-ფორმების ჯამი 1-ფორმაა და ფუნქციის ნამრავლი 1-ფორმაზე ისევ 1- ფორმაა იმავე წერტილში. ე.ი $\land^1(R^3)$ არის $\binom{3}{1}=3$ განზომილებიანი ვექტორული სივრცე. ანალოგიურად  $\land^2(R^3)$ ბაზისით $dxdy,dxdz,dydz$ და $\land^3(R^3)$ ბაზისით $dxdydz$. 
	
	რ.დ.გ
	
	\begin{example} 
		\label{dxdy=dudv}
		ვთქვათ $f,g$ ორი ცვლადის ნამდვილი ფუნქციებია. აჩვენეთ
		\begin{equation*}
			df\land dg=\left|
			\begin{array}{cc}
				\frac{\partial f}{\partial x} & \frac{\partial f}{\partial y}\\ 
				\frac{\partial g}{\partial x} & \frac{\partial g}{\partial y}\\
			\end{array}
			\right|dxdy.
		\end{equation*}
	\end{example}
	
	\begin{example}
		ნებისმიერი 1-ფორმებისთვის $\phi_{i}=\sum_{j=1}^3f_{ij}dx_j$ აჩვენეთ
		\begin{equation*}
			\phi_1\land \phi_2 \land \phi_3=\left|
			\begin{array}{ccc}
				f_{11} & f_{12} & f_{13}\\ 
				f_{21}& f_{22} & f_{23} \\
				f_{31} & f_{32} & f_{33}\\
			\end{array}
			\right|dx_1dx_2dx_3.
		\end{equation*}
	\end{example}
	
	\bigskip
	
	\begin{theorem} ვთქვათ $\phi$ და $\psi$ 1-ფორმებია. 
		
		1. $d(f\phi)=df\land \phi+fd\phi$;
		
		2. $d(\phi\land \psi)=d\phi\land \psi-\phi\land d\psi$.
	\end{theorem}
	
	1 უფრო ადვილია და დამოუკიდებლად ვცადოთ.
	ვაჩვენოთ 2. საკმარისია განვიხილოთ $\phi=fdx$ და $\psi=gdy$ და ვნახოთ საიდან გაჩნდა მინუს ნიშანი. ზოგადი შემთხვევა ასეთი ფორმების ნამრავლების წრფივი კომბინაციაა. 
	
	\begin{align*}
		d(\phi\land \psi) &= \\
		& d(fg)\land dxdy=(df g+fdg)\land dxdy=\\
		& df g\land dxdy+fdg\land dxdy=\\
		& g\frac{\partial f}{\partial z}dz\land dxdy+f\frac{\partial g}{\partial z}dz\land dxdy=\\
		& (g\frac{\partial f}{\partial z}+f\frac{\partial g}{\partial z}) dxdydz.
	\end{align*}
	მეორეს მხრივ
	\begin{align*}
		d \phi\land \psi &= \frac{\partial f}{\partial z} dz\land dx \land gdy=g\frac{\partial f}{\partial z}\land dxdydz;\\
		\phi \land d\psi &= fdx\land \frac{\partial g}{\partial z}dzdy=-f\frac{\partial g}{\partial z}dxdydz.
	\end{align*}
	
	რ.დ.გ.
	
	\bigskip
	
	$w$ დიფერენციალურ ფორმას ეწოდება ჩაკეტილი თუ $dw=0$. თუ $w$ არის რაიმე ფორმის დიფერენციალი $w=d \phi $, მაშინ $w$-ს ეწოდება ზუსტი. შეიძლება ფორმა იყოს ჩაკეტილი მაგრამ არ იყოს ზუსტი. სამაგიეროდ 
	ნებისმიერი ზუსტი ფორმა ჩაკეტილია, ანუ სამართლიანია
	
	\begin{theorem} 
		ნებისმიერი $w$ დიფერენციალური ფორმისათვის
		$$d(dw)=0.$$
	\end{theorem}
	
	\subsection{დიფერენციალური ფორმები და ვექტორული ველები} 
	\label{fields-forms} 
	
	ვთქვათ $x,y,z$ სტანდარტული კოორდინატებია $R^3$-ში.
	ჩვენი მიზანია განვსაზღვროთ ასახვები 
	
	\begin{align*}
		&T_0:0\text{ -ფორმები} \to \text{ფუნქციები } R^3\text{-ში}\\
		&T_1:1\text{ -ფორმები} \to \text{ვექტორული ველები } R^3\text{-ში}\\
		&T_2:2\text{ -ფორმები} \to \text{ვექტორული ველები } R^3\text{-ში}\\
		&T_3:3\text{ -ფორმები} \to \text{ფუნქციები } R^3\text{-ში}.\\
	\end{align*}
	
	$T_0,T_1,T_3$-ის განსაზღვრებები ადვილია, $T_2$-ის განსაზღვრება ცოტა მოსაზრებას ითხოვს.
	
	$T_0$ იგივური ასახვაა.
	
	ვიცით, რომ $1$-ფორმა საზოგადოდ ასე ჩაიწერება $w=f_1dx+f_2dy+f_3dz$. განვსაზღვროთ
	$$
	T_1(w)=(f_1,f_2,f_3).
	$$
	
	$T_3$-ის განსაზღვრებაც ვხადია: თუ 
	$w=fdxdydz$ 3-ფორმაა, მაშინ
	$$
	T_3(w)=f.
	$$
	შემდეგ, $T_2$-ის ფორმულა ცოტა უცნაურია თუ $w_2=f_1dxdy+f_2dxdz+f_3dydz$:
	$$
	T_2(w)=(f_3,-f_2,f_1)
	$$
	
	ეს ასე შეიძლება აიხსნას: $R^3$-ის ყოველ წერტილში 2-ფორმების სივრცე $\land^2(R^3)$ სამგანზომილებიანი ნამდვილი ვექტორული სივრცეა ბაზისით $dxdy, dxdz, dydz$. მისი იზომორფულია 1-ფორმების $\land^1(R^3)$ ვექტორული სივრცე ბაზისით $dx, dy, dz$. ეს იზომორფიზმი ასეა მოწყობილი:: სივრცის ყოველ წერტილში 
	1-ფორმებს განვიხილავთ როგორც წრფივ ასახვებს 2-ფორმებიდან $R$-ში, ანუ $\land^1(R^3)$ არის $\land^2(R^3)$ ორადული სივრცე.
	ამ იზომორფისმით 
	$$ 
	w_1(w_2)=T_3(w_1\land w_2)
	$$
	ანუ $dxdy\to dz$, $dxdz\to -dy$, $dydz\to dx$. ეს არის მოტივაცია $T_2$-ასახვის ფორმულის.
	
	გავიხსენოთ 
	$$
	\mathrm{grad}(f)=\left(\frac{\partial f}{\partial x},\frac{\partial f}{\partial y}, \frac{\partial f}{\partial z}\right);
	$$
	\begin{equation*}
		\mathrm{curl}(f_1,f_2,f_3)=\left|
		\begin{array}{ccc}
			i & j & k\\
			\frac{\partial }{\partial x} & \frac{\partial }{\partial y}&\frac{\partial }{\partial z}\\ 
			f_1 & f_2 & f_3\\
		\end{array}
		\right|;
	\end{equation*}
	$$
	\mathrm{div}(f_1,f_2,f_3)=\frac{\partial f_1}{\partial x} + \frac{\partial f_2}{\partial y}+\frac{\partial f_3}{\partial z}.
	$$
	
	გვაქვს თეორემა
	
	\begin{theorem}
		ვთქვათ $w_k$ არის $R^3$-ში $k$-ფორმა. მაშინ
		\begin{align*}
			&T_1(dw_0)=\mathrm{grad}(T_0(w_0));\\
			&T_2(dw_1)=\mathrm{curl}(T_1(w_1));\\
			&T_3(dw_2)=\mathrm{div}(T_2(w_2)).\\
		\end{align*}
	\end{theorem}
	
	\begin{example}
		$f$ და $g$ დიფერენცირებადი ფუნქციებია $R^2$-ზე. იპოვეთ $df\wedge dg$.
	\end{example}
	
	\begin{example}
		განვიხილოთ ანალიზური ფუნქციები $f_1,\,\,f_2:R^2\to R$ 
		\begin{equation*}
			f_1(x_1,x_2)=x_1+x_2;\,\,\,\,f_2(x_1,x_2)=x^2_1+x^2_2-1.
		\end{equation*}
		
		ა) იპოვეთ $df_1$,\,\, $df_2$,\,\, $df_1\wedge df_2$.
		
		ბ) ამოხსენით განტოლებათა სისტემა
		$$df_1\wedge df_2=0,\,\,\,x^2_1+x^2_2-1=0.$$
	\end{example}
	
	\bigskip
	\bigskip
	
	\subsection{დიფერენციალური ფორმები ზედაპირებზე} 
	
	ვთქვათ $f:M\to R$ არის რაიმე ნამდვილი ფუნქცია $M$ ზედაპირზე და 
	$${\bf x}={\bf x}(u,v):D\to M,\,\,\,(u,v)\in D\subset R^2$$ 
	არის საკოორდინატო რუქა. კომპოზიციას 
	$$D\xrightarrow{{\bf x}} M \xrightarrow{f}R,$$
	ანუ, $f({\bf x}(u,v))$-ს ეწოდება 
	$f$-ის კოორდინატული გამოსახვა; ეს კომპოზიცია უკვე არის ორი ნამდვილი ცვლადის 
	ნამდვილი მნიშვნელობის ფუნქცია
	$$R^2\ni (u,v)\to f({\bf x}(u,v))\in R,$$
	და კალკულუსიდან ვიცით მისი კერძო წარმოებულების და დიფერენციალის განსაზღვრება. 
	
	\begin{definition} 
		$f:M\to R$ ფუნქციას ეწოდება დიფერენცირებადი, თუ მისი ყველა საკოორდინატო გამოსახვა დიფერენცირებადია ჩვეულებრივი ევკლიდური გაგებით.
	\end{definition} 
	
	\bigskip
	
	ასევ, თუ გვაქვ სრაიმე $n$ ნამდვილი ცვლადის $F:R^n \to M$ ფუნცია, თითოეული რუქა $\bf{x}$ იძლევა $F$-ის საკოორდინატო გამოსახვას 
	$$
	R^n\xrightarrow{F} M\xrightarrow{{\bf x}^{-1}} D,
	$$
	ანუ
	$${\bf x}^{-1}(F):R^n\to D.$$
	ცხადია ეს კომპოზია განისაზღვრება $R^n$–ის იმ ღია $O$ ქვესიმრავლეზე, რომ $F(O)\subset x(D$. 
	
	\begin{definition} $F:R^n \to M$ ფუნცია არის დიფერენცირებადი, თუ მისი ყველა საკოორდინატო გამოსახვა დიფერენცირებადია
		ჩვეულებრივი გაგებით.
	\end{definition}
	
	კერძოდ განვიხილოთ $n=1$ შემთხვევა, ანუ წირი 
	$$
	\alpha:I\to M
	$$ 
	რომლის სახე ძევს ${\bf x}(D)\subset M$-ში. მაშინ არსებობს საკოორდინატო ფუნქციები 
	$a_1(t),a_2(t)$, ისე რომ, $\alpha(t)=(a_1(t),a_2(t))$. მართლაც ავიღოთ
	$$(a_1(t),a_2(t)):\bf{x}^{-1}(\alpha(t)).$$ 
	
	როცა ზედაპირი $R^3$ ევკლიდურ სივრცეში ძევს მაშინ ყოველი რუქა არის ასახვა $R^2\to R^3$, ხოლო ზედაპირზე განსაზღვრული ვეტორ-ფუნქცია შეიძლება განვიხილოთ როფორც ფუნქცია $R^3\to R^n$. ისე რომ ამ შემთხვევებში დიფერენცირებადობა "ჩვეულებრივია". 
	
	\begin{definition}
		აღვნიშნოთ \textbf{დიფერენცირებადი ფუნქციების სიმრავლე} $M\subset R^3$-ზე
			$\mathcal{F}(M)$-ით. მხები ვექტორების სიმრავლე ზედაპირის $p$-წერტილში ავღნიშნოთ $T_p$-თი.\textbf{ დიფერენცირებადი ვექტორული ველების ფუნქცია} (რომელიც ზედაპირის ყოველ წერტილს შეაუსაბამებს რაიმე მხებ ვექტორს ამ წერტილში ) სიმრავლე აღინიშნება $\mathcal{V}(M)$-ით.
	\end{definition}
	
	დიფერენციალური ფორმები ზედაპირებზე განისაზღვრება ევკლიდურ სივრცეში დიფერენციალური ფორმების პირდაპირი ანალოგიით:
	დიფერენციალური $m$ ფორმა $M\in R^n$–ზე არის ანტისიმეტრიული (ორი არგუმენტის გადასმისას ნიშანს იცვლის) მრავალწრფივი (ყოველი კომპონენტის მიმართ წრფივი)
	ფუნქცია
	$$T_p(M) \times T_p(M) \times \cdots \times T_p(M)\to R,$$
	რომელიც $M$-ზე დიფერენცირებადი ვექტორულ ი ველების $m$-ეულებს აქცევს $M$-ზე ფუნქციებად.
	
	მაგალითად, თუ $\eta$ არის $2$ ფორმა $M$-ზე 
	$$\eta:T_p(M) \times T_p(M)\to R,$$
	და $V,W \in \mathcal{V} (M)$, ჩვენ მოვითხოვთ, რომ $\eta(V, W)\in F(M)$.
	ფორმების წრფივი კომბინაციები და გარე ნამრავლები განისაზღვრება, როგორც ადრე:
	
	\begin{definition}[$0$-ფორმის დიფერენციალი.]
		თუ გვაქვს $f \in F(M)$, განვსაზღვროთ $1$-ფორმა $M$-ზე, ანუ დიფერენციალი $df$ ასე 
		$$df(v_p) = v_p[f],\,\,\, v_p \in T_p(M).$$
	\end{definition}
	
	თუ მოცემულია რუქა $\mathbf{x}: D \to M$, საჭიროა რუქის საკოორდინატო დიფერენცირებადი ფუნქციები $\tilde{u}$, $\tilde{v}$ და 1-ფორმები $d\tilde{u}, d\tilde{v}$, იმის ანალოგიურად რაც სიბრტყეზე გვქონა $u,v,du,dv$. 
	
	განვსაზღვროთ
	$$
	\tilde{u}: \mathbf{x}(D)\to R \text{ \,\,\,და \,\,\,} \tilde{v}:\mathbf{x}(D) \to R
	$$
	ისე, რომ 
	$\tilde{u}(\mathbf{x} (u, v))=u$ და $\tilde{v} (\mathbf{x}(u, v)) = v$.
	გავიხსენოთ, რომ $\mathbf{x} (u_0, v_0) = p$ წერტილში მოცემული $T_p(M)$ მხები სიბრტყის ბაზისი მცემულია წყვილით: 
	$$(\mathbf{x}_u(u_0, v_0), \mathbf{x}_v (u_0, v_0)) = (\mathbf{x}_u \circ\mathbf{x}^{-1}(p), \mathbf{x}_v \circ \mathbf{x}^{-1}(p)).$$ 
	
	ამ საფუძველზე 
	გამოვთვლით
	$$
	d\tilde{u}(\mathbf{x}_u) = \mathbf{x}_u [\tilde{u}] = \frac{\partial(\tilde{u} \circ x)}{u}=1
	$$
	და
	$$
	d\tilde{u}(\mathbf{x}_v) = \mathbf{x}_u [\tilde{v}] = \frac{\partial(\tilde{u} \circ x)}{v}=0.
	$$
	ანალოგიურად $d\tilde{v}(\mathbf{x}_v) =1$ და $d\tilde{v}(\mathbf{x}_u) =0$.
	
	ამრიგად $d\tilde{u}, d\tilde{v}$ ფორმები არის $\mathbf{x}_u, \mathbf{x}_v$ ბაზისის ორადული და სტანტარტული გზით მტკიცდება შემდეგი
	
	\begin{theorem}
		ვთქვათ $\mathbf{x}: D \to M$ არის $M\subset R^3$ ზედაპირის რუქა.
		
		(ა) ყოველი $1$-ფორმის $\eta $-ს ჩაწერა $\mathbf{x}(D)$-ზე ასე შეიძლება (ცალსახად) 
		$$\eta=fd\tilde{u} + gd\tilde{v},$$
		კერძოდ დიფერენცირებადი ფუნქციებით $f = \eta(\mathbf{x}_u)$ და $g = \eta (\mathbf{x}_v)$.
		
		(ბ) ყოველი $2$-ფორმის $\psi$-ს ჩაწერა $\mathbf{x}(D)$-ზე ასე შეიძლება (ცალსახად) 
		$$
		\psi= h d\tilde{u}\wedge d\tilde{v},
		$$
		კერძოდ დიფერენცირებადი ფუნქციით $h= \psi(\mathbf{x}_u,\mathbf{x}_v)$.
	\end{theorem}
	
	\begin{corollary}
		ზედაპირზე მოცემული ფუნქციისთვის
		$$
		df=\frac{\partial f}{\partial u}d\tilde{u}+\frac{\partial f}{\partial v} d\tilde{v},
		$$
		სადაც $\frac{\partial f}{\partial u}=df({{\bf x}_u})$ და 
		$\frac{\partial f}{\partial v}=df({{\bf x}_v})$.
	\end{corollary}
	
	\begin{definition}[$1$-ფორმის გარე წარმოებული]
		ვთქვათ $\eta $ არის $1$ ფორმა $M$ ზედაპირზე. ჩვენ განვსაზღვრავთ $2$-ფორმას
		$d\eta$ ჯერ $\eta=fd\tilde{u}+gd\tilde{v}$ ფორმის გამოსახვით 
		ყოველ $\mathbf{x}: D \to M$ რუქაზე და შემდეგ იმის გამოცხადებით, რომ 
		$$d\eta = df \wedge d\tilde{u} + dg \wedge d\tilde{v}$$
		ყოველ $p\in \mathbf{x}(D)$ წერტილში მხებ სიბრტყეზე $T_p (M)$. ბოლოს უნდა დამტკიცდეს,
		რომ ეს განსაზღვრება არ არის დამოკიდებული კონკრეტული რუქის არჩევანზე.
	\end{definition}
	
	რაგდან $d\eta$ სრულიად განისაზღვრება მისი მნიშვნელობით $(\mathbf{x}_u, \mathbf{x}_v)$ წყვილზე, ვისარგებლოთ განსაზღვრების ფორმილით. მივიღებთ შემდეგ სასარგებლო ფორმულას.
	
	\begin{theorem}[$1$-ფორმის დიფერენციალის ფორმულა]
		\begin{align*}
			d\eta(\mathbf{x}_u, \mathbf{x}_v)= 
			\frac{\partial}{\partial u}(\eta(\mathbf{x}_v))-\frac{\partial}{\partial v}(\eta(\mathbf{x}_u)). 
		\end{align*}
	\end{theorem}
	
	დამტკიცება.
	\begin{align*}
		d\eta(\mathbf{x}_u, \mathbf{x}_v) &= \\ 
		& df(\mathbf{x}_u)d\tilde{u}(\mathbf{x}_v)-df(\mathbf{x}_v)d\tilde{u}(\mathbf{x}_u) 
		+ dg(\mathbf{x}_u)d\tilde{v}(\mathbf{x}_v)- dg(\mathbf{x}_v)d\tilde{v}(\mathbf{x}_u)=\\
		& df(\mathbf{x}_u) \cdot 0-df(\mathbf{x}_v) \cdot 1 + dg(\mathbf{x}_u) \cdot 1 - dg(\mathbf{x}_v) \cdot 0=\\
		& \left( \frac{\partial g}{\partial u}-\frac{\partial f}{\partial v} \right)=\\
		& \frac{\partial}{\partial u}(\eta(\mathbf{x}_v))-\frac{\partial}{\partial v}(\eta(\mathbf{x}_u)). 
	\end{align*}
	
	\textbf{შენიშვნა}. ვინაიდან თითოეულ $T_p(M)$-ს აქვს ორი ვექტორისგან შემდგარი ბაზისი, ყველა $\geq 3$-ფორმა $R^3$-ის ზედაპირებზე ნულია.
	
	\medskip
	
	შემდეგი ორ თეორემის მტკიცება ევკლიდეს სივრციდან მოდის.
	
	\begin{theorem}
		ვთქვათ $M \subset R^3$ ზედაპირია,
		$f,g \in \mathcal{F}(M)$, და ვთქვათ $\eta$ და $\sigma$ $1$-ფორმებია $M$-ზე.
		მაშინ
		$$
		d(fg) = g df+fdg,\,\,\, d(f \eta) = (df)\wedge \eta +f d\eta,\,\,\, d(f +g) = df +dg,\,\,\, d(\eta+\sigma) = d\eta+d\sigma.
		$$
	\end{theorem}
	
	\begin{theorem}
		ნებისმიერი $\eta $ დიფერენციალური ფორმისთვის $M \subset R^3$ ზედაპირზე ორჯერ დიფერენცირება ნულია, ანუ
		$d^2\eta=d(d\eta)=0$.
	\end{theorem}
	
	შემდეგი თეორემა ეხება ზედაპირზე და გარემომცველ ევკლიდური სივრცეში $1$-ფორმების გარე დიფერენციალის შენათხმებულობას. 
	
	\begin{theorem} ვთქვათ $\eta$ $1$-ფორმაა $R^3$ სივრცეში და $\sigma$ $1$-ფორმა $M\subset R^3$ ზედაპირზე, ისე რომ ზედაპირის ნებისმიერი $v_p$ მხები ვექტორისთვის 
		$\eta(v_p) = \sigma(v_p)$. მაშინ ზედაპირის ნებისმიერი მხები ვექტორების წყვილისთვის
		$$
		d\eta(v_p)(v_p,w_p) = d\sigma(v_p,w_p). 
		$$
	\end{theorem}
	
	\subsection*{დიფერენციალური   ფორმის  პულბექი} 
	
	\begin{definition}
		ვთქვა გვაქვს ორი ზედაპირი, $M$ ზედაპირი მოცემულია $M=\mathbf{x}(D)$ რუქით ხოლო $N$ ზედაპირი $N=\mathbf{y}(D')$ რუქით. ზედაპირების ასახვას $F: M \to N$ ეწოდება დიფერენცირებადი, თუ კომპოზიცია 
		$$\mathbf{y}^{-1}\circ F\circ \mathbf{x}:D\to D'$$
		დიფერენცირებადია, ჩვეულებრივი აზრით, როგორც $D\subset R^2$-ის ასახვა $D'\subset R^2$-ში. 
		
		
		$F$-ასახვის დიფერენციალი ნებისმიერ $p\in M$ წერტილში არის წრფივი ასახვა 
		$$F_*:T_p\to T_{F(p)},$$
		ანუ $F_*$ ასახვა $M$ ზედაპირის $p$ წერტილში ყოველ მხებ $v_p$ ვექტორს გადაიყვანს $N$-ზედაპირის $F(p)$ წერტილში $w_{F(p)}$ მხებ ვექტორში ასე:
		$$v_p=\alpha'(t)|_{t=0} \to (F(\alpha(t)))'|_{t=0},$$
		სადაც $\alpha(0)=p$, ხოლო კომპოზიცია $\beta=F(\alpha(t))$ განიხილება უკვე როგორც წირი $N$ ზედაპირზე, $\beta(0)=F(p)$. 
	\end{definition}
	
	\begin{definition}
		განვიხილოთ $R^3$ სივრცის ზედაპირების დიფერენცირებადი ასახვა $F:M\to N$.
		არსებობს კონსტრუქცია, ($N$-ზე მოცემული დიფერენციალური ფორმის ე.წ. პულბექი $F$ ასახვით) რომლითაც $N$-ზე მოცემული დიფერენციალური ფორმით აიგება მისი "პულბექი" ფორმა $M$-ზე ასე:
		
		a) 0-ფორმის (ანუ ფუნქციის) პულბექი: 
		
		თუ $f \in \mathcal{F}(N)$, განვსაზღვროთ $f$-ის პულბექი $F^*(f)\in \mathcal{F}(M)$ ტოლობით $$F^*(f)(p)=f(F(p)).$$
		ანუ, $N$-ზე მოცემული $f$ ფუნქციის პულბექი $F^*(f)$ არის ფუნქცია $M$-ზე, რომლის მნიშვნელობა $p\in M$ წერტილში იგივეა რაც $f$-ის მნიშვნელობა $F(p)$ წერტილში.
		
		ბ) თუ $\eta$ $1$-ფორმაა $N$-ზე, მისი პულბექი $F^*(\eta)$ არის $1$-ფორმა $M$-ზე, და განისაზღვრება ტოლობით 
		$$F^*(\eta)(v_p)=\eta(F_*(v_p)), \,\,\,\forall v_p \in T_p(M).$$
		
		ანუ, $N$-ზე 1-ფორმის $F^*(\eta)$ პულბექი არის 1-ფორმა $M$-ზე, რომლის მნიშვნელობა $M$-ის ნებისმიერ მხებ $v_p$ ვექტორზე იგივეა, რაც $\eta$-ს მნიშვნელობა $N$-ის იმ მხებ ვექტორზე, სადაც გადადის $v_p$ ვექტორი $F$ ფუნქციის დიფერენციალით.
		
		გ) თუ $\sigma$ $2$-ფორმაა $N$-ზე, მისი პულბექი $F^*(\sigma)$ არის $2$-ფორმა $M$-ზე, და განისაზღვრება ტოლობით 
		$$F^*(\sigma)(v_p,w_p)=\sigma(F_*(v_p),F_*(w_p)), \,\,\,\forall v_p,w_p \in T_p(M).$$
	\end{definition}
	
	\begin{theorem}
		ვთქვათ $F : M \to N$ ზედაპირების დიფერენცირებადი სახვაა 
		და $\eta$ და $\sigma $ ფორმებია $N$-ზე. მაშინ
		\begin{align*} 
			&(a) F^*(\eta+\sigma)=F^*(\eta)+F^*(\sigma);\\
			&(b) F^*(\eta \wedge \sigma ) = F^*(\eta) \wedge F^*(\sigma);\\
			&(c) F^*(d\eta) = dF^*(\eta).
		\end{align*}
	\end{theorem}
	
	\subsection{ფორმების  ინტეგრება}
	
	\begin{definition}
		ვთქვათ $\phi$ $1$-ფორმაა $M$-ზე, $\alpha:[a,b]\to M$ რაიმე წირის ნაწილია. მაშინ $\phi$-დან ინტეგრალი $\alpha$-ს გასწვრივ ეწოდება
		$$
		\int_{\alpha}\phi=\int_{[a,b]}\alpha^*(\phi)=\int_a^b \phi(\alpha'(t))dt
		$$
	\end{definition}
	აქ $\alpha^*(\phi)$ არის $\phi$-ფორმის პულბექი $[a,b]$ მონაკვეთზე. [a.b]-ზე ყოველ $1$ ფორმას აქვს სახე $f(t)dt$, სადაც ბულბექის განსაზღვრებით $$f(t)=\alpha^*(\phi)(U_1)=\phi(\alpha'(t)).$$
	
	ვთქვათ ${\bf F}=P(x,y,z)i+Q(x,,z)j+R(x,,y,z)k$ გლუვი ვექტორული ველია $R^3$-ში.
	\ref{fields-forms} პუნქტის მიხედვით
	${ \bf F}$-ს შეესაბამება 1-ფორმა $\phi=Pdx+Qdy+Rdz$ სივრცეში. ამ ფორმიდან წირითი ინტეგრალი $\alpha(t)=(x(t),y(t),z(t))$ წირის გასწვრივ არის
	$$
	\int_{\alpha}\phi=\int_{\alpha}Pdx+Qdy+Rdz=\int_a^b Px'(t)dt+Qy'(t)dt+Rz'(t)dt=
	\int_{\alpha} F\cdot d\alpha.
	$$
	
	\begin{definition}
		ვთქვათ $\sigma $ $2$-ფორმაა $M$-ზე, 
		$$ D=\{(x,y)|a\leq x \leq b, c\leq y \leq d \} =[a,b]\times [c,d]$$ მართკუთხედია სიბრტყეზე, ხოლო 
		${\bf x}:D \to M $ მართკუთხედის სახეა ზედაპირზე. ინტეგრალი 
		$\sigma $-დან $R$-ზე ეწოდება
		$$
		\int_{\bf x}\sigma=\iint_D{\bf x}^*(\sigma)=\int_a^b\int_c^d\sigma({\bf x}_u,{\bf x}_v)dudv.
		$$
	\end{definition}
	აქ ${\bf x}^*(\sigma)$ არის $\sigma$-ფორმის პულბექი $R$ მართკუთხედზე.
	
	\begin{theorem}
		[სტოქსის ზოგადი თეორემა] თუ $R$ არის $p$-განზომილებიანი რეგიონი და $\eta$ არის $(p-1)$-ფორმა , მაშინ $d\eta$-დან ინტეგრალი $R$-ზე
		უდრის $\eta$-დან ინტეგრალს $R$-ის საზღვარზე. 
	\end{theorem}
	
	ამ თეორემის გამოყენების მიზნით განვიხილოთ ისევ ${\bf F}=Pi+Qj+Rk$ ვექტორული ველი და \ref{fields-forms} პუნქტის მიხედვით მისი შესაბამისი $2$ ფორმა
	$$\sigma= Pdx\wedge dy-Qdx\wedge dz+Rdy\wedge dz$$
	მისი გარე დიფერენციალი $d\sigma$ უდრის
	$$d(Pdx\wedge dy-Qdx\wedge dz+Rdy \wedge dz) 
	=\cdots =\left(\frac{\partial P}{\partial x}+\frac{\partial Q}{\partial y}+\frac{\partial R}{\partial z}\right)dx\wedge dy\wedge dz.
	$$
	(შევამოწმოთ!) ამიტომ მივიღეთ 3-ფორმა რომლის კომპონენტი 
	(ანუ ისევ \ref{fields-forms} პუნქტის მიხედვით შესაბამისი $0$ ფორმა ) არის 
	\begin{equation}
		\label{divergence}
		\left(\frac{\partial P}{\partial x}+\frac{\partial Q}{\partial y}+\frac{\partial R}{\partial z}\right)=\mathrm{div}(Pi+Qj+Rk),
	\end{equation}
	თავდაპირველი ${\bf F}$ ვექტორული ველის დივერგენცია. 
	
	\eqref{dxdy=dudv} მაგალითით გვაქვს
	
	\begin{equation*}
		dx\land dy=\left|
		\begin{array}{cc}
			\frac{\partial x}{\partial u} & \frac{\partial x}{\partial v}\\ 
			\frac{\partial y}{\partial u} & \frac{\partial y}{\partial v}\\
		\end{array}
		\right|dudv,
	\end{equation*}
	
	\begin{equation*}
		dx\land dz=\left|
		\begin{array}{cc}
			\frac{\partial x}{\partial u} & \frac{\partial x}{\partial v}\\ 
			\frac{\partial z}{\partial u} & \frac{\partial z}{\partial v}\\
		\end{array}
		\right|dudv,
	\end{equation*}
	
	\begin{equation*}
		dy\land dz=\left|
		\begin{array}{cc}
			\frac{\partial y}{\partial u} & \frac{\partial y}{\partial v}\\ 
			\frac{\partial z}{\partial u} & \frac{\partial z}{\partial v}\\
		\end{array}
		\right|dudv,
	\end{equation*}
	
	ამიტომ
	\begin{align*}
		&\int_{\bf x}Pdx\wedge dy-Qdx\wedge dz+Rdy \wedge dz) = \\
		&\int_D P\left|
		\begin{array}{cc}
			\frac{\partial y}{\partial u} & \frac{\partial y}{\partial v}\\ 
			\frac{\partial z}{\partial u} & \frac{\partial z}{\partial v}\\
		\end{array}
		\right|dudv-
		\int_D Q\left|
		\begin{array}{cc}
			\frac{\partial x}{\partial u} & \frac{\partial x}{\partial v}\\ 
			\frac{\partial z}{\partial u} & \frac{\partial z}{\partial v}\\
		\end{array}
		\right|dudv+
		\int_D R\left|
		\begin{array}{cc}
			\frac{\partial y}{\partial u} & \frac{\partial y}{\partial v}\\ 
			\frac{\partial z}{\partial u} & \frac{\partial z}{\partial v}\\
		\end{array}
		\right|dudv=\\
		&\int_D  {\bf F} \cdot \frac{{\bf x}_u\times {\bf x}_v}{\left| {\bf x}_u \times {\bf x}_v \right| } \left| {\bf x}_u \times {\bf x}_v \right|dudv=\\
		&\int_{{\bf x}}  {\bf F} \cdot {\bf n} \, dS.
	\end{align*}
	
	\bigskip
	
	სტოკსის ზოგადი თეორემით და \eqref{divergence} ტოლობით თუ $T$ სივრცით სხეულს (რეგიონს) აქვს $S$ საზღვარი, მაშინ ${\bf F}$ ვექტორული ველისთვის გვაქვს
	
	\begin{equation}
		\label{div}
		\int_S  {\bf F} \cdot {\bf n} \, dS=\int_T \mathrm{div}{\bf F} dxdydz.
	\end{equation}
	
	ეს ტოლობა ცნობილია { \bf დივერგენციის თეორემის } სახელით და მნიშვნელოვან როლს თამაშობს ფიზიკაში.
	
	\begin{example}
		ვთქვათ გვაქვს ვექტორული ველი ${\bf F}=(2x,3y,z^2)$ და განვიხილოთ კუბი, 
		საკოორდინატ ო სიბრტყეების პარალელური გვერდებით და მოპირდაპირე კუთხეებით $(0,0,0)$ და $(1,1,1)$. ჩვენ გამოვთვლით დივერგენციის თეორემის ორივე ინტეგრალს.. 
	\end{example}
	სამჯერადი ინტეგრალი ადვილია:
	$$
	\int_0^1\int_0^1\int_0^1(2+3+2z)dxdydz=6
	$$
	ზედაპირული ინტეგრალი უნდა გაიყოს ექვს ნაწილად კუბის 6 წახნაგისთვის. ერთი წახნაგია
	არის ${\bf x}=(u,v,0)$, ანუ $z=0$, $0\leq u,v\leq 1$. ამიტომ ${\bf x}_u=(1,0,0)$, ${\bf x}_v=(0,1,0)$,  ${\bf x}_u\times {\bf x}_v=(0,0,1)$.
	
	ჩვენ გვჭირდება ეს ეს ბოლო ვექტორი ორიენტიაციით ქვევით (კუბიდან გარეთ), ამიტომ ავიღოთ $(0,0,-1)$ და შესაბამისი ინტეგრალი არის 
	$$
	\int_0^1\int_0^1(-z^2)dudv=\int_0^1\int_0^1 0\, dudv=0.
	$$
	
	სხვა წახნაგია $y=1$, ანუ ${\bf x}=(u,1,v)$,  ამიტომ  ${\bf x}_u=(1,0,0)$, ${\bf x}_v=(0,0,1)$,  ${\bf x}_u\times {\bf x}_v=(0,-1,0)$. ჩვენ გვჭირდება ბოლო ვექტორის ორიენტაცია $y$ ის დადებითი მიმართულებით , ანუ უნდა ავიღოთ $(0,1,0)$
	და გვექნება
	$$
	\int_0^1\int_0^1 3ydudv=\int_0^1\int_0^1 3\, dudv=3.
	$$
	კუბის გადარჩენი ოტხი წახნაგის შესაბამისი ინტეგრალები ტოლია, 0,0,2 და 1. (შეამოწმეთ!) ისე, რომ ჯამში გვექნება 6, რაც ეთანხმება \eqref{div} დივერგენციის თეორემას.
	
	\bigskip
	
	\begin{example}
		ვთქვათ გვაქვს ვექტორული ველი ${\bf F}=(x^3,y^3,z^2)$ და განვიხილოთ ცილინდრული სხეული $x^2+y^2\leq 9$, $0\leq z\leq 2$.  
	\end{example}
	ცილინდრული კოორდინატების გამოყენებით სამჯერადი ინტეგრალი ტოლია 
	$$
	\int_0^{2\pi}\int_0^3\int_0^2(3r^2+2z)dzdrd\theta=279\pi.
	$$
	ზედაპირისთვის გვჭირდება სამი ინტეგრალი. ცილინდრის ზედა ფუძე შეიძლება ასე წარმოვადგინოთ ${\bf x}=(v\cos u,v\sin u,2)$, ხოლო ${\bf x}_u\times {\bf x}_v=(0,0,-v)$. ეს ვექტორი მიმართულია ცილიდრის შიგნით, და არაა გარე ნორმალი. ამიტომ მას ვცვლით ვექტორით $(0,0,v)$ და გვაქვს
	$$
	\int_0^{2\pi}\int_0^3(v^3\cos^3 u,v^3\sin^3 u ,4)\cdot (0,0,v)dvdu=\int_0^{2\pi}\int_0^3 4vdvdu=36\pi.
	$$
	ცილინდრის ქვედა ფუძეა ${\bf x}=(v\cos u,v\sin u,0)$, ხოლო ${\bf x}_u\times {\bf x}_v=(0,0,-v)$ რაც მიმართულია ცილიდრის გარეთ და გვაქვს
	$$
	\int_0^{2\pi}\int_0^3(v^3\cos^3 u,v^3\sin^3 u ,0)\cdot (0,0,-v)dvdu=\int_0^{2\pi}\int_0^3 0\,dvdu=0.
	$$
	ცილინდრის გვერდითი ზედაპირია 
	$${\bf x}=(3\cos u,3\sin u,v),$$
	ხოლო ${\bf x}_u\times {\bf x}_v=(3\cos u,3 \sin u,0)$ რაც მიმართულია ცილიდრის გარეთ და გვაქვს
	\begin{align*}
		\int_0^{2\pi}\int_0^3(27\cos^3 u,& 27^3\sin^3 u ,v^2)\cdot (3\cos u,3\sin u,0)dvdu=\\
		&\int_0^{2\pi}\int_0^3 (81\cos^4 u+81\sin^4 u)dvdu=243\pi.
	\end{align*}
	ამრიგად სრული ზედაპირუ ინტეგრალი ტოლია $36\pi+0+243\pi=279\pi$, რაც თავიდან გამოთვლილ სივრცით ინტეგრალ უდრის და ამიტომ ეთანხმება დივერგენციის თეორემას.
	
	\bigskip
	
	\subsection{სავარჯიშოები} 
	ქვემოთ მოცემულ პირობებში გამოთვალეთ \eqref{div} დივეგენციის თეორემის ორივე ინტეგრალი,
	
	\begin{example}
		ვექტორული ველია ${\bf F}=(3x,y^3,-2z^2)$ და რეგიონი შემოსაზღვრულია 
		$$x^2+y^2=9, \,z=0, \, z=5 .$$
	\end{example}
	პასუხი: ორივე ტოლია $\frac{-45\pi}{4}$
	
	\begin{example}
		ვექტორული ველია ${\bf F}=(x^2,y^2,z^2)$, ხოლო რეგიონი 
		$$0\leq x\leq a,\, 0\leq y\leq b, \,0\leq z \leq c.$$
	\end{example}
	პასუხი: ორივე ტოლია $a^2bc+ab^2c+abc^2$.
	
	
	\begin{example}
		ვექტორული ველია ${\bf F}=(2xy,3xy,ze^{x+y})$, ხოლო რეგიონი 
		$$0\leq x\leq 1,\, 0\leq y\leq 1, \,0\leq z \leq 1.$$
	\end{example}
	პასუხი: 3.
	
	\begin{example}
		ვექტორული ველია ${\bf F}=(x^3,y^3,z^3)$, ხოლო რეგიონი 
		$$x^2+y^2+z^2\leq 4.$$
	\end{example}
	პასუხი: $\frac{384\pi}{5}$.
	
	\begin{example}
		ვექტორული ველია ${\bf F}=(\sqrt{x^2+y^2+z^2},\sqrt{x^2+y^2+z^2},\sqrt{x^2+y^2+z^2})$, ხოლო რეგიონი 
		$$0\leq z\leq{1-x^2-y^2}.$$
	\end{example}
	პასუხი: $\frac{\pi}{3}$.
	
	\begin{example}
		ვექტორული ველია ${\bf F}=(xy^2,yz,x^2z)$, ხოლო რეგიონი 
		$$x^2+y^2\leq 1,\,0\leq z\leq 4 .$$
	\end{example}
	პასუხი: $ 10\pi$.
	
	
	
	
	
	
	
	
	\section{მრავალნაირობები}
	
	\subsection{ტოპოლოგიური მრავალნაირობები}
	
	\begin{definition}
		ჰაუსდორფის ტოპოლოგიურ $M$ სივრცეს თვლადი ბაზისით, რომელიც ლოკალურად $\mathbb{R}^m$-ის ჰომეომორფულია, ეწოდება
		$m$ განზომილებიანი ტოპოლოგიური მრავალნაირობა.
	\end{definition}
	იგულისხმება, რომ ევკლიდურ სივრცეს $\mathbb{R}^m$-ს აქვს მეტრიკული ტოპოლოგია. ლოკალურად ჰომეომორფულობა ნიშნავს, რომ $\forall p\in M$ წერტილისთვის არსებობს მისი მიდამო $U\ni p$, ღია სიმრავლე $U'\subseteq \mathbb{R}^m$ და ჰომეომორფიზმი $\phi:U\to U'$.
	
	$\phi$-ს ეწოდება რუქა, ხოლო $U$-ს რუქის არე. რუქების კოლექციას $\{\phi_{\alpha}: U_{\alpha} \to U'_{\alpha}\}$, ისე რომ $\bigcup U_{\alpha}=M$, ეწოდება ატლასი.
	
	\subsection{დიფერენცირებადი მრავალნაირობები}
	
	მრავალნაირობის უმარტივესი მაგალითი არის ევკლიდური სივრცე $\mathbb{R}^n$. ჩვენ ვიცით რა არის $F:\mathbb{R}^m\to \mathbb{R}$ ფუნქციის დიფერენციალი. თუ მოვინდომებთ დიფერენციალური აღრიცხვის გადატანას მრავალნაირობებზე, მაგალითად, განისაზღვროს მრავალნაირობების $M\to N$ ასახვის დიფერენციალი, პირველ რიგში უნდა განისაზღვროს $f:M\to \mathbb{R}$ ასახვის დიფერენციალი. თუ გვაქვს $a\in U\subset M$ წერტილისთვის რუქა
	$\phi:U\to U'\subseteq\mathbb{R}^m$ მაშინ ბუნებრივია განვიხილოთ კომპოზიცია
	\begin{equation}
		\label{comp}
		f\circ \phi^{-1}:U'\subset \mathbb{R}^m\to \mathbb{R}.
	\end{equation}
	$f$-ფუნქცია დიფერენცირებადია $a\in M$ წერტილში, თუ \eqref{comp} კომპოზიცია დიფერენცირებადია $\phi(a)$ წერტილში და
	\[
	Df(a)=D(f\circ \phi^{-1})(\phi(a)).
	\]
	
	\begin{figure}[h]
		\centering
		\begin{tikzpicture}
			% Functions
			\path[->] (0.8, 0) edge [bend right] node[left, xshift=-2mm] {$\phi$} (-1, -2.9);
			\draw[white,fill=white] (0.06,-0.57) circle (.15cm);
			\path[->] (-0.7, -3.05) edge [bend right] node [right, yshift=-3mm] {$\phi^{-1}$} (1.093, -0.11);
			\draw[white, fill=white] (0.95,-1.2) circle (.15cm);
			
			\path[->] (4.2, 0) edge [bend left] node[right, xshift=2mm] {$f$} (6.2, -2.8);
			\draw[white, fill=white] (4.54,-0.12) circle (.15cm);
			
			% Manifold
			\draw[smooth cycle, tension=0.4, fill=white, pattern color=brown, pattern=north west lines, opacity=0.7] plot coordinates{(2,2) (-0.5,0) (3,-2) (5,1)} node at (3,2.3) {$M$} node at (1.5,0.3) {$. a$};
			
			% Subsets
			\draw[smooth cycle, pattern color=orange, pattern=crosshatch dots]
			plot coordinates {(1,0) (1.5, 1.2) (2.5,1.3) (2.6, 0.4)}
			node [label={[label distance=-0.3cm, xshift=-2cm, fill=white]:$U$}] {};
			
			% R^m space
			\node at (-2.4,-5) {$U'$};
			\node at (-3,-3) {$\mathbb{R}^m$};
			
			% Arrow to R
			\draw[->] (5, -3.85) -- node[midway, above]{$\mathbb{R}$} (7.5, -3.85);
			
			% Sets in R^m
			\fill[pattern color=orange, pattern=crosshatch dots, smooth cycle] plot coordinates{(-2, -4.5) (-2, -3.2) (-0.8, -3.2) (-0.8, -4.5)};
			\draw[smooth cycle] plot coordinates{(-2, -4.5) (-2, -3.2) (-0.8, -3.2) (-0.8, -4.5)};
			
		\end{tikzpicture}
		\caption{რუქა $\phi:U\to U'$. დიფერენცირებადი ასახვა $f:M\to \mathbb{R}$}
	\end{figure}
	
	მრავალნაირობაზე განსაზღვრული ფუნქციის დიფერენცირებადობა არ უნდა იყოს დამოკიდებული რუქაზე. ანუ თუ გვაქვს ორი რუქა $\phi_1, \phi_2$,
	\[
	f \circ \phi_1^{-1}= (f \circ \phi_2^{-1}) \circ (\phi_2 \circ \phi_1^{-1}).
	\]
	ამიტომ უნდა მოვითხოვოთ ეს ორი რუქა იყოს თავსებადი, ანუ კომპოზიცია $\phi_2 \circ \phi_1^{-1}$, ე.წ. \textbf{რუქების გადასვლის ფუნქცია}, იყოს დიფერენცირებადი.
	
	განვიხილოთ მრავალნაირობა $(M,A)$ დიფერენცირებადი $A$ ატლასით, ანუ ატლასის რუქების გადასვლის ფუნქციები დიფერენცირებადია.
	
	\begin{figure}[h]
		\centering
		\begin{tikzpicture}
			% Functions i
			\path[->] (0.8, 0) edge [bend right] node[left, xshift=-2mm] {$\phi_i$} (-1, -2.9);
			\draw[white,fill=white] (0.06,-0.57) circle (.15cm);
			\path[->] (-0.7, -3.05) edge [bend right] node [right, yshift=-3mm] {$\phi^{-1}_i$} (1.093, -0.11);
			\draw[white, fill=white] (0.95,-1.2) circle (.15cm);
			
			% Functions j
			\path[->] (5.8, -2.8) edge [bend left] node[midway, xshift=-5mm, yshift=-3mm] {$\phi^{-1}_j$} (3.8, -0.35);
			\draw[white, fill=white] (4,-1.1) circle (.15cm);
			\path[->] (4.2, 0) edge [bend left] node[right, xshift=2mm] {$\phi_j$} (6.2, -2.8);
			\draw[white, fill=white] (4.54,-0.12) circle (.15cm);
			
			% Manifold
			\draw[smooth cycle, tension=0.4, fill=white, pattern color=brown, pattern=north west lines, opacity=0.7] plot coordinates{(2,2) (-0.5,0) (3,-2) (5,1)} node at (3,2.3) {$M$};
			
			% Subsets
			\draw[smooth cycle, pattern color=orange, pattern=crosshatch dots]
			plot coordinates {(1,0) (1.5, 1.2) (2.5,1.3) (2.6, 0.4)}
			node [label={[label distance=-0.3cm, xshift=-2cm, fill=white]:$U_i$}] {};
			\draw[smooth cycle, pattern color=blue, pattern=crosshatch dots]
			plot coordinates {(4, 0) (3.7, 0.8) (3.0, 1.2) (2.5, 1.2) (2.2, 0.8) (2.3, 0.5) (2.6, 0.3) (3.5, 0.0)}
			node [label={[label distance=-0.8cm, xshift=.75cm, yshift=1cm, fill=white]:$U_j$}] {};
			
			% First Axis
			\draw[thick, ->] (-3,-5) -- (0, -5) node [label=above:{$\phi_i(U_i)$}] {};
			\draw[thick, ->] (-3,-5) -- (-3, -2) node [label=right:{$\mathbb{R}^m$}] {};
			
			% Arrow from i to j
			\draw[->] (0.2, -3.85) -- node[midway, above, text width=3cm, align=center]{$\Psi_{ij}=\phi_{j}\rvert_{U_i\cap U_j}\circ \phi^{-1}_i$} (4.8, -3.85);
			
			% Second Axis
			\draw[thick, ->] (5, -5) -- (8, -5) node [label=above:{$\phi_j(U_j)$}] {};
			\draw[thick, ->] (5, -5) -- (5, -2) node [label=right:{$\mathbb{R}^m$}] {};
			
			% Sets in R^m
			\fill[pattern color=orange, pattern=crosshatch dots, smooth cycle] plot coordinates{(-2, -4.5) (-2, -3.2) (-0.8, -3.2) (-0.8, -4.5)};
			\draw[smooth cycle] plot coordinates{(-2, -4.5) (-2, -3.2) (-0.8, -3.2) (-0.8, -4.5)};
			
			\fill[pattern color=blue, pattern=crosshatch dots, smooth cycle] plot coordinates{(7, -4.5) (7, -3.2) (5.8, -3.2) (5.8, -4.5)};
			\draw[smooth cycle] plot coordinates{(7, -4.5) (7, -3.2) (5.8, -3.2) (5.8, -4.5)};
			
		\end{tikzpicture}
		\caption{ატლასის გადასვლის ფუნქციები}
	\end{figure}
	
	აქ გავიხსენოთ, რომ გადასვლის ფუნქცია, როგორც
	\[\Psi_{ij}=(g_1,\dots, g_m):\mathbb{R}^m\to \mathbb{R}^m\]
	ასახვა არის $C^k$ კლასის (ან გლუვია, ანუ არის $C^{\infty}$ კლასის) თუ ფუნქციებს $g_1(x_1,\dots,x_m), \dots, g_m(x_1,\dots,x_m)$ გააჩნია $k$ რიგის (შესაბამისად ნებისმიერი რიგის) უწყვეტი კერძო წარმოებულები. ამრიგად გლუვი ატლასის გადასვლის ფუნქციები დიფეომორფიზმებია.
	
	ცხადია ერთი და იგივე მრავალნაირობას შეიძლება არაერთი ატლასი ჰქონდეს. ხშირად საჭიროა რუქების საკმარისი რაოდენობა, ამიტომ შემოდის მაქსიმალური ატლასის ცნება.
	
	\begin{definition}
		ვთქვათ მოცემული გვაქვს $(M,A)$ მრავალნაირობა გლუვი ატლასით. $A$-ატლასით განისაზღვრება $M$-ის მაქსიმალური ატლასი $\tilde{A}$:
		
		$\tilde{A}$ არის $M$-ზე $A$ ატლასთან თავსებადი რუქების კოლექცია $\big\{\psi:V\to V' \big\}$, ანუ $A$-ს ყველა რუქებისთვის $\big\{ \phi:U\to U' \big\}$ გადასვლის ფუნქციები $\big\{\phi \circ \psi^{-1} \big\}$ დიფეომორფიზმებია.
	\end{definition}
	
	$A$ ატლასთან თავსებადი რუქების კოლექცია ადგენს ექვივალენტობის კლასს. განსაზღვრებით ამ კლასის მაქსიმალური ატლასი არის $\tilde{A}$. ეს ნიშნავს, რომ მოცემულ კლასში არ არსებობს $\tilde{A}$-ს მკაცრად მომცველი ატლასი.
	
	\begin{definition}
		მოცემული ნატურალური $n,k \geq 1$ რიცხვებისთვის ან $k=\infty$, $n$-განზომილებიანი $C^k$ მრავალნაირობა არის ტოპოლოგიური მრავალნაირობა, $M$ და მასზე $C^k$ $n$-ატლასების $\bar{A}$ ექვივალენტობის კლასი.
		ნებისმიერ $A$ ატლასს $\bar{A}$ ექვივალენტობის კლასიდან ეწოდება დიფერენცირებადი სტრუქტურა $M$-ზე. როცა $k=\infty$ ვიტყვით, რომ $M$ არის გლუვი მრავალნაირობა.
	\end{definition}
	
	ხშირად $M$ სიმრავლე არაა თავიდანვე რაიმე კარგი ტოპოლოგიით აღჭურვილი. ამ შემთხვევებში მოსახერხებელია მრავალნაირობის შემდეგი განსაზღვრება.
	
	\begin{definition}
		\label{M-set}
		მოცემული ა $M$ სიმრავლე და ნატურალური $n,k \geq 1$ რიცხვები ან $k=\infty$. $M$-ზე $C^k$ $n$-ატლასი $A$ არის $\{(\phi_{\alpha},U_{\alpha})\}$ რუქების კოლექცია ისე, რომ
		
		(1) ყოველი $U_{\alpha}$ არის $M$-ის ქვესიმრავლე და $\phi_{\alpha}:U_{\alpha}\to \phi_{\alpha}(U_{\alpha})$ არის ბიექცია $\phi_{\alpha}(U_{\alpha})\subseteq \mathbb{R}^n$ ღია ქვესიმრავლეზე.
		
		(2) კოლექცია $\{U_{\alpha}\}$ ფარავს $M$-ს, ე.ი.,
		\[ M=\bigcup_{\alpha} U_{\alpha}; \]
		
		(3) როცა $U_{\alpha}\cap U_{\beta}\neq \emptyset$, გადასვლის ფუნქციები $\Psi_{\alpha \beta}$ არის $C^k$-დიფეომორფიზმები. როცა $k=\infty$ ატლასს ეწოდება გლუვი.
	\end{definition}
	
	\begin{example}
		ა) რა ტოპოლოგია ჩნდება $M$-სიმრავლეზე \eqref{M-set} განსაზღვრებით?
		
		ბ) დამატებით კიდევ რა უნდა მოვითხოვოთ რუქის სიმრავლეების ტერმინებში იმისათვის, რომ (როგორც მრავალნაირობის წინა განსაზღვრებაში) ტოპოლოგია $M$-ზე იყოს თვლადი ბაზისით (თვლადობის მეორე აქსიომა) და იყოს ჰაუსდორფის?
	\end{example}
	
	\bigskip
	
	მრავალნაირობის მაგალითებში უკეთ გასარკვევად სასარგებლოა ჯერ განვიხილოთ "არამაგალითები".
	\begin{example}
		$M$ იყოს ბრტყელი წირი, რომლის განტოლებაა
		\[ y^2=x^2-x^3, \]
		ანუ სიბრტყის ქვესიმრავლე
		\[ M=\big\{(x,y)\in \mathbb{R}^2 \mid y^2=x^2-x^3\big\}. \]
		
		ვაჩვენოთ, რომ $M$ ტოპოლოგიური მრავალნაირობა არ არის.
	\end{example}
	საქმე ისაა, რომ $M$-ს $(0, 0)=O$ სათავეში აქვს თვითგადაკვეთა. დავუშვათ რომ $M$ არის ტოპოლოგიური მრავალნაირობა. მაშინ არსებობს სათავის ბმული მიდამო, ანუ სათავის მომცველი ბმული ღია სიმრავლე
	$U \subseteq M$, (კერძოდ, ქვესივრცის ტოპოლოგიით ეს არის
	სიბრტყის საკმარისად პატარა ღია დისკის თანაკვეთა $M$-თან) და აგრეთვე არსებობს რუქა $\phi: U \to U'$, სადაც $U'$ არის $\mathbb{R}$-ის რაღაც ბმული ღია ქვესიმრავლე (ანუ ღია ინტერვალი). ეხლა შევნიშნოთ, რომ $U\setminus\{O\}$ შედგება ოთხი არაბმული კომპონენტისგან, ხოლო
	$U'\setminus\{\phi(O)\}$ შედგება ორი არაბმული კომპონენტისაგან. ეს ეწინააღმდეგება იმ ფაქტს, რომ $\phi$ არის ჰომეომორფიზმი. ე.ი. ჩვენი დაშვება მცდარია და მივიღეთ რ.დ.გ.
	
	\begin{figure}[h]
		\centering
		\begin{tikzpicture}
			\begin{axis}[
				axis x line=middle,
				axis y line=middle,
				axis equal,
				xlabel = {$x$},
				ylabel = {$y$},
				restrict y to domain = -2:2,
				restrict x to domain = -1:1.2,
				]
				\addplot [domain = 0:1, samples=200, smooth, blue] ({x}, {sqrt(x^2-x^3)});
				\addplot [domain = 0:1, samples=200, smooth, blue] ({x}, {-sqrt(x^2-x^3)});
			\end{axis}
		\end{tikzpicture}
		\caption{კვანძოვანი კუბიკა $M:y^2=x^2-x^3$}
	\end{figure}
	
	\begin{example}
		განვიხილოთ $y^2=x^3$ განტოლებით მოცემული წირი $M$. ამ წირის პარამეტრიზაციაა $r(t)=(t^2,t^3)$. $M$ როგორც აბსტრაქტული (და არა სიბრტყეში ჩადგმული) მრავალნაირობა შეიძლება განვიხილოთ ერთადერთი რუქით $\phi:M\to \mathbb{R}$, ისე რომ $\phi^{-1}(t)=(t^2,t^3)$. რაოდენ საკვირველიც არ უნდა იყოს $M$ ამ ატლასით არის გლუვი მრავალნაირობა, თუმცა არაა სიბრტყის რეგულარული წირი, რადგან სათავეში მხები ვექტორი $r'(t)=(2t,3t^2)$ ნულოვანია $t=0$-თვის.
	\end{example}
	
	\begin{figure}[h]
		\centering
		\begin{tikzpicture}
			\begin{axis}[
				axis x line=middle,
				axis y line=middle,
				axis equal,
				xlabel = {$x$},
				ylabel = {$y$},
				restrict y to domain = -2:2,
				restrict x to domain = -0.5:2,
				]
				\addplot [domain = 0:2, samples=200, smooth, red] ({x}, {x^1.5});
				\addplot [domain = 0:2, samples=200, smooth, red] ({x}, {-x^1.5});
			\end{axis}
		\end{tikzpicture}
		\caption{ბასრი კუბიკა $M: y^2=x^3$}
	\end{figure}
	
	\begin{example}
		ნამდვილი პროექციული სივრცე $\mathbb{RP}^n$.
	\end{example}
	
	პროექციული სივრცის ატლასის განსაზღვრისთვის მოსახერხებელია $\mathbb{RP}^n$ განვიხილოთ, როგორც $\mathbb{R}^{n+1}\setminus\{0\}$-ის ვექტორების ეკვივალენტობის კლასების სიმრავლე
	\[ u\sim v \iff v=\lambda u, \text{ რაღაც } \lambda\neq 0\in \mathbb{R}. \]
	ნებისმიერი მოცემული წერტილისთვის $p=[x_1,\dots,x_{n+1}]\in \mathbb{RP}^n$, კომპონენტებს $(x_1,\dots,x_{n+1})$ ეწოდება $p$ წერტილის ერთგვაროვანი კოორდინატები და აღინიშნება ასე $p=(x_1:\dots:x_{n+1})$.
	
	მეორეს მხრივ $\mathbb{RP}^n$ შეგვიძლია წარმოვიდგინოთ, როგორც  $S^n$ სფეროს ფაქტორ  სივრცე, სადაც ექვივალენტობის კლასები არის ანტიპოდალური $-x\sim x$ წერტილები.
	
	გავიხსენოთ ზოგადი ტოპოლოგიიდან, როგორ შემოდის ტოპოლოგია $X/\sim $ ფაქტორ სივრცეზე.
	
	ადვილი საჩვენებელია, რომ ჩვენს შემთხვევაში ეს ყველაფერი კარგადაა, ამიტომ ისღა დაგვრჩენია, რომ განვსაზღვროთ პროექციული სივრცის ატლასი.
	
	განვიხილოთ ყოველი $1\leq i\leq n+1$ და შესაბამისი რუქა $(U_i,\phi_i)$ იყოს
	\begin{align*}
		&U_i=\{(x_1:\dots:x_{n+1})\in \mathbb{RP}^n \mid x_i\neq 0 \},\\
		&\phi_i:(x_1:\dots:x_{n+1})\to \left(\frac{x_1}{x_i}, \dots, \frac{x_{i-1}}{x_i}, \frac{x_{i+1}}{x_i}, \dots, \frac{x_{n+1}}{x_i}\right).
	\end{align*}
	აქ მარჯვნივ $i$-ური კომპონენტი გამოტოვებულია.
	
	\subsection{მხები ვექტორები, მხები სივრცე და კომხები სივრცე}
	
	ვთქვათ $M$ არის $n$ განზომილების $C^k$ მრავალნაირობა, $k \geq 1$. ჩვენი მიზანია განვსაზღვროთ $M$ მრავალნაირობის $p$ წერტილში მხები სივრცე, $T_p(M)$ და მისი ორადული კომხები სივრცე, $T_p^*(M)$.
	
	\begin{definition}
		\label{eqviv}
		ვთქვათ $M$ არის $n$ განზომილებიანი $C^k$ მრავალნაირობა. ყოველი $p\in M$ წერტილისთვის ამ წერტილზე გამავალ ორ $C^1$ წირს
		\[ \gamma_1: (-\epsilon_1,\epsilon_1) \to M \text{ და } \gamma_2: (-\epsilon_2, \epsilon_2) \to M, \quad \gamma_1(0)=\gamma_2(0)=p, \]
		ეწოდება \textbf{ექვივალენტური}, თუ არსებობს $p$ წერტილის მომცველი რუქა $(U,\phi)$, რომ
		\[ (\phi \circ\gamma_1)'(0)=(\phi \circ\gamma_2)'(0). \]
	\end{definition}
	
	\bigskip
	
	საჩვენებელია, რომ ეს განსაზღვრება არაა დამოკიდებული რუქის არჩევაზე.
	
	\begin{definition}[მხები ვექტორი. ვერსია 1]
		ვთქვათ $M$ არის $n$ განზომილებიანი $C^k$ მრავალნაირობა. ნებისმიერ $p\in M$ წერტილში მხები ვექტორი ეწოდება ამ წერტილზე გამავალი $C^1$ წირების ექვივალენტობის კლასს \ref{eqviv} განსაზღვრების მიხედვით. $p\in M$ წერტილში ყველა მხები ვექტორების სიმრავლეს ეწოდება მხები სივრცე და აღინიშნება $T_pM$-ით.
	\end{definition}
	
	\begin{definition}[მხები ვექტორი. ვერსია 2]
		ვთქვათ $M$ არის $n$ განზომილებიანი $C^k$ მრავალნაირობა, $k\geq 1$ და $(U,\phi)$ არის $p\in M$ წერტილის მომცველი რუქა. $p$ წერტილში $M$-ის მხები ვექტორი არის $p$ წერტილში-გაწარმოება გერმების $\mathcal{O}^{(k)}_{M,p}$ რგოლზე.
		
		დაწვრილებით: მხები ვექტორი არის ისეთი $v$ წრფივი ასახვა გერმების რგოლზე
		\[ v:\mathcal{O}^{(k)}_{M,p}\to \mathbb{R}, \quad v:[f]=v(f) \]
		რომელიც აკმაყოფილებს ლაიბნიცის წესს
		\[ v(fg)=v(f)g(p)+f(p)v(g). \]
	\end{definition}
	
	
	
	\bigskip
	
	
	
	\bigskip
	
	
	
	
	
		
	\section{ლის ჯგუფები და ლის ალგებრები}
		
		
		
		
		
	
		
		ლის ჯგუფებისა და ლის ალგებრების თეორია XIX საუკუნის მეორე ნახევარში ჩაისახა და დღესდღეობით თანამედროვე მათემატიკისა და თეორიული ფიზიკის ერთ-ერთ ქვაკუთხედს წარმოადგენს. ამ თეორიის შექმნა ნორვეგიელი მათემატიკოსის, \textbf{სოფუს ლის (Sophus Lie, 1842–1899)} სახელს უკავშირდება.
		
		ლის ფუნდამენტური იდეა იყო, შეექმნა თეორია \textbf{დიფერენციალური განტოლებებისთვის}, რომელიც იქნებოდა იმის ანალოგი, რაც \textbf{ევარისტ გალუამ (Évariste Galois)} შექმნა ალგებრული განტოლებებისთვის. გალუამ გამოიყენა \textbf{სასრული ჯგუფები} (რომლებიც აღწერენ დისკრეტულ სიმეტრიებს) პოლინომიური განტოლებების ამოსახსნელობის საკითხის შესასწავლად. ანალოგიურად, ლის სურდა გამოეყენებინა \textbf{უწყვეტი გარდაქმნების ჯგუფები} დიფერენციალური განტოლებების სიმეტრიების გასაანალიზებლად და მათი ამოხსნის მეთოდების მოსაძებნად. ეს უწყვეტი ჯგუფები, რომლებიც თავად გლუვ მრავალნაირობებს წარმოადგენენ, დღეს ლის ჯგუფების სახელითაა ცნობილი.
		
		ლის გენიალურმა ინსტრუმენტმა — \textbf{ინფინიტეზიმალურმა (უსასრულოდ მცირე) გარდაქმნებმა} — ის მიიყვანა ობიექტამდე, რომელსაც დღეს \textbf{ლის ალგებრას} ვუწოდებთ. ლის ალგებრა არის ვექტორული სივრცე, რომელიც აღწერს ლის ჯგუფის ლოკალურ სტრუქტურას ერთეულოვან ელემენტთან ახლოს. აღმოჩნდა, რომ ლის ჯგუფის რთული, არაწრფივი სტრუქტურის შესახებ უამრავი ინფორმაცია მის გაცილებით მარტივ, წრფივ ანალოგში — ლის ალგებრაშია კოდირებული.
		
		თეორიის განვითარებაში უმნიშვნელოვანესი როლი შეასრულეს სხვა მათემატიკოსებმაც:
		
		\begin{itemize}
			\item \textbf{ფელიქს კლაინმა (Felix Klein)} თავის ცნობილ \textit{ერლანგენის პროგრამაში} (1872) გეომეტრია განსაზღვრა, როგორც გარკვეული გარდაქმნების ჯგუფის მიმართ ინვარიანტული თვისებების შესწავლა. ამ იდეამ ლის ჯგუფების როლი გეომეტრიაში კიდევ უფრო მნიშვნელოვანი გახადა.
			
			\item \textbf{ვილჰელმ კილინგმა (Wilhelm Killing)}, რომელიც ლისგან დამოუკიდებლად მუშაობდა, 1888–1890 წლებში შეასრულა ფუნდამენტური სამუშაო \textbf{მარტივი ლის ალგებრების კლასიფიკაციაზე}. მიუხედავად იმისა, რომ მის ნაშრომს თავდაპირველად სათანადო აღიარება არ მოჰყოლია (თვით სოფუს ლისაც კი სკეპტიკურად უყურებდა), სწორედ კილინგის იდეები დაედო საფუძვლად ლის ალგებრების სტრუქტურულ თეორიას.
			
			\item XX საუკუნის დასაწყისში ფრანგმა მათემატიკოსმა \textbf{ელი კარტანმა (Élie Cartan)} განაზოგადა, დახვეწა და მკაცრად დაასაბუთა კილინგისა და ლის შედეგები. კარტანმა საბოლოო სახე მისცა მარტივი ლის ალგებრების კლასიფიკაციას, განავითარა წარმოდგენათა თეორია და დაადგინა ღრმა კავშირი ლის ჯგუფებსა და დიფერენციალურ გეომეტრიას შორის, განსაკუთრებით სიმეტრიული სივრცეების თეორიაში. სწორედ კარტანს ეკუთვნის ცნობილი თეორემა, რომლის თანახმადაც ზოგადი წრფივი ჯგუფის $GL(n, \mathbb{R})$ ნებისმიერი ჩაკეტილი ქვეჯგუფი თავად ლის ჯგუფია, რაც მატრიცული ჯგუფების შესწავლას განსაკუთრებით ამარტივებს.
			
			\item \textbf{ჰერმან ვაილმა (Hermann Weyl)} გადამწყვეტი როლი ითამაშა ლის ჯგუფების თეორიის ფიზიკაში, კერძოდ, \textbf{კვანტურ მექანიკასა} და \textbf{ფარდობითობის თეორიაში} დანერგვაში. მან აჩვენა, რომ ფიზიკური სისტემების სიმეტრიები (მაგალითად, ბრუნვები, ლორენცის გარდაქმნები) ბუნებრივად აღიწერება ლის ჯგუფებით, ხოლო ფიზიკური სიდიდეები — მათი ალგებრების წარმოდგენებით.
		\end{itemize}
		
		დღეს ლის თეორია მათემატიკის თითქმის ყველა დარგში გამოიყენება — ალგებრული ტოპოლოგიიდან და რიცხვთა თეორიიდან დაწყებული, რობოტოტექნიკითა და კომპიუტერული ხედვით დამთავრებული. განსაკუთრებით დიდი როლი მას \textbf{ელემენტარული ნაწილაკების ფიზიკაში} აკისრია, სადაც ბუნების ფუნდამენტური ურთიერთქმედებები აღიწერება კალიბრაციული ლის ჯგუფებით, როგორიცაა $U(1)$, $SU(2)$ და $SU(3)$.
		
	
	

		
	
		
		აქ ჩვენი მიზანია შესავალი ისეთ ობიექტებზე, როგორიცაა მრავალნაირობები,
		ლის ჯგუფები და ლის ალგებრები, რაც ჩვენი ძირითადი შესწავლის ობიექტებია.
		ამ მშვენიერ სამყაროში შესაღწევად კარგი გზაა
		თამაში ცხად მატრიცულ ჯგუფებთან, როგორიცაა $\mathbb{R}^2$-ში (ან $\mathbb{R}^3$-ში) ბრუნვების ჯგუფი
		და ექსპონენციალური ასახვა. რეალურად, ანტისიმეტრიული $2\times 2$ $B$ მატრიცის ექსპონენციალის $A = e^B$ გამოთვლის შემდეგ
		და დაკვირვებით, რომ ეს არის ბრუნვის
		მატრიცა, და ანალოგიურად $3\times 3$ ანტისიმეტრიული $B$ მატრიცისთვის, ადამიანი იწყებს ფიქრს, რომ რაღაც ღრმა და მნიშვნელოვანი ხდება. ანალოგიურად, მას შემდეგ რაც აღმოჩნდა, რომ ყველა ნამდვილი შებრუნებადი $n\times n$
		მატრიცა $A$ შეიძლება ჩაიწეროს როგორც $A = RP$, სადაც $R$ არის ორთოგონალური მატრიცა, ხოლო $P$ არის დადებითად განსაზღვრული სიმეტრიული მატრიცა, და რომ $P$ შეიძლება ჩაიწეროს როგორც $P = e^S$, რაღაც $S$ სიმეტრიული მატრიცისთვის,
		ადამიანი იწყებს ექსპონენციალური ასახვის დაფასებას.
		
		იმის ნაცვლად, რომ განვსაზღვროთ ლის ჯგუფები სრული ზოგადობით, ჩვენ ჯერ განვიხილავთ წრფივ ლის ჯგუფებს, რადგან
		კარტანის ცნობილი შედეგით, $GL(n,\mathbb{R})$ ზოგადი წრფივი ჯგუფის ჩაკეტილი ქვეჯგუფი არის მრავალნაირობა და შესაბამისად, ლის ჯგუფი. ამ გზით, ლის ალგებრა
		"გამოითვლება", ლის ჯგუფის ნეიტრალურ ელემენტში (ერთეულში) მხები ვექტორების გამოყენებით $t\to A(t)$ ფორმის წირებზე, სადაც $A(t)$ არის მატრიცა.
		
		მოცემული $n\times n$ ნამდვილი ან კომპლექსური მატრიცისათვის $A=(a_{ij})$ ჩვენ გვინდა განვსაზღვროთ ექსპონენტა $e^A$ როგორც მწკრივის ჯამი
		\[
		e^A=I_n+\sum_{p\geq 1}\frac{A^p}{p!}=\sum_{p\geq 0}\frac{A^p}{p!},
		\]
		თუ მივიღებთ, რომ $A^0=I_n$ ერთეულოვანი მატრიცაა.
		
		აქ საჩვენებელია, რომ განსაზღვრება კორექტულია, ანუ მარჯვენა მხარეში მწკრივი კრებადია.
		ეს გამომდინარეობს შემდეგი ლემიდან (რომელიც ინდუქციით მტკიცდება. ვცადოთ დამოუკიდებლად) და იმ ფაქტიდან რომ ბანახის სივრცეში მწკრივის აბსოლუტური კრებადობა იწვევს კრებადობას.
		
		\begin{lemma}
			ვთქვათ $A=(a_{ij})$ ნამდვილი ან კომპლექსური $n\times n$ მატრიცაა და
			\[
			\mu=\max \{|a_{ij}| \mid 1\leq i,j\leq n \}.
			\]
			თუ $A^p=(a^{(p)}_{ij})$, მაშინ $|a^{(p)}_{ij}|\leq (n\mu)^p$ ნებისმიერი $1\leq i,j\leq n$.
			შედეგად $n^2$ რაოდენობის მწკრივი
			\[
			\sum_{p\geq 0}\frac{a^{(p)}_{ij}}{p!}
			\]
			აბსოლუტურად კრებადია და
			\[
			\sum_{p\geq 0}\frac{A^p}{p!}
			\]
			მწკრივი კარგადაა განსაზღვრული.
		\end{lemma}
		
		\begin{example}
			ვიპოვოთ ანტისიმეტრიული
			\[
			A=\begin{pmatrix}
				0 & -\theta\\
				\theta & 0
			\end{pmatrix}
			\]
			მატრიცის ექსპონენტა.
		\end{example}
		
		\medskip
		
		თუ გავითვალისწინებთ
		\[
		\begin{pmatrix}
			0 & -\theta\\
			\theta & 0
		\end{pmatrix}
		= \theta
		\begin{pmatrix}
			0 & -1\\
			1 & 0
		\end{pmatrix}
		\text{ და }
		\begin{pmatrix}
			0 & -\theta\\
			\theta & 0
		\end{pmatrix}^2
		= -\theta^2
		\begin{pmatrix}
			1 & 0\\
			0 & 1
		\end{pmatrix},
		\]
		და ავღნიშნავთ
		\[
		J=\begin{pmatrix}
			0 & -1\\
			1 & 0
		\end{pmatrix},
		\]
		ინდუქციით გვექნება
		\begin{align*}
			&A^{4n}=\theta^{4n}I_2,\\
			&A^{4n+1}=\theta^{4n+1}J,\\	
			&A^{4n+2}=-\theta^{4n+2}I_2,\\
			&A^{4n+3}=-\theta^{4n+3}J,
		\end{align*}
		ამიტომ
		\[
		e^A=I_2+\frac{\theta}{1!}J-\frac{\theta^2}{2!}I_2-\frac{\theta^3}{3!}J+
		\frac{\theta^4}{4!}I_2+\frac{\theta^5}{5!}J-\frac{\theta^6}{6!}I_2-
		\frac{\theta^7}{7!}J+\cdots.
		\]
		
		შეგვიძლია გადავალაგოთ (რადგან ამით ჯამი არ შეიცვლება ლემის გამო) და გვექნება
		\begin{align*}
			e^A &= \left(1-\frac{\theta^2}{2!}+\frac{\theta^4}{4!}-\frac{\theta^6}{6!}+\cdots\right)I_2+
			\left(\frac{\theta}{1!}-\frac{\theta^3}{3!}+\frac{\theta^5}{5!}-
			\frac{\theta^7}{7!}+\cdots\right)J.
		\end{align*}
		
		აქ ადვილი ამოსაცნობია $\cos$ და $\sin$ მწკრივები, ამიტომ
		\[
		e^A=\cos\theta I_2+\sin\theta J,
		\]
		ანუ
		\[
		e^A=\begin{pmatrix}
			\cos\theta & -\sin\theta\\
			\sin\theta & \cos\theta
		\end{pmatrix}.
		\]
		
		ამრიგად გვაქვს $\theta$ კუთხით მობრუნების მატრიცა! ეს ზოგადი ფაქტია: თუ $A$ არის ანტისიმეტრიული მატრიცა, მაშინ $e^A$ არის ორთოგონალური მატრიცა დეტერმინანტით 1, ანუ მობრუნების მატრიცა; პირიქითაც, ნებისმიერ მობრუნების მატრიცას აქვს ასეთი ფორმა.
		
		\subsection{ლის ჯგუფები $GL(n,\mathbb{R})$, $SL(n,\mathbb{R})$, $O(n)$, $SO(n)$,\\ ლის ალგებრები $\mathfrak{gl}(n,\mathbb{R})$, $\mathfrak{sl}(n, \mathbb{R})$, $\mathfrak{o}(n)$, $\mathfrak{so}(n)$\\ და ექსპონენციალური ასახვა}
		
		პირველ რიგში, გავიხსენოთ რამდენიმე ძირითადი ფაქტი და განსაზღვრება. ნამდვილი შებრუნებადი $n \times n$ მატრიცების სიმრავლე
		ქმნის ჯგუფს გამრავლების მიმართ, რომელიც აღინიშნება $GL(n,\mathbb{R})$-ით. მისი ქვესიმრავლე, დეტერმინანტით
		+1 არის $GL(n,\mathbb{R})$-ის ქვეჯგუფი, რომელიც აღინიშნება $SL(n,\mathbb{R})$-ით.
		ასევე ადვილია იმის შემოწმება, რომ ნამდვილი $n \times n$ ორთოგონალური მატრიცების სიმრავლე ქმნის ჯგუფს
		გამრავლების მიმართ, და აღინიშნება $O(n)$-ით. $O(n)$-ის ქვესიმრავლე, რომელიც შედგება იმ მატრიცებისგან, რომლებსაც აქვთ
		დეტერმინანტი +1, არის ქვეჯგუფი, რომელიც აღინიშნება $SO(n)$-ით. ჩვენ $SO(n)$ მატრიცებს
		ასევე ვუწოდებთ ბრუნვის მატრიცებს. მარტივად შეგვიძლია შევამოწმოთ, რომ ნამდვილი $n \times n$
		ნულოვანი კვალის (დიაგონალის ელემენტების ჯამი) მქონე მატრიცები ქმნიან ვექტორულ სივრცეს შეკრების მიმართ. ანალოგიურად ანტისიმეტრიული მატრიცების სიმრავლისთვის.
		
		\begin{definition}
			ჯგუფს $GL(n,\mathbb{R})$ ეწოდება ზოგადი წრფივი ჯგუფი და მის ქვეჯგუფს
			$SL(n,\mathbb{R})$ ეწოდება სპეციალური წრფივი ჯგუფი. ორთოგონალური მატრიცების $O(n)$ ჯგუფს ეწოდება
			ორთოგონალური ჯგუფი და მის ქვეჯგუფს $SO(n)$ ეწოდება სპეციალური ორთოგონალური ჯგუფი (ან
			ბრუნვების ჯგუფი). ნულოვანი კვალის მქონე ნამდვილი $n\times n$ მატრიცების ვექტორული სივრცე აღინიშნება $\mathfrak{sl}(n,\mathbb{R})$,
			ხოლო ნამდვილი $n \times n$ ანტისიმეტრიული მატრიცების ვექტორული სივრცე აღინიშნება $\mathfrak{so}(n)$-ით.
		\end{definition}
		
		\textbf{შენიშვნა.} აღნიშვნა $\mathfrak{sl}(n,\mathbb{R})$ და $\mathfrak{so}(n)$ საკმაოდ უცნაურია და იმსახურებს გარკვეულ განმარტებას.
		ჯგუფები $GL(n,\mathbb{R})$, $SL(n,\mathbb{R})$, $O(n)$ და $SO(n)$ უფრო მეტია, ვიდრე უბრალოდ ჯგუფები. ისინიც არიან
		ტოპოლოგიური ჯგუფები, რაც ნიშნავს, რომ ისინი ტოპოლოგიური სივრცეებია (განიხილება, როგორც $\mathbb{R}^{n^2}$-ის ქვესივრცეები) და რომ გამრავლება და შებრუნებული ოპერაციები უწყვეტია (ფაქტობრივად, გლუვი). გარდა ამისა, ისინი გლუვი ნამდვილი მრავალნაირობებია.
		
		ასეთ ობიექტებს \textbf{ლის ჯგუფები} ეწოდება. ნამდვილი
		ვექტორული სივრცეები $\mathfrak{sl}(n,\mathbb{R})$ და $\mathfrak{so}(n)$ არის ის, რასაც \textbf{ლის ალგებრას} უწოდებენ.
		ალგებრის სტრუქტურა მოცემულია ე.წ.
		\textbf{ლის ფრჩხილებით}, რომელიც განისაზღვრება როგორც
		\[
		[A,B]=AB-BA.
		\]
		
		ლის ალგებრა ასოცირდება ლის ჯგუფებთან: ლის ალგებრა არის
		ლის ჯგუფის მხები სივრცე ერთეულოვან ელემენტში, ანუ ყველა მხები ვექტორების სივრცე
		$I_n$-ში (ამ შემთხვევაში). გარკვეული გაგებით, ლის ალგებრა აღწევს ლის ჯგუფის „ლინეარიზაციას“. ექსპონენციალური ასახვა არის ასახვა ლის ალგებრიდან ლის ჯგუფში, მაგალითად,
		\[
		\exp:\mathfrak{so}(n)\to SO(n),
		\]
		ან
		\[
		\exp:\mathfrak{sl}(n,\mathbb{R})\to SL(n,\mathbb{R}).
		\]
		
		ექსპონენციალური ასახვა ხშირად იძლევა ლის ჯგუფის ელემენტების უფრო მარტივი
		ობიექტებით, ლის ალგებრის ელემენტებით, პარამეტრიზაციის საშუალებას.
		
		\subsection{ლის ჯგუფების განსაზღვრება}
		
		\begin{definition}
			ტოპოლოგიური ჯგუფი $G$, არის ტოპოლოგიური სივრცე, რომელიც ასევე არის ჯგუფი, ისე რომ ჯგუფის ოპერაცია
			\[ \cdot : G \times G \to G,\quad (x,y) \to x \cdot y \]
			და აგრეთვე შებრუნებული ელემენტის აღება
			\[ \text{inv}: G \to G, \quad x \to x^{-1},\]
			როგორც ასახვები, უწყვეტია. აქ $G \times G$ განიხილება, როგორც ტოპოლოგიური სივრცე ნამრავლის ტოპოლოგიით. ამბობენ, რომ ასეთი ტოპოლოგია თავსებადია ჯგუფურ ოპერაციებთან და მას ჯგუფურ ტოპოლოგიას უწოდებენ.
			
			\textbf{ლის ჯგუფი} ეწოდება ისეთ ტოპოლოგიურ ჯგუფს, რომელიც არის დიფერენცირებადი მრავალნაირობა და ჯგუფური ოპერაციები $\cdot$ და $\text{inv}$ დიფერენცირებადი ასახვებია.
		\end{definition}
		
		\subsection{ლის ჯგუფების მოქმედება ტოპოლოგიურ სივრცეებზე}
		
		$G$ ტოპოლოგიური ჯგუფის მოქმედება (მარჯვნიდან "მიწერა") $X$ სივრცეზე არის ასახვა
		\[\mu: X \times G \to X,\quad (x,g)\to xg\]
		ისე, რომ სრულდება ორი პირობა:
		\begin{enumerate}
			\item $x \cdot 1= x, \quad \forall x\in X,$
			\item $x(g_1g_2) = (xg_1)g_2, \quad \forall x\in X, \, \forall g_1, g_2 \in G.$
		\end{enumerate}
		
		ყურადღება მიაქციეთ, რომ ასეთი მოქმედების გათვალისწინებით, ყოველი ელემენტი $g$ მოქმედებს როგორც ჰომეომორფიზმი, ვინაიდან მას აქვს შებრუნებული ასახვა, კერძოდ $g^{-1}$ ელემენტით მოქმედება. ამრიგად, ჯგუფის მოქმედება $\mu$ განსაზღვრავს ასახვას
		\[\mu : G\to \mathrm{Homeo}(X),\]
		სადაც $\mathrm{Homeo}(X)$ აღნიშნავს $X$-ის ჰომეომორფიზმების ჯგუფს. მოთხოვნა, რომ $\mu$ ასახვა იყოს ჯგუფური ჰომომორფიზმი, ექვივალენტურია ზემოთ ჩამოთვლილი ორი (1), (2) პირობის.
		
		კერძო შემთხვევაში, როცა $G$ ლის ჯგუფია და $X=M$ გლუვი მრავალნაირობა, ვიტყვით, რომ $G$ მარჯვნიდან გლუვად მოქმედებს $M$-ზე, თუ $\mu$ არის ჯგუფური ჰომომორფიზმი
		\[\mu : G\to \mathrm{Diffeo}(M),\]
		სადაც $\mathrm{Diffeo}(M)$ აღნიშნავს $M$-ის დიფეომორფიზმების ჯგუფს.
		
		$X$ იყოს სივრცე მარჯვენა $G$-მოქმედებით. მოცემული $x\in X$ წერტილისთვის სიმრავლეს
		\[xG =\{xg \mid g\in G\} \subset X\]
		ეწოდება $x$-ის $G$-ორბიტა. $x$-ის იზოტროპიის ქვეჯგუფი, $\mathrm{Iso}(x)\subset G$, განისაზღვრება
		როგორც იმ ელემენტების სიმრავლე, რომლებიც $x$-ზე ტრივიალურად მოქმედებენ, ანუ $x$-ს უძრავად ტოვებენ:
		\[\mathrm{Iso}(x) = \{g\in G \mid xg = x\}.\]
		
		გაითვალისწინეთ, რომ ასახვა
		$G\to xG$,
		რომელიც $g$-ს ასახავს $xg$-ში, განსაზღვრავს ჰომეომორფიზმს მოსაზღვრე კლასების სივრციდან ორბიტაში
		\[X/\mathrm{Iso}(x)\xrightarrow{\cong}xG \subset X.\]
		
		მოქმედება $X$ სივრცეზე ეწოდება ტრანზიტული, თუ სივრცე $X$ არის ერთი წერტილის ორბიტა,
		$X = xG$. გავითვალისწინოთ, რომ თუ $X = x_0G$ რაიმე $x_0\in X$-ისთვის, მაშინ $X = xG$ ნებისმიერი $x\in X$-ისთვის.
		
		ტრანზიტულობის პირობა უდრის იმის თქმას, რომ ნებისმიერი ორი წერტილისთვის $x_1, x_2 \in X$ არსებობს
		ელემენტი $g \in G$ ისეთი, რომ $x_1 = x_2g$. ბოლოს შევნიშნოთ, რომ თუ $G$-ტრანზიტულად მოქმედებს $X$-ზე, მაშინ
		ზემოხსენებული დისკუსია იზოტროპიის ქვეჯგუფების შესახებ გულისხმობს, რომ არსებობს ქვეჯგუფი $H <G$ და ჰომეომორფიზმი
		\[G/H\xrightarrow{\cong} X.\]
		
		რა თქმა უნდა, თუ $X$ არის გლუვი, $G$ არის ლის ჯგუფი და მოქმედება გლუვია, მაშინ ზემოთ მოცემული ასახვა იქნება დიფეომორფიზმი.
		
		ჯგუფურ მოქმედებას ეწოდება (ფიქსირებული წერტილებისგან) თავისუფალი, თუ ყოველი $x$ წერტილის იზოტროპიული ჯგუფები ტრივიალურია,
		\[\mathrm{Iso}(x) = \{1\},\quad \forall x\in X.\]
		
		სხვაგვარად რომ ვთქვათ, მოქმედება თავისუფალია $\iff$ თუ
		განტოლება
		$xg = x$ სრულდება მხოლოდ როცა
		$g = 1 \in G$.
		
		თუ $g \in G$ ელემენტისთვის, ფიქსირებული წერტილების სიმრავლეს ავღნიშნავთ ასე
		\[\mathrm{Fix}(g) = \{x \in X \mid xg = x \},\]
		მაშინ მოქმედება თავისუფალია
		$\iff$ $\forall g\neq 1 \in G$ გვაქვს $\mathrm{Fix}(g) = \emptyset$.
		
		\section{ვექტორული   ფიბრაციები}
		
		
	ვექტორული ფიბრაცია თანამედროვე გეომეტრიისა და ტოპოლოგიის ერთ-ერთი ყველაზე ფუნდამენტური კონცეფციაა. ის საშუალებას გვაძლევს, განვაზოგადოთ ერთი ვექტორული სივრცის ცნება და ვისაუბროთ ვექტორული სივრცეების "უწყვეტ ოჯახზე", რომელიც პარამეტრიზებულია სხვა, როგორც წესი, მრუდე სივრცით.
			
	ინტუიციურად, ვექტორული ფიბრაცია შეგვიძლია წარმოვიდგინოთ, როგორც საბაზო სივრცის (მაგალითად, სფეროს ან წრეწირის) თითოეულ წერტილზე "დამაგრებული" ვექტორული სივრცე, რომელსაც "ფიბრი" ეწოდება. ამის კლასიკური მაგალითია \textbf{მხები ფიბრაცია} ($TM$). თუ ავიღებთ სფეროს, როგორც საბაზო სივრცეს, მაშინ მის თითოეულ წერტილზე შეგვიძლია განვიხილოთ ამ წერტილში გავლებული მხები სიბრტყე. ყველა ამ მხები სიბრტყის ერთობლიობა, რომლებიც "გლუვად" იცვლება წერტილიდან წერტილში, ქმნის სფეროს მხებ ფიბრაციას.
			
	ფორმალურად, ვექტორული ფიბრაცია შედგება:
	\begin{itemize}
				\item \textbf{ტოტალური სივრცისგან ($E$)} – ეს არის ყველა ფიბრის გაერთიანება.
				\item \textbf{საბაზო სივრცისგან ($B$)} – სივრცე, რომლის წერტილებზეც ფიბრებია "დამაგრებული".
				\item \textbf{პროექციის ასახვისგან ($p: E \to B$)} – უწყვეტი ასახვა, რომელიც $E$-ს თითოეულ წერტილს (რომელიც რომელიმე ფიბრს ეკუთვნის) შეუსაბამებს $B$-ს იმ წერტილს, რომელზეც ეს ფიბრია დამაგრებული.
			\end{itemize}
			
			ყველაზე მნიშვნელოვანი პირობა, რომელიც ვექტორულ ფიბრაციას განასხვავებს ფიბრების უბრალო ერთობლიობისგან, არის \textbf{ლოკალური ტრივიალურობა}. ეს ნიშნავს, რომ მიუხედავად იმისა, რომ გლობალურად ფიბრაცია შეიძლება "გრეხილი" იყოს, საბაზო სივრცის ნებისმიერი საკმარისად პატარა $U$ არისთვის, მისი თავზე არსებული ფიბრაციის ნაწილი $p^{-1}(U)$ გამოიყურება, როგორც ჩვეულებრივი დეკარტული ნამრავლი $U \times V$, სადაც $V$ სტანდარტული ვექტორული სივრცეა.
			
			ამის საუკეთესო მაგალითებია \textbf{ცილინდრი} და \textbf{მობიუსის ლენტი}. ორივე მათგანის საბაზო სივრცე არის წრეწირი, ხოლო ფიბრები – მონაკვეთები (ერთგანზომილებიანი ვექტორული სივრცეები). ცილინდრი \textbf{ტრივიალური} ფიბრაციაა, რადგან ის გლობალურადაც კი $S^1 \times I$ ნამრავლის იზომორფულია. მობიუსის ლენტი კი \textbf{არატრივიალურია} – მიუხედავად იმისა, რომ მისი ნებისმიერი პატარა ნაწილი ჩვეულებრივ ბრტყელ ოთხკუთხედს ჰგავს, გლობალურად მას აქვს "გრეხა", რომლის გასწორებაც შეუძლებელია.
			
			ვექტორული ფიბრაციები ქმნიან იმ ჩარჩოს, რომელიც საშუალებას გვაძლევს, გამოვიყენოთ წრფივი ალგებრის მეთოდები მრავალნაირობებზე, შევისწავლოთ დიფერენციალური განტოლებები მრუდე სივრცეებზე და განვსაზღვროთ ისეთი ცნებები, როგორიცაა კოვარიანტული წარმოებული. მათ ფუნდამენტური როლი აკისრიათ არა მხოლოდ სუფთა მათემატიკაში, არამედ თეორიულ ფიზიკაშიც, განსაკუთრებით კალიბრაციულ თეორიებსა და ზოგად ფარდობითობაში.
			
		
		
		\begin{definition}
			ვთქვათ $B$ ბმული ტოპოლოგიური სივრცეა ფიქსირებული $b\in B$ წერტილით.
			\[p:E\to B\]
			ასახვას ეწოდება \textbf{ლოკალურად ტრივიალური ფიბრაცია} $F$ ფიბრით, თუ სრულდება პირობები:
			\begin{enumerate}
				\item $p^{-1}(b)=F$;
				\item $p:E\to B$ არის ზე-ასახვა (სიურექცია);
				\item $\forall x\in B$ არსებობს მიდამო $U_x$ და ფიბრის შემნახავი ჰომეომორფიზმი
				\[\Psi_{U_x}:p^{-1}(U_x)\to U_x \times F,\]
				ანუ შემდეგი დიაგრამა კომუტაციურია:
			\end{enumerate}
			\[
			\begin{tikzcd}
				p^{-1}(U_x) \arrow[r, "\Psi_{U_x}"] \arrow[d, "p"]
				& U_x \times F \arrow[d, "proj" ] \\
				U_x \arrow[r, "="]
				& U_x.
			\end{tikzcd}
			\]
		\end{definition}
		
		\bigskip
		
		\textbf{ტრივიალური ფიბრაცია}
			ეს არის დეკარტული ნამრავლის პროექცია პირველ მამრავლზე  $p:B\times F \to B$.
		
		
		\textbf{დაფარვა}
			ვთქვათ $S^1\subset \mathbb{C}$ ერთეულოვანი წრეწირია ფიქსირებული წერტილით $1\in S^1$. განვიხილოთ $f_n : S^1\to S^1$ ასახვა, ფორმულით $f_n(z)=z^n$. მაშინ $f_n$ არის ლოკალურად ტრივიალური ფიბრაცია, რომლის ფიბრია $n$-ელემენტიანი სიმრავლე, კერძოდ $n$ ხარისხის ფესვები 1-დან. სასრულფიბრიან ლოკალურად ტრივიალურ ფიბრაციას დაფარვა ეწოდება.
		
		
		მოცემული $p: E\to B$ ფიბრაციისთვის $F$ ფიბრით, ხშირად $B$-ს ეწოდება \textbf{ფიბრაციის საბაზო სივრცე} (ბაზა), $E$-ს კი \textbf{ფიბრაციის ტოტალური სივრცე}. ეს ინფორმაცია ავღნიშნოთ ოთხეულით $(p,F,E,B)$.
		
		\begin{definition} 
			ფიბრაციების ასახვა $\Phi:(p_1,F_1,E_1,B_1)\to (p_2,F_2,E_2,B_2)$
			ანუ "მორფიზმი" არის ასახვათა წყვილი
			\[
			\bar{\phi}:E_1\to E_2,\quad \phi:B_1\to B_2,
			\]
			ისე რომ კომუტაციურია დიაგრამა
			\[
			\begin{tikzcd}
				E_1 \arrow[r, "\bar{\phi}"] \arrow[d, "p_1"]
				& E_2 \arrow[d, "p_2" ] \\
				B_1 \arrow[r, "\phi"]
				& B_2.
			\end{tikzcd}
			\]
			\medskip
			
			ამასთან მოითხოვება ფიბრაციების ბაზების ფიქსირებული წერტილების შენახვა, ანუ თუ $b_1\in B_1$, $b_2\in B_2$ ფიქსირებული წერტილებია, მაშინ $\phi(b_1)=b_2$.
			
			შევნიშნოთ რომ ფიბრაციების ასახვა განსაზღვრავს ფიბრების უწყვეტ ასახვას $\phi_0:F_1\to F_2$.
			
			ფიბრაციების $\Phi $ მორფიზმს ეწოდება \textbf{იზომორფიზმი}, თუ არსებობს შექცეული მორფიზმი
			\[\Phi^{-1}:(p_2,F_2,E_2,B_2)\to(p_1,F_1,E_1,B_1),\]
			ისე რომ
			\[\Phi \circ\Phi^{-1}=\Phi^{-1}\circ\Phi=1.\]
		\end{definition}
		
\subsection*{ვექტორული ფიბრაციები}
		
		\begin{definition}
			$n$ განზომილებიანი ვექტორული ფიბრაცია $k$ ველის მიმართ არის ლოკალურად ტრივიალური ფიბრაცია
			\[p:E\to B\]
			რომლის ფიბრი არის $n$-განზომილებიანი $k$-ვექტორული სივრცე $V$, რომელიც აკმაყოფილებს დამატებით პირობას, რომ ლოკალური ტრივიალიზაცია
			\[ \psi:p^{-1}(U)\to U\times V \]
			ინდუცირებს იზომორფიზმს ფიბრებზე. ანუ $\forall x\in U$ წერტილისთვის $\psi$ ასახვის შეზღუდვა ფიბრზე
			\[ \psi\rvert_{p^{-1}(x)}:p^{-1}(x)\to \{x\}\times V \]
			არის ვექტორული სივრცეების იზომორფიზმი.
			
			ვექტორული  ფიბრაციების ასახვა (ანუ მორფიზმი) $\phi:\zeta_1\to\zeta_2$ არის ფიბრაციების ასახვა, რომელიც დამატებით $k$-წრფივი ასახვაა ნებისმიერ ფიბრზე შეზღუდვისას.
		\end{definition}
	
	
	
		
		\subsection*{ოპერაციები ვექტორულ ფიბრაციებზე}
		
		ისევე როგორც ვექტორული სივრცეებისგან შეგვიძლია ახალი ვექტორული სივრცეების აგება (მაგალითად, პირდაპირი ჯამის ან ტენზორული ნამრავლის მეშვეობით), ასევე შეგვიძლია არსებული ვექტორული ფიბრაციებისგან ახლების კონსტრუირება. ეს ოპერაციები, როგორც წესი, სრულდება "ფიბრის დონეზე" (fiber-wise), ანუ ოპერაცია ტარდება თითოეულ ფიბრზე ცალ-ცალკე, შემდეგ კი მიღებული შედეგები ერთიანდება ახალ ტოტალურ სივრცეში, რომელსაც შესაბამისი ტოპოლოგია ენიჭება.
		
		სანამ ოპერაციებს განვიხილავდეთ, უნდა შემოვიტანოთ ორი ფუნდამენტური ცნება: იზომორფიზმი და პულბექი.
		
		\subsection*{ვექტორული ფიბრაციების იზომორფიზმი}
		
		ორი ვექტორული ფიბრაცია ითვლება "ერთნაირად", თუ ისინი იზომორფულია.
		
		\begin{definition}
			ვთქვათ $p_1: E_1 \to B$ და $p_2: E_2 \to B$ ორი ვექტორული ფიბრაციაა ერთი და იგივე $B$ ბაზაზე. \textbf{ფიბრაციების იზომორფიზმი} არის ფიბრაციების ასახვა $\Phi = (\bar{\phi}, \phi)$, სადაც $\phi = \id_B$ (ანუ ასახვა ბაზაზე იგივურია), ხოლო $\bar{\phi}: E_1 \to E_2$ არის ჰომეომორფიზმი, რომელიც ყოველ ფიბრზე $p_1^{-1}(b)$ ინდუცირებს ვექტორული სივრცეების წრფივ იზომორფიზმს $p_2^{-1}(b)$-ზე.
			
			თუ ასეთი იზომორფიზმი არსებობს, ვამბობთ, რომ $E_1$ და $E_2$ იზომორფული ფიბრაციებია და ვწერთ $E_1 \cong E_2$.
		\end{definition}
		
		განსაკუთრებით მნიშვნელოვანია ტრივიალური ფიბრაციის ცნება. ფიბრაცია $p:E \to B$ არის \textbf{ტრივიალური}, თუ ის $B \times V$ სახის ტრივიალური ფიბრაციის იზომორფულია, სადაც $V$ ვექტორული სივრცეა.
		
		\subsection*{ვექტორული ფიბრაციის პულბექი (Pullback)}
		
		პულბექი არის უმნიშვნელოვანესი კონსტრუქცია, რომელიც საშუალებას გვაძლევს, "გადმოვიტანოთ" ვექტორული ფიბრაცია ერთი ბაზიდან მეორეზე უწყვეტი ასახვის მეშვეობით.
		
		\begin{definition}
			ვთქვათ მოცემულია ვექტორული ფიბრაცია $p: E \to B$ და უწყვეტი ასახვა $f: B' \to B$. \textbf{$E$ ფიბრაციის პულბექი} $f$ ასახვით არის ახალი ვექტორული ფიბრაცია $f^*E \to B'$, რომლის ტოტალური სივრცე განისაზღვრება შემდეგნაირად:
			\[ f^*E = \{ (b', e) \in B' \times E \mid f(b') = p(e) \}. \]
			პროექცია $p': f^*E \to B'$ მოიცემა, როგორც $p'(b', e) = b'$.
		\end{definition}
		
		ინტუიციურად, ახალი ფიბრაციის $f^*E$ ფიბრი $b' \in B'$ წერტილის თავზე არის "ასლი" ძველი ფიბრაციის $E$ ფიბრისა $f(b') \in B$ წერტილის თავზე. ანუ, $(f^*E)_{b'} \cong E_{f(b')}$.
		
		პულბექი საშუალებას გვაძლევს განვსაზღვროთ ოპერაციები ფიბრაციებზე, რომლებიც სხვადასხვა ბაზისზეა მოცემული, მათი ერთ ბაზაზე "გადმოტანით". თუმცა, თუ ფიბრაციები თავიდანვე ერთსა და იმავე ბაზაზეა, ოპერაციები უფრო მარტივად განიმარტება.
		
		\subsection*{ძირითადი ოპერაციები ვექტორულ ფიბრაციებზე}
		
		ვთქვათ $p_1: E_1 \to B$ და $p_2: E_2 \to B$ ორი ვექტორული ფიბრაციაა ერთსა და იმავე $B$ ბაზისზე.
		
		\subsubsection{უიტნის ჯამი (Whitney Sum)}
		
		ეს არის ვექტორული სივრცეების პირდაპირი ჯამის ანალოგი.
		
		\begin{definition}
			$E_1$ და $E_2$ ფიბრაციების \textbf{უიტნის ჯამი} არის ვექტორული ფიბრაცია $E_1 \oplus E_2 \to B$, რომლის ტოტალური სივრცე არის ფიბრების პირდაპირი ჯამების დიზიუნქტური გაერთიანება:
			\[ E_1 \oplus E_2 = \bigsqcup_{b \in B} (E_1)_b \oplus (E_2)_b. \]
			მიღებული ფიბრაციის რანგია $\rank(E_1 \oplus E_2) = \rank(E_1) + \rank(E_2)$.
		\end{definition}
		
		\subsubsection{ტენზორული ნამრავლი}
		
		ეს კონსტრუქცია ვექტორული სივრცეების ტენზორული ნამრავლის განზოგადებაა.
		
		\begin{definition}
			$E_1$ და $E_2$ ფიბრაციების \textbf{ტენზორული ნამრავლი} არის ვექტორული ფიბრაცია $E_1 \otimes E_2 \to B$, რომლის ტოტალური სივრცეა:
			\[ E_1 \otimes E_2 = \bigsqcup_{b \in B} (E_1)_b \otimes (E_2)_b. \]
			მიღებული ფიბრაციის რანგია $\rank(E_1 \otimes E_2) = \rank(E_1) \cdot \rank(E_2)$.
		\end{definition}
		
		\subsubsection{სხვა მნიშვნელოვანი კონსტრუქციები}
		
		ზემოთხსენებული ოპერაციების გარდა, არსებობს სხვა სტანდარტული კონსტრუქციებიც:
		
		\begin{itemize}
			\item \textbf{ორადული ფიბრაცია (Dual Bundle):} $p: E \to B$ ფიბრაციისთვის, მისი ორადული ფიბრაცია $E^* \to B$ არის ფიბრაცია, რომლის ფიბრი თითოეულ $b \in B$ წერტილში არის $(E_b)^*$ (ორადული ვექტორული სივრცე). მისი რანგი $E$-ს რანგის ტოლია.
			
			\item \textbf{ჰომორფიზმების ფიბრაცია (Hom Bundle):} ორი ფიბრაციისთვის $E_1, E_2 \to B$, ფიბრაცია $\Hom(E_1, E_2) \to B$ აიგება ფიბრების დონეზე წრფივი ასახვების სივრცეებით $\Hom((E_1)_b, (E_2)_b)$. მართებულია იზომორფიზმი: $\Hom(E_1, E_2) \cong E_1^* \otimes E_2$.
			
			\item \textbf{გარე ხარისხები (Exterior Powers):} $p: E \to B$ ფიბრაციისთვის, $k$-ური გარე ხარისხის ფიბრაცია $\Lambda^k(E) \to B$ აგებულია $\Lambda^k((E_b))$ ფიბრებით. ეს კონსტრუქცია განსაკუთრებით მნიშვნელოვანია დიფერენციალურ გეომეტრიაში დიფერენციალური ფორმების შესწავლისას.
			
			\item \textbf{დეტერმინანტის ფიბრაცია (Determinant Bundle):} ეს არის $n$-ური გარე ხარისხი $\Lambda^n(E)$, სადაც $n = \rank(E)$. ეს არის წრფივი ფიბრაცია (ანუ 1-განზომილებიანი ფიბრებით), რომლის თვისებები მრავალნაირობის მნიშვნელოვან ტოპოლოგიურ ინფორმაციას ატარებს.
		\end{itemize}
		

	
	
	
	
	
		
	
	
	
	
	
	
	
	
	
	
	
	
	
	
	
	
	
	
	
	
	
	
	
	
		
		\newpage
		%======================================================================
		\part{გლობალური გეომეტრია და ტოპოლოგია}
		%======================================================================
		\addcontentsline{toc}{section}{გლობალური გეომეტრია და ტოპოლოგია}
		
		\subsection*{გაუს-ბონეს თეორემა}
		ეს არის დიფერენციალური გეომეტრიის ერთ-ერთი უმნიშვნელოვანესი შედეგი, რომელიც ერთმანეთთან აკავშირებს ზედაპირის \textbf{გეომეტრიასა} (სიმრუდე) და \textbf{ტოპოლოგიას} (ფორმა).
		
		\begin{definition}[ეილერის მახასიათებელი]
			კომპაქტური ზედაპირისთვის $S$, მისი \textbf{ეილერის მახასიათებელი} $\chi(S)$ არის ტოპოლოგიური ინვარიანტი. თუ ზედაპირი დავყოფთ პოლიგონებად, მაშინ $\chi = V - E + F$, სადაც $V$ წვეროების, $E$ წიბოების და $F$ წახნაგების რაოდენობაა.
			\begin{itemize}
				\item სფეროსთვის $\chi(S^2)=2$.
				\item ტორისთვის $\chi(T^2)=0$.
				\item $g$ ნახვრეტის მქონე ტორისთვის $\chi = 2-2g$.
			\end{itemize}
		\end{definition}
		
		\begin{theorem}[გაუს-ბონეს თეორემა]
			$S$ იყოს კომპაქტური, ორიენტირებადი რეგულარული ზედაპირი. მაშინ:
			\[
			\iint_S K \, dA = 2\pi \chi(S)
			\]
		\end{theorem}
		\begin{proof}[დამტკიცების იდეა]
			დამტკიცება საკმაოდ რთულია. ერთი მიდგომაა ზედაპირის დაფარვა გეოდეზიური სამკუთხედებით. თითოეული სამკუთხედისთვის მტკიცდება ლოკალური გაუს-ბონეს ფორმულა:
			\[ \iint_T K\,dA = (\alpha+\beta+\gamma) - \pi \]
			სადაც $\alpha, \beta, \gamma$ სამკუთხედის შიდა კუთხეებია. ამ ფორმულების შეკრებით მთელ ტრიანგულაციაზე და ეილერის ფორმულის გამოყენებით, მიიღება გლობალური შედეგი.
		\end{proof}
		
		\begin{corollary}
			თუ კომპაქტურ ზედაპირზე შესაძლებელია მეტრიკის შემოღება, რომლის გაუსის სიმრუდე ყველგან დადებითია ($K>0$), მაშინ ზედაპირი ტოპოლოგიურად სფეროს ეკვივალენტურია ($\chi > 0$). თუ $K<0$, მაშინ $\chi<0$ (მაგ. ორი ნახვრეტის მქონე ტორი). თუ $K=0$, მაშინ $\chi=0$ (ტორი).
		\end{corollary}
		
		\subsection{მინიმალური ზედაპირები}
		\begin{definition}[მინიმალური ზედაპირი]
			ზედაპირს ეწოდება \textbf{მინიმალური}, თუ მისი საშუალო სიმრუდე იგივურად ნულის ტოლია ($H=0$).
		\end{definition}
		ეს ნიშნავს, რომ $k_1 = -k_2$ (გარდა ბრტყელი წერტილებისა). ასეთი ზედაპირები ლოკალურად ამცირებენ ფართობს.
		\begin{example}[კატენოიდი და ჰელიკოიდი]
			\textbf{კატენოიდი} არის $y=\cosh x$ წირის $x$ ღერძის გარშემო ბრუნვით მიღებული ზედაპირი.
			\textbf{ჰელიკოიდი} არის ხვეული კიბის ფორმის ზედაპირი.
			ეს ორი ზედაპირი ლოკალურად იზომეტრიულია.
		\end{example}
		
		\subsection{განზოგადებები: რიმანის გეომეტრია}
		რიმანის გეომეტრია არის დიფერენციალური გეომეტრიის განზოგადება $n$-განზომილებიან მრავალნაირობებზე.
		\begin{itemize}
			\item \textbf{მრავალნაირობა} არის სივრცე, რომელიც ლოკალურად ევკლიდურ სივრცეს ჰგავს.
			\item \textbf{რიმანის მეტრიკა} არის პირველი ფუნდამენტური ფორმის ანალოგი, რომელიც თითოეულ მხებ სივრცეში სკალარულ ნამრავლს განსაზღვრავს.
			\item \textbf{სიმრუდე} განისაზღვრება რიმანის სიმრუდის ტენზორის საშუალებით, რომელიც გაუსის სიმრუდის განზოგადებაა.
		\end{itemize}
		რიმანის გეომეტრია არის მათემატიკური ენა, რომელზეც დაწერილია აინშტაინის \textbf{ფარდობითობის ზოგადი თეორია}, სადაც ოთხგანზომილებიანი სივრცე-დროის სიმრუდე აღიქმება როგორც გრავიტაცია.
		
		\subsection{სავარჯიშოები}
		\begin{exercise}
			გამოიყენეთ გაუს-ბონეს თეორემა, რომ იპოვოთ $R$ რადიუსის სფეროზე მდებარე გეოდეზიური სამკუთხედის კუთხეების ჯამი, თუ მისი ფართობია $A$.
		\end{exercise}
		\begin{exercise}
			დაამტკიცეთ, რომ კომპაქტური მინიმალური ზედაპირი $\mathbb{R}^3$-ში არ არსებობს.
		\end{exercise}
	
	
	\newpage
	

	
	\part{დანართი}
	
	
	\section{ზოგადი ტოპოლოგიის განსაზღვრებები}
	
	
	
	\subsection{ტოპოლოგიური სივრცის განსაზღვრება ღია და ჩაკეტილი სიმრავლეებით}
	
	\begin{definition}
		\label{open}
		ვთქვათ $X$ რაიმე სიმრავლეა; $\Omega$ არის $X$-ის ქვესიმრავლეების რაიმე კოლექცია, ისეთი რომ:
		
		ა) თვითონ $X$ და ცარიელი სიმრავლე $\emptyset$ ეკუთვნიან $\Omega$-ს;
		
		ბ) $\Omega $-ს ნებისმიერი ქვეკოლექციის ელემენტების გაერთიანება ეკუთვნის $\Omega$-ს.
		
		(ანუ თუ ამოვარჩევთ $\Omega $-დან ნებისმიერი რაოდენობის ელემენტებს, რადგან ეს ელემენტები $X$ ქვესიმრავლეებია, შეგვიძლია ავიღოთ მათი გაერთიანება. მოითხოვება, რომ მიღებული სიმრავლე ისევ $\Omega$ კოლექციაში იყოს.)
		
		გ) $\Omega$-ს ნებისმიერი \textbf{სასრული} ქვეკოლექციის ელემენტების თანაკვეთა ეკუთვნის $\Omega$-ს.
		
		(ანუ თუ ნებისმიერად ამოვარჩევთ $\Omega $-დან \textbf{სასრული} რაოდენობის ელემენტებს და ავიღებთ მათ თანაკვეთას, მიღებული სიმრავლე ისევ $\Omega$ კოლექციაში იყოს.)
		
		მაშინ ვამბობთ, რომ მოცემულია ტოპოლოგიური სივრცე, როგორც წყვილი $\big( X,\Omega \big)$.
	\end{definition} 
	
	\medskip
	
	როცა ეს გაუგებრობას არ გამოიწვევს და გასაგები იქნება რომელ კოლექციაზეა საუბარი, ჩვენ $\Omega$-ს სიმბოლოს გამოვტოვებთ. 
	
	თუ $X$ ტოპოლოგიური სივრცეა $\Omega$ ტოპოლოგიით, მაშინ ვიტყვით რომ $X$-ის ქვესიმრავლე $U$ \textbf{ღიაა} $X$-ში თუ $U$ ეკუთვნის $\Omega$ კოლექციას. ამ ტერმინოლოგიით ჩვენ შეგვიძლია ვთქვათ, რომ ტოპოლოგიური სივრცე არის $X$ სიმრავლე თავის ღია სიმრავლეების კოლექციასთან ერთად, ისე რომ $X$ და $\emptyset$ ორივე ღიაა, და ისე რომ ნებისმიერი რაოდენობის ღია სიმრავლეების გაერთიანება და ნებისმიერი სასრული რაოდენობის ღია სიმრავლეების თანაკვეთა ისევ ღიაა.
	
	\begin{example} 
		განსაზღვრეთ სხვადასხვა ტოპოლოგიები სამელემენტიან სიმრავლეზე $X=\{ a,b,c \} $. მოიყვანეთ ქვესიმრავლეების კოლექცია, რომელიც არაა ტოპოლოგია.
	\end{example}
	
	$X$-სიმრავლეზე ყველა ქვესიმრავლეების კოლექცია გვაძლევს \textbf{დისკრეტულ ტოპოლოგიას}. თუ ღია სიმრავლეებად გამოცხადებულია მხოლოდ $X$ და $\emptyset$, გვაქვს ე.წ. \textbf{ტრივიალური ტოპოლოგია}.
	
	\begin{example} 
		ვთქვათ $X$ არის სხივი $[0, \infty )$ და $\Omega$ კოლექციაში შედიან $\emptyset, X$ და სხივები $[a,+\infty )$ $a\geq 0$-თვის. აჩვენეთ რომ $(X,\Omega)$ ტოპოლოგიური სივრცეა.
	\end{example}
	
	\begin{example} 
		ვთქვათ $X$ არის სიბრტყე და $\Omega$ კოლექციაში შედიან $\emptyset, X$ და ღია ბირთვები ცენტრით სათავეში. არის თუ არა $(X,\Omega)$ ტოპოლოგიური სივრცე?
	\end{example}
	
	\begin{example} 
		$X$ შედგება ოთხი ელემენტისგან: $X = \{a, b, c, d\}$. შემდეგი კოლექციებიდან რომელი აკმაყოფილებს ტოპოლოგიური სტრუქტურის აქსიომებს:
		
		(1) \,\,\,$\emptyset, X, \{a\}, \{b\}, \{a, c\}, \{a, b, c\}, \{a, b\};$
		
		(2) \,\,\, $\emptyset, X, \{a\}, \{b\}, \{a, b\}, \{b, d\};$
		
		(3) \,\,\, $\emptyset, X, \{a, c, d\}, \{b, c, d\}$.	
	\end{example}
	
	\subsection{ტოპოლოგიური სივრცეების სხვა მაგალითები}
	
	\fbox{\begin{minipage}{38em}
			ძალიან მნიშვნელოვანი მაგალითია ნამდვილ რიცხვთა წრფე $X=\mathbb{R}$, ხოლო ღია სიმრავლეები $(a,b)$, $a,b \in \mathbb{R}$ ინტერვალების ნებისმიერი გაერთიანებები. ამ ტოპოლოგიურ სტრუქტურას $\mathbb{R}$-ზე ეწოდება სტანდარტული ან კანონიკური ტოპოლოგია (შეამოწმეთ აქსიომები).
	\end{minipage}}
	
	\begin{example} 
		$X=\mathbb{R}$, ხოლო ღია სიმრავლეები კი $\mathbb{R}$ ყველა უსასრულო სიმრავლეები. არის ეს ტოპოლოგიური სივრცე?
	\end{example}
	
	\begin{example} 
		$X=\mathbb{R}$, ხოლო ღია სიმრავლეები კი $\mathbb{R}$ ყველა სასრული სიმრავლეების დამატებები. არის ეს ტოპოლოგიური სივრცე?
	\end{example}
	
	\begin{example} 
		$(X, \Omega)$ ტოპოლოგიური სივრცეა. განვიხილოთ $Y$ სიმრავლე რომელიც მიიღება $X$-ისგან ერთი ელემენტის დამატებით. არის თუ არა
		\[ \{ \{a \} \cup U | U \in \Omega \} \cup \{ \emptyset \} \]
		ტოპოლოგიური სტრუქტურა $Y$-ზე?
	\end{example}
	
	\begin{example} 
		არის თუ არა კოლექცია $\{\emptyset, \{0\}, \{0, 1\}\}$ ტოპოლოგიური სტრუქტურა $\{0, 1\}$-ზე ?
	\end{example}
	
	\bigskip
	
	\fbox{\begin{minipage}{38em}
			გავიხსენოთ რომ $(X, \Omega)$ ტოპოლოგიური სივრცისთვის, $X$-ის ელემენტები წერტილებია, $\Omega$-ს ელემენტები კი $X$-ის ღია სიმრავლეები.
	\end{minipage}}
	
	\begin{definition} 
		$X$ ტოპოლოგიური სივრცის \textbf{ჩაკეტილი} სიმრავლე ეწოდება $X$-ის რაიმე ღია სიმრავლის დამატებას.
	\end{definition}
	
	გავიხსენოთ სიმრავლეებზე დე მორგანის ფორმულები. ვთქვათ $\Gamma$ არის $X$ სიმრავლის ქვესიმრავლეების ნებისმიერი კოლექციაა, მაშინ გვაქვს
	\[
	X \setminus \bigcup_{A\in \Gamma} A =\bigcap_{A\in \Gamma}(X \setminus A ).
	\]
	\[
	X \setminus \bigcap_{A\in \Gamma} A =\bigcup_{A\in \Gamma}(X \setminus A ).
	\]
	
	\fbox{\begin{minipage}{35em}
			დე მორგანის ფორმულების გამოყენებით
			აჩვენეთ, რომ ტოპოლოგიის განსაზღვრება $X$-სიმრავლეზე შეიძლება არა მხოლოდ ღია სიმრავლებით, არამედ ჩაკეტილი სიმრავლეებითაც.
	\end{minipage}}
	
	\bigskip
	
	სხვა სიტყვებით, განსაზღვრება \ref{open} ექვივალენტურია შემდეგი განსაზღვრების
	
	\begin{definition} 
		ტოპოლოგიური სივრცე არის წყვილი $(X, \Gamma)$, სადაც $\Gamma$ არის $X$-სიმრავლიდან არჩეული "ჩაკეტილი" ქვესიმრავლეთა კოლექცია და სრულდება აქსიომები:
		
		\medskip
		ა) $X$ და $\emptyset$ ჩაკეტილია;
		
		\medskip
		ბ) თუ $\{ X_{i}\}_{i \in I}$ ჩაკეტილ ქვესიმრავლეთა ნებისმიერი კოლექციაა, მაშინ $\bigcap_{i \in I} X_{i}$ ჩაკეტილია;
		
		\medskip
		გ) თუ $\{ X_{i}\}_{i=1}^n$ ჩაკეტილ ქვესიმრავლეთა სასრული კოლექციაა, მაშინ $\bigcup_{i=1}^n X_{i}$ ჩაკეტილია.
	\end{definition}	
	
	\bigskip
	
	მაგალითით აჩვენეთ, რომ გ) აქსიომაში სასრულობის მოთხოვნა არსებითია.
	
	მითითება: მოიყვანეთ წრფის ღია ინტერვალის, მაგალითად $(-1,1)$-ის დაფარვა მონაკვეთებით. 	
	
	\bigskip
	
	\fbox{\begin{minipage}{38em}
			ფრანგ მათემატიკოსების (N. Bourbaki) მიბაძვით, ჩვენთვის სივრცის რაიმე წერტილის \textbf{მიდამო} არის ნებისმიერი სიმრავლე რომელიც მოიცავს ამ წერტილს.
	\end{minipage}}
	
	\subsubsection{ტოპოლოგიის ბაზისი}
	
	\fbox{\begin{minipage}{38em}
			ხშირად ტოპოლოგია განისაზღვრება არა ღია სიმრავლეების სრული კოლექციის აღწერით, არამედ გაცილებით მცირე ქვეკოლექციის საშუალებით.
	\end{minipage}}
	
	\begin{definition}
		$X$-სიმრავლეზე რაიმე $\Omega$ ტოპოლოგიის ბაზისი $\mathcal{B}$ არის $X$-სიმრავლის ქვესიმრავლეების, ე.წ. საბაზისო ელემენტების კოლექცია, ისე რომ $\Omega$ ტოპოლოგიით $X$-სიმრავლის ნებისმიერი არაცარიელი ღია სიმრავლე არის $\mathcal{B}$-ს ელემენტების გაერთიანება.
	\end{definition}
	
	\fbox{\begin{minipage}{25em}
			როდისაა რაიმე კოლექცია ბაზისი?
	\end{minipage}}
	
	\begin{definition}
		\label{base}
		$X$-სიმრავლეზე რაიმე ტოპოლოგიის ბაზისი $\mathcal{B}$ არის $X$-სიმრავლის ქვესიმრავლეების, ე.წ. საბაზისო ელემენტების კოლექცია, მაშინ და მხოლოდ მაშინ როცა:
		
		(1) ყოველი $x\in X$ ელემენტისთვის არსებობს ერთი მაინც საბაზისო ელემენტი $B$ რომელიც შეიცავს $x$-ს.
		
		(2) თუ $x$ ეკუთვნის ორი საბაზისო ელემენტის $B_1$ და $B_2$ თანაკვეთას, მაშინ არსებობს საბაზისო ელემენტი $B_3$, რომელიც შეიცავს $x$-ს და $B_3 \subset B_1 \cap B_2$.
	\end{definition}
	
	ორი საბაზისო ელემენტის თანაკვეთა საზოგადოდ არაა საბაზისო ელემენტი!
	
	თუ $\mathcal{B}$ აკმაყოფილებს ამ ორ პირობას, მაშინ განისაზღვრება $\mathcal{B}$-თი წარმოქმნილი ტოპოლოგია $\mathcal{T}$ შემდეგნაირად:
	$U \subset X$ სიმრავლე არის ღია, თუ $\forall x \in U$ არსებობს საბაზისო ელემენტი $ B \in \mathcal{B}$, ისე რომ $x\in B$ და $B\subset U$.
	
	\textbf{შევამოწმოთ} რომ ასე ბაზისი მართლა განსაზღვრავს ტოპოლოგიას.
	
	% \fbox{\begin{minipage}{38em}
	% ...
	%\end{minipage}}
	
	\subsubsection{ერთ სიმრავლეზე განსხვავებული ტოპოლოგიური სტრუქტურების იერარქია}
	
	\begin{definition} 
		ვთქვათ ერთი და იგივე სიმრავლეზე გვაქვს ორი ტოპოლოგიური სტრუქტურა. ვიტყვით რომ $(X,\Omega_1)$ უფრო ნაზია ანუ უფრო მდიდარი ვიდრე $(X, \Omega_2)$, თუ $\Omega_2 \subset \Omega_1$.
		სხვა სიტყვებით ღარიბი ტოპოლოგიური სტრუქტურის ყოველი ღია სიმრავლე მით უმეტეს ღიაა მდიდარ ტოპოლოგიაში. ცხადია შედარება ყოველთვის არაა შესაძლებელი, ანუ ნებისმიერი ორი ტოპოლოგიური სტრუქტურა არაა სადარი $\iff$ ($\Omega_1 \not\subset\Omega_2$ და $\Omega_2 \not\subset\Omega_1$).
	\end{definition}
	
	\subsection{ახალი ტოპოლოგიური სტრუქტურების აგების გზები}
	\subsubsection{ქვესივრცის ტოპოლოგია}
	
	\begin{definition}[კონსტრუქცია]
		ვთქვათ $(X,\Omega)$ ტოპოლოგიური სივრცეა და $A\subset X$ ქვესიმრავლე. მაშინ არსებობს ტოპოლოგია $(A,\Omega_A)$, სადაც $A$-ს ღია სიმრავლეებად გამოცხადებულია $X$-ის ღია სიმრავლეების თანაკვეთები $A$ სიმრავლესთან. ფორმალურად ეს ასე ჩაიწერება:
		\[
		\Omega_A=\{A\cap V \mid V \in \Omega \}.
		\]
		წყვილს $(A,\Omega_A)$ ეწოდება $(X,\Omega)$ ტოპოლოგიური სივრცის ქვესივრცე. კოლექცია $\Omega_A$ არის ქვესივრცის ტოპოლოგია (ან ფარდობითი ტოპოლოგია, ან ინდუცირებული ტოპოლოგია).
	\end{definition}
	(შეამოწმეთ ტოპოლოგიის აქსიომები).
	
	\begin{example}
		არის თუ არა ნახევრად ღია ინტერვალი $[0, 1)$ ღია სიმრავლე $[0, 2]$ სეგმენტში, როგორც წრფის ქვესივრცეში.
	\end{example}
	
	\fbox{\begin{minipage}{38em}
			სიმრავლის თვისება იყოს ღია (ან ჩაკეტილი) ფარდობითია. ანუ, ქვესივრცის ღია (ჩაკეტილი) სიმრავლე საზოგადოდ არაა ღია (ჩაკეტილი) მომცველ სივრცეში. მაგალითად, ერთადერთი ღია სიმრავლე $\mathbb{R}$-ში, რომელიც აგრეთვე ღიაა $\mathbb{R}^2$-ში არის ცარიელი სიმრავლე $\emptyset$.
	\end{minipage}}
	
	\begin{example}
		თუმცა მართალია შემდეგი: ღია ქვესივრცის ღია ქვესიმრავლე აგრეთვე ღიაა მომცველ სივრცეში. ანუ თუ $A$ ღიაა $X$-ში და $U\in \Omega_A$, მაშინ $U\in \Omega$.
	\end{example}
	
	\begin{example}
		ანალოგიურად, ჩაკეტილი ქვესივრცის ჩაკეტილი ქვესიმრავლე აგრეთვე ჩაკეტილია მომცველ სივრცეში.
	\end{example}
	
	\subsubsection{ნამრავლის ტოპოლოგია $X\times Y$-ზე}
	
	\begin{definition}
		ვთქვათ $X$ და $Y$ ტოპოლოგიური სივრცეებია. \textbf{ნამრავლის ტოპოლოგია} $X \times Y$ დეკარტულ ნამრავლზე განისაზღვრება მისი $\mathcal{B}$ ბაზისით, როგორც $U\times V$ ფორმის სიმრავლეების კოლექციით, სადაც $U\subset X$ და $V\subset Y$ ღია სიმრავლეებია.
	\end{definition}	
	
	\begin{example}
		გამოვიყენოთ განსაზღვრება \eqref{base} და შევამოწმოთ რომ კოლექცია $\mathcal{B}$ მართლაც ბაზისია.
	\end{example}
	
	\begin{example}
		ვაჩვენოთ რომ კოლექცია $\mathcal{B}$ არაა ტოპოლოგია $X \times Y$ დეკარტულ ნამრავლზე.
	\end{example}
	
	\begin{theorem}
		ვთქვათ $\mathcal{B}$ არის $X$-ზე ტოპოლოგიის ბაზისი და $\mathcal{C}$ არის $Y$-ზე ტოპოლოგიის ბაზისი, მაშინ კოლექცია
		\[
		\mathcal{D}=\{B\times C \mid B\in \mathcal{B},\, C\in \mathcal{C}\}
		\]
		არის $X \times Y$-ზე ტოპოლოგიის ბაზისი.
	\end{theorem}
	(დაამტკიცეთ)
	
	\subsubsection{დალაგების ტოპოლოგია}
	
	ვთქვათ $X$ მარტივად დალაგებული სიმრავლეა. ეს ნიშნავს რომ $X$-ზე არსებობს მარტივი დალაგების მიმართება $<$, ე.ი.
	\[
	\forall x, y \in X, x\neq y \implies x<y \text{ ან } y<x.
	\]
	არსებობს ნაწილობრივი დალაგების მიმართებაც, მაგალითად ქვესიმრავლის $\subset$ მიმართება სიმრავლეებში.
	
	$X$ მარტივად დალაგებული სიმრავლეში მოცემული ელემენტებისათვის
	$a<b$ განისაზღვრება შემდეგი ღია, ნახევრად ღია, და ჩაკეტილი ინტერვალები
	\begin{align*}
		&(a,b)=\{x\in X \mid a<x<b \},\\
		&[a,b)=\{x\in X \mid a\leq x<b \},\\
		&(a,b]=\{x\in X \mid a< x\leq b \},\\
		&[a,b]=\{x\in X \mid a\leq x\leq b \}.
	\end{align*}
	
	\begin{definition} 
		ვთქვათ $X$ მარტივად დალაგებული სიმრავლეა. განვიხილოთ $\mathcal{B}$ კოლექცია, $X$ სიმრავლის ყველა შემდეგი ტიპის ქვესიმრავლეების:
		
		(1) ყველა ღია ინტერვალები $(a,b)$;
		
		(2) ყველა ნახევრად ღია ინტერვალები $[a_0,b)$, თუ არსებობს მინიმალური ელემენტი $a_0 \in X$;
		
		(3) ყველა ნახევრად ღია ინტერვალები $(a,b_0]$, თუ არსებობს მაქსიმალური ელემენტი $b_0 \in X$.
		
		$\mathcal{B}$ კოლექცია არის ბაზისი $X$-ზე ე.წ. \textbf{დალაგების ტოპოლოგიის}.
	\end{definition}
	
	\begin{example}
		რიცხვით წრფეზე დალაგების ტოპოლოგია ემთხვევა სტანდარტულ ტოპოლოგიას.
	\end{example}
	
	\begin{example}
		სიბრტყის წერტილებზე, როგორც ნამდვილ რიცხვთა წყვილებზე, არსებობს მარტივი დალაგების მიმართება, ე.წ. ლექსიკოგრაფიული დალაგება. ამ დალაგების დროს სიბრტყის ორ წერტილს ასე ვადარებთ: ჯერ ვადარებთ პირველი კოორდინატების მიხედვით, ხოლო თუ პირველი კოორდინატები ტოლია, ვადარებთ მეორე კოორდინატების მიხედვით. დახატეთ სიბრტყეზე ამ დალაგების მიხედვით ღია, ნახევრად ღია და ჩაკეტილი ინტერვალები.
	\end{example}
	
	\begin{example}
		ვთქვათ $Y$ არის $X$-ის ქვესივრცე და $A$ არის $Y$-ის ქვესიმრავლე. აჩვენეთ რომ $A$-ზე ორი ტოპოლოგია, ერთი როგორც $Y$-ის ქვესიმრავლის და მეორე როგორც $X$-ის ქვესიმრავლის, ერთი და იგივეა.
	\end{example}
	
	\begin{example}
		ვთქვათ ორი ტოპოლოგიდან $X$-ზე $\mathcal{T}$ უფრო მდიდარია ვიდრე $\mathcal{T}'$. რა შეიძლება ითქვას შესაბამის ქვესიმრავლის ტოპოლოგიებზე $Y\subset X$ სიმრავლეზე?
	\end{example}
	
	\begin{example}
		ასახვას $f:X\to Y$ ეწოდება ღია, თუ ნებისმიერი $U\subset X$ ღია სიმრავლისათვის $f(U)$ ღიაა $Y$-ში. აჩვენეთ რომ პროექციები
		\[
		\pi_1:X\times Y \to X : \,\,\, \pi_1(x,y)=x,
		\]
		\[
		\pi_2:X\times Y \to Y : \,\,\, \pi_2(x,y)=y
		\]
		ღია ასახვებია.
	\end{example}
	
	\subsection{ფაქტორ-ტოპოლოგია}
	
	\subsubsection{ფაქტორ-ასახვა}
	ფაქტორ-ასახვა $p:X \to Y$ ეწოდება ისეთ ასახვას, რომელიც არის ზე-ასახვა (სიურექცია) და $V$ ღიაა $Y$-ში, თუ $p^{-1}(V)$ ღიაა $X$-ში.
	\begin{itemize}
		\item სიურექციული ასახვა $p$ არის ფაქტორ-ასახვა $\iff$ ($V$ ჩაკეტილია $Y$-ში, თუ $p^{-1}(V)$ ჩაკეტილია $X$-ში).
		\item $p$ არის ფაქტორ-ასახვა, თუ ის არის სიურექცია, უწყვეტი და ასახავს ღია გაჯერებულ სიმრავლეებს ღია სიმრავლეებში, სადაც $A$-ს $X$-ში უწოდებენ გაჯერებულს, თუ ეს არის $Y$-ის რაიმე ქვესიმრავლის წინარესახე.
		\item ფაქტორ-ასახვა არაა აუცილებლად ღია ან ჩაკეტილი, ფაქტორ-ასახვა რომელიც ღიაა არ არის აუცილებლად ჩაკეტილი და პირიქით.
	\end{itemize}
	ფაქტორ-ასახვების თვისებების ჩამოთვლამდე შევნიშნოთ, რომ ფაქტორ-ასახვის შეზღუდვა ქვედომენზე შეიძლება არ იყოს ფაქტორ-ასახვა, მაშინაც კი, თუ ის კვლავ სიურექციაა (და უწყვეტი).
	\begin{itemize}
		\item თუ $p:X \to Y$ არის ფაქტორ-ასახვა და $A$ არის გაჯერებული $p$-ის მიმართ, მაშინ თუ $A$ ღიაა ან ჩაკეტილია, ან $p$ ღიაა ან დახურულია, მაშინ $p|A:A\to p(A)$ არის ფაქტორ-ასახვა.
		\item თუ $p:X\to Y$ არის ღია ფაქტორ-ასახვა და $A\subset X$ ღიაა, მაშინ $p|A:A\to p(A)$ არის ღია ფაქტორ-ასახვა.
		\item ორი ფაქტორ-ასახვის კომპოზიცია არის ფაქტორ-ასახვა.
		\item ორი ფაქტორ-ასახვის ნამრავლი შეიძლება არ იყოს ფაქტორ-ასახვა. თუ ორივე ფაქტორ-ასახვა ღიაა, მაშინ ნამრავლი არის ღია ფაქტორ-ასახვა.
		\item თუ უწყვეტ ფუნქციას აქვს უწყვეტი მარჯვენა შებრუნებული, მაშინ ეს არის ფაქტორ-ასახვა.
		\item რეტრაქცია არის ფაქტორ-ასახვა. $X$-ის $A$ ქვესიმრავლეზე რეტრაქცია არის უწყვეტი ფაქტორ-ასახვა, $f:X\to A$ ისეთი, რომ თუ $x\in A$, $f(x)=x$.
	\end{itemize}
	
	\subsubsection{ფაქტორ-ტოპოლოგია და ფაქტორ-სივრცე}
	თუ $X$ არის სივრცე და $p:X\to Y$ არის სიურექცია, მაშინ $Y$-ზე არის ზუსტად ერთი ტოპოლოგია, ისეთი, რომ $p$ არის ფაქტორ-ასახვა. ეს არის ფაქტორ-ტოპოლოგია $Y$-ზე, რომელიც გამოწვეულია $p$-ით.
	
	დავუშვათ, $X^*$ არის $X$ სივრცის დაყოფა ფაქტორ-ტოპოლოგიით, რომელიც გამოწვეულია $p:X\to X^*$, სადაც $p(x)=A$ ისეთი, რომ $x\in A$, მაშინ $X^*$ ეწოდება $X$-ის ფაქტორ-სივრცეს.
	
	ფაქტორ-სივრცეზე შეიძლება ვიფიქროთ, როგორც სივრცეზე სადაც ხდება სხვადასხვა ნაწილების „შეწებება“. მაგალითად, თუ მონაკვეთის ბოლოებს შევაწებებთ ერთმანეთს, მივიღებთ ფაქტორ-სივრცეს წრეწირს.
	
	\bigskip
	\bigskip
	
	\subsection{სიმრავლის შიგა, გარე და საზღვრის წერტილები}
	
	\begin{definition}
		ვთქვათ $X$ ტოპოლოგიური სივრცეა, $A\subset X$ ქვესიმრავლე და $b\in X$. $b$ წერტილს ეწოდება:
		
		$A$-სიმრავლის \textbf{შიგა (interior) წერტილი}, თუ $b$ წერტილს აქვს მიდამო, რომელიც შედის $A$-ში;
		
		$A$-სიმრავლის \textbf{გარე (exterior) წერტილი}, თუ $b$ წერტილს აქვს მიდამო, რომელიც არ კვეთს $A$-ს (ანუ შედის $A$-ს დამატებაში $X \setminus A$-ში);
		
		$A$-სიმრავლის \textbf{საზღვრის (boundary) წერტილი}, თუ $b$ წერტილის ნებისმიერი მიდამო კვეთს $A$-ს და $A$-ს დამატებას.
	\end{definition}
	
	\subsubsection{სიმრავლის შიგანი (ინტერიორი)}
	
	\begin{definition}
		$X$-ტოპოლოგიური სივრცის $A$ ქვესიმრავლის \textbf{შიგანი} არის $X$-ის უდიდესი ღია სიმრავლე, რომელიც შედის $A$-ში, ანუ ღია $X$-ის სიმრავლე, რომელიც მოიცავს $A$-ს ყველა სხვა ღია ქვესიმრავლეს. $A$ ქვესიმრავლის შიგანი აღინიშნება ასე $\mathrm{Int}(A)$ ან უფრო დაწვრილებით $\mathrm{Int}_X A$.
	\end{definition}
	
	შიგანის თვისებებია:
	
	ტოპოლოგიური სივრცის ყოველ ქვესიმრავლეს აქვს შიგანი. ეს არის ყველა იმ ღია სიმრავლეების გაერთიანება, რომლებიც შედიან ამ სიმრავლეში.
	
	ტოპოლოგიური სივრცის ქვესიმრავლის შიგანი არის მისი ყველა შიგა წერტილების სიმრავლე.
	
	ტოპოლოგიური სივრცის ქვესიმრავლე ღიაა მაშინ და მხოლოდ მაშინ, როცა ეს სიმრავლე ემთხვევა თავის შიგანს.
	
	\begin{example}
		აჩვენეთ, რომ $\mathrm{Int} [0,1)=(0,1)$;
	\end{example}	
	
	\bigskip
	
	\subsection{კომპაქტური სივრცეები}
	
	\begin{definition}
		\label{compact}
		$X$ ტოპოლოგიური სივრცის ქვესიმრავლეების $\mathcal{A}$ კოლექციას ეწოდება $X$-ის დაფარვა, თუ $\mathcal{A}$-ს ელემენტების გაერთიანება ემთხვევა $X$-ს. დაფარვას ეწოდება ღია, თუ $\mathcal{A}$ კოლექციის ელემენტები ღია სიმრავლეებია.
		
		$X$ ტოპოლოგიური სივრცეს ეწოდება კომპაქტური, თუ მისი ნებისმიერი ღია დაფარვა $\mathcal{A}$ შეიცავს სასრულ ქვედაფარვას (ანუ $\mathcal{A}$-დან შეიძლება ავირჩიოთ სასრული რაოდენობის ღია სიმრავლეები, რომელთა გაერთიანება აგრეთვე ემთხვევა $X$-ს).
	\end{definition}
	
	\begin{example}
		\label{line-non-compact}
		არაკომპაქტური სივრცის მაგალითია $\mathbb{R}$ რიცხვითი წრფე: განვიხილოთ წრფის ღია დაფარვა ინტერვალებით $\mathcal{A}=\{(n,n+2) \mid n\in \mathbb{Z}\}$. ამ კოლექციიდან ერთი ელემენტიც რომ ამოვაკლოთ, მაგალითად $(0,2)$, რიცხვით წრფეზე დარჩება დაუფარავი წერტილი 1.
	\end{example}
	
	\begin{example}
		\label{0,1/n-compact}
		კომპაქტური სივრცის მაგალითია შემდეგ წერტილთა სიმრავლე რიცხვით წრფეზე:
		\[X=\{0\}\cup \{1/n \mid n\in \mathbb{Z}_{+}\}.\]
		განვიხილოთ წრფის ნებისმიერი $\mathcal{A}$ ღია დაფარვა ინტერვალებით. ამ დაფარვიდან ამოვირჩიოთ 0-ის მიდამო $U_0$. ცხადია $X$ წერტილთა მიმდევრობის ზღვარი არის 0, ამიტომ $U_0$ ღია სიმრავლე შეიცავს $X$-ის ყველა წერტილს, გარდა სასრული რაოდენობისა. $\mathcal{A}$-კოლექციიდან შევარჩიოთ ამ სასრული რაოდენობა წერტილების მიდამოები $U_1,\dots,U_m$. მივიღებთ $X$-ის ღია ქვედაფარვას
		$U_0,U_1,\dots,U_m$.
	\end{example}
	
	\subsubsection{კომპაქტურ სივრცეთა ნამრავლი}
	
	\begin{theorem}
		\label{product-of-compacts}
		კომპაქტურ სივრცეთა სასრული რაოდენობის ნამრავლი კომპაქტურია.
	\end{theorem}
	\begin{proof}
		საკმარისია დავამტკიცოთ, რომ ორი კომპაქტური სივრცის ნამრავლი კომპაქტურია (რატომ?).
		
		გვაქვს კომპაქტური სივრცეები $X$ და $Y$. დავაფიქსიროთ წერტილი $x_0 \in X$ და განვიხილოთ "ფენა" $x_0 \times Y$. თუ $U_{x_0}$ არის $x_0$ წერტილის რაიმე მიდამო, მაშინ $U_{x_0}\times Y$-ს ეწოდება "მილი" ანუ "მილისებრი მიდამო" $x_0 \times Y$ ფენის გარშემო.
		
		დამტკიცების პირველი ნაბიჯი არის შემდეგი: ნებისმიერი ღია სიმრავლე $N\subset X\times Y$ რომელიც მოიცავს $x_0 \times Y$ ფენას, მოიცავს აგრეთვე რაიმე მილს მის გარშემო. $X\times Y$-სივრცის ნებისმიერი ღია სიმრავლე და მათ შორის $N$ დაიფარება $U \times V$ სახის სიმრავლეებით (რადგან ესენია ბაზისის ელემენტები დეკარტული ნამრავლის ტოპოლოგიაში). რადგან $x_0 \times Y$ ფენა კომპაქტურია (ეს იგივე $Y$-ის კოპიოა) შეგვიძლია ავარჩიოთ მისი სასრული ქვედაფარვა
		\[U_1 \times V_1, U_2 \times V_2, \dots, U_n \times V_n.\]
		აქ ვიგულისხმოთ, რომ ყოველი $V_j$ არის $x_0$ წერტილის მიდამო (წინააღმდეგ შემთხვევაში გადავაგდოთ). ავიღოთ
		\[W=U_1\cap U_2 \cap \dots \cap U_n\]
		და მილის როლში $W\times Y$. მართლაც ამ მილის ნებისმიერი წერტილი $x\times y$ ეკუთვნის რომელიღაც $U_i\times V_i \subset N$: კერძოდ მეორე ინდექსს შევარჩევთ ისე, რომ $y\in V_i$, ხოლო პირველ ინდექსად ნებისმიერი ვარგა, რადგან $x\in W$ და ამიტომ ბოლო ტოლობით $x\in U_j$ ნებისმიერი ინდექსისთვის.
		
		მეორე ნაბიჯი. ავიღოთ $X\times Y$-ის ნებისმიერი ღია დაფარვა $\mathcal{A}$. ნებისმიერი ფენა $x\times Y$ კომპაქტურია და ის დაიფარება სასრული ქვედაფარვით $A_1, \dots, A_k$. ამ სიმრავლეების გაერთიანება $x\times Y$ ფენასთან ერთად მოიცავს მის რაიმე მილისებრ მიდამოს $W_x\times Y$.
		
		$\bigcup W_x$ არის $X$-ის დაფარვა და მისი კომპაქტურობის გამო შეგვიძლია ავარჩიოთ სასრული ქვედაფარვა $W_1, \dots, W_m$. ამრიგად $X\times Y$-ს ფარავს სასრული რაოდენობის მილი და თითოეული მილისთვის გვაქვს სასრული დაფარვა. რ.დ.გ.
	\end{proof}
	
	\subsubsection{$\mathbb{R}^n$-ევკლიდური სივრცის კომპაქტური სიმრავლეები}
	
	\begin{definition}
		\label{bounded}
		$\mathbb{R}^n$ ევკლიდური სივრცის $A$ ქვესიმრავლეს ეწოდება შემოსაზღვრული, თუ $\exists M>0$, ისეთი რომ, $\forall x,y \in A$ წერტილებს შორის მანძილი $d(x,y)\leq M$.
		აქ $d$ სტანდარტული ევკლიდური მანძილია: თუ $x=(x_1,\dots,x_n)$, $y=(y_1,\dots,y_n)$, მაშინ
		\[d(x,y)=\sqrt{(x_1-y_1)^2+\dots+(x_n-y_n)^2}.\]
	\end{definition}
	
	არსებობს აგრეთვე სხვა, ე.წ. კვადრატული მეტრიკა $\mathbb{R}^n$-ში:
	\[\rho(x,y)=\max\{|x_1-y_1|,\dots, |x_n-y_n|\}.\]
	ადვილი საჩვენებელია, რომ
	\[ \rho(x,y)\leq d(x,y) \leq \sqrt{n}\rho(x,y). \]
	
	\begin{theorem}
		\label{closed+bounded=compact}
		$\mathbb{R}^n$ ევკლიდური სივრცის ქვესიმრავლე კომპაქტურია მაშინ და მხოლოდ მაშინ, თუ ის ჩაკეტილია და შემოსაზღვრული ევკლიდური $d$ მეტრიკით, ან კვადრატული $\rho$ მეტრიკით.
	\end{theorem}
	\begin{proof}
		ბოლო უტოლობით ცხადია, რომ სიმრავლის $d$ მეტრიკით და $\rho$ მეტრიკით შემოსაზღვრულობა ექვივალენტურია.
		
		დავუშვათ $A\subset \mathbb{R}^n$ სიმრავლე კომპაქტურია. რადგან $\mathbb{R}^n$ ჰაუსდორფის სივრცეა, ამიტომ $A$ ჩაკეტილია. რატომაა შემოსაზღვრულიც? განვიხილოთ $\mathbb{R}^n$-ის დაფარვა ღია ბირთვების კოლექციით, ცენტრით სათავეში და რადიუსით $m$
		\[\{B(0,m) \mid m\in \mathbb{Z}_+\}.\]
		რადგან $A$ კომპაქტურია, მას დაფარავს რაიმე სასრული ქვეკოლექცია, ამიტომ
		\[ A\subset B(0,M) \]
		რაიმე $M$ რადიუსისთვის. ამიტომ $\forall x,y \in A$ წერტილებს შორის მანძილი $\rho(x,y)\leq 2M$.
		
		პირიქით, ვთქვათ $A$ ჩაკეტილია და შემოსაზღვრული. საკმარისია ვაჩვენოთ, რომ $A$ ქვესიმრავლეა რაღაც n-განზომილებიანი კუბის $[-P,P]^n$. მაშინ $A$ იქნება კომპაქტური, როგორც კომპაქტური სიმრავლის (კუბის, როგორც მონაკვეთების ნამრავლის) ჩაკეტილი ქვესიმრავლე. ამისათვის დავაფიქსიროთ $A$ სიმრავლის რაიმე წერტილი $x_0$. ვთქვათ $\rho(x_0,0)=b$. რახან $A$ შემოსაზღვრულია, $\rho(x,y)\leq N$, $\forall x,y \in A$. მეტრიკის სამკუთხედის წესით
		\[ \rho(x,0)\leq N+b, \quad \forall x\in A, \]
		ამიტომ თუ ავღნიშნავთ $P=N+b$, მაშინ $A$ სიმრავლე $[-P,P]^n$ კუბის ქვესიმრავლეა. რ.დ.გ.
	\end{proof}
	
	\subsection{სავარჯიშოები ტოპოლოგიის საკითხებზე}
	
	\begin{example} \label{ex:1-s}
		აჩვენეთ, რომ ტოპოლოგიის განსაზღვრება $X$-სიმრავლეზე შეიძლება არა მხოლოდ აქსიომებით ღია სიმრავლეებისთვის, არამედ შემდეგი აქსიომებითაც (ჩაკეტილი სიმრავლეებისთვის):
		
		ა) $X$ და $\emptyset$ ჩაკეტილია;
		
		ბ) თუ $\{ X_{i}\}_{i \in I}$ ჩაკეტილ ქვესიმრავლეთა ნებისმიერი კოლექციაა, მაშინ $\bigcap_{i \in I} X_{i}$ ჩაკეტილია;
		
		გ) თუ $\{ X_{i}\}_{i=1}^n$ ჩაკეტილ ქვესიმრავლეთა სასრული კოლექციაა, მაშინ $\bigcup_{i=1}^n X_{i}$ ჩაკეტილია.
		
		მითითება: განსაზღვრეთ ჩაკეტილი სიმრავლე, როგორც ღია სიმრავლის დამატება.
		
		მაგალითით აჩვენეთ, რომ გ) აქსიომაში სასრულობის მოთხოვნა არსებითია.
		
		მითითება: მოიყვანეთ წრფის ღია ინტერვალის, მაგალითად $(-1,1)$-ის დაფარვა მონაკვეთებით.
	\end{example}
	
	\begin{example} \label{ex:2-s}
		მოცემულ $X=\{a,b,c\}$ სიმრავლეზე შემოიტანეთ სამი განსხვავებული ტოპოლოგია.
	\end{example}
	
	\begin{example} \label{ex:3-s-hausdorff}
		$X=\{a,b,c\}$ სიმრავლეზე შემოიტანეთ ტოპოლოგია, რომელიც არაა ჰაუსდორფის.
	\end{example}
	
	\begin{example} \label{ex:4-s}
		აჩვენეთ, რომ თუ $X$ ქვესივრცეა $Y$-ში და $A$ ქვესიმრავლეა $Y$-ში, მაშინ $A$ სიმრავლეზე ინდუცირებული ორი ტოპოლოგია, როგორც $Y$-ის და $X$-ის ქვესიმრავლე, ერთმანეთს ემთხვევა.
	\end{example}
	
	\begin{example} \label{ex:5-s}
		განვიხილოთ $Y=[-1,1]$, როგორც $\mathbb{R}$ წრფის ქვესიმრავლე. შემდეგი სიმრავლეებიდან რომელია ღია $Y$-ში? რომელია ღია $\mathbb{R}$-ში?
		\begin{align*}
			& A=\{x \mid \tfrac{1}{2} < x < 1 \};\\
			& B=\{x \mid \tfrac{1}{2} < x \leq 1 \};\\
			& C=\{x \mid \tfrac{1}{2} \leq x < 1 \};\\
			& D=\{x \mid \tfrac{1}{2} \leq x \leq 1 \}.
		\end{align*}
	\end{example}
	
	\begin{example} \label{ex:6-s}
		იპოვეთ $\mathbb{R}^2$ სიბრტყის შემდეგი ქვესიმრავლეების საზღვარი და ინტერიორი:
		\begin{align*}
			& A=\{(x,y) \mid y=0\};\\
			& B=\{(x,y) \mid x\geq 0, y\neq 0\};\\
			& C=A\cup B;\\
			& D=\{(x,y) \mid x \text{ რაციონალურია}\};\\
			& E=\{(x,y) \mid 0< x^2-y^2\leq 1\};\\
			& F=\{(x,y) \mid x\neq 0, y\leq 1/x \}.
		\end{align*}
	\end{example}
	
	\begin{example} \label{ex:7-s}
		აჩვენეთ, რომ $f:\mathbb{R}\to \mathbb{R}$ ფუნქციისათვის უწყვეტობის $\epsilon-\delta$ განსაზღვრებიდან გამომდინარეობს ღია სიმრავლით განსაზღვრება.
	\end{example}
	
	\begin{example} \label{ex:8-s1}
		ვთქვათ $f:X\to Y$ უწყვეტი ფუნქციაა და $x$ არის $A\subset X$ სიმრავლის ზღვრული წერტილი. მართალია თუ არა, რომ აუცილებლად $f(x)$ წერტილი $f(A)$ სიმრავლის ზღვრული წერტილია?
	\end{example}
	
	\begin{example} \label{ex:8-s2}
		ვთქვათ $X$ და $X'$ ერთიდაიგივე სიმრავლეა ორი სხვადასხვა ტოპოლოგიით $\mathcal{T}$ და $\mathcal{T}'$. ვთქვათ $i$ იგივური ფუნქციაა.
		
		(1) აჩვენეთ რომ $i$ უწყვეტია $\iff$ $\mathcal{T}$ მდიდარია $\mathcal{T}'$-ზე;
		
		(2) აჩვენეთ რომ $i$ ჰომეომორფიზმია $\iff$ $\mathcal{T}'=\mathcal{T}$.
	\end{example}
	
	\begin{example} \label{ex:40-s}
		დაამტკიცეთ, რომ შემდეგი ტოპოლოგიური სივრცეები ჰომეომორფულია. დამტკიცებაში მოიყვანეთ ფორმულები ან გეომეტრიული აგებები.
		
		ა) წრეწირი $S^1$ და სამკუთხედი $\Delta ^1$;
		
		ბ) წრე და ბრტყელი სამკუთხედი;
		
		გ) სფერო $S^2$ და ტეტრაედრის ზედაპირი $\Delta^2$;
		
		დ) კუბის ზედაპირი და ტეტრაედრის ზედაპირი;
		
		ე) ღია ინტერვალი $(0,1)$ და ღია ინტერვალი $(0,10)$;
		
		ვ) ნებისმიერი ორი ღია ინტერვალი;
		
		ზ) ღია ინტერვალი და წრფე;
		
		თ) ღია წრე და სიბრტყე;
		
		ი) ღია რგოლი და სიბრტყე ერთი წერტილის გარეშე (გახვრეტილი სიბრტყე).
	\end{example}
	
	\begin{example} \label{ex:3-s-homeo}
		დაამტკიცეთ, რომ ევკლიდური სივრცის შემდეგი გარდაქმნა არის ჰომეომორფიზმი.
		
		ა) პარალელური გადატანა $\vec{p}$ ვექტორით;
		
		ბ) შეკუმშვა $\vec{p}$ ვექტორის მიმართულებით $\lambda$ კოეფიციენტით;
		
		გ) მობრუნება კოორდინატთა სათავის გარშემო $\phi$ კუთხით;
		
		დ) წრფივი გარდაქმნა, რომელიც მოცემულია არაგადაგვარებული $A$ მატრიცით.
	\end{example}
	
	\newpage
	\section{განსაზღვრებები კალკულუსიდან}
	
	სავარჯიშოების სახით გავიხსენოთ შემდეგი მასალა კალკულუსიდან.
	
	\subsubsection{ნამდვილი ფუნქციები $\mathbb{R}^3$-ზე}
	
	\bigskip
	
	\begin{example} \label{calc:5-s}
		$f = x^2y$ და $g = y\sin z$. გამოსახეთ შემდეგი ფუნქციები $x, y, z$ ცვლადებში: 
		ა) $fg^2$. \quad 
		ბ) $\frac{\partial f}{\partial x} g+\frac{\partial g}{\partial y} f$. \quad
		გ) $\frac{\partial^2(fg)}{\partial y\partial z}$. \quad 
		დ) $\frac{\partial}{\partial y}(\sin f)$.
	\end{example}
	
	\begin{example} \label{calc:6-s}
		იპოვეთ $f = x^2y - y^2z$ ფუნქციის მნიშვნელობა წერტილში:
		
		ა) $(1,1,1)$.\quad
		ბ) $(a,1,1-a)$.\quad
		გ) $(3,-1,1/2)$.\quad
		დ) $(t,t^2,t^3)$.
	\end{example}
	
	\begin{example} \label{calc:7-s}
		იპოვეთ $\frac{\partial f}{\partial x }$, თუ
		
		ა) $f = x\sin(xy) + y\cos(xz)$.\quad
		ბ) $f = \sin g$, $g = e^h$, $h = x^2 + y^2 + z^2$.
	\end{example}
	
	\begin{example} \label{calc:8-s}
		თუ $g_1,g_2,g_3$ ნამდვილი ცვლადის ფუნქციებია $\mathbb{R}^3$-ზე, მაშინ $f=h(g_1,g_2,g_3)$ არის ფუნქცია, ისეთი რომ $f(\mathbf{p})=h(g_1(\mathbf{p}),g_2(\mathbf{p}),g_3(\mathbf{p}))$, $\forall \mathbf{p} \in \mathbb{R}^3$. იპოვეთ $\frac{\partial f}{\partial x}$, თუ $h=x^2-yz$ და
		
		ა) $f = h(x + y, y^2, x + z)$.\quad
		ბ) $f = h(e^z, e^{x+y}, e^x)$.\quad
		გ) $f = h(x, -x, x)$.
	\end{example}
	
	\subsubsection{ნამდვილი ცვლადის ვექტორ-ფუნქცია}
	
	\begin{example} \label{calc:9-s}
		აჩვენეთ რომ $\alpha(t)=(3t,3t^2,2t^3)$ ვექტორ-ფუნქციის წარმოებული $\alpha'(t)$ მუდმივ კუთხეს ადგენს $y=0$, $z=x$ წრფესთან.
	\end{example}
	
	\begin{example} \label{calc:10-s}
		გააწარმოეთ ვექტორის სიგრძე $|\alpha(t)|$. გამოთვალეთ $|\alpha(t)|'$, თუ $\alpha(t)$ არის შემდეგი ვექტორი
		
		ა) $(e^t, e^{-t}, t^3)$; \quad
		ბ) $(\sin^3t, \cos(2t), 1-t^2)$; \quad 
		გ) $(1-\sin(2t), \cos^2(2t), -1/t)$.
	\end{example}
	
	\begin{example} \label{calc:11-s}
		გააწარმოეთ სკალარული ნამრავლი $\alpha(t)\cdot \beta(t)$.
	\end{example}
	
	\begin{example} \label{calc:12-s}
		დაამტკიცეთ, რომ თუ $t$ პარამეტრის რაღაც შუალედში $|\alpha(t)|$ მუდმივია, მაშინ $\alpha(t)\bot \alpha'(t)$. ახსენით ეს გეომეტრიულად. მართალია თუ არა ამ დებულების შებრუნებული?
	\end{example}
	
	\begin{example} \label{calc:13-s}
		გააწარმოეთ ვექტორული ნამრავლი $\alpha(t)\times \beta(t)$.
	\end{example}
	
	\subsubsection{ვექტორული ველი $\mathbb{R}^3$ სივრცეზე}
	\index{ვექტორული ველი}
	სივრცეზე მოცემული ვექტორული ველი სივრცის ყოველ წერტილს შეუსაბამებს ვექტორს ამ წერტილში. ქვემოთ $v_p$ აღნიშნავს $p$ წერტილში მოდებულ $v$ ვექტორს. სტანდარტული ბაზისის ვექტორები $U_1=(1,0,0)$, $U_2=(0,1,0)$, $U_3=(0,0,1)$ გვაძლევენ მუდმივ ვექტორულ ველებს $U_1(p)=(1,0,0)_p$, $U_2(p)=(0,1,0)_p$, $U_3(p)=(0,0,1)_p$, $\forall p\in \mathbb{R}^3$.
	
	\medskip
	
	\begin{example} \label{calc:14-s}
		ვთქვათ $v=(-2, 1, -1)$ და $w=(0, 1, 3)$.
		(ა) ჩაწერეთ $p$ წერტილში ვექტორი $3v_p - 2w_p$ როგორც $U_1(p)$, $U_2(p)$, $U_3(p)$ ვექტორების წრფივი კომბინაცია.
		(ბ) დახატეთ $p = (1, 1, 0)$ წერტილში მოდებული შემდეგი ვექტორები $v_p$, $w_p$, $-2v_p$, $v_p + w_p$.
	\end{example}
	
	\begin{example} \label{calc:15-s}
		მოცემულია ვექტორული ველები $V = xU_1 + yU_2$ და $W = 2x^2U_2 - U_3$. გამოთვალეთ ვექტორული ველი $W - xV$ და იპოვეთ მისი მნიშვნელობა წერტილში $p = (-1, 0, 2)$.
	\end{example}
	
	\begin{example} \label{calc:16-s}
		გამოსახეთ ვექტორული ველი $V$ სტანდარტული ფორმით $\sum_{i=1}^3 v_i U_i$, თუ
		
		ა) $2z^2U_1 = 7V + xyU_3$.
		
		ბ) $V(p) = (p_1, p_3 - p_1, 0)_p$ ნებისმიერ $p$ წერტილში.
		
		გ) $V = 2(xU_1 + yU_2) - x(U_1 - y^2U_3)$.
		
		დ) ყოველ $p$ წერტილში, $V(p)$ არის ვექტორი $(p_1, p_2, p_3)$ წერტილიდან $(1 + p_1, p_2p_3, p_2)$ წერტილამდე.
		
		ე) ყოველ $p$ წერტილში, $V(p)$ არის ვექტორი $p$ წერტილიდან სათავემდე.
	\end{example}
	
	\begin{example} \label{calc:17-s}
		მოცემულია $V = y^2U_1 - x^2U_3$ და $W = x^2U_1 - zU_2$ ვექტორული ველები. იპოვეთ ისეთი $f$ და $g$ ფუნქციები, რომ ვექტორული ველი $fV + gW$ შეიძლება გამოისახოს მხოლოდ $U_2$ და $U_3$-ის ტერმინებში.
	\end{example}
	
	\begin{example} \label{calc:18-s}
		მოცემულია $V_1 = U_1 - xU_3$, $V_2 = U_2$ და $V_3 = xU_1 + U_3$ ვექტორული ველები.
		
		ა) დაამტკიცეთ, რომ $V_1(p)$, $V_2(p)$, $V_3(p)$ ვექტორები წრფივად დამოუკიდებელია $\mathbb{R}^3$ სივრცის ნებისმიერ წერტილში.
		
		ბ) გამოსახეთ $xU_1 + yU_2 + zU_3$ ვექტორული ველი როგორც $V_1$, $V_2$, $V_3$ ველების წრფივი კომბინაცია.
	\end{example}
	
	\begin{example} \label{calc:19-s}
		იპოვეთ შემდეგი ვექტორული ველის დივერგენცია, როტორი, ნაკადი.
	\end{example}
	
	\subsubsection{მიმართულებით წარმოებული}
	\index{მიმართულებით წარმოებული}
	
	\begin{example}\label{calc:20-s}
		იპოვეთ ფუნქციის მიმართულებითი წარმოებული $V_p[f]$, თუ $f=x^2yz; \, P=(1, 1, 0);  \, V=(1, 0, -3).$
		
		\bigskip
		\textbf{ამოხსნა I ხერხი:}
		$P+tV=(1, 1, 0)+t(1, 0, -3)=(1+t, 1, -3t)$;
		
		$f(P+tV)=x^2yz=(1+t)^2 \cdot 1 \cdot (-3t) =-3t(1+2t+t^2)=-3t-6t^2-3t^3 $;
		
		$ \frac{d}{dt}(f(P+tV))\big|_{t=0} =(-3-12t-9t^2)\big|_{t=0}=-3$.
		
		\bigskip
		\textbf{ამოხსნა II ხერხი:}
		\[
		\mathrm{grad} f \big|_P =\left(\frac{\partial f}{\partial x};\frac{\partial f}{\partial y}; \frac{\partial f}{\partial z}\right)\bigg|_P = (2xyz, x^2z, x^2y)\big|_{(1, 1, 0)} = (0, 0, 1)
		\]
		\[
		V_p[f]= \mathrm{grad} f \big|_P \cdot V = (0, 0, 1) \cdot (1, 0, -3) = -3.
		\]
	\end{example}
	
	\begin{example} \label{calc:21-s}
		იპოვეთ ფუნქციის მიმართულებითი წარმოებული $V_p[f]$, თუ
		
		ა) $f=x^2z^2; \, P=(1, 2, -3);  \, V=(1, 0, -2).$
		
		ბ) $f=x^3y^2z; \, P=(0, -1, 4);  \, V=(1, -3, 2);$
		
		გ) $f=e^x \cos y; \, P=(2, 0, -1);  \, V=(2, -1, 3);$
	\end{example}
	
	\begin{example} \label{calc:22-s}
		მოცემულია $p$ წერტილში მოდებული ვექტორი $v_p\in \mathbb{R}^3$, სადაც
		$v = (2, -1, 3)$ და $p = (2, 0, -1)$. იპოვეთ მიმართულებითი წარმოებული $v_p[f]$, თუ ა) $f = y^2z$.\,\,\, ბ) $f = x^7$. \,\,\,გ) $f = e^x\cos y$.
	\end{example}
	
	\begin{example} \label{calc:23-s}
		ვთქვათ $V = y^2U_1 - xU_3$ და $f = xy$, $g = z^3$. გამოთვალეთ ფუნქციები
		
		ა) $V[f]$.\quad ბ) $V[g]$.\quad გ) $V[fg]$. \quad დ) $fV[g] - gV[f]$.\quad ე) $V[f^2 + g^2]$.\quad ვ) $V[V[f]]$.
	\end{example}
	
	\begin{example} \label{calc:24-s}
		დაამტკიცეთ იგივობა $V = \sum_{i=1}^3 V[x_i]U_i$, სადაც $x_1$, $x_2$, $x_3$ საკოორდინატო ფუნქციებია.
	\end{example}
	
	\begin{example} \label{calc:25-s}
		დაამტკიცეთ, რომ თუ $V[f] = W[f]$ ნებისმიერი $f$ ფუნქციისთვის $\mathbb{R}^3$-ზე, მაშინ $V = W$.
	\end{example}
	
	\bigskip
	\newpage
	
	\subsection{ტეილორის თეორემა მრავალი ცვლადის ფუნქციისთვის}
	
	\begin{theorem}
		განვიხილოთ $n$ ცვლადის ფუნქცია $g:\mathbb{R}^n\to \mathbb{R}$. ვთქვათ ფუნქციას აქვს $k+1$ რიგის უწყვეტი კერძო წარმოებულები რაიმე ჩაკეტილ ბირთვზე $B(p)$ ცენტრით $p\in \mathbb{R}^n$. მაშინ სამართლიანია ტეილორის თეორემა ნაშთი წევრის ინტეგრალური ფორმით
		\[
		g(x)=g(p)+\sum_{1 \leq |\alpha| \leq k} \frac{1}{\alpha!}(D^{\alpha}g)(p)(x-p)^{\alpha} + \sum_{|\alpha|=k+1}\frac{k+1}{\alpha!}(x-p)^{\alpha}\int_0^1(1-t)^k(D^{\alpha}g)((1-t)p+tx)dt.
		\]
		აქ გამოყენებულია სიმბოლოები
		\begin{align*}
			&\alpha=(\alpha_1,\alpha_2,\dots,\alpha_n);\\
			&|\alpha|=\alpha_1+\alpha_2+\dots+\alpha_n;\\
			&\alpha!=\alpha_1!\alpha_2!\dots\alpha_n!;\\
			&\binom{j}{\alpha}=
			\binom{j}{\alpha_1\alpha_2\dots\alpha_n}=
			\frac{j!}{\alpha_1!\alpha_2!\dots\alpha_n!};\\
			&x=(x_1,x_2,\dots,x_n), \,\,\, p=(p_1,p_2,\dots,p_n);\\
			&(x-p)^{\alpha}=(x_1-p_1)^{\alpha_1}(x_2-p_2)^{\alpha_2}\dots(x_n-p_n)^{\alpha_n};\\
			&D^{\alpha}g=\frac{\partial^{|\alpha|}g}{\partial x_1^{\alpha_1}\partial x_2^{\alpha_2}\dots\partial x_n^{\alpha_n}}.
		\end{align*}
	\end{theorem}
	
	ჩვენ გვინდა დამტკიცება დავიყვანოთ $[0,1]$ მონაკვეთზე ერთი ცვლადის $f$ ფუნქციაზე $k+1$ რიგის უწყვეტი წარმოებულებით
	\begin{equation}
		\label{f-k}
		f(1)=f(0)+\frac{1}{1!}f'(0)+\dots+\frac{1}{k!}f^{(k)}(0)+
		\frac{1}{k!}\int_0^1(1-t)^kf^{(k+1)}(t)dt
	\end{equation}
	და შემდეგ გადავიდეთ მრავალი ცვლადის ფუნქციაზე
	\begin{equation}
		\label{g}	
		g(u(t))=f(t), \quad u(t)=(1-t)p+tx.
	\end{equation}
	
	ფორმულა \eqref{f-k} მტკიცდება ინდუქციით, დაწყებული $k=0$ შემთხვევით.
	\begin{align*}
		f(1) &= f(0)+\int_0^1 f'(t)dt \quad \text{(ნაწილობითი ინტეგრება)} \\
		&= f(0) + \left[-(1-t)f'(t)\right]\bigg|_0^1 - \int_0^1 \left(-(1-t)f''(t)\right)dt \\
		&= f(0)+f'(0)+\int_0^1(1-t)f''(t)dt.
	\end{align*}
	
	\eqref{f-k} და \eqref{g} ტოლობებით გვაქვს
	\begin{equation}
		\label{f^k}	
		g(x)=f(1)=f(0)+\frac{1}{1!}f'(0)+\dots+\frac{1}{k!}f^{(k)}(0)+ \frac{1}{k!}\int_0^1(1-t)^kf^{(k+1)}(t)dt
	\end{equation}
	რთული ფუნქციის გაწარმოებით გვაქვს
	\begin{align*}
		f^{(j)}(t) &= \frac{d^{j}}{dt^j}g(u(t)) = \frac{d^{j}}{dt^j}g((1-t)p+tx) \\
		&= \sum_{|\alpha|=j}\binom{j}{\alpha}(D^{\alpha}g)((1-t)p+tx)(x-p)^{\alpha}.
	\end{align*}
	
	ამ ფორმულით შესაბამისი წევრები ჩავსვათ \eqref{f^k} ტოლობაში და გავითვალისწინოთ, რომ
	\[\frac{1}{j!}\binom{j}{\alpha}=\frac{1}{\alpha!}.\]
	გვაქვს რ.დ.გ.
	
	კერძოდ, როცა $k=1$, $x=q$ გვაქვს
	\begin{equation}
		\label{k=1}
		g(q)=g(p)+\sum_{i=1}^{n}\frac{\partial g}{\partial x_i}(p)(q_i-p_i) + \sum_{i,j=1}^{n}(q_i-p_i)(q_j-p_j)\int_0^1(1-t)\frac{\partial^2 g}{\partial x_i\partial x_j}((1-t)p+tq)dt.
	\end{equation}
	
	\bigskip
	\newpage
	
	\section{ასახვები $\mathbb{R}^n \to \mathbb{R}^m$}
	
	განვიხილოთ ასახვები $\mathbb{R}^n$ ევკლიდური სივრციდან $\mathbb{R}^m$–ში.
	
	მაგალითად, თუ $n=3$, $m=1$, მაშინ გვაქვს სამი ნამდვილი ცვლადის ნამდვილი ფუნქცია, თუ $n=1$ და $m=3$, გვაქვს ერთი ნამდვილი ცვლადის ვექტორ-ფუნქცია, ანუ წირი $\mathbb{R}^3$-ში.
	
	ძირითადი დაკვირვება $F:\mathbb{R}^n\to \mathbb{R}^m$ ასახვის შესახებ არის ის, რომ $F$ მთლიანად აღიწერება $\mathbb{R}^n$–ზე განსაზღვრული ნამდვილი საკოორდინატო ფუნქციებით $f_1,\dots,f_m$, ანუ
	\[F(x_1,\dots,x_n)=(f_1(x_1,\dots,x_n),\dots,f_m(x_1,\dots,x_n)).\]
	
	მაგალითად ელიფსი მოიცემა ასახვით $r(t)=(a\cos t, b\sin t)$. მეორე რიგის ზედაპირის, ელიფსოიდის, პარამეტრიზაციაა
	\[r(\theta, \phi)=(a\sin\theta\cos\phi, b\sin\theta\sin\phi, c\cos\theta).\]
	
	\index{საკოორდინატო ფუნქციები}
	\index{გლუვი ასახვა}
	\index{დიფერენცირებადი ასახვა}
	
	\begin{definition}
		\label{tangent-map-def}
		მოცემული $F:\mathbb{R}^n \to \mathbb{R}^m$ ფუნქციის მნიშვნელობა $p\in \mathbb{R}^n$ წერტილში არის $F(p) \in \mathbb{R}^m$, ამიტომ
		\[ F(p)=(f_1(p), f_2(p), \dots, f_m(p)), \]
		სადაც $f_1, f_2, \dots, f_m$ არის $\mathbb{R}^n$–ზე ნამდვილი ფუნქციები. მათ $F$-ის საკოორდინატო ფუნქციები ეწოდება. \textbf{$F$-ს ეწოდება დიფერენცირებადი} თუ თითოეულ $f_i$-ს გააჩნია კერძო წარმოებულები ყოველი $n$ ცვლადის მიმართ. $F$ \textbf{ფუნქციის მხები ასახვა} ანუ \textbf{დიფერენციალი} $F_*$ არის წრფივი ასახვა $F_*:\mathbb{R}^n \to \mathbb{R}^m$, ისე რომ
		\[ F_*(v)=v[F]=(v[f_1], v[f_2], \dots, v[f_m]) \text{ არის მხები ვექტორი } F(p)\text{-ში}. \]
	\end{definition}
	
	\textbf{მაგალითად} $F:\mathbb{R}^2\to \mathbb{R}^3$, $F=(x,y,xy)$, $p=(1,2)$, $v=(4,5)$;
	\begin{align*}
		&p+tv=(p_1+tv_1,p_2+tv_2)=(1+4t,2+5t);\\
		&F(p+tv)=(f_1(p+tv),f_2(p+tv),f_3(p+tv))=(1+4t,2+5t,(1+4t)(2+5t));\\
		&F_*(v)=(v[x],v[y],v[xy])=(4,5,13).
	\end{align*}
	
	\textbf{$F$-ს ეწოდება გლუვი} თუ თითოეულ $f_i$-ს გააჩნია ნებისმიერი რიგის კერძო წარმოებულები ყოველი $n$ ცვლადის მიმართ.
	
	გავიხსენოთ წრფივი ალგებრიდან, რომ წრფივი ასახვა და მისი მატრიცა რაიმე ბაზისში განისაზღვრება ამ ასახვის მნიშვნელობებით საბაზისო ვექტორებზე. კერძოდ, ჩვენს შემთხვევაში:
	\begin{align*}
		& F_*(U_1)=\left(\frac{\partial f_1}{\partial x_1},\frac{\partial f_2}{\partial x_1}, \dots,\frac{\partial f_m}{\partial x_1}\right);\\
		& \dots \\
		& F_*(U_n)=\left(\frac{\partial f_1}{\partial x_n},\frac{\partial f_2}{\partial x_n}, \dots, \frac{\partial f_m}{\partial x_n}\right);
	\end{align*}
	
	ამიტომ $F_*$-ის მატრიცა, რომელსაც $F$ ასახვის იაკობის მატრიცა ეწოდება, $U_1,U_2,\dots,U_n$ ბაზისში არის $m\times n $ მატრიცა:
	\[
	J=
	\begin{bmatrix}
		\frac{\partial f_1}{\partial x_1} & \frac{\partial f_1}{\partial x_2} & \dots & \frac{\partial f_1}{\partial x_n} \\
		\frac{\partial f_2}{\partial x_1} & \frac{\partial f_2}{\partial x_2} & \dots
		& \frac{\partial f_2}{\partial x_n} \\
		\vdots & \vdots & \ddots & \vdots \\
		\frac{\partial f_m}{\partial x_1} & \frac{\partial f_m}{\partial x_2}
		& \dots & \frac{\partial f_m}{\partial x_n} \\
	\end{bmatrix}.
	\]
	
	\medskip
	
	\begin{definition}
		\label{regularmap}
		$F:\mathbb{R}^n\to \mathbb{R}^m$, $n\leq m$ \textbf{ასახვას ეწოდება რეგულარული}, იმ პირობით, რომ ნებისმიერ $p$ წერტილში მისი მხები ასახვა $F_*$ არის ჩადგმა (ინექცია), ანუ $F_*(v_p)=0 \implies v_p = 0$, ან რაც იგივეა, $p$ წერტილში $F$-ის (დიფერენციალის) იაკობის მატრიცას აქვს რანგი $n$.
	\end{definition}
	\index{იაკობის მატრიცა}
	\index{რეგულარული ასახვა}
	
	ეხლა განვიხილოთ შემთხვევა $m\leq n$.
	
	\begin{definition}
		\label{submersion}
		ვთქვათ $U\subset \mathbb{R}^n$. გლუვ ასახვას $F:U\to \mathbb{R}^m$, $m\leq n$ ეწოდება ჩატვირთვა (submersion) $p\in U$ წერტილში, თუ დიფერენციალი $F_*:U\to \mathbb{R}^m$ არის ზე-ასახვა (სიურექცია), $p$-ს ეწოდება რეგულარული წერტილი; წინააღმდეგ შემთხვევაში $p$-ს ეწოდება კრიტიკული წერტილი. თუ $F$ ასახვა ყოველ წერტილში ჩატვირთვაა, მას ჩატვირთვა ეწოდება. სხვა სიტყვებით, $F$ ასახვა ჩატვირთვაა, თუ მისი დიფერენციალის რანგი მუდმივია და უდრის $m$-ს. კრიტიკული $p$ წერტილში ფუნქციის $F(p)$ მნიშვნელობას კრიტიკული მნიშვნელობა ეწოდება, ხოლო რეგულარულ წერტილში ფუნქციის მნიშვნელობას - რეგულარული მნიშვნელობა.
	\end{definition}
	\index{კრიტიკული წერტილი}
	\index{ჩატვირთვა}
	
	\begin{remark}
		ზოგიერთ წიგნში კრიტიკულ წერტილს უწოდებენ წერტილს, სადაც იაკობიანის მატრიცის რანგი არაა მაქსიმალური. ეს არ ეთანხმება ჩვენს განსაზღვრებებს. კერძოდ, როცა $m\leq n$, მაშინ წინააღმდეგობა არაა აქ მოყვანილ \ref{submersion} განსაზღვრებასთან. მაგრამ განსაზღვრება \ref{regularmap}-ის შემთხვევაში ყველა წერტილი კრიტიკული გამოდის, თუნდაც რანგი მაქსიმალური იყოს და უდრიდეს $n$-ს.
	\end{remark}
	
	\begin{example}
		მაგალითად, როცა გვაქვს გლუვი ასახვა $F:\mathbb{R}\to \mathbb{R}$, მაშინ $p$ წერტილის რეგულარობა ნიშნავს $F'(p)\neq 0$. მართლაც, თუ წერტილი კრიტიკულია
		\[dF=0\implies F'(x)dx=0\implies F'(x)=0.\]
	\end{example}
	
	\begin{example}
		გვაქვს გლუვი ასახვა $F:\mathbb{R}^3\to \mathbb{R}$, მაშინ $p$ წერტილის რეგულარობა ნიშნავს, რომ $F_x, F_y, F_z$ კერძო წარმოებულები ერთდროულად ყველა ნული არაა. მართლაც, $F$-ის დიფერენციალის მატრიცა $e_1,e_2,e_3\in \mathbb{R}^3$ სტანდარტულ ბაზისში არის
		\[J_F = [dF(e_1)\, dF(e_2) \,dF(e_3)]=[F_x \,F_y \,F_z].\]
	\end{example}
	
	\bigskip
	
	\subsection{ვექტორული ველის კოვარიანტული წარმოებული}
	კოვარიანტული წარმოებული არის მიმართულებითი წარმოებულის განზოგადება.
	
	ვთქვათ გვაქვს ორჯერ უწყვეტად დიფერენცირებადი რეგულარული ჩადგმა
	\[ \Phi:\mathbb{R}^d\supset U\to \mathbb{R}^n, \]
	ისე რომ, ყოველ $\Phi(p)\in \Phi(U)$ წერტილში მხები სივრცე მოჭიმულია ვექტორებზე
	\[ \{\frac{\partial \Phi}{\partial x^i}\bigg|_p, \quad i=1,\dots,d \} \]
	და სკალარული ნამრავლი $\Phi(U)$-ზე ნასესხებია მისი მომცველი $\mathbb{R}^n$ სივრცისგან.
	
	$\Phi(U)$-ზე მოცემული $V=\sum_j v^j\frac{\partial \Phi}{\partial x^j}$ მხები ვექტორული ველისთვის გვაქვს
	\[ \frac{\partial V}{\partial x^i}=\sum_j\left(\frac{\partial v^j}{\partial x^i}\frac{\partial \Phi}{\partial x^j}+v^j\frac{\partial^2 \Phi}{\partial x^i\partial x^j}\right). \]
	ფორმულის მეორე შესაკრები ვექტორი არ ძევს მხებ სიბრტყეში, ამიტომ დაიშლება მხებ და ნორმალურ კომპონენტებად. გამოვიყენოთ ქრისტოფელის სიმბოლოები $\Gamma_{ij}^k$:
	\[ \Gamma_{ij}^k=\frac{\partial e_j}{\partial x^i}\cdot e_k, \quad \text{სადაც } e_m=\frac{\partial \Phi}{\partial x^m}. \]
	შევნიშნოთ, რომ $\frac{\partial^2 \Phi}{\partial x^i\partial x^j}$ ვექტორის კომპონენტი მხები სივრციდან (ორთოგონალური პროექცია) არის
	\[\sum_k \Gamma_{ij}^k \frac{\partial \Phi}{\partial x^k},\]
	ამიტომ გვაქვს დაშლა მხებ და ნორმალურ კომპონენტებად
	\[\sum_j v^j\frac{\partial^2 \Phi}{\partial x^i\partial x^j}=\sum_{j,k}v^j \Gamma_{ij}^k \frac{\partial \Phi}{\partial x^k}+\vec{n}.\]
	
	ამრიგად, შემოვიღოთ
	\begin{definition}
		ზედაპირზე მოცემული $V$ ვექტორული ველის კოვარიანტული წარმოებული $\nabla_{e_i} =\nabla_i$, როგორც ვექტორული ველის ზედაპირის მხების $i$-ური საბაზისო ვექტორის მიმართულებით წარმოებული, ორთოგონალურად დაპროექტებული ზედაპირის მხებ სივრცეზე.
		\begin{equation}
			\nabla_i V := \left(\frac{\partial V}{\partial x^i}\right)^\top =
			\sum_{k} \left(\frac{\partial v^k}{\partial x^i} + \sum_j v^j \Gamma_{ij}^k \right)\frac{\partial \Phi}{\partial x^k}.
		\end{equation}
	\end{definition}
	
	\bigskip
	
	\subsection{მრუდწირული კოორდინატები $\mathbb{R}^n$-ში}
	განვიხილოთ რაიმე მიდამოში დეკარტის კოორდინატთა სისტემა ბაზისით $e_1, \dots, e_n$ და კოორდინატებით $x^1, \dots, x^n$. ვთქვათ მოცემული გვაქვს ახალი \textbf{მრუდწირული კოორდინატები}, ანუ
	\[y^1=y^1(x^1, \dots,x^n), \dots, y^n=y^n(x^1,\dots,x^n)\]
	უწყვეტი ფუნქციები, რომელთაც გააჩნიათ უწყვეტი კერძო წარმოებულები. გარდა ამისა, ვგულისხმობთ, რომ ამ ფუნქციებს გააჩნიათ უწყვეტი შებრუნებული ფუნქციები
	\[x^1=x^1(y^1, \dots,y^n), \dots, x^n=x^n(y^1,\dots,y^n),\]
	რომელთაც აგრეთვე აქვთ უწყვეტი კერძო წარმოებულები იგივე მიდამოში.
	
	განვიხილოთ $r=(x^1,\dots,x^n)$ რადიუს-ვექტორის კერძო წარმოებულები მრუდწირული კოორდინატების მიმართ.
	\begin{equation}
		\label{e-E}
		E_i=\frac{\partial r}{\partial y^i}=\frac{\partial}{\partial y^i}\left(\sum_{j=1}^n e_j x^j\right) = \sum_{j=1}^n e_j \frac{\partial x^j}{\partial y^i}.
	\end{equation}
	აქ უნდა გვახსოვდეს აინშტაინის შეთანხმება ჯამზე, ის რომ $e_j$ მუდმივია და $\frac{\partial x^j}{\partial y^i}$ კერძო წარმოებულები ადგენენ მატრიცას, რომელიც გადაუგვარებელია, რადგან $y^i=y^i(x^1, \dots,x^n)$, $i=1, \dots, n$ შებრუნებადია.
	
	თუ ავღნიშნავთ
	\begin{equation}
		\label{S,T}
		S^j_i=\frac{\partial x^j}{\partial y^i},\quad T^m_k=\frac{\partial y^m}{\partial x^k},
	\end{equation}
	გვექნება, რომ $T$ მატრიცა $S$ მატრიცის შებრუნებულია, ანუ
	\begin{equation}
		\label{S=T}
		\sum_{j=1}^n T^m_j S^j_i = \delta^m_i, \quad \sum_{m=1}^n S^i_m T^m_k = \delta^i_k.
	\end{equation}
	აქ $\delta^i_j$ სიმბოლოს კრონეკერის სიმბოლო ეწოდება
	\begin{align*}
		\delta^i_j=
		\begin{cases} 1 & \text{თუ } i=j \\
			0 & \text{თუ } i\neq j.
		\end{cases}
	\end{align*}
	
	ვექტორებს $E_1, \dots, E_n$ ეწოდება მრუდწირული კოორდინატების მხები ვექტორები $r=(x^1, \dots, x^n)$ წერტილში.
	
	\bigskip
	\textbf{მაგალითად}, $\mathbb{R}^3$-ში განვიხილოთ \textbf{სფერული კოორდინატები} $y=(y^1,y^2,y^3)=(r,\theta,\phi)$, ისე რომ
	\begin{equation}
		\label{spherical}
		x^1=r\sin\theta\cos\phi,\quad x^2=r\sin\theta\sin\phi,\quad x^3=r\cos\theta.
	\end{equation}
	საბაზისო ვექტორებია:
	\begin{align*}
		& E_1(y) = \sin\theta\cos\phi\,e_1 + \sin\theta\sin\phi\,e_2 + \cos\theta\,e_3, \\
		& E_2(y) = r(\cos\theta\cos\phi\,e_1 + \cos\theta\sin\phi\,e_2 - \sin\theta\,e_3), \\
		& E_3(y) = r\sin\theta(-\sin\phi\,e_1 + \cos\phi\,e_2).
	\end{align*}
	გამოვთვალოთ $g_{ij}(y)=E_i(y)\cdot E_j(y)$:
	\begin{align}
		& g_{11}(y)=E_1(y)\cdot E_1(y)=1, \\
		& g_{22}(y)=E_2(y)\cdot E_2(y)=r^2,\\
		& g_{33}(y)=E_3(y)\cdot E_3(y)=r^2\sin^2\theta,\\
		& g_{ij}(y)=E_i(y)\cdot E_j(y)=0, \text{ თუ } i\neq j.
	\end{align}
	\medskip
	
	ვექტორის კომპონენტების ზემოთ მოყვანილ გარდაქმნის ფორმულებს აქვს შემდეგი შედეგი: 
	$a=(a^1,a^2,a^3)$ ვექტორის კომპონენტები გარდაიქმნება $J$ იაკობის მატრიცით 
	\[ \bar{a}^{i}=\sum_{j=1}^n\frac{\partial \bar{x}^{i}}{\partial x^{j}}a^{j}. \]
	

	\subsubsection{ქრისტოფელის სიმბოლოები}
	თუ თითოეულ $E_i$ მხებ ვექტორს გავაწარმოებთ თითოეული მრუდწირული $y^j$ კოორდინატის მიმართ, მივიღებთ $n^2$ რაოდენობის ვექტორებს $\frac{\partial E_i}{\partial y^j}$. გამოვსახოთ ეს ვექტორები $(E_1,\dots,E_n)$ ბაზისში
	\begin{equation}
		\label{Gamma}
		\frac{\partial E_i}{\partial y^j}=\sum_{k=1}^n E_k\Gamma _{ij}^k.
	\end{equation}
	აქ აინშტაინის შეთანხმების თანახმად, ჯამი იგულისხმება. $\Gamma_{ij}^k$ კოეფიციენტებს ეწოდება \textbf{აფინური ბმულობა, ან მეორე ტიპის ქრისტოფელის სიმბოლოები}.
	
	\textbf{აფინური ბმულობის პირველი ფორმულა:}
	\begin{equation}
		\label{Gamma-ij}
		\Gamma^p_{ij} = \sum_{k=1}^n \frac{\partial y^p}{\partial x^k} \frac{\partial^2 x^k}{\partial y^i \partial y^j}.
	\end{equation}
	ამრიგად, აფინური ბმულობა სიმეტრიულია ქვედა ინდექსების მიმართ.
	\medskip
	
	\textbf{აფინური ბმულობის მეორე ფორმულა} მიიღება ასე:
	\begin{equation}
		\label{Gamma-ij-alt}
		\Gamma^p_{ij} = - \sum_{k=1}^n \frac{\partial x^k}{\partial y^i} \frac{\partial^2 y^p}{\partial y^j \partial x^k}.
	\end{equation}
	


	

		
	
\printindex
\end{document}







